% !TeX program = xelatex
% Compile twice: xelatex -interaction=nonstopmode main.tex
\documentclass[UTF8,10.5pt,a4paper,oneside,fontset=fandol,scheme=plain]{ctexbook}
% ==================== Art-Book Editorial Preamble for AI Series Volume 2 ====================
\usepackage[top=24mm,bottom=24mm,left=23mm,right=23mm,headheight=22pt,headsep=7mm]{geometry}
\usepackage{amsmath,amssymb,mathtools,amsthm}
\usepackage{booktabs,tabularx,array,multirow,longtable}
\usepackage{graphicx}
\usepackage{xcolor}
\usepackage{tikz}
\usepackage{pgfplots}
\pgfplotsset{compat=1.18}
\usetikzlibrary{arrows.meta,positioning,shapes.geometric,calc,decorations.pathreplacing,backgrounds}
\usepackage[most]{tcolorbox}
\usepackage{fancyhdr}
\usepackage{listings}
\usepackage{enumitem}
\usepackage{hyperref}

% ==================== Curated Gallery Color Palette ====================
\definecolor{ink}{HTML}{1B2430}        % Deep charcoal-slate ink
\definecolor{teal}{HTML}{0B5555}       % Refined verdigris teal
\definecolor{tealDark}{HTML}{083B3B}   % Deep gallery spine teal
\definecolor{tealLight}{HTML}{F2F8F8}  % Delicate ice-teal parchment
\definecolor{tealLine}{HTML}{A3C6C5}   % Soft architectural rule
\definecolor{tealMist}{HTML}{DCEBEB}   % Watermark chapter numeral tint
\definecolor{navy}{HTML}{1D3B5E}       % Muted Prussian indigo
\definecolor{navyLight}{HTML}{F2F5FA}  % Soft indigo tint
\definecolor{ochre}{HTML}{9E4A22}      % Burnished terracotta-ochre
\definecolor{gold}{HTML}{B97A24}       % Fine burnished amber-gold accent
\definecolor{ochreLight}{HTML}{FCF5F0} % Warm parchment tint
\definecolor{plum}{HTML}{5B3A6B}       % Muted editorial plum
\definecolor{plumLight}{HTML}{F8F4FB}  % Soft plum tint
\definecolor{slate}{HTML}{526071}      % Muted architectural gray
\definecolor{codebg}{HTML}{F7F9FB}     % Clean light code canvas
\definecolor{codeborder}{HTML}{D4DCE5}

\hypersetup{
  unicode=true,
  colorlinks=true,
  urlcolor=teal,
  linkcolor=teal,
  citecolor=teal,
  pdftitle={天书中的智能 · 第2卷：AI概率论、数理统计、信息论与首个回归/分类闭环},
  pdfauthor={天书中的智能：人工智能的七卷手记}
}

% ==================== Editorial Chapter Openers & Typography ====================
\ctexset{
  contentsname={目\quad 录},
  chapter={
    name={第,章},
    number=\arabic{chapter},
    format=\LARGE\sffamily\bfseries\color{tealDark},
    nameformat=\small\sffamily\bfseries\color{gold},
    aftername=\hfill{\fontsize{38}{38}\selectfont\sffamily\bfseries\color{tealMist}\ifnum\value{chapter}<10 0\fi\thechapter}\par\vspace{-4mm},
    aftertitle=\par\vspace{2.8mm}{\noindent\color{gold}\rule{26mm}{1.5pt}}{\color{tealLine}\rule{\dimexpr\linewidth-26mm\relax}{0.55pt}}\par,
    beforeskip=6pt,
    afterskip=10pt
  },
  section={
    format=\Large\sffamily\bfseries\color{teal},
    beforeskip=14pt plus 3pt minus 2pt,
    afterskip=6pt plus 1pt
  },
  subsection={
    format=\large\sffamily\bfseries\color{ink},
    beforeskip=10pt plus 2pt minus 1pt,
    afterskip=4pt plus 1pt
  },
  subsubsection={
    format=\normalsize\sffamily\bfseries\color{navy},
    beforeskip=7pt,
    afterskip=2pt
  }
}
\renewcommand{\lstlistingname}{代码}
\renewcommand{\figurename}{图}
\renewcommand{\tablename}{表}

% Artisanal Chapter Epigraph macro (2 args: quote + right-aligned attribution)
\newcommand{\chapEpigraph}[2]{%
  \vspace{-2mm}%
  \begin{tcolorbox}[
    enhanced,
    colback=ochreLight!45!white,
    colframe=ochreLight!45!white,
    borderline west={1.8pt}{0pt}{gold},
    boxrule=0pt,
    arc=0.8mm,
    left=4mm,right=3.5mm,top=2mm,bottom=2mm,
    before skip=0pt,after skip=8pt
  ]
    {\kaishu\normalsize\color{ink!85}“#1”}\par\vspace{1mm}%
    \hfill{\footnotesize\sffamily\color{slate}——\,#2}%
  \end{tcolorbox}%
}

% Exercise tier header macro
\newcommand{\exTier}[2]{%
  \par\vspace{3.8mm}\noindent
  \colorbox{tealLight}{\strut\;\sffamily\bfseries\footnotesize\color{tealDark}#1\;}\;%
  {\small\sffamily\bfseries\color{slate}#2}\par\nopagebreak\vspace{1.2mm}%
}
\newcommand{\exHint}[1]{\par\smallskip\noindent{\footnotesize\sffamily\color{navy}\textbf{【提示】}\,#1}}

\setlength{\parindent}{0pt}
\setlength{\parskip}{5pt plus 1.5pt minus 0.5pt}
\linespread{1.26}
\setlist[itemize]{leftmargin=1.8em,topsep=3pt,itemsep=2.5pt,parsep=1pt}
\setlist[enumerate]{leftmargin=2.0em,topsep=3pt,itemsep=2.5pt,parsep=1pt}
\setlength{\tabcolsep}{5.5pt}
\renewcommand{\arraystretch}{1.22}

% ==================== Running Headers & Footers ====================
\pagestyle{fancy}
\fancyhf{}
\fancyhead[L]{\small\sffamily\color{slate}天书中的智能 \;\textcolor{gold}{$\cdot$}\; 第 2 卷}
\fancyhead[R]{\small\sffamily\color{slate}\nouppercase{\leftmark}}
\fancyfoot[C]{\small\sffamily\color{slate}\textcolor{gold}{$\diamond$}\quad\textcolor{tealDark}{\thepage}\quad\textcolor{gold}{$\diamond$}}
\renewcommand{\headrulewidth}{0.4pt}
\renewcommand{\headrule}{\hbox to\headwidth{\color{tealLine}\leaders\hrule height \headrulewidth\hfill}}
\renewcommand{\chaptermark}[1]{\markboth{第 \thechapter\ 章\quad #1}{}}

\fancypagestyle{plain}{
  \fancyhf{}
  \fancyfoot[C]{\small\sffamily\color{slate}\textcolor{gold}{$\diamond$}\quad\textcolor{tealDark}{\thepage}\quad\textcolor{gold}{$\diamond$}}
  \renewcommand{\headrulewidth}{0pt}
}

% ==================== Theorem & Pedagogical Box Environments ====================
\theoremstyle{definition}
\newtheorem{theorem}{定理}[chapter]
\newtheorem{lemma}[theorem]{引理}
\newtheorem{proposition}[theorem]{命题}
\newtheorem{corollary}[theorem]{推论}
\newtheorem{definition}[theorem]{定义}
\newtheorem{example}[theorem]{例}
\newtheorem{remark}[theorem]{注记}
\renewcommand{\proofname}{\textbf{\sffamily\color{teal}证明}}

\newtcolorbox{thmBox}[1][]{
  enhanced,breakable,
  colback=tealLight,
  colframe=tealLine,
  borderline west={2.4pt}{0pt}{teal},
  boxrule=0.45pt,
  arc=1.5mm,
  left=4mm,right=3.8mm,top=2.8mm,bottom=2.8mm,
  before skip=6.5pt,after skip=6.5pt,
  #1
}

\newtcolorbox{defBox}[1][]{
  enhanced,breakable,
  colback=navyLight,
  colframe=navy!30,
  borderline west={2.4pt}{0pt}{navy},
  boxrule=0.45pt,
  arc=1.5mm,
  left=4mm,right=3.8mm,top=2.8mm,bottom=2.8mm,
  before skip=6.5pt,after skip=6.5pt,
  #1
}

\newtcolorbox{insightBox}[1]{
  enhanced,breakable,
  colback=white,
  colframe=navy!35,
  borderline west={2.4pt}{0pt}{gold},
  coltitle=navy,
  fonttitle=\sffamily\bfseries\small,
  attach title to upper={\quad},
  title={【直觉与下游坐标】#1},
  boxrule=0.45pt,
  arc=1.2mm,
  left=4mm,right=3.8mm,top=2.6mm,bottom=2.6mm,
  before skip=6.5pt,after skip=6.5pt
}

\newtcolorbox{recallBox}[1]{
  enhanced,breakable,
  colback=codebg,
  colframe=codeborder,
  borderline west={2.2pt}{0pt}{slate},
  coltitle=slate,
  fonttitle=\sffamily\bfseries\small,
  attach title to upper={\quad},
  title={【前置速忆】#1},
  boxrule=0.4pt,
  arc=1.2mm,
  left=3.8mm,right=3.5mm,top=2.4mm,bottom=2.4mm,
  before skip=5.5pt,after skip=5.5pt
}

\newtcolorbox{failBox}[1]{
  enhanced,breakable,
  colback=ochreLight,
  colframe=ochre!35,
  borderline west={2.4pt}{0pt}{ochre},
  coltitle=ochre,
  fonttitle=\sffamily\bfseries\small,
  attach title to upper={\quad},
  title={【反例、数值陷阱与排错】#1},
  boxrule=0.45pt,
  arc=1.2mm,
  left=4mm,right=3.8mm,top=2.6mm,bottom=2.6mm,
  before skip=6.5pt,after skip=6.5pt
}

\newtcolorbox{extBox}[1]{
  enhanced,breakable,
  colback=plumLight,
  colframe=plum!30,
  borderline west={2.4pt}{0pt}{plum},
  coltitle=plum,
  fonttitle=\sffamily\bfseries\small,
  attach title to upper={\quad},
  title={【选读拓展 · 初读可跳过】#1},
  boxrule=0.45pt,
  arc=1.2mm,
  left=4mm,right=3.8mm,top=2.6mm,bottom=2.6mm,
  before skip=6.5pt,after skip=6.5pt
}

\newtcolorbox{chapIntro}[1]{
  enhanced,
  colback=tealLight,
  colframe=tealLine,
  borderline west={2.8pt}{0pt}{tealDark},
  boxrule=0.45pt,
  arc=1.5mm,
  left=4.2mm,right=4mm,top=3mm,bottom=3mm,
  before skip=2pt,after skip=10pt,
  title={\sffamily\bfseries\small\color{tealDark}#1},
  colbacktitle=tealLight,
  attach boxed title to top left={yshift=-1mm,xshift=2mm},
  boxed title style={empty,size=minimal}
}

% ==================== Code Listing Style ====================
\lstdefinestyle{pycode}{
  language=Python,
  basicstyle=\footnotesize\ttfamily\color{ink},
  keywordstyle=\bfseries\color{teal},
  commentstyle=\itshape\color{slate},
  stringstyle=\color{ochre},
  backgroundcolor=\color{codebg},
  frame=single,
  rulecolor=\color{codeborder},
  framesep=5pt,
  xleftmargin=4pt,
  xrightmargin=4pt,
  numbers=left,
  numberstyle=\tiny\ttfamily\color{slate!80},
  numbersep=6pt,
  columns=fullflexible,
  keepspaces=true,
  showstringspaces=false,
  breaklines=true,
  aboveskip=6.5pt,
  belowskip=6.5pt
}
\lstset{style=pycode}

% ==================== Math Shortcuts ====================
\newcommand{\R}{\mathbb{R}}
\newcommand{\C}{\mathbb{C}}
\newcommand{\N}{\mathbb{N}}
\newcommand{\E}{\mathbb{E}}
\newcommand{\Prob}{\mathbb{P}}
\newcommand{\Var}{\operatorname{Var}}
\newcommand{\Cov}{\operatorname{Cov}}
\newcommand{\Bias}{\operatorname{Bias}}
\newcommand{\DKL}{D_{\mathrm{KL}}}
\newcommand{\col}{\operatorname{col}}
\newcommand{\row}{\operatorname{row}}
\newcommand{\nul}{\operatorname{null}}
\newcommand{\rank}{\operatorname{rank}}
\newcommand{\diag}{\operatorname{diag}}
\newcommand{\tr}{\operatorname{tr}}
\newcommand{\spanop}{\operatorname{span}}
\newcommand{\argmin}{\operatorname*{arg\,min}}
\newcommand{\argmax}{\operatorname*{arg\,max}}
\newcommand{\cond}{\kappa}
\newcommand{\norm}[1]{\left\|#1\right\|}
\newcommand{\abs}[1]{\left|#1\right|}
\newcommand{\inner}[2]{\left\langle #1,\, #2 \right\rangle}
\newcommand{\had}{\odot}
\newcommand{\indep}{\perp\!\!\!\perp}

\makeatletter
\setlength{\@fpsep}{16pt}
\makeatother


\begin{document}
\color{ink}

% ==================== Cover Page (Gallery Art-Book Edition) ====================
\frontmatter
\thispagestyle{empty}
\begin{tikzpicture}[remember picture, overlay]
  % Architectural deep teal spine + burnished gold hairline + vertical monograph label
  \fill[tealDark] (current page.north west) rectangle ([xshift=11mm]current page.south west);
  \fill[gold] ([xshift=11mm]current page.north west) rectangle ([xshift=12.2mm]current page.south west);
  \node[rotate=90, anchor=center, font=\sffamily\bfseries\scriptsize, color=tealMist, opacity=0.85]
    at ([xshift=5.5mm]current page.west) {INTELLIGENCE FROM THE BOOK \quad $\cdot$ \quad VOLUME II};

  % Bespoke fine-art geometric probability & KL-divergence composition (top-right)
  \begin{scope}[shift={([xshift=-95mm, yshift=-20mm]current page.north east)}, scale=0.48, transform shape]
    % Soft concentric 2D Gaussian contour ellipses
    \begin{scope}[rotate=22, opacity=0.14]
      \draw[navy, line width=0.9pt] (0.5, 0.4) ellipse (2.6cm and 1.15cm);
      \draw[navy, line width=0.75pt] (0.5, 0.4) ellipse (1.85cm and 0.80cm);
      \draw[navy, line width=0.6pt] (0.5, 0.4) ellipse (1.1cm and 0.48cm);
    \end{scope}
    % Translucent bimodal vs unimodal KL curves
    \draw[slate!40, thin] (-2.8,-0.6) -- (2.8,-0.6);
    \fill[teal, opacity=0.09, domain=-2.7:2.7, samples=90]
      (-2.7,-0.6) -- plot (\x, {-0.6 + 1.85*(0.55*exp(-1.1*(\x+1.05)*(\x+1.05)) + 0.55*exp(-1.1*(\x-1.05)*(\x-1.05)))}) -- (2.7,-0.6) -- cycle;
    \draw[teal, opacity=0.35, line width=1.2pt, domain=-2.7:2.7, samples=90]
      plot (\x, {-0.6 + 1.85*(0.55*exp(-1.1*(\x+1.05)*(\x+1.05)) + 0.55*exp(-1.1*(\x-1.05)*(\x-1.05)))});
    \draw[gold, opacity=0.45, line width=1.1pt, dashed, domain=-2.7:2.7, samples=90]
      plot (\x, {-0.6 + 1.35*exp(-0.26*\x*\x)});
    \draw[ochre, opacity=0.42, line width=1.0pt, dotted, domain=-0.6:2.7, samples=60]
      plot (\x, {-0.6 + 1.55*exp(-1.15*(\x-1.05)*(\x-1.05))});
  \end{scope}
\end{tikzpicture}
\vspace*{1mm}
\noindent{\small\sffamily\bfseries\color{teal}天书中的智能 \;\textcolor{gold}{$\cdot$}\; 人工智能的七卷手记 \hfill 第 2 卷 \quad 2026 年 11 月版}
\vspace{1.5mm}
\noindent{\color{gold}\rule{30mm}{1.8pt}}{\color{tealLine}\rule{\dimexpr\linewidth-30mm\relax}{0.7pt}}
\vspace{6.5mm}

{\fontsize{24}{32}\selectfont\sffamily\bfseries\color{ink}AI 概率论、数理统计、\\信息论与首个回归/分类闭环}\par
\vspace{2.5mm}
{\large\sffamily\color{teal}用概率、熵与散度理解噪声与学习，亲手完成纯 NumPy 端到端模型闭环}\par
\vspace{4.5mm}
{\small\color{slate}对标 Stanford CS109 / MIT 6.3700 / CMU 36-401 / 北大概率建模，承接第 1 卷线代与微积分}\par

\vspace{5.5mm}

\begin{tcolorbox}[
  enhanced,
  colback=tealLight,
  colframe=tealLine,
  arc=2mm,
  boxrule=0.6pt,
  left=4.5mm,right=4.5mm,top=3mm,bottom=3mm
]
\textbf{\sffamily\color{tealDark}本卷定位：补齐 AI 专属概率统计底座，打通首个完整监督学习闭环}\par\vspace{1.2mm}
\small
传统考研概率偏重二维积分技巧，而人工智能真正高频依赖的是一套完全不同的概率语言。本卷承接第 1 卷的 NumPy、谱定理、矩阵微分与梯度下降，用 6 章打通“不确定性建模 $\to$ 信息论损失 $\to$ 端到端分类实战”：
\begin{itemize}
  \item \textbf{直击 AI 高频概率工具}：条件独立与 Explaining Away、多元高斯几何与条件分布、Jensen 与 Hoeffding 浓度不等式、蒙特卡洛采样与重要性采样。
  \item \textbf{从零讲透信息论核心}：熵、条件熵、互信息、交叉熵与 KL 散度（含正向 vs 反向 KL 行为差异），一次性打通分类损失、语言模型困惑度、决策树、VAE 与 DPO 的共同语言。
  \item \textbf{统一统计估计与回归/分类}：严格推导 MLE $\equiv$ 最小二乘、MAP $\equiv$ 岭/Lasso、Bernoulli MLE $\equiv$ 二元交叉熵 $\equiv$ 最小 KL 散度，并用 Bootstrap 量化模型置信区间。
  \item \textbf{端到端无框架实战 + 48 题四层精解}：纯 NumPy 实现防溢出逻辑/Softmax 回归、防数据泄漏流水线与完整评测指标，所有曲线与数据表均由本机实测生成。
\end{itemize}
\end{tcolorbox}

\vspace{4.5mm}

\noindent\textbf{\sffamily\color{tealDark}第 2 卷章节导览}
\vspace{1.5mm}

{\footnotesize\renewcommand{\arraystretch}{1.15}
\begin{tabularx}{\linewidth}{@{}p{15mm}p{47mm}X@{}}
\toprule
\textbf{章节} & \textbf{标题} & \textbf{核心能力、下游用途与代码实战} \\
\midrule
序\quad 章 & 本卷坐标与八校考点覆盖对照 & 为什么不能只靠应试概率、八校对标表与本卷知识依赖 DAG \\
第 1 章 & 概率建模、条件独立与 Bayes 推断 & 先验/似然/后验更新、罕见病反例、独立 vs 条件独立、朴素贝叶斯雏形 \\
第 2 章 & AI 核心概率分布与多元高斯 & Bernoulli/Categorical/Beta、多元高斯马氏距离、仿射封闭性与条件高斯 \\
第 3 章 & 极限定理、浓度不等式与蒙特卡洛 & Jensen 不等式、Markov/Chebyshev/Hoeffding 样本量界、CLT 与重要性采样 \\
第 4 章 & 信息论：熵、交叉熵与 KL 散度 & Shannon 熵、互信息、Gibbs 非负性、正向/反向 KL 差异与困惑度 \\
第 5 章 & 统计估计：MLE、MAP 与 Bootstrap & 偏差-方差分解、MLE 过拟合反例、Beta-Bernoulli 共轭与非参数置信区间 \\
第 6 章 & 纯 NumPy 回归/分类端到端闭环 & MLE/MAP、Sigmoid+MSE 反例、BCE/Softmax 稳定实现、防泄漏与实测扫参 \\
附\quad 录 & 全卷 48 道分层习题完整精解 & 第 1--6 章 A/B/C/D 四层习题分步推导、反例验算与代码自检基准 \\
\bottomrule
\end{tabularx}}

\vfill
\noindent{\color{tealLine}\rule{\linewidth}{0.5pt}}
{\footnotesize\color{slate}配套实验脚本 \texttt{code/vol2\_experiments.py} 与实测数据随本卷提供。上一卷：第 1 卷；下一卷：第 3 卷《经典人工智能：搜索、博弈、概率图与序列决策》。}

% ==================== Table of Contents ====================
\clearpage
\tableofcontents

% ==================== Main Matter ====================
\mainmatter
\chapter*{序章：本卷在七卷本中的坐标与八校考点对照}
\addcontentsline{toc}{chapter}{序章：本卷在七卷本中的坐标与八校考点对照}
\markboth{序章\quad 本卷坐标与八校考点对照}{}

\begin{chapIntro}{本章定位}
第 1 卷已经把 Python/NumPy 数值编程、AI 线性代数与多元微积分装配完毕。本卷（第 2 卷）专攻 AI 底座的另一半——\textbf{概率论、数理统计与信息论}，并在第 6 章将两卷数理工具合流，完成无框架纯 NumPy 的线性回归与逻辑/Softmax 回归端到端实战闭环。
\end{chapIntro}

\section*{0.1\quad 为什么不能只靠传统应试概率学 AI？}
\addcontentsline{toc}{section}{0.1\quad 为什么不能只靠传统应试概率学 AI？}

传统工科概率论（及考研数学一概率部分）侧重二维连续联合密度积分、顺序统计量与抽样分布查表，但现代人工智能论文与代码里天天出现的却是另一批核心结构：
\begin{enumerate}
  \item \textbf{高维结构化分布}：不是两维手算积分，而是用第 1 卷的实对称谱定理看懂\textbf{多元高斯分布 $\mathcal{N}(\mu,\Sigma)$} 的椭球几何、仿射封闭性与条件更新；
  \item \textbf{信息论语言}：没有\textbf{熵、交叉熵、KL 散度与互信息}，就无法真正理解分类损失函数、决策树信息增益、语言模型困惑度、VAE 的 ELBO 与大模型对齐（DPO/RLHF）中的 KL 惩罚；
  \item \textbf{有限样本保证与重采样}：通过\textbf{Jensen 与 Hoeffding 浓度不等式}量化经验误差逼近真实泛化误差所需的样本量，并用\textbf{蒙特卡洛与 Bootstrap} 在不假设正态分布时给出模型指标的置信区间；
  \item \textbf{概率建模直通监督学习}：看清最小二乘即高斯 MLE、岭回归即高斯 MAP、逻辑回归交叉熵即 Bernoulli MLE 与最小 KL 散度，并严格防范数据泄漏与不平衡准确率陷阱。
\end{enumerate}

\section*{0.2\quad 第 2 卷八校核心考点覆盖对照表}
\addcontentsline{toc}{section}{0.2\quad 第 2 卷八校核心考点覆盖对照表}

\begin{center}
\footnotesize
\begin{tabularx}{\linewidth}{@{}p{34mm}p{34mm}p{16mm}X@{}}
\toprule
\textbf{八校核心考点} & \textbf{主要对标课程} & \textbf{本卷位置} & \textbf{下游直接用途} \\
\midrule
条件独立、Explaining Away 与 Bayes 更新 & Stanford CS109 / 北大通班 & 第 1 章 & 第 3 卷贝叶斯网络、HMM 与朴素贝叶斯 \\
Bernoulli / Categorical / Beta / 多元高斯 & MIT 6.3700 / CS109 & 第 2 章 & 分类建模、第 4 卷 GMM-EM、第 5 卷 VAE 重参数化 \\
Jensen、Markov/Chebyshev/Hoeffding 界 & Berkeley CS70 / 清华 & 第 3 章 & 第 4 章 KL 非负性、第 4 卷泛化界与 EM 下界 \\
蒙特卡洛估计与重要性采样 & MIT 6.3700 / CMU & 第 3 章 & 第 3 卷图采样、第 6 卷 PPO 策略梯度 \\
熵、互信息、交叉熵、正/反向 KL 散度 & CS109 / CS229 / 北大 & 第 4 章 & 分类损失、语言模型困惑度、决策树、VAE/DPO \\
MSE 偏差-方差、MLE/MAP、共轭与 Bootstrap & CMU 36-401 / 南大概统 & 第 5 章 & 第 4 卷正则化与泛化、第 7 卷实验置信区间 \\
纯 NumPy 线性/逻辑/Softmax 回归与评测 & CS229 / CMU 07-280 & 第 6 章 & 阶段一验收里程碑，全系列监督学习基线 \\
\bottomrule
\end{tabularx}
\end{center}

\section*{0.3\quad 第 2 卷知识阶梯（零先修倒挂）}
\addcontentsline{toc}{section}{0.3\quad 第 2 卷知识阶梯（零先修倒挂）}

\begin{center}
\resizebox{0.98\linewidth}{!}{%
\begin{tikzpicture}[
  node distance=6.5mm and 7.5mm,
  ch/.style={
    rectangle, rounded corners=1.8mm, draw=teal, thick, fill=tealLight,
    text width=34mm, minimum height=14.5mm, align=center, font=\sffamily\footnotesize, text=ink, inner sep=4pt
  },
  goal/.style={
    rectangle, rounded corners=1.8mm, draw=ochre, very thick, fill=ochreLight,
    text width=37mm, minimum height=14.5mm, align=center, font=\sffamily\footnotesize, text=ink, inner sep=4pt
  },
  arr/.style={-{Stealth[length=2.2mm]}, thick, teal, preaction={draw=white,line width=2.5pt, arrows=-}}
]
  \node[ch] (c1) {\textbf{第 1 章}\\概率建模、条件独立\\与 Bayes 推断};
  \node[ch, right=of c1] (c2) {\textbf{第 2 章}\\AI 核心分布与\\多元高斯几何};
  \node[ch, right=of c2] (c3) {\textbf{第 3 章}\\Jensen、浓度不等式\\与蒙特卡洛};

  \node[ch, below=9mm of c1] (c4) {\textbf{第 4 章}\\信息论：熵、交叉熵\\与正/反向 KL 散度};
  \node[ch, below=9mm of c2] (c5) {\textbf{第 5 章}\\统计估计：MLE、MAP\\共轭先验与 Bootstrap};
  \node[goal, below=9mm of c3] (c6) {\textbf{\color{ochre}第 6 章（底座大合流）}\\纯 NumPy 线性/逻辑/\\Softmax 回归全闭环};

  \draw[arr] (c1) -- (c2);
  \draw[arr] (c2) -- (c3);
  \draw[arr] (c1.south) -- (c4.north);
  \draw[arr] ([xshift=-9mm]c3.south) -- ([xshift=9mm]c4.north);
  \draw[arr] (c2) -- (c5);
  \draw[arr] (c3) -- (c5);
  \draw[arr] (c4) -- (c5);
  \draw[arr] (c4.south) -- ++(0,-4.5mm) -| (c6.south);
  \draw[arr] (c5) -- (c6);
\end{tikzpicture}}
\end{center}

\chapter[概率建模、条件独立与 Bayes 推断]{概率建模、条件独立与 Bayes 推断}
\label{ch:ch1}

\chapEpigraph{概率论归根结底，只是化作计算的常识。}{皮埃尔-西蒙·拉普拉斯（Pierre-Simon Laplace），《概率的哲学随笔》}

\begin{chapIntro}{本章定位：从确定性线性代数走向不确定性世界}
在第 1 卷中，我们建立了人工智能的确定性几何底座：内积、投影、实对称阵谱定理、半正定矩阵（PSD）与多元梯度下降。然而，真实世界的 AI 任务几乎从不“干净”——传感器读数带噪声、图像会被遮挡、医学检测存在假阳性、自然语言的下一词本身就具有多种合理续写。只靠确定性方程 $Ax = b$，模型无法区分“数据本身的随机波动”与“模型还没学到的规律”，更无法回答“我对这次预测有多大把握”。

本章将用四步搭好 AI 概率建模的第一层地基，并直接对接后续卷次：
\begin{enumerate}
  \item 从条件概率与\textbf{链式法则}出发，看清大语言模型（GPT）的自回归生成为什么不是工程巧合，而是概率链式法则的精确化身；
  \item 通过\textbf{Bayes 定理}建立“先验 $\times$ 似然 $\to$ 后验”的动态证据更新观，并用罕见病检测反例击碎直觉误区，顺势搭出你的第一个概率分类器雏形（朴素贝叶斯与对数空间防下溢）；
  \item 彻底厘清\textbf{独立 vs 条件独立}的本质分野，借“共因结构（Fork）”与“汇合结构（Collider / Explaining Away）”两枚反例，打通第 3 卷贝叶斯网络的核心神经；
  \item 将离散/连续随机变量、期望线性性、全期望定律与\textbf{协方差矩阵 $\Sigma \succeq 0$} 闭环起来，与第 1 卷的实对称阵谱定理无缝对接。
\end{enumerate}
\end{chapIntro}

% =========================================================================
\section{概率公理速忆、条件概率与链式法则}
\label{sec:ch1-axioms-chain}

\begin{recallBox}{样本空间、事件与 Kolmogorov 概率公理}
记随机试验所有可能结果的集合为\textbf{样本空间} $\Omega$，其子集 $A \subseteq \Omega$ 称为\textbf{事件}。概率函数 $\Prob$ 为每个（可测）事件赋予一个实数，满足三条 Kolmogorov 公理：（1）\textbf{非负性} $\Prob(A) \ge 0$；（2）\textbf{规范性} $\Prob(\Omega) = 1$；（3）\textbf{可列可加性}：若事件序列 $A_1, A_2, \dots$ 两两互斥（$A_i \cap A_j = \varnothing,\; i\neq j$），则 $\Prob\!\left(\bigcup_{i=1}^\infty A_i\right) = \sum_{i=1}^\infty \Prob(A_i)$。由此立得补事件公式 $\Prob(A^c) = 1 - \Prob(A)$、单调性（$A \subseteq B \implies \Prob(A) \le \Prob(B)$）与加法公式 $\Prob(A \cup B) = \Prob(A) + \Prob(B) - \Prob(A \cap B)$。
\end{recallBox}

\subsection{条件概率：获得新信息后的样本空间坍缩}

概率论真正成为人工智能武器的那一刻，并不在静态的 $\Prob(A)$，而在\textbf{观察到新证据 $B$ 之后，我们对事件 $A$ 的信念如何更新}。

先看一个极小的具体场景。假设一款邮件过滤器面对的收件箱中共有 $100$ 封邮件：其中 $40$ 封是垃圾邮件（事件 $S$），$60$ 封是正常邮件。进一步统计发现，全部 $100$ 封邮件里有 $36$ 封包含关键词“免费”（事件 $F$），而这 $36$ 封里恰好有 $30$ 封是垃圾邮件（事件 $S \cap F$）。在还没打开邮件前，它属于垃圾邮件的先验概率是 $\Prob(S) = 40/100 = 0.40$；但一旦系统扫到标题里出现了“免费”，原本的 $100$ 封邮件样本空间 $\Omega$ 就瞬间\textbf{坍缩}为满足条件 $F$ 的那 $36$ 封邮件，垃圾邮件所占比例随之重新归一化为 $30/36 \approx 0.833$。

这正是条件概率的全部几何直觉：把世界缩小到已发生的子集 $B$，再把 $B$ 内部的概率总和拉回 $1$。

\begin{defBox}
\begin{definition}[条件概率（Conditional Probability）]
\label{def:ch1-cond-prob}
设 $A, B \subseteq \Omega$ 为事件且 $\Prob(B) > 0$。在已知事件 $B$ 发生的条件下，事件 $A$ 发生的\textbf{条件概率}定义为
\begin{equation}
  \Prob(A \mid B) := \frac{\Prob(A \cap B)}{\Prob(B)}. \label{eq:ch1-cond-prob}
\end{equation}
\end{definition}
\end{defBox}

注意一个容易被忽略却极其有用的事实：当固定条件事件 $B$（$\Prob(B) > 0$）时，映射 $A \mapsto \Prob(A \mid B)$ 本身就是一个完全合法的概率测度，满足全部三条 Kolmogorov 公理——例如 $\Prob(A^c \mid B) = 1 - \Prob(A \mid B)$，且对互斥事件 $A_1, A_2$ 有 $\Prob(A_1 \cup A_2 \mid B) = \Prob(A_1 \mid B) + \Prob(A_2 \mid B)$。凡是对普通概率 $\Prob(\cdot)$ 成立的恒等式，只要在每一项的条件栏里统一加上“$\mid B$”，依然全部成立。

\subsection{乘法法则与概率链式法则}

将定义式 \eqref{eq:ch1-cond-prob} 两边同乘 $\Prob(B)$，便得到\textbf{乘法法则}：$\Prob(A \cap B) = \Prob(B)\,\Prob(A \mid B) = \Prob(A)\,\Prob(B \mid A)$。将其反复嵌套到 $n$ 个事件上，就得到整个人工智能序列建模最重要的恒等式之一。

\begin{thmBox}
\begin{theorem}[概率链式法则（Chain Rule of Probability）]
\label{thm:ch1-chain-rule}
设事件 $A_1, A_2, \dots, A_n$ 满足 $\Prob(A_1 \cap A_2 \cap \dots \cap A_{n-1}) > 0$，则它们的联合概率可精确分解为条件概率的连乘积：
\begin{equation}
  \Prob(A_1 \cap A_2 \cap \dots \cap A_n)
  = \Prob(A_1)\,\Prob(A_2 \mid A_1)\,\Prob(A_3 \mid A_1 \cap A_2)\cdots\Prob(A_n \mid A_1 \cap \dots \cap A_{n-1}).
  \label{eq:ch1-chain-rule}
\end{equation}
简记联合事件为逗号分隔形式，即 $\Prob(A_1, \dots, A_n) = \prod_{t=1}^n \Prob(A_t \mid A_1, \dots, A_{t-1})$（约定 $t=1$ 时条件为空）。
\end{theorem}
\end{thmBox}

\begin{proof}
对右边按定义 \eqref{eq:ch1-cond-prob} 展开，得到一个逐项相消的望远镜乘积（telescoping product）：
\[
  \Prob(A_1)\cdot \frac{\Prob(A_1 \cap A_2)}{\Prob(A_1)}\cdot \frac{\Prob(A_1 \cap A_2 \cap A_3)}{\Prob(A_1 \cap A_2)}\cdots \frac{\Prob(A_1 \cap \dots \cap A_n)}{\Prob(A_1 \cap \dots \cap A_{n-1})}
  = \Prob(A_1 \cap \dots \cap A_n).
\]
由于假设 $\Prob(A_1 \cap \dots \cap A_{n-1}) > 0$，由概率单调性可知所有中间分母 $\Prob(A_1 \cap \dots \cap A_{t-1}) \ge \Prob(A_1 \cap \dots \cap A_{n-1}) > 0$，每一步除法均有意义。
\end{proof}

\begin{insightBox}{为什么 GPT 的“预测下一个词元”能给整篇文章建模？}
假设词表大小为 $V = 50{,}000$，要直接建模一段长为 $T = 100$ 个词元（token）的文本 $(w_1, w_2, \dots, w_T)$ 的联合分布，可能句子总数高达 $V^T = 50{,}000^{100}$——远远超过可观测宇宙的原子数量，绝无可能建表存储。

然而，\textbf{不需要任何近似假设}，概率链式法则 \eqref{eq:ch1-chain-rule} 允许我们把整段文本的联合概率精确重写为：
\[
  \Prob(w_1, w_2, \dots, w_T) = \prod_{t=1}^T \Prob(w_t \mid w_1, w_2, \dots, w_{t-1}) = \prod_{t=1}^T \Prob(w_t \mid w_{<t}).
\]
现代自回归大语言模型（Autoregressive LLM，如 GPT 系列）所做的唯一一件事，就是用一个带参数 $\theta$ 的 Transformer 神经网络来拟合条件分布 $\Prob_\theta(w_t \mid w_{<t})$（第 5 卷）。因此，“给定上文预测下一个词元（Next-Token Prediction）”并非权宜之计，而是概率链式法则在语言建模上的\textbf{无损精确分解}。
\end{insightBox}

\subsection{全概率公式：分情况讨论的严格语言}

实际问题中，我们往往很难直接算出某事件 $E$ 的总概率，但只要把隐藏原因拆成若干互斥且完备的情形 $H_1, \dots, H_k$，在每种情形下算 $\Prob(E \mid H_i)$ 却非常容易。

\begin{thmBox}
\begin{theorem}[全概率公式（Law of Total Probability）]
\label{thm:ch1-total-prob}
设事件族 $\{H_1, H_2, \dots, H_k\}$ 构成样本空间 $\Omega$ 的一个\textbf{划分}（partition），即两两互斥（$H_i \cap H_j = \varnothing,\; i\neq j$）且覆盖全空间（$\bigcup_{i=1}^k H_i = \Omega$），同时满足 $\Prob(H_i) > 0$。则对任意事件 $E \subseteq \Omega$，有
\begin{equation}
  \Prob(E) = \sum_{i=1}^k \Prob(E \cap H_i) = \sum_{i=1}^k \Prob(E \mid H_i)\,\Prob(H_i). \label{eq:ch1-total-prob}
\end{equation}
\end{theorem}
\end{thmBox}

\begin{proof}
由划分性质，事件 $E$ 可拆分为互斥并集 $E = E \cap \Omega = E \cap \left(\bigcup_{i=1}^k H_i\right) = \bigcup_{i=1}^k (E \cap H_i)$。对两两互斥的 $E \cap H_i$ 应用公理 3（可加性），再用乘法法则 $\Prob(E \cap H_i) = \Prob(E \mid H_i)\Prob(H_i)$ 即证。
\end{proof}

\begin{example}[两座晶圆厂的芯片良品率]
\label{ex:ch1-chip-factory}
某 AI 加速卡的芯片由 A 厂供应 $70\%$（$\Prob(H_A)=0.7$）、B 厂供应 $30\%$（$\Prob(H_B)=0.3$）。已知 A 厂芯片次品率为 $1\%$（$\Prob(E\mid H_A)=0.01$），B 厂次品率为 $4\%$（$\Prob(E\mid H_B)=0.04$）。随机抽检一颗芯片，其次品总概率为加权平均：
\[
  \Prob(E) = \Prob(E\mid H_A)\Prob(H_A) + \Prob(E\mid H_B)\Prob(H_B) = 0.01\times 0.7 + 0.04\times 0.3 = 0.019\;(1.9\%).
\]
若进一步反问：“现在已经测出这颗芯片是次品，它来自 B 厂的概率有多大？”这便是从“由因推果”逆转到“由果溯因”的 Bayes 推断。
\end{example}

% =========================================================================
\section{Bayes 定理：从先验、似然到后验的证据更新}
\label{sec:ch1-bayes}

在科学与人工智能中，正向的“生成方向”通常容易描述：已知患者患病 $H$，化验呈阳性的概率 $\Prob(E\mid H)$ 可以通过临床试验测出来；已知邮件是垃圾邮件 $y=1$，里面出现“中奖”的概率 $\Prob(x\mid y=1)$ 可以统计出来。但真正要做决策时，摆在我们面前的总是\textbf{反问题}——只看到了化验单 $E$ 或邮件内容 $x$，想反推背后的隐藏假设 $H$ 或标签 $y$。Bayes 定理正是翻转条件方向的桥梁。

\begin{thmBox}
\begin{theorem}[Bayes 定理（Bayes' Theorem）]
\label{thm:ch1-bayes}
设 $\{H_1, \dots, H_k\}$ 是样本空间 $\Omega$ 的一个划分（$\Prob(H_i)>0$），事件 $E$ 满足 $\Prob(E)>0$，则对任意假设 $H_i$，有
\begin{equation}
  \Prob(H_i \mid E)
  = \frac{\Prob(E \mid H_i)\,\Prob(H_i)}{\Prob(E)}
  = \frac{\Prob(E \mid H_i)\,\Prob(H_i)}{\sum_{j=1}^k \Prob(E \mid H_j)\,\Prob(H_j)}.
  \label{eq:ch1-bayes}
\end{equation}
\end{theorem}
\end{thmBox}

\begin{proof}
由乘法法则的对称性，$\Prob(H_i \mid E)\Prob(E) = \Prob(H_i \cap E) = \Prob(E \mid H_i)\Prob(H_i)$。两边除以 $\Prob(E) > 0$ 得第一等号，再将分母用全概率公式 \eqref{eq:ch1-total-prob} 展开即得第二等号。
\end{proof}

Bayes 定理的代数推导虽然只有两行，但它建立了一整套贯穿全系列（第 2、3、5、7 卷）的推断哲学。表 \ref{tab:ch1-bayes-terms} 总结了公式中四个核心对象的术语与职责。

\begin{table}[htbp]
  \centering
  \caption{Bayes 推断的四个核心组件及其在机器学习中的对应角色}
  \label{tab:ch1-bayes-terms}
  \small
  \begin{tabularx}{\linewidth}{@{}p{28mm}p{26mm}X@{}}
    \toprule
    \textbf{名称} & \textbf{数学记号} & \textbf{含义与机器学习对应角色} \\
    \midrule
    \textbf{先验（Prior）} & $\Prob(H_i)$ & 看到证据 $E$ 之前对假设 $H_i$ 的初始信念；对应第 5 章参数估计中的先验分布 $p(\theta)$（即正则化项）。 \\
    \textbf{似然（Likelihood）} & $\Prob(E \mid H_i)$ & 假设 $H_i$ 成立时，产生当前观测证据 $E$ 的匹配程度（注意：作为 $H_i$ 的函数，它\textbf{不一定}对 $i$ 求和为 $1$）。 \\
    \textbf{边际证据（Evidence）} & $\Prob(E)$ & 对所有可能假设加权平均后的观测总概率，充当\textbf{归一化常数}，与具体 $H_i$ 无关。 \\
    \textbf{后验（Posterior）} & $\Prob(H_i \mid E)$ & 吸收证据 $E$ 后更新得到的信念分布，是最终做分类与决策的直接依据。 \\
    \bottomrule
  \end{tabularx}
\end{table}

在比较不同假设 $H_i$ 孰优孰劣时，由于分母 $\Prob(E)$ 对所有 $i$ 都相同，我们常常写作简洁的比例形式：
\[
  \underbrace{\Prob(H_i \mid E)}_{\text{后验}} \;\propto\; \underbrace{\Prob(E \mid H_i)}_{\text{似然}} \times \underbrace{\Prob(H_i)}_{\text{先验}}.
\]
更重要的是，Bayes 更新天然支持\textbf{序贯学习（Sequential Updating）}：若依次获得两份在给定假设 $H$ 下条件独立的证据 $E_1, E_2$，则
\[
  \Prob(H \mid E_1, E_2) \;\propto\; \Prob(E_1, E_2 \mid H)\Prob(H)
  = \Prob(E_2 \mid H)\,\bigl[\Prob(E_1 \mid H)\Prob(H)\bigr]
  \;\propto\; \Prob(E_2 \mid H)\,\Prob(H \mid E_1).
\]
也就是说，\textbf{吸收第一份证据后的“今天后验”，直接变成迎接第二份证据的“明天先验”}！下面这个经典反直觉案例把这一点展现得淋漓尽致。

\begin{example}[罕见病筛查中的基础率谬误与序贯二次复检]
\label{ex:ch1-rare-disease}
设某罕见病在人群中的患病率（先验基础率）为 $\Prob(D) = 0.001$（即 $0.1\%$，千分之一）。现有一款精度看起来“极高”的检测试剂盒：
\begin{itemize}
  \item \textbf{灵敏度（Sensitivity / 真阳性率）}：$\Prob(+\mid D) = 0.99$（真患者有 $99\%$ 概率检出阳性）；
  \item \textbf{特异度（Specificity / 真阴性率）}：$\Prob(-\mid \neg D) = 0.98$，故\textbf{假阳性率（FPR）}为 $\Prob(+\mid \neg D) = 1 - 0.98 = 0.02$。
\end{itemize}
某人无明显症状参加普查，第一次检测竟呈阳性（$+_1$）。他真正患病的后验概率是多少？直觉常让人脱口而出“接近 $98\%$”，但 Bayes 定理给出了截然不同的答案。

\textbf{（1）自然频数直觉与精确计算（首次阳性）}：想象对 $100{,}000$ 人做普查。其中真正患病者约 $100$ 人，试剂盒能检出 $100\times 0.99 = 99$ 例\textbf{真阳性}；而未患病的健康人多达 $99{,}900$ 人，即便假阳性率只有区区 $2\%$，也会制造出 $99{,}900\times 0.02 = 1{,}998$ 例\textbf{假阳性}！健康人基数太大，假阳性人数足足是真阳性的 $20$ 倍。代入式 \eqref{eq:ch1-bayes} 精确计算：
\begin{align*}
  \Prob(D \mid +_1)
  &= \frac{\Prob(+_1 \mid D)\Prob(D)}{\Prob(+_1 \mid D)\Prob(D) + \Prob(+_1 \mid \neg D)\Prob(\neg D)} \\
  &= \frac{0.99 \times 0.001}{0.99 \times 0.001 + 0.02 \times 0.999}
  = \frac{0.00099}{0.00099 + 0.01998}
  = \frac{0.00099}{0.02097} \approx 0.04721\;(4.72\%).
\end{align*}
也就是说，即便 $99\%$ 灵敏度、$98\%$ 特异度的检测报了阳性，此人真正患病的概率仍\textbf{仅有 $4.72\%$}！忽视极低先验 $\Prob(D)$ 而把 $\Prob(+\mid D)$ 错当成 $\Prob(D\mid +)$，在统计学中称为\textbf{基础率谬误（Base Rate Fallacy）}。

\textbf{（2）第二次条件独立复检阳性（序贯更新）}：医生当然不会就此下结论，而是换用原理独立的第二套试剂盒复检（假设给定真实健康状态 $D$ 或 $\neg D$ 时，两次检测结果条件独立），结果再次呈阳性（$+_2$）。此时把第一次的后验 $p_1 = 0.047210$ 当作新的先验，再套一次 Bayes 定理：
\[
  \Prob(D \mid +_1, +_2)
  = \frac{0.99 \times p_1}{0.99 \times p_1 + 0.02 \times (1 - p_1)}
  \approx \frac{0.046738}{0.046738 + 0.019056}
  \approx 0.71037\;(71.04\%).
\]
如果第三次独立复检依然呈阳性，后验概率将进一步飙升至 $99.18\%$！用\textbf{几率（Odds）形式}看更加清晰：阳性似然比 $\mathrm{LR}_+ = \frac{\Prob(+\mid D)}{\Prob(+\mid \neg D)} = \frac{0.99}{0.02} = 49.5$，每出现一次独立阳性，患病几率 $\frac{\Prob(D)}{\Prob(\neg D)}$ 就乘以 $49.5$：
\[
  \frac{1}{999} \xrightarrow{\;\times 49.5\;} \frac{49.5}{999} \approx 0.04955\;(4.72\%)
  \xrightarrow{\;\times 49.5\;} \frac{49.5^2}{999} \approx 2.4527\;(71.04\%)
  \xrightarrow{\;\times 49.5\;} \frac{49.5^3}{999} \approx 121.41\;(99.18\%).
\]
以上数值均与配套脚本 \texttt{code/vol2\_experiments.py} 实测输出完全一致。
\end{example}

\begin{figure}[htbp]
  \centering
  \begin{tikzpicture}[font=\footnotesize\sffamily]
    % -------- Left Panel: Probability Tree for 1st Test --------
    \begin{scope}[xshift=0cm, yshift=0cm]
      \node[font=\small\sffamily\bfseries, color=tealDark] at (3.2, 3.97) {(a) 首次检测的概率树分解};

      % Root node
      \node[circle, draw=ink, fill=codebg, inner sep=2pt, font=\scriptsize\sffamily] (root) at (0.3, 1.55) {普查};

      % Level 1: D and ~D
      \node[rounded corners=2pt, draw=ochre, fill=ochreLight, minimum width=1.45cm, minimum height=0.6cm, align=center] (D) at (2.7, 2.55) {患病 $D$};
      \node[rounded corners=2pt, draw=teal, fill=tealLight, minimum width=1.45cm, minimum height=0.6cm, align=center] (nD) at (2.7, 0.55) {健康 $\neg D$};

      % Level 2: + and -
      \node[rounded corners=2pt, draw=ochre, thick, fill=ochreLight, minimum width=1.95cm, minimum height=0.52cm, align=center] (Dp) at (5.75, 3.00) {真阳性 $D\cap +$};
      \node[rounded corners=2pt, draw=slate, fill=codebg, minimum width=1.95cm, minimum height=0.52cm, align=center] (Dm) at (5.75, 2.10) {假阴性 $D\cap -$};
      \node[rounded corners=2pt, draw=ochre, thick, fill=ochreLight, minimum width=1.95cm, minimum height=0.52cm, align=center] (nDp) at (5.75, 1.00) {假阳性 $\neg D\cap +$};
      \node[rounded corners=2pt, draw=slate, fill=codebg, minimum width=1.95cm, minimum height=0.52cm, align=center] (nDm) at (5.75, 0.10) {真阴性 $\neg D\cap -$};

      % Edges Level 1
      \draw[-{Stealth[length=2mm]}, ochre, thick] (root) -- (D) node[midway, above left, xshift=3pt, yshift=-1pt, color=ochre] {$0.001$};
      \draw[-{Stealth[length=2mm]}, teal, thick] (root) -- (nD) node[midway, below left, xshift=3pt, yshift=1pt, color=teal] {$0.999$};

      % Edges Level 2
      \draw[-{Stealth[length=2mm]}, ochre, thick] (D.east) -- (Dp.west) node[midway, above, yshift=1pt, color=ochre] {$0.99$};
      \draw[-{Stealth[length=2mm]}, slate] (D.east) -- (Dm.west) node[midway, below, yshift=-1pt, color=slate] {$0.01$};
      \draw[-{Stealth[length=2mm]}, ochre, thick] (nD.east) -- (nDp.west) node[midway, above, yshift=1pt, color=ochre] {$0.02$};
      \draw[-{Stealth[length=2mm]}, slate] (nD.east) -- (nDm.west) node[midway, below, yshift=-1pt, color=slate] {$0.98$};

      % Joint probabilities on right
      \node[anchor=west, color=ochre, font=\scriptsize\sffamily\bfseries] at (6.82, 3.00) {$0.00099$};
      \node[anchor=west, color=slate, font=\scriptsize\sffamily] at (6.82, 2.10) {$0.00001$};
      \node[anchor=west, color=ochre, font=\scriptsize\sffamily\bfseries] at (6.82, 1.00) {$0.01998$};
      \node[anchor=west, color=slate, font=\scriptsize\sffamily] at (6.82, 0.10) {$0.97902$};

      % Bottom summary note
      \node[anchor=north, font=\scriptsize\sffamily, color=ink, align=center, fill=codebg, draw=codeborder, rounded corners=1.5pt, inner sep=3.5pt] at (3.7, -0.45)
        {阳性总概率 $\Prob(+)=0.00099+0.01998=0.02097$，真阳性仅占 $4.72\%$};
    \end{scope}

    % -------- Right Panel: Sequential Bayes Posterior Bar Chart --------
    \begin{scope}[xshift=9.2cm, yshift=0.1cm]
      \node[font=\small\sffamily\bfseries, color=tealDark] at (2.8, 3.87) {(b) 序贯复检下的后验信念跃升};

      % Axes
      \draw[-{Stealth[length=2mm]}, ink] (0, 0) -- (5.9, 0);
      \draw[-{Stealth[length=2mm]}, ink] (0, 0) -- (0, 3.15) node[above, font=\scriptsize\sffamily] {$\Prob(D\mid \text{证据})$};

      % Horizontal gridlines
      \foreach \y/\lbl in {0.7/25\%, 1.4/50\%, 2.1/75\%, 2.8/100\%} {
        \draw[dashed, gray!30] (0, \y) -- (5.65, \y);
        \node[left, font=\scriptsize\sffamily, color=slate] at (-0.05, \y) {\lbl};
      }

      % Bars: scale 100% -> 2.8cm
      % 0.10% -> 0.0028 (draw tiny visible bar 0.04cm)
      \fill[slate!50, draw=slate] (0.5, 0) rectangle (1.3, 0.05);
      \node[above, font=\scriptsize\sffamily\bfseries, color=slate] at (0.9, 0.08) {$0.10\%$};
      \node[below, font=\scriptsize\sffamily, align=center] at (0.9, -0.08) {先验\\[-1pt]无检测};

      % 4.72% -> 0.132cm
      \fill[ochre!65, draw=ochre] (1.85, 0) rectangle (2.65, 0.15);
      \node[above, font=\scriptsize\sffamily\bfseries, color=ochre] at (2.25, 0.18) {$4.72\%$};
      \node[below, font=\scriptsize\sffamily, align=center] at (2.25, -0.08) {首次\\[-1pt]阳性 $+_1$};

      % 71.04% -> 1.989cm
      \fill[teal!70, draw=tealDark] (3.2, 0) rectangle (4.0, 1.99);
      \node[above, font=\scriptsize\sffamily\bfseries, color=tealDark] at (3.6, 2.02) {$71.04\%$};
      \node[below, font=\scriptsize\sffamily, align=center] at (3.6, -0.08) {二次\\[-1pt]阳性 $+_2$};

      % 99.18% -> 2.777cm
      \fill[navy!75, draw=navy] (4.55, 0) rectangle (5.35, 2.78);
      \node[above, font=\scriptsize\sffamily\bfseries, color=navy] at (4.95, 2.80) {$99.18\%$};
      \node[below, font=\scriptsize\sffamily, align=center] at (4.95, -0.08) {三次\\[-1pt]阳性 $+_3$};
    \end{scope}
  \end{tikzpicture}
  \caption{例 \ref{ex:ch1-rare-disease} 罕见病筛查的概率树与序贯 Bayes 更新。（a）极低基础率 $0.1\%$ 导致假阳性分支（$0.01998$）远大于真阳性分支（$0.00099$），首次阳性后验仅 $4.72\%$；（b）在条件独立复检下，每次阳性都将患病几率放大 $49.5$ 倍，第二次即跃升至 $71.04\%$，第三次逼近 $99.18\%$。}
  \label{fig:ch1-bayes-update}
\end{figure}

\subsection{首个概率分类器雏形：朴素贝叶斯与对数空间计算}

把例 \ref{ex:ch1-rare-disease} 中的“患病与否 $D$”换成“邮件是否为垃圾邮件 $y \in \{0, 1\}$”，把“两次独立检测”换成“邮件是否包含 $d$ 个特征词 $x = (x_1, x_2, \dots, x_d) \in \{0, 1\}^d$”，我们就直接得到了机器学习历史上最经典也最轻量的分类器之一——\textbf{朴素贝叶斯（Naive Bayes）}。

如果直接用 Bayes 定理算后验 $\Prob(y=1 \mid x_1, \dots, x_d) \propto \Prob(y=1)\,\Prob(x_1, \dots, x_d \mid y=1)$，立刻会撞上一堵参数爆炸墙：$d$ 个二值词特征的联合分布 $\Prob(x_1, \dots, x_d \mid y)$ 需要估计 $2^d - 1$ 个参数，当词表 $d = 10{,}000$ 时完全无法学习。

朴素贝叶斯做了一个大胆而“朴素”的简化假设：\textbf{给定类别标签 $y$ 时，各个词特征 $x_1, \dots, x_d$ 之间条件独立}。于是高维联合似然立刻拆成一维标量概率的连乘积：
\begin{equation}
  \Prob(y = c \mid x_1, \dots, x_d)
  \;\propto\; \Prob(y = c) \prod_{j=1}^d \Prob(x_j \mid y = c), \qquad c \in \{0, 1\}.
  \label{eq:ch1-naive-bayes}
\end{equation}
参数量从指数级的 $O(2^d)$ 暴降为线性级的 $O(d)$！

\begin{example}[3 词垃圾邮件朴素贝叶斯手算]
\label{ex:ch1-nb-mini}
设垃圾邮件先验 $\Prob(y=1)=0.4$，正常邮件先验 $\Prob(y=0)=0.6$。考察三个词特征 $x_1=\text{``免费''}$、$x_2=\text{``中奖''}$、$x_3=\text{``论文''}$ 的条件概率：
\begin{center}
\small
\begin{tabular}{cccc}
  \toprule
  类别 $y$ & $\Prob(x_1=1\mid y)$（免费） & $\Prob(x_2=1\mid y)$（中奖） & $\Prob(x_3=1\mid y)$（论文） \\
  \midrule
  垃圾邮件 $y=1$ & $0.80$ & $0.70$ & $0.05$ \\
  正常邮件 $y=0$ & $0.10$ & $0.05$ & $0.60$ \\
  \bottomrule
\end{tabular}
\end{center}
现收到一封邮件，含“免费”和“中奖”、不含“论文”，即 $x = (1, 1, 0)$。由式 \eqref{eq:ch1-naive-bayes} 分别计算未归一化后验分子：
\begin{align*}
  \text{垃圾分子：}\quad & \Prob(y=1)\,\Prob(x_1=1\mid 1)\,\Prob(x_2=1\mid 1)\,\Prob(x_3=0\mid 1) = 0.4 \times 0.80 \times 0.70 \times 0.95 = 0.2128, \\
  \text{正常分子：}\quad & \Prob(y=0)\,\Prob(x_1=1\mid 0)\,\Prob(x_2=1\mid 0)\,\Prob(x_3=0\mid 0) = 0.6 \times 0.10 \times 0.05 \times 0.40 = 0.0012.
\end{align*}
归一化后，该邮件为垃圾邮件的后验概率为 $\Prob(y=1\mid x) = \frac{0.2128}{0.2128 + 0.0012} \approx 0.9944$（$99.44\%$）。
\end{example}

\begin{failBox}{高维连乘的浮点下溢陷阱与“零概率一票否决”}
把例 \ref{ex:ch1-nb-mini} 放到真实的 $d = 2{,}000$ 词邮件分类中，直接按式 \eqref{eq:ch1-naive-bayes} 连乘会瞬间触发两大经典工程灾难：
\begin{enumerate}
  \item \textbf{浮点下溢（Floating-Point Underflow）}：假设每个词概率约为 $0.01$，连乘 $1{,}000$ 项的结果是 $0.01^{1000} = 10^{-2000}$，远远低于 IEEE-754 双精度浮点数（\texttt{float64}）能表示的最小正数约 $2.2\times 10^{-308}$。计算机会把两类的分子\textbf{同时截断为精确的 \texttt{0.0}}，归一化除法 $0.0 / 0.0$ 直接炸出 \texttt{NaN}！
  \textbf{标准修复}：永远在\textbf{对数空间（Log-Space）}中把连乘转化为连加，计算对数得分
  \[
    s_c = \log \Prob(y=c) + \sum_{j=1}^d \log \Prob(x_j \mid y=c), \qquad c \in \{0, 1\},
  \]
  若只需类别预测，直接取 $\argmax_c s_c$；若需精确后验概率，先减去 $m = \max_c s_c$ 再做指数归一化（即 \textbf{Log-Sum-Exp 稳定技巧}，第 6 章将系统展开）：$\Prob(y=c\mid x) = \frac{e^{s_c - m}}{\sum_{c'} e^{s_{c'} - m}}$。
  \item \textbf{零概率一票否决}：若训练集里某罕见词从未在正常邮件中出现，频率估计会给出 $\Prob(x_j=1\mid y=0)=0$，导致连乘积瞬间归零（或 $\log 0 = -\infty$），抹杀其余 $1{,}999$ 个词的全部证据。解决办法是在计数分子分母上各加伪计数（\textbf{Laplace 平滑}，第 5 章将证明它恰是 Beta/Dirichlet 先验下的后验估计）。
\end{enumerate}
\end{failBox}

% =========================================================================
\section{独立 vs 条件独立：AI 概率图的灵魂（Explaining Away）}
\label{sec:ch1-indep-condindep}

上一节朴素贝叶斯的核心假设是“给定邮件标签后，各词特征条件独立”。但许多初学者的第一个直觉会立刻提出质疑：“‘免费’和‘中奖’在现实中明明经常一起出现，怎么能说它们条件独立呢？”

这个追问直指 AI 概率建模最核心、也最容易混淆的概念边界——\textbf{边际独立（Independence）}与\textbf{条件独立（Conditional Independence）}是两件截然不同的事，它们之间\textbf{没有任何单向或双向的蕴含关系}！

\begin{defBox}
\begin{definition}[独立性与条件独立性]
\label{def:ch1-indep}
设 $A, B, C$ 为事件（且 $\Prob(C) > 0$）：
\begin{enumerate}
  \item 若满足
  \begin{equation}
    \Prob(A \cap B) = \Prob(A)\,\Prob(B), \label{eq:ch1-indep}
  \end{equation}
  则称 $A$ 与 $B$ \textbf{相互独立}（或\textbf{边际独立}），记作 $A \indep B$。当 $\Prob(B)>0$ 时，它等价于 $\Prob(A \mid B) = \Prob(A)$（知道 $B$ 不改变 $A$ 的概率）。
  \item 若在给定 $C$ 发生的条件下满足
  \begin{equation}
    \Prob(A \cap B \mid C) = \Prob(A \mid C)\,\Prob(B \mid C), \label{eq:ch1-cond-indep}
  \end{equation}
  则称 $A$ 与 $B$ \textbf{给定 $C$ 条件独立}，记作 $A \indep B \mid C$。当 $\Prob(B \cap C)>0$ 时，它等价于 $\Prob(A \mid B, C) = \Prob(A \mid C)$（一旦已知 $C$，再告知 $B$ 不会提供关于 $A$ 的任何新信息）。
\end{enumerate}
\end{definition}
\end{defBox}

同样地，对随机变量 $X, Y, Z$，$X \indep Y \mid Z$ 意味着对一切取值 $(x, y, z)$（满足 $p(z)>0$），均有 $p(x, y \mid z) = p(x \mid z)\,p(y \mid z)$。下面用两枚彻底算清的离散反例，把“边际相关却条件独立”与“边际独立却条件相关”两种相反机制钉进直觉里。

\subsection{反例一：共因结构（Fork $A \leftarrow C \rightarrow B$）——边际相关，条件独立}

回到“免费”（$A$）与“中奖”（$B$）的例子。设 $C \in \{0, 1\}$ 表示邮件是否为垃圾邮件，且 $\Prob(C=1) = \Prob(C=0) = 0.5$。假设给定邮件类型 $C$ 后，作者独立决定是否写入两个词：
\[
  \Prob(A=1\mid C=1) = 0.8,\quad \Prob(B=1\mid C=1) = 0.8;\qquad
  \Prob(A=1\mid C=0) = 0.1,\quad \Prob(B=1\mid C=0) = 0.1,
\]
且严格满足条件独立性 $A \indep B \mid C$（即给定 $C$ 时联合概率等于条件概率之积）。由此可写出 $(A, B, C)$ 的完整联合概率，并对 $C$ 求和得到 $(A, B)$ 的边际联合分布表（表 \ref{tab:ch1-fork-table}）：

\begin{table}[htbp]
  \centering
  \caption{共因结构 $A \leftarrow C \rightarrow B$ 下的 $(A, B)$ 边际联合概率表（由条件独立模型对 $C$ 求和所得）}
  \label{tab:ch1-fork-table}
  \small
  \begin{tabular}{cccc}
    \toprule
    & $B = 1$（含“中奖”） & $B = 0$（不含“中奖”） & 边际行和 $\Prob(A)$ \\
    \midrule
    $A = 1$（含“免费”） & $0.5\times 0.64 + 0.5\times 0.01 = \mathbf{0.325}$ & $0.5\times 0.16 + 0.5\times 0.09 = 0.125$ & $0.450$ \\
    $A = 0$（不含“免费”） & $0.5\times 0.16 + 0.5\times 0.09 = 0.125$ & $0.5\times 0.04 + 0.5\times 0.81 = 0.425$ & $0.550$ \\
    \midrule
    边际列和 $\Prob(B)$ & $0.450$ & $0.550$ & $1.000$ \\
    \bottomrule
  \end{tabular}
\end{table}

立刻检验边际独立性：
\[
  \Prob(A=1, B=1) = 0.325 \;\neq\; 0.450 \times 0.450 = 0.2025 = \Prob(A=1)\,\Prob(B=1)!
\]
不仅如此，$\Prob(B=1 \mid A=1) = 0.325 / 0.450 \approx 0.722$，远高于先验 $\Prob(B=1) = 0.450$！为什么？因为当我们不知道 $C$ 时，看到“免费”会大幅提高“这是垃圾邮件”的后验概率，从而间接拉高“中奖”出现的概率——信息通过隐藏共因 $C$ 从 $A$ 流向了 $B$。而一旦我们\textbf{直接观测并固定了共因 $C$}（比如已知这就是垃圾邮件），$A \leftarrow C \rightarrow B$ 这条信息通道就被阻断，$A$ 和 $B$ 立刻恢复条件独立。这就是朴素贝叶斯之所以能合理建模相关词共现的数学秘密！

\subsection{反例二：汇合结构（Collider $A \rightarrow C \leftarrow B$）——边际独立，条件相关（Explaining Away）}

反过来更令人拍案叫绝：两个原本风马牛不相及的独立事件，一旦观测到它们的\textbf{共同结果}，就会立刻产生强烈的负相关！

\begin{example}[两枚独立硬币与辩解效应（Explaining Away）]
\label{ex:ch1-collider-coins}
独立抛掷两枚均匀硬币，记结果为 $X_1, X_2 \in \{0, 1\}$（正面为 $1$，反面为 $0$），再定义它们的共同结果（和）$S := X_1 + X_2 \in \{0, 1, 2\}$。四个等可能样本点列表如下：
\begin{center}
\small
\begin{tabular}{cccccc}
  \toprule
  样本点 $(X_1, X_2)$ & $(0, 0)$ & $(1, 0)$ & $(0, 1)$ & $(1, 1)$ \\
  \midrule
  联合概率 $\Prob(X_1, X_2)$ & $1/4$ & $1/4$ & $1/4$ & $1/4$ \\
  观测结果 $S = X_1 + X_2$ & $0$ & $\mathbf{1}$ & $\mathbf{1}$ & $2$ \\
  \bottomrule
\end{tabular}
\end{center}
\begin{itemize}
  \item \textbf{边际上严格独立（$X_1 \indep X_2$）}：$\Prob(X_1=1, X_2=1) = \frac{1}{4} = \frac{1}{2}\times\frac{1}{2} = \Prob(X_1=1)\Prob(X_2=1)$。
  \item \textbf{给定共同结果 $S = 1$ 后强相关（$X_1 \not\indep X_2 \mid S$）}：一旦被告知 $S = 1$，样本空间坍缩为 $\{(1,0), (0,1)\}$，二者条件概率各为 $1/2$。于是 $\Prob(X_1=1\mid S=1) = 1/2$，$\Prob(X_2=1\mid S=1) = 1/2$，但联合条件概率却为
  \[
    \Prob(X_1=1, X_2=1 \mid S=1) = 0 \;\neq\; \frac{1}{2}\times\frac{1}{2} = \frac{1}{4}!
  \]
  更直观地看：在已知 $S=1$ 的前提下，只要我告诉你第一枚硬币是正面（$X_1=1$），你就 $100\%$ 笃定第二枚必是反面（$X_2=0$）——原本完全独立的两枚硬币变成了\textbf{完全负相关}（条件协方差 $\Cov(X_1, X_2 \mid S=1) = -1/4$）。
\end{itemize}
\end{example}

在人工智能与认知科学中，这种汇合结构（Collider）诱导的条件相关性有一个极富画面感的名字——\textbf{辩解效应（Explaining Away）}。想象你早晨推开门看到花园草地是湿的（共同结果 $W=1$），这可能源于两个原本相互独立的独立原因：昨晚下过雨（$R=1$）或自动洒水器开了（$S=1$）。在没看到草地时，天气和洒水器独立（$R \indep S$）；可一旦看到草地湿了（观测到 $W=1$），两个原因就开始“竞争解释同一份证据”——倘若你紧接着抬头发现邻居屋顶也在滴水（确认 $R=1$），“昨晚下雨”已经充分解释了草地为什么湿，于是“洒水器开过”的后验概率便应声下降（被下雨“辩解掉”了）。

\begin{figure}[htbp]
  \centering
  \begin{tikzpicture}[
    node distance=1.6cm,
    rv/.style={circle, draw=ink, thick, minimum size=0.92cm, font=\small\sffamily\bfseries, fill=white},
    obs/.style={circle, draw=tealDark, thick, minimum size=0.92cm, font=\small\sffamily\bfseries, fill=tealLight},
    obsbad/.style={circle, draw=ochre, thick, minimum size=0.92cm, font=\small\sffamily\bfseries, fill=ochreLight},
    edge/.style={-{Stealth[length=2.3mm]}, thick, ink}
  ]
    % -------- Panel (a): Fork --------
    \begin{scope}[xshift=0cm]
      \node[font=\small\sffamily\bfseries, color=tealDark] at (0, 2.35) {(a) 共因结构（Fork）：观测 $C$ 阻断通道};

      \node[obs] (C1) at (0, 1.30) {$C$};
      \node[rv] (A1) at (-2.1, -0.10) {$A$};
      \node[rv] (B1) at (2.1, -0.10) {$B$};

      \draw[edge] (C1) -- (A1);
      \draw[edge] (C1) -- (B1);

      % Blocked path indicator
      \draw[dashed, teal, thick] (A1) to[bend right=18] node[midway, below, font=\scriptsize\sffamily\bfseries, color=tealDark] {通道被阻断：$A \indep B \mid C$} (B1);
      \node[font=\scriptsize\sffamily, color=tealDark, align=center] at (0, 1.30) {};
      \node[above right, font=\scriptsize\sffamily, color=tealDark] at (0.48, 1.35) {已观测（阴影）};
      \node[below, font=\scriptsize\sffamily, color=slate, align=center] at (0, -1.05) {例：垃圾邮件 $C$ $\to$ 词 $A, B$\\未观测 $C$ 时 $A \not\indep B$，给定 $C$ 后独立};
    \end{scope}

    % Divider line
    \draw[tealLine, dashed] (3.65, -1.65) -- (3.65, 2.55);

    % -------- Panel (b): Collider --------
    \begin{scope}[xshift=7.3cm]
      \node[font=\small\sffamily\bfseries, color=ochre] at (0, 2.35) {(b) 汇合结构（Collider）：观测 $C$ 打开通道};

      \node[obsbad] (C2) at (0, -0.10) {$C$};
      \node[rv] (A2) at (-2.1, 1.30) {$A$};
      \node[rv] (B2) at (2.1, 1.30) {$B$};

      \draw[edge] (A2) -- (C2);
      \draw[edge] (B2) -- (C2);

      % Opened path indicator
      \draw[ochre, thick, <->, dashed] (A2) to[bend left=18] node[midway, above, font=\scriptsize\sffamily\bfseries, color=ochre] {辩解效应打通关联：$A \not\indep B \mid C$} (B2);
      \node[right, font=\scriptsize\sffamily, color=ochre] at (0.52, -0.10) {已观测（阴影）};
      \node[below, font=\scriptsize\sffamily, color=slate, align=center] at (0, -1.05) {例：独立硬币 $X_1, X_2 \to$ 和 $S$\\未观测 $S$ 时 $X_1 \indep X_2$，给定 $S=1$ 后负相关};
    \end{scope}
  \end{tikzpicture}
  \caption{有向概率图中的两种对偶基本结构（阴影节点代表已被观测到的条件变量）。（a）共因结构 $A \leftarrow C \rightarrow B$ 中，观测共因 $C$ 会切断 $A, B$ 的统计联系；（b）汇合结构 $A \rightarrow C \leftarrow B$ 中，观测共同结果 $C$ 反而会打开 $A, B$ 之间的信息通道，引发 Explaining Away（辩解效应）。}
  \label{fig:ch1-fork-collider}
\end{figure}

\begin{insightBox}{直通第 3 卷贝叶斯网络与因果推断（$d$-分离与对撞偏倚）}
图 \ref{fig:ch1-fork-collider} 中的共因结构、汇合结构，再加上一条直链 $A \to C \to B$（观测 $C$ 同样阻断 $A$ 与 $B$，如马尔可夫链），共同构成了\textbf{有向图模型（贝叶斯网络）}仅有的三种原子路径。第 3 卷将在此基础上建立完整的 $d$-分离（$d$-separation）判定准则，只需看图就能瞬间读出高维分布中的所有条件独立关系。此外，在数据分析与大模型评估中，如果不慎按汇合节点做数据过滤（例如只统计“进入顶级高校的学生”里智商与努力程度的相关性），就会人为制造出虚假负相关，称为\textbf{对撞偏倚（Collider Bias / Berkson 悖论）}。
\end{insightBox}

% =========================================================================
\section{随机变量、联合分布、期望线性性与协方差矩阵}
\label{sec:ch1-rv-cov}

到目前为止，我们都在用抽象事件 $A, B$ 讨论概率。要把概率论塞进 NumPy 向量和矩阵里，就需要把每个样本点 $\omega \in \Omega$ 映射为实数或实向量——这就是\textbf{随机变量（Random Variable）} $X: \Omega \to \R$ 与\textbf{随机向量} $X: \Omega \to \R^d$。

\subsection{PMF、PDF、CDF 与联合/边际/条件分布}

\begin{itemize}
  \item \textbf{离散随机变量}由\textbf{概率质量函数（PMF）} $p_X(x) := \Prob(X = x)$ 刻画，满足 $p_X(x) \ge 0$ 与 $\sum_x p_X(x) = 1$。
  \item \textbf{连续随机变量}由\textbf{概率密度函数（PDF）} $p_X(x) \ge 0$ 刻画，满足 $\int_{-\infty}^\infty p_X(x)\,dx = 1$，事件概率通过积分给出：$\Prob(a \le X \le b) = \int_a^b p_X(x)\,dx$。
  \item 无论离散还是连续，\textbf{累积分布函数（CDF）}统一定义为 $F_X(x) := \Prob(X \le x)$，它是单调非减、右连续且满足 $F_X(-\infty)=0,\; F_X(+\infty)=1$ 的函数；对连续变量有 $p_X(x) = F_X'(x)$。
\end{itemize}

\begin{failBox}{PDF 是“密度”而非概率——它可以大于 $1$！}
连续随机变量取任意单个精确点的概率恒为零：$\Prob(X = x_0) = \int_{x_0}^{x_0} p_X(x)\,dx = 0$。因此 $p_X(x)$ 代表的是“单位长度上的概率密度”，而非概率本身，\textbf{完全可以严格大于 $1$}！例如区间 $[0, 0.2]$ 上的均匀分布 $\mathrm{Uniform}(0, 0.2)$，其密度为 $p_X(x) = 5 > 1$（因为 $5 \times 0.2 = 1$）。同样，第 4 章在计算连续分布的微分熵时，由于 $p_X(x)$ 可以大于 $1$，会出现“微分熵可为负数”的独特现象。
\end{failBox}

当同时考察多个随机变量 $(X, Y)$ 时，我们用\textbf{联合分布} $p_{X,Y}(x,y)$ 描述它们共同的规律，并通过两种基本操作在维度间自由穿梭（以下为离散求和形式，连续情形将 $\sum$ 换为 $\int$ 即可）：
\begin{align*}
  \textbf{边际化（Marginalization / 消去冗余变量）：}\; & p_X(x) = \sum_y p_{X,Y}(x, y), \\
  \textbf{条件化（Conditioning）：}\; & p_{X\mid Y}(x \mid y) = \frac{p_{X,Y}(x, y)}{p_Y(y)}.
\end{align*}
从工程视角看，“边际化”就是对你不关心的隐变量把所有可能取值求和/积分掉；例如例 \ref{ex:ch1-rare-disease} 中计算分母 $\Prob(+)$ 时，本质上就是在对隐藏患病状态 $D \in \{0, 1\}$ 做边际化。

\subsection{期望的线性性与全期望定律}

随机变量的\textbf{数学期望（Expectation）}定义为按概率加权的平均值：离散时 $\E[X] := \sum_x x\,p_X(x)$，连续时 $\E[X] := \int_{-\infty}^\infty x\,p_X(x)\,dx$（前提是绝对收敛）。更一般地，对任意函数 $g$，有无意识统计学家法则（LOTUS）：$\E[g(X)] = \sum_x g(x)\,p_X(x)$（连续情形为积分）。

在全部概率计算中，最强大、最省力的两条恒等式莫过于\textbf{期望线性性}与\textbf{全期望定律}。

\begin{thmBox}
\begin{theorem}[期望的线性性与全期望定律（Tower Property）]
\label{thm:ch1-expectation-laws}
设 $X, Y$ 为期望存在的随机变量，$a, b, c \in \R$ 为任意常数，则：
\begin{enumerate}
  \item \textbf{期望线性性}（\textbf{不要求 $X$ 与 $Y$ 独立！}）：
  \begin{equation}
    \E[aX + bY + c] = a\,\E[X] + b\,\E[Y] + c. \label{eq:ch1-linearity}
  \end{equation}
  \item \textbf{全期望定律（迭代期望公式）}：记 $\E[X \mid Y]$ 为以 $Y$ 为随机自变量的条件期望函数，则
  \begin{equation}
    \E[X] = \E_Y\bigl[\E[X \mid Y]\bigr]. \label{eq:ch1-tower}
  \end{equation}
\end{enumerate}
\end{theorem}
\end{thmBox}

\begin{proof}
我们在离散有限（或绝对收敛）情形下给出清晰证明（连续情形仅需将求和换为积分并用 Fubini 定理交换积分次序）。
\begin{enumerate}
  \item 利用联合分布与边际化公式：
  \begin{align*}
    \E[aX + bY + c]
    &= \sum_x \sum_y (ax + by + c)\,p_{X,Y}(x, y) \\
    &= a \sum_x x \Bigl(\sum_y p_{X,Y}(x,y)\Bigr) + b \sum_y y \Bigl(\sum_x p_{X,Y}(x,y)\Bigr) + c \sum_{x,y} p_{X,Y}(x,y) \\
    &= a \sum_x x\,p_X(x) + b \sum_y y\,p_Y(y) + c = a\,\E[X] + b\,\E[Y] + c.
  \end{align*}
  注意全程只用到了联合分布的边际和，\textbf{完全没有假设} $p_{X,Y}(x,y) = p_X(x)p_Y(y)$。
  \item 展开右侧外层对 $Y$ 的期望，再利用 $p_Y(y)\,p_{X\mid Y}(x\mid y) = p_{X,Y}(x,y)$ 交换求和次序：
  \begin{align*}
    \E_Y\bigl[\E[X \mid Y]\bigr]
    &= \sum_y \Bigl(\sum_x x\,p_{X\mid Y}(x \mid y)\Bigr) p_Y(y)
     = \sum_x x \sum_y p_{X,Y}(x, y)
     = \sum_x x\,p_X(x) = \E[X]. \qedhere
  \end{align*}
\end{enumerate}
\end{proof}

期望线性性为何如此惊人？试想一个小问题：把 $n$ 份写有姓名的试卷完全随机地发给 $n$ 位同学，平均有多少人刚好拿到自己的试卷？每个人拿到自己试卷的指示变量 $I_k \in \{0, 1\}$ 彼此显然\textbf{不独立}（若前 $n-1$ 人都拿对了，第 $n$ 人必然也拿对），但借助线性性，完全不用算错排概率，只需一行：
\[
  \E\Bigl[\sum_{k=1}^n I_k\Bigr] = \sum_{k=1}^n \E[I_k] = \sum_{k=1}^n \frac{1}{n} = 1!
\]
在机器学习中，训练损失的经验均值 $\frac{1}{N}\sum_{i=1}^N \ell_i$、随机梯度下降（SGD）的梯度无偏性 $\E[\nabla \ell_{\text{batch}}] = \nabla L(\theta)$，以及第 7 卷强化学习的 Bellman 方程（全期望定律），背后全靠这两条恒等式支撑。

\subsection{方差、协方差与半正定协方差矩阵 $\Sigma \succeq 0$}

期望只刻画分布的“重心”，要描述数据的离散程度与变量间的线性协同关系，需要二阶矩。定义随机变量 $X$（均值 $\mu_X = \E[X]$）的\textbf{方差}、\textbf{标准差}以及 $X, Y$ 间的\textbf{协方差}分别为：
\[
  \Var(X) := \E\bigl[(X - \mu_X)^2\bigr] = \E[X^2] - (\E[X])^2, \qquad
  \sigma_X := \sqrt{\Var(X)},
\]
\[
  \Cov(X, Y) := \E\bigl[(X - \mu_X)(Y - \mu_Y)\bigr] = \E[XY] - \E[X]\E[Y].
\]
直接展开 $(X + Y - \mu_X - \mu_Y)^2$ 并用期望线性性，立刻得到\textbf{和的方差公式}：
\begin{equation}
  \Var(aX + bY) = a^2\Var(X) + b^2\Var(Y) + 2ab\,\Cov(X, Y). \label{eq:ch1-var-sum}
\end{equation}
只有当 $\Cov(X, Y) = 0$（例如 $X \indep Y$）时，交叉项才消失，方差才满足可加性 $\Var(X+Y) = \Var(X) + \Var(Y)$。

\begin{failBox}{反例：不相关（$\Cov(X,Y)=0$）绝不等于独立（$X \indep Y$）！}
若 $X \indep Y$，则 $\E[XY] = \E[X]\E[Y]$，必有 $\Cov(X,Y)=0$；但\textbf{逆命题不成立}！取 $X$ 等概率取 $\{-1, 0, 1\}$（各占 $1/3$），并令 $Y := X^2$。显然 $Y$ 完全由 $X$ 决定，二者强相关，但 $\E[X]=0$ 且 $\E[XY] = \E[X^3] = \frac{1}{3}(-1 + 0 + 1) = 0$，于是
\[
  \Cov(X, Y) = \E[XY] - \E[X]\E[Y] = 0 - 0 = 0!
\]
协方差只能捕捉\textbf{线性}相关性，对对称的非线性依赖（如抛物线 $Y=X^2$）完全“失明”。（唯一著名的例外是第 2 章的多元联合高斯分布：对联合高斯向量，零协方差与独立性等价。）
\end{failBox}

现在把目光升维到 $d$ 维随机特征向量 $X = (X_1, \dots, X_d)^T \in \R^d$，其均值向量记为 $\mu := \E[X] = (\E[X_1],\allowbreak \dots,\allowbreak \E[X_d])^T \in \R^d$。把任意两维之间的协方差排成一个 $d\times d$ 矩阵，就得到 AI 数据分析中最核心的矩阵对象——\textbf{协方差矩阵（Covariance Matrix）}：
\begin{equation}
  \Sigma := \Cov(X) := \E\bigl[(X - \mu)(X - \mu)^T\bigr] \in \R^{d\times d}, \qquad \Sigma_{ij} = \Cov(X_i, X_j). \label{eq:ch1-cov-mat}
\end{equation}
它与第 1 卷第 3 章的半正定矩阵（PSD）立刻产生完美合流！

\begin{thmBox}
\begin{proposition}[协方差矩阵的对称半正定性与线性变换律]
\label{prop:ch1-cov-psd}
设 $X \in \R^d$ 为具有有限二阶矩的随机向量，均值 $\mu = \E[X]$，协方差矩阵 $\Sigma = \Cov(X)$。则：
\begin{enumerate}
  \item $\Sigma$ 必为\textbf{实对称半正定矩阵}（$\Sigma = \Sigma^T$ 且 $\Sigma \succeq 0$）；
  \item 对任意确定性矩阵 $A \in \R^{m\times d}$ 与向量 $b \in \R^m$，随机向量 $Y = AX + b$ 的均值与协方差矩阵为
  \begin{equation}
    \E[AX + b] = A\mu + b, \qquad \Cov(AX + b) = A\,\Sigma\,A^T. \label{eq:ch1-cov-affine}
  \end{equation}
\end{enumerate}
\end{proposition}
\end{thmBox}

\begin{proof}
（1）对称性由 $\Sigma_{ij} = \Cov(X_i, X_j) = \Cov(X_j, X_i) = \Sigma_{ji}$ 直接得证。对任意确定性向量 $v \in \R^d$，利用期望对确定性矩阵/向量的线性性：
\[
  v^T \Sigma v
  = v^T \E\bigl[(X - \mu)(X - \mu)^T\bigr] v
  = \E\bigl[v^T(X - \mu)(X - \mu)^T v\bigr]
  = \E\Bigl[\bigl(v^T(X - \mu)\bigr)^2\Bigr]
  = \Var(v^T X) \ge 0.
\]
因为平方的期望永不为负，故 $\Sigma \succeq 0$。

（2）均值式 $\E[AX + b] = A\E[X] + b = A\mu + b$ 由期望线性性逐分量成立。从而 $Y - \E[Y] = A(X - \mu)$，代入协方差矩阵定义：
\[
  \Cov(Y) = \E\bigl[A(X-\mu)(X-\mu)^T A^T\bigr] = A\,\E\bigl[(X-\mu)(X-\mu)^T\bigr] A^T = A\Sigma A^T. \qedhere
\]
\end{proof}

\begin{insightBox}{与第 1 卷谱定理、PCA 和第 2 章多元高斯的无缝合流}
命题 \ref{prop:ch1-cov-psd} 的证明里藏着一句含金量极高的等式：$v^T \Sigma v = \Var(v^T X)$！它意味着：把随机数据 $X$ 投影到任意单位方向 $v$（$\|v\|_2 = 1$）上，所得一维标量的方差恰好是协方差矩阵的 Rayleigh 商 $v^T \Sigma v$。因此，第 1 卷第 4 章的\textbf{主成分分析（PCA）}——寻找使投影方差最大的正交主轴——在概率意义下正是对实对称半正定矩阵 $\Sigma$ 做谱分解 $\Sigma = Q\Lambda Q^T$，特征值 $\lambda_i \ge 0$ 就是第 $i$ 个主成分方向上的精确方差！下一章我们将看到，正是这个 $\Sigma$ 决定了多元高斯分布的椭球几何。
\end{insightBox}

\subsection{纯 NumPy 实战：序贯 Bayes 更新、对数空间朴素贝叶斯与协方差验证}

代码 \ref{lst:ch1-bayes-cov} 用不到 40 行纯 NumPy 同时验证了本章三个核心机制：（1）罕见病序贯 Bayes 更新（复现 $4.72\%$ 与 $71.04\%$）；（2）对数空间 Log-Sum-Exp 朴素贝叶斯（在高维 $d=1000$ 词下对比直接连乘下溢与对数空间稳定性）；（3）经验协方差矩阵的半正定性验证（对应命题 \ref{prop:ch1-cov-psd}）。

\begin{lstlisting}[caption={序贯 Bayes 更新、对数空间防下溢与经验协方差矩阵半正定性验证}, label={lst:ch1-bayes-cov}]
import numpy as np

# 1. 罕见病序贯 Bayes 更新（对齐 vol2_experiments.py）
prior, sens, fpr = 0.001, 0.99, 0.02
p = prior
for step in range(1, 4):
    p = (sens * p) / (sens * p + fpr * (1.0 - p))
    print(f"After test {step} (+): posterior = {p:.5f} ({p*100:.2f}%)")

# 2. 高维 Naive Bayes：直接连乘下溢 vs 对数空间 Log-Sum-Exp
d = 1000
p_feat_spam = np.full(d, 0.012)  # 每词概率约 0.012
p_feat_ham  = np.full(d, 0.010)
raw_spam = 0.4 * np.prod(p_feat_spam)  # 远小于 float64 最小正数 2.2e-308
raw_ham  = 0.6 * np.prod(p_feat_ham)
print(f"Raw product underflow: spam={raw_spam}, ham={raw_ham}")

log_scores = np.array([
    np.log(0.4) + np.sum(np.log(p_feat_spam)),
    np.log(0.6) + np.sum(np.log(p_feat_ham)),
])
m = np.max(log_scores)
probs = np.exp(log_scores - m) / np.sum(np.exp(log_scores - m))
print(f"Log-space posterior: P(spam|x)={probs[0]:.6f}, P(ham|x)={probs[1]:.6f}")

# 3. 经验协方差矩阵半正定性与仿射变换律验证
rng = np.random.default_rng(42)
X = rng.normal(size=(5000, 3))
A = np.array([[2.0, -1.0, 0.5], [0.0, 1.5, 1.0]])  # 2x3 线性映射
Y = X @ A.T + np.array([1.0, -2.0])
S_x = np.cov(X, rowvar=False)
S_y = np.cov(Y, rowvar=False)
eigvals = np.linalg.eigvalsh(S_y)
print(f"min eigval(S_y) = {eigvals.min():.6f} >= 0")
print(f"||S_y - A S_x A^T||_F = {np.linalg.norm(S_y - A @ S_x @ A.T):.2e}")
\end{lstlisting}

在本机运行上述脚本，可以看到：\texttt{raw\_spam} 与 \texttt{raw\_ham} 双双下溢成精确的 \texttt{0.0}，而对数空间版本稳定输出精确后验 $P(\text{spam}\mid x) = 1.000000$；同时经验协方差矩阵的最小特征值严格非负，且仿射变换误差 $\|S_y - A S_x A^T\|_F < 10^{-15}$（机器精度级完全吻合）。

% =========================================================================
\section{本章习题}
\label{sec:ch1-exercises}

\exTier{A. 概念辨析与反例查错}{合上书独立作答，精准定位独立性与概率概念边界}

\begin{enumerate}[label=\textbf{1.\arabic*}, leftmargin=2.4em]
  \item \textbf{（独立、互斥与条件独立辨析）} 判断以下四个命题正误；对错误命题给出最小反例，对正确命题简述理由：
  \begin{enumerate}[label=(\alph*), leftmargin=1.8em]
    \item 若两事件 $A, B$ 互斥（$A\cap B = \varnothing$）且 $\Prob(A)>0,\; \Prob(B)>0$，则 $A$ 与 $B$ 必相互独立；
    \item 若随机变量 $X, Y$ 满足 $X \indep Y$，则对任意第三个随机变量 $Z$，必有 $X \indep Y \mid Z$；
    \item 连续随机变量的概率密度函数 $p_X(x)$ 在任意点处恒满足 $0 \le p_X(x) \le 1$；
    \item 对任意有限期望的随机变量 $X, Y$，即使二者强相关，仍恒有 $\E[X + Y] = \E[X] + \E[Y]$。
  \end{enumerate}

  \item \textbf{（协方差为零与线性相关性陷阱）} 设随机变量 $X$ 满足 $\Prob(X=-2)=\Prob(X=0)=\Prob(X=2)=\frac{1}{3}$，令 $Y = |X|$。请验证 $\Cov(X, Y) = 0$，但 $X$ 与 $Y$ 并不独立；并举例说明为什么“$\Var(X+Y) = \Var(X) + \Var(Y)$”也不能反推出 $X \indep Y$。
\end{enumerate}

\exTier{B. 小例手算与数值直觉}{亲手打通 Bayes 更新、朴素贝叶斯与辩解效应计算}

\begin{enumerate}[label=\textbf{1.\arabic*}, leftmargin=2.4em, resume]
  \item \textbf{（两厂芯片溯源与检测后验）} 在例 \ref{ex:ch1-chip-factory} 中（A 厂供货 $70\%$、次品率 $1\%$；B 厂供货 $30\%$、次品率 $4\%$）：
  \begin{enumerate}[label=(\arabic*), leftmargin=1.8em]
    \item 已知随机抽检的一颗芯片是次品，求它来自 B 厂的后验概率 $\Prob(H_B \mid E)$；
    \item 若改用一台并非完美的检测仪筛查（对真次品报警率 $95\%$，对良品误报率 $2\%$），求随机一颗芯片触发报警的总概率，以及报警芯片确实为次品的后验概率。
  \end{enumerate}

  \item \textbf{（草地湿润中的 Explaining Away 精确手算）} 设“昨晚下雨” $R \in \{0,1\}$ 与“洒水器开启” $S \in \{0,1\}$ 边际独立，先验概率分别为 $\Prob(R=1) = 0.2$、$\Prob(S=1) = 0.1$。草地湿润 $W \in \{0,1\}$ 满足：
  \begin{align*}
    \Prob(W=1\mid R=1, S=1) &= 0.99, & \Prob(W=1\mid R=1, S=0) &= 0.80, \\
    \Prob(W=1\mid R=0, S=1) &= 0.90, & \Prob(W=1\mid R=0, S=0) &= 0.00.
  \end{align*}
  手算以下三个概率并量化辩解效应的幅度：
  \begin{enumerate}[label=(\arabic*), leftmargin=1.8em]
    \item 仅观测到草地湿润时，洒水器开启的后验概率 $\Prob(S=1 \mid W=1)$；
    \item 既观测到草地湿润、又确认昨晚下过雨时，洒水器开启的后验概率 $\Prob(S=1 \mid W=1, R=1)$；
    \item 对比 $\Prob(S=1)$、$\Prob(S=1\mid W=1)$ 与 $\Prob(S=1\mid W=1, R=1)$ 的变化方向，用一句话解释其含义。
  \end{enumerate}
\end{enumerate}

\exTier{C. 实用推导与性质证明}{AI 概率建模与后续卷次反复使用的核心恒等式}

\begin{enumerate}[label=\textbf{1.\arabic*}, leftmargin=2.4em, resume]
  \item \textbf{（全方差定律：Law of Total Variance）} 设随机变量 $X, Y$ 具有有限二阶矩，证明著名的\textbf{方差分解恒等式（全方差定律）}：
  \[
    \Var(X) = \E\bigl[\Var(X \mid Y)\bigr] + \Var\bigl(\E[X \mid Y]\bigr),
  \]
  其中 $\Var(X \mid Y) := \E[(X - \E[X \mid Y])^2 \mid Y]$。并解释两项的直观含义（类内残差方差 + 类间均值方差）。
  \exHint{利用 $\Var(X) = \E[X^2] - (\E[X])^2$ 与全期望定律 $\E[X^2] = \E[\E[X^2\mid Y]]$，在中间加上并减去 $\E\bigl[(\E[X\mid Y])^2\bigr]$，再整理成两项即可。}

  \item \textbf{（Log-Odds 形式下的朴素贝叶斯线性分类器本质）} 考虑二分类朴素贝叶斯模型 \eqref{eq:ch1-naive-bayes}，其中特征 $x = (x_1, \dots, x_d)^T \in \{0, 1\}^d$，记 $p_j := \Prob(x_j=1 \mid y=1) \in (0,1)$，$q_j := \Prob(x_j=1 \mid y=0) \in (0,1)$，先验 $\pi := \Prob(y=1) \in (0,1)$。证明其对数后验几率（log-odds）恰好是特征 $x$ 的\textbf{线性函数}：
  \[
    \log \frac{\Prob(y=1 \mid x)}{\Prob(y=0 \mid x)} = w^T x + b,
  \]
  并写出权重向量 $w \in \R^d$ 与偏置 $b \in \R$ 的显式表达式（说明为何第 6 章的逻辑回归与朴素贝叶斯拥有完全相同的线性决策边界形式）。
  \exHint{将 $\Prob(x_j \mid y=1)$ 写成 $p_j^{x_j}(1-p_j)^{1-x_j}$，取对数比值后将含 $x_j$ 的项与常数项分离。}
\end{enumerate}

\exTier{D. 编程实战与可自检实验}{用纯 NumPy 跑出概率机制并通过数值断言}

\begin{enumerate}[label=\textbf{1.\arabic*}, leftmargin=2.4em, resume]
  \item \textbf{（高维对数空间朴素贝叶斯与 Laplace 平滑实验）} 编写纯 NumPy 函数实现二值特征朴素贝叶斯的训练（含 Laplace 平滑系数 $\alpha \ge 0$：$\hat{p}_{jc} = \frac{N_{jc} + \alpha}{N_c + 2\alpha}$）与对数空间 Log-Sum-Exp 预测：
  \begin{enumerate}[label=(\arabic*), leftmargin=1.8em]
    \item 生成 $N_{\text{tr}}=400$ 条训练样本、$N_{\text{te}}=400$ 条测试样本、特征维度 $d=4000$ 的模拟二分类文本数据（类别 $1$ 的前 $40$ 个关键词出现概率为 $0.35$、其余为 $0.05$；类别 $0$ 所有词出现概率为 $0.05$）；
    \item 对比 $\alpha = 0$（无平滑）与 $\alpha = 1.0$（Laplace 平滑）在测试集出现训练集未见词时的表现（是否产生 \texttt{NaN}/\texttt{-inf}），并对比直接概率连乘与对数空间预测；
    \item 输出测试集准确率，并断言所有预测后验概率之和与 $1$ 的最大偏差小于 $10^{-12}$。
  \end{enumerate}
  \par\smallskip\noindent{\footnotesize\sffamily\color{tealDark}\textbf{【自检标准】}\,在 $d=4000$ 维度下直接概率连乘必须复现大面积浮点下溢（超过半数测试样本两类分子同时为 \texttt{0.0}）；对数空间 + Laplace 平滑（$\alpha=1.0$）版本须无任何 \texttt{NaN}，所有样本满足 $|\Prob(y=0\mid x)+\Prob(y=1\mid x)-1| < 10^{-12}$，且测试准确率高于 $90\%$。}

  \item \textbf{（蒙特卡洛验证 Explaining Away 与条件协方差）} 用 \texttt{np.random.default\_rng(2026)} 独立采样 $N = 200{,}000$ 组习题 1.4 中的 $(R, S, W)$：
  \begin{enumerate}[label=(\arabic*), leftmargin=1.8em]
    \item 经验计算边际相关系数 $\mathrm{Corr}(R, S)$ 与经验概率 $\hat{\Prob}(S=1 \mid W=1)$、$\hat{\Prob}(S=1 \mid W=1, R=1)$；
    \item 进一步仅在子集 $W = 1$ 内计算条件相关系数 $\mathrm{Corr}(R, S \mid W=1)$，定量展示“原本无关的两变量，经由观测共同结果 $W=1$ 后变为显著负相关”。
  \end{enumerate}
  \par\smallskip\noindent{\footnotesize\sffamily\color{tealDark}\textbf{【自检标准】}\,边际相关系数 $|\mathrm{Corr}(R, S)| < 0.01$；条件后验估计与习题 1.4 理论解析值偏差均小于 $0.005$；给定 $W=1$ 后的条件相关系数必须严格为负（落在 $[-0.86, -0.80]$ 区间内）。}
\end{enumerate}

\chapter[AI 核心概率分布与多元高斯]{AI 核心概率分布与多元高斯分布}
\label{ch:ch2}

\chapEpigraph{我几乎不知道还有什么东西，能像误差频率定律那样令人惊叹：它在最狂野的混沌之中，悄然主宰着宇宙的秩序。}{弗朗西斯·高尔顿（Francis Galton），《自然遗传》}

\begin{chapIntro}{本章定位：认识人工智能的“概率乐高积木”}
传统概率教材常按历史顺序罗列十余种分布，却很少回答“神经网络和机器学习算法到底在用哪几种，分别用在哪一层”。事实上，从逻辑回归、Softmax 分类器、大语言模型采样，到变分自编码器（VAE）、高斯过程（GP）与卡尔曼滤波，现代 AI 真正高频依赖的概率分布屈指可数，但对它们的\textbf{几何结构与矩阵运算法则}要求极深。

本章将为你组装一套精炼而完整的 AI 分布工具箱，并将第 1 卷的谱定理与 Cholesky 分解彻底变现为概率几何直觉：
\begin{enumerate}
  \item 系统掌握四大离散主力分布（\textbf{Bernoulli、Binomial、Categorical、Poisson}）与四大连续主力分布（\textbf{Uniform、Exponential、Beta、一维高斯}），明确每一种分布在 AI 中的真实岗位，并用图形直观看见 Beta 分布如何刻画“概率的概率”；
  \item 以第 1 卷的\textbf{实对称阵谱定理 $\Sigma = Q\Lambda Q^T$} 为手术刀，彻底解剖\textbf{多元高斯分布 $\mathcal{N}(\mu, \Sigma)$} 的椭球几何与\textbf{马氏距离（Mahalanobis Distance）}，区分各向同性、对角与全协方差三类结构；
  \item 掌握多元高斯的\textbf{仿射变换封闭性}，顺势打通\textbf{Cholesky 采样}与第 5 卷深度生成模型的\textbf{重参数化技巧（Reparameterization Trick）}；
  \item 建立多元高斯的\textbf{边际与条件分布公式}（主线配具体二维算例与实测切片，选读框提供优雅的分块矩阵 Schur 补推导），为线性回归、卡尔曼滤波与高斯过程铺平道路。
\end{enumerate}
\end{chapIntro}

% =========================================================================
\section{离散主力分布：Bernoulli、Binomial、Categorical 与 Poisson}
\label{sec:ch2-discrete}

当模型的输出是类别标签、词表索引或事件计数时，我们需要离散概率分布来建模。表 \ref{tab:ch2-discrete} 汇总了 AI 中最常出现的四个离散主力分布及其核心指标，随后逐一剖析它们的推导与工程角色。

\begin{table}[htbp]
  \centering
  \caption{人工智能四大离散主力分布速查表}
  \label{tab:ch2-discrete}
  \footnotesize
  \begin{tabularx}{\linewidth}{@{}p{23mm}p{20mm}p{33mm}p{15mm}p{17mm}X@{}}
    \toprule
    \textbf{分布名称} & \textbf{记号与支撑} & \textbf{概率质量函数（PMF）} & $\E[X]$ & $\Var(X)$ & \textbf{AI 核心岗位} \\
    \midrule
    \textbf{Bernoulli} & $\mathrm{Bern}(p)$ $x\in\{0,1\}$ & $p^x(1-p)^{1-x}$ & $p$ & $p(1-p)$ & 二分类标签（逻辑回归 / BCE）、Dropout 随机掩码 \\
    \textbf{Binomial} & $\mathrm{Bin}(n,p)$ $k\in\{0,\dots,n\}$ & $\binom{n}{k}p^k(1-p)^{n-k}$ & $np$ & $np(1-p)$ & 测试集答对题数、A/B 测试点击次数统计 \\
    \textbf{Categorical} & $\mathrm{Cat}(\pi)$ $y\in\{e_1,\dots,e_K\}$ & $\prod_{k=1}^K \pi_k^{y_k}$ & $\pi$（向量） & $\diag(\pi)-\pi\pi^T$ & 多分类 Softmax 输出、大语言模型下一词元分布 \\
    \textbf{Poisson} & $\mathrm{Pois}(\lambda)$ $k\in\N$ & $e^{-\lambda}\dfrac{\lambda^k}{k!}$ & $\lambda$ & $\lambda$ & 单位时间请求/点击计数建模、泊松回归 \\
    \bottomrule
  \end{tabularx}
\end{table}

\subsection{Bernoulli 与 Binomial：从单次试验到独立求和}

\textbf{（1）Bernoulli 分布} $\mathrm{Bern}(p)$（$p \in [0, 1]$）描述最简单的“是/否”随机事件：$X=1$ 的概率为 $p$，$X=0$ 的概率为 $1-p$。它的 PMF 可以写成极其紧凑的指数形式：
\begin{equation}
  p(x) = p^x (1 - p)^{1 - x}, \qquad x \in \{0, 1\}. \label{eq:ch2-bern}
\end{equation}
这个写法并非炫技——第 5、6 章对它取负对数 $-\log p(x) = -\bigl[x\log p + (1-x)\log(1-p)\bigr]$ 时，\textbf{二元交叉熵损失（Binary Cross-Entropy, BCE）}会原封不动地跳出来！直接由定义立得期望与方差：
\[
  \E[X] = 1\cdot p + 0\cdot(1-p) = p, \qquad
  \E[X^2] = 1^2\cdot p + 0^2\cdot(1-p) = p \;\implies\; \Var(X) = p - p^2 = p(1 - p).
\]
观察方差函数 $p(1-p)$：在 $p=0$ 或 $p=1$（结果确定）时方差为 $0$；在 $p=0.5$（最不确定）时方差达到最大值 $1/4$。在深度学习中，\textbf{Dropout} 正则化层正是按 $\mathrm{Bern}(1 - p_{\text{drop}})$ 为每个神经元生成 $0/1$ 随机掩码。

\textbf{（2）Binomial 分布} $\mathrm{Bin}(n, p)$ 描述 $n$ 次\textbf{独立同分布（i.i.d.）} $\mathrm{Bern}(p)$ 试验的总成功次数 $S_n = \sum_{i=1}^n X_i$。若用组合恒等式直接求 $\sum_{k=0}^n k\binom{n}{k}p^k(1-p)^{n-k}$ 相当繁琐，但借助第 1 章的\textbf{期望线性性}与独立变量方差可加性，只需两行即可秒杀：
\[
  \E[S_n] = \sum_{i=1}^n \E[X_i] = np, \qquad
  \Var(S_n) = \sum_{i=1}^n \Var(X_i) = np(1 - p).
\]
例如，若一个分类器在单条独立测试样本上的真实准确率为 $p = 0.85$，在 $n = 180$ 条测试样本上答对的题数就服从 $\mathrm{Bin}(180, 0.85)$，均值 $153$ 题，标准差 $\sqrt{180\times 0.85\times 0.15} \approx 4.79$ 题（对应准确率标准误约 $2.66\%$，第 5 章 Bootstrap 实测恰为 $2.63\%$！）。

\subsection{Categorical 与 Poisson：多分类 Softmax 与稀疏计数}

\textbf{（3）Categorical 分布（范畴分布）} $\mathrm{Cat}(\pi_1, \dots, \pi_K)$ 是 Bernoulli 分布向 $K$ 个互斥类别的自然推广，其中概率向量 $\pi = (\pi_1, \dots, \pi_K)^T$ 满足 $\pi_k \ge 0$ 且 $\sum_{k=1}^K \pi_k = 1$（位于 $K-1$ 维概率单纯形 $\Delta^{K-1}$ 上）。在机器学习中，我们通常把类别标签表示为\textbf{独热向量（One-Hot Vector）} $y \in \{0, 1\}^K$（仅真实类别所在分量为 $1$，其余为 $0$），从而 PMF 可统一写作：
\begin{equation}
  p(y) = \prod_{k=1}^K \pi_k^{y_k}, \qquad -\log p(y) = -\sum_{k=1}^K y_k \log \pi_k. \label{eq:ch2-cat}
\end{equation}
请牢记式 \eqref{eq:ch2-cat} 右边那一项——它就是贯穿所有多分类神经网络与大语言模型训练的\textbf{多类交叉熵损失}！当把 $n$ 次独立的 $\mathrm{Cat}(\pi)$ 试验向量相加，就得到\textbf{多项分布（Multinomial）}，常用于词袋模型（Bag-of-Words）统计文本词频。

\textbf{（4）Poisson 分布} $\mathrm{Pois}(\lambda)$（速率参数 $\lambda > 0$）描述固定时间或空间窗口内稀有独立事件的发生次数（如每秒 API 请求数、单位面积缺陷数），PMF 为 $p(k) = e^{-\lambda}\frac{\lambda^k}{k!}$（$k=0, 1, 2, \dots$）。它是二项分布 $\mathrm{Bin}(n, \lambda/n)$ 在 $n \to \infty$、事件发生概率极小而平均次数保持为 $\lambda$ 时的极限。利用泰勒级数 $e^\lambda = \sum_{k=0}^\infty \frac{\lambda^k}{k!}$ 容易验证：
\[
  \E[X] = \sum_{k=1}^\infty k e^{-\lambda}\frac{\lambda^k}{k!} = \lambda e^{-\lambda}\sum_{k=1}^\infty \frac{\lambda^{k-1}}{(k-1)!} = \lambda, \qquad
  \E[X(X-1)] = \lambda^2 \;\implies\; \Var(X) = \lambda.
\]
期望与方差严格相等（$\E[X] = \Var(X) = \lambda$）是识别数据是否近似服从 Poisson 分布的快捷试金石。

\begin{insightBox}{现代神经网络不是在“猜答案”，而是在输出分布的参数}
初学者常误以为分类神经网络输出的是一个硬邦邦的类别编号（如“猫”）。实际上，无论是二分类的 Sigmoid 输出 $p_\theta(x) \in (0,1)$，还是多分类与大语言模型的 Softmax 输出 $\pi_\theta(x) \in \Delta^{K-1}$，网络输出的都是\textbf{概率分布的参数}——分别定义了一个 $\mathrm{Bern}(p_\theta(x))$ 或 $\mathrm{Cat}(\pi_\theta(x))$ 分布！只有输出合法的概率分布，我们才能在第 4--6 章用极大似然估计（MLE）和交叉熵损失对其进行端到端梯度训练，并在推理时量化预测置信度。
\end{insightBox}

% =========================================================================
\section{连续主力分布：Uniform、Exponential、Beta 与一维高斯}
\label{sec:ch2-continuous}

当需要建模连续参数、等待时间、先验概率权重或真实世界的连续测量值时，我们转向四大连续主力分布（表 \ref{tab:ch2-continuous}）。

\begin{table}[htbp]
  \centering
  \caption{人工智能四大连续主力分布速查表}
  \label{tab:ch2-continuous}
  \footnotesize
  \begin{tabularx}{\linewidth}{@{}p{21mm}p{19mm}p{33mm}p{14mm}p{26mm}X@{}}
    \toprule
    \textbf{分布名称} & \textbf{记号与支撑} & \textbf{概率密度函数（PDF）} & $\E[X]$ & $\Var(X)$ & \textbf{AI 核心岗位} \\
    \midrule
    \textbf{Uniform} & $\mathcal{U}(a,b)$ $x\in[a,b]$ & $\dfrac{1}{b-a}$ & $\dfrac{a+b}{2}$ & $\dfrac{(b-a)^2}{12}$ & 神经网络权重初始化（Xavier/He）、随机超参搜索 \\
    \textbf{Exponential} & $\mathrm{Exp}(\lambda)$ $x\ge 0$ & $\lambda e^{-\lambda x}$ & $\dfrac{1}{\lambda}$ & $\dfrac{1}{\lambda^2}$ & 等待时间/生存分析、Gumbel-Max 与重采样构造 \\
    \textbf{Beta} & $\mathrm{Beta}(\alpha,\beta)$ $x\in(0,1)$ & $\dfrac{x^{\alpha-1}(1-x)^{\beta-1}}{B(\alpha,\beta)}$ & $\dfrac{\alpha}{\alpha+\beta}$ & $\frac{\alpha\beta}{(\alpha+\beta)^2(\alpha+\beta+1)}$ & 概率 $p$ 的先验分布、Beta-Bernoulli 共轭、Thompson 采样 \\
    \textbf{1D Gaussian} & $\mathcal{N}(\mu,\sigma^2)$ $x\in\R$ & $\dfrac{1}{\sqrt{2\pi}\sigma}\exp\!\left(-\dfrac{(x-\mu)^2}{2\sigma^2}\right)$ & $\mu$ & $\sigma^2$ & 回归噪声、权重 $L_2$ 先验、激活标准化与中心极限 \\
    \bottomrule
  \end{tabularx}
\end{table}

\subsection{Uniform 与 Exponential：随机初始化与无记忆性}

\textbf{（1）Uniform 分布} $\mathcal{U}(a, b)$ 在区间 $[a, b]$ 上均匀铺开概率质量，是所有计算机伪随机数生成器（\texttt{rng.random()} 产出 $\mathcal{U}(0, 1)$）的源头。在深度学习中，为了防止深层网络前向信号爆炸或消失，Xavier 初始化把线性层权重从 $\mathcal{U}(-r, r)$ 采样，并利用方差公式 $\Var = r^2 / 3$ 精确控制每层输出方差保持稳定。

\textbf{（2）Exponential 分布} $\mathrm{Exp}(\lambda)$（$\lambda > 0$）描述 Poisson 过程中相邻两次事件的时间间隔，其生存函数（右尾概率）极简：$\Prob(X > t) = \int_t^\infty \lambda e^{-\lambda x}\,dx = e^{-\lambda t}$。由此立得其独一无二的\textbf{无记忆性（Memoryless Property）}：对任意 $s, t \ge 0$，
\[
  \Prob(X > s + t \mid X > s) = \frac{\Prob(X > s + t)}{\Prob(X > s)} = \frac{e^{-\lambda(s+t)}}{e^{-\lambda s}} = e^{-\lambda t} = \Prob(X > t).
\]
也就是说，如果一台服务器的无故障间隔服从指数分布，那么它“已经平稳运行了 $s$ 小时”这一事实，并不会让它未来还能撑过 $t$ 小时的概率变高或变低。

\subsection{Beta 分布：刻画“关于概率的不确定性”}

传统工科概率课几乎不讲 Beta 分布，但在贝叶斯机器学习、推荐系统探索-利用（Multi-Armed Bandit）和第 5 章的 MAP 估计中，Beta 分布是不可或缺的主角。为什么需要它？因为当我们想给一枚硬币的正面概率 $p \in (0, 1)$ 赋予一个\textbf{先验概率分布}时，支撑集必须恰好落在开区间 $(0, 1)$ 上！

\begin{defBox}
\begin{definition}[Beta 分布]
\label{def:ch2-beta}
设形状参数 $\alpha > 0,\; \beta > 0$。若连续随机变量 $X \in (0, 1)$ 的概率密度函数为
\begin{equation}
  p(x \mid \alpha, \beta) = \frac{1}{B(\alpha, \beta)}\, x^{\alpha - 1}(1 - x)^{\beta - 1}, \qquad x \in (0, 1), \label{eq:ch2-beta}
\end{equation}
其中归一化常数（Beta 函数）$B(\alpha, \beta) := \int_0^1 t^{\alpha-1}(1-t)^{\beta-1}\,dt = \frac{\Gamma(\alpha)\Gamma(\beta)}{\Gamma(\alpha+\beta)}$（对正整数有 $\Gamma(n) = (n-1)!$），则称 $X \sim \mathrm{Beta}(\alpha, \beta)$。其均值与方差分别为：
\begin{equation}
  \E[X] = \frac{\alpha}{\alpha + \beta}, \qquad
  \Var(X) = \frac{\alpha\beta}{(\alpha+\beta)^2(\alpha+\beta+1)}. \label{eq:ch2-beta-moments}
\end{equation}
\end{definition}
\end{defBox}

把式 \eqref{eq:ch2-beta} 的核 $x^{\alpha-1}(1-x)^{\beta-1}$ 与二项分布似然 $p^k(1-p)^{n-k}$ 放在一起看，你会立刻恍然大悟：\textbf{$\alpha - 1$ 和 $\beta - 1$ 恰好扮演了“先验伪计数（pseudo-counts）”的角色}——相当于在正式做实验前，我们已经在脑海里预设看到了 $\alpha - 1$ 次成功和 $\beta - 1$ 次失败！总强度 $\alpha + \beta$ 越大，密度峰越尖锐，代表先验信念越笃定（方差按 $\approx 1/(\alpha+\beta)$ 缩小）。

图 \ref{fig:ch2-beta-shapes} 绘制了四组典型参数的 $\mathrm{Beta}(\alpha, \beta)$ 密度曲线，让你一眼看清它极强的形状可塑性：
\begin{itemize}
  \item $\mathrm{Beta}(1, 1)$：$p(x)=1$，退化为 $\mathcal{U}(0,1)$，代表“完全无知的均匀先验”；
  \item $\mathrm{Beta}(2, 2)$：$p(x)=6x(1-x)$，对称钟形，柔和偏好 $p$ 靠近 $0.5$ 而避开极端值 $0$ 和 $1$；
  \item $\mathrm{Beta}(5, 2)$：$p(x)=30x^4(1-x)$，明显右偏（均值 $5/7 \approx 0.714$，峰值在 $x=0.8$），表示先验相信成功率较高；
  \item $\mathrm{Beta}(0.5, 0.5)$：$p(x)=\frac{1}{\pi\sqrt{x(1-x)}}$，呈 U 形（Jeffreys 无信息先验），认为 $p$ 更可能接近两个极端。
\end{itemize}

\begin{figure}[htbp]
  \centering
  \begin{tikzpicture}
    \begin{axis}[
      width=0.80\linewidth,
      height=5.6cm,
      xlabel={概率参数 $x \in (0, 1)$},
      ylabel={概率密度 $p(x \mid \alpha, \beta)$},
      xmin=0, xmax=1,
      ymin=0, ymax=3.15,
      xtick={0, 0.2, 0.4, 0.6, 0.8, 1.0},
      ytick={0, 0.5, 1.0, 1.5, 2.0, 2.5, 3.0},
      grid=major,
      grid style={dashed, gray!25},
      legend pos=north west,
      legend style={at={(0.5,1.035)}, anchor=south, cells={anchor=west}, font=\scriptsize, fill=white, draw=codeborder, rounded corners=1.5pt},
      tick label style={font=\footnotesize},
      label style={font=\small\sffamily}
    ]
      % Beta(1,1) = 1
      \addplot[color=slate, thick, dashed, domain=0:1, samples=40] {1.0};
      \addlegendentry{$\mathrm{Beta}(1, 1)$：均匀无知先验}

      % Beta(2,2) = 6 x (1-x)
      \addplot[color=teal, very thick, domain=0:1, samples=100] {6*x*(1-x)};
      \addlegendentry{$\mathrm{Beta}(2, 2)$：对称中庸先验（峰值 $1.5$）}

      % Beta(5,2) = 30 x^4 (1-x)
      \addplot[color=navy, very thick, domain=0:1, samples=120] {30*(x^4)*(1-x)};
      \addlegendentry{$\mathrm{Beta}(5, 2)$：偏高先验（均值 $5/7$，峰值 $x=0.8$）}

      % Beta(0.5, 0.5) = 1 / (pi * sqrt(x*(1-x)))
      \addplot[color=ochre, thick, dotted, domain=0.015:0.985, samples=140] {1.0 / (3.14159265 * sqrt(x*(1-x)))};
      \addlegendentry{$\mathrm{Beta}(0.5, 0.5)$：U 形极端偏好先验}

      % Peak annotation for Beta(5,2)
      \addplot[only marks, mark=*, mark size=2pt, color=navy] coordinates {(0.8, 2.4576)};
      \node[anchor=east, font=\scriptsize\sffamily, color=navy, fill=white, fill opacity=0.85, text opacity=1, inner sep=1.5pt] at (axis cs:0.75, 2.68) {众数 $x=0.8$};
      \draw[-{Stealth[length=1.8mm]}, color=navy, thin] (axis cs:0.75, 2.60) -- (axis cs:0.793, 2.48);
    \end{axis}
  \end{tikzpicture}
  \caption{四组典型参数下的 $\mathrm{Beta}(\alpha, \beta)$ 概率密度曲线。仅调节 $(\alpha, \beta)$ 两个正数，就能在 $(0,1)$ 区间上平滑地表达“完全无知”、“中庸对称”、“偏向高成功率”乃至“偏向极端值”等任意先验信念；第 5 章的 Beta-Bernoulli 共轭更新将直接在这条曲线族内部闭环演化。}
  \label{fig:ch2-beta-shapes}
\end{figure}

\subsection{一维高斯分布：标准化与 68--95--99.7 经验法则}

一维高斯分布（正态分布）$\mathcal{N}(\mu, \sigma^2)$（均值 $\mu \in \R$，标准差 $\sigma > 0$）是连续世界最核心的基准分布：
\[
  p(x \mid \mu, \sigma^2) = \frac{1}{\sqrt{2\pi}\,\sigma} \exp\!\left(-\frac{(x - \mu)^2}{2\sigma^2}\right), \qquad x \in \R.
\]
对任意 $X \sim \mathcal{N}(\mu, \sigma^2)$，做\textbf{标准化（Z-score 变换）} $Z := \frac{X - \mu}{\sigma}$ 总能将其拉回标准正态分布 $\mathcal{N}(0, 1)$（反过来 $X = \mu + \sigma Z$）。记忆三条极其实用的置信区间概率（68--95--99.7 法则），能让你一眼心算模型误差的合理波动范围：
\[
  \Prob(|X-\mu|\le \sigma) \approx 68.27\%, \qquad
  \Prob(|X-\mu|\le 2\sigma) \approx 95.45\%, \qquad
  \Prob(|X-\mu|\le 3\sigma) \approx 99.73\%.
\]
高斯分布之所以在 AI 与统计学中占据统治地位，源于三大支柱：第 3 章将证明大量独立小扰动之和趋近高斯（\textbf{中心极限定理}）；第 4 章将证明在给定均值与方差的所有连续分布中，高斯具有最大熵（\textbf{最大熵原理}）；而接下来两节将揭示它在高维空间中无可替代的解析优雅性。

% =========================================================================
\section{多元高斯分布 $\mathcal{N}(\mu, \Sigma)$：谱定理下的椭球几何与马氏距离}
\label{sec:ch2-mvg-geom}

现实中的 AI 样本几乎从不是孤立的一维标量，而是存在复杂协同关系的 $d$ 维特征向量（如图像隐变量、机器人状态向量、金融资产收益）。把一维高斯的标量方差 $\sigma^2 > 0$ 升级为对称正定协方差矩阵 $\Sigma \succ 0$，便得到高维连续概率建模的绝对支柱——\textbf{多元高斯分布（Multivariate Gaussian, MVG）}。

\begin{defBox}
\begin{definition}[多元高斯分布 $\mathcal{N}(\mu, \Sigma)$]
\label{def:ch2-mvg}
设均值向量 $\mu \in \R^d$，协方差矩阵 $\Sigma \in \R^{d\times d}$ 为\textbf{实对称正定矩阵}（$\Sigma = \Sigma^T \succ 0$）。若随机向量 $X \in \R^d$ 具有概率密度函数
\begin{equation}
  p(x \mid \mu, \Sigma)
  = \frac{1}{(2\pi)^{d/2}(\det\Sigma)^{1/2}}
    \exp\!\left(-\frac{1}{2}(x - \mu)^T \Sigma^{-1} (x - \mu)\right), \qquad x \in \R^d,
  \label{eq:ch2-mvg-pdf}
\end{equation}
则称 $X$ 服从 $d$ 维多元高斯分布，记作 $X \sim \mathcal{N}(\mu, \Sigma)$。
\end{definition}
\end{defBox}

一眼看去，式 \eqref{eq:ch2-mvg-pdf} 像是一团令人生畏的矩阵指数，其实它完全可以从最简单的标准各向同性高斯 $Z = (Z_1, \dots, Z_d)^T \sim \mathcal{N}(0, I_d)$ 生长出来！因为 $Z_1, \dots, Z_d$ 独立同分布于 $\mathcal{N}(0,1)$，其联合密度就是一维密度的连乘：
\[
  p_Z(z) = \prod_{i=1}^d \frac{1}{\sqrt{2\pi}} e^{-z_i^2 / 2} = \frac{1}{(2\pi)^{d/2}} \exp\!\left(-\frac{1}{2} z^T z\right).
\]
现在用对称正定平方根 $\Sigma^{1/2}$（第 1 卷第 3 章）对 $Z$ 做仿射拉伸与平移：$X = \mu + \Sigma^{1/2} Z$，反解得 $z = \Sigma^{-1/2}(x - \mu)$。根据多元微积分的变量代换公式，体积拉伸雅可比行列式为 $|\det(\frac{\partial x}{\partial z})| = |\det(\Sigma^{1/2})| = (\det\Sigma)^{1/2}$，而指数上的 $z^T z$ 恰好变成
\[
  z^T z = (x - \mu)^T \Sigma^{-1/2} \Sigma^{-1/2} (x - \mu) = (x - \mu)^T \Sigma^{-1} (x - \mu)!
\]
这正是式 \eqref{eq:ch2-mvg-pdf} 中归一化因子 $(\det\Sigma)^{-1/2}$ 与二次型 $(x-\mu)^T\Sigma^{-1}(x-\mu)$ 的全部来历！

\subsection{第 1 卷谱定理的巅峰变现：椭球等高面与马氏距离}

现在请出第 1 卷第 3 章的王牌定理——\textbf{实对称矩阵谱定理}！由于 $\Sigma \succ 0$，它必存在正交特征分解
\[
  \Sigma = Q \Lambda Q^T = \sum_{i=1}^d \lambda_i q_i q_i^T,
\]
其中正交矩阵 $Q = [q_1, \dots, q_d]$（$Q^T Q = I_d$）的列向量是单位正交主轴，$\Lambda = \diag(\lambda_1, \dots, \lambda_d)$ 的特征值全部严格为正（$\lambda_i > 0$），且其逆矩阵即为 $\Sigma^{-1} = Q \Lambda^{-1} Q^T$。

将中心化向量 $x - \mu$ 投影到特征向量构成的新正交坐标系下，记新坐标为 $y := Q^T(x - \mu)$（第 $i$ 个分量 $y_i = q_i^T(x - \mu)$）。立刻把二次型彻底解耦：
\begin{equation}
  d_M^2(x, \mu) := (x - \mu)^T \Sigma^{-1} (x - \mu)
  = (x - \mu)^T Q \Lambda^{-1} Q^T (x - \mu)
  = y^T \Lambda^{-1} y
  = \sum_{i=1}^d \frac{y_i^2}{\lambda_i}.
  \label{eq:ch2-mahalanobis-decouple}
\end{equation}

\begin{thmBox}
\begin{proposition}[多元高斯的椭球几何与马氏距离]
\label{prop:ch2-mvg-ellipsoid}
设 $X \sim \mathcal{N}(\mu, \Sigma)$（$\Sigma = Q\Lambda Q^T \succ 0$）。定义点 $x$ 到分布中心 $\mu$ 的\textbf{马氏距离（Mahalanobis Distance）}为
\[
  d_M(x, \mu) := \sqrt{(x - \mu)^T \Sigma^{-1} (x - \mu)} = \norm{\Sigma^{-1/2}(x - \mu)}_2.
\]
则：
\begin{enumerate}
  \item 概率密度 $p(x)$ 仅通过 $d_M^2(x, \mu)$ 单调递减地依赖于 $x$；
  \item 等密度曲面 $\{x \in \R^d : d_M(x, \mu) = c\}$（$c > 0$）是以 $\mu$ 为中心、以特征向量 $q_1, \dots, q_d$ 为主轴方向、各主轴半长为 $c\sqrt{\lambda_1}, \dots, c\sqrt{\lambda_d}$ 的\textbf{$d$ 维超椭球}。
\end{enumerate}
\end{proposition}
\end{thmBox}

马氏距离在异常检测、度量学习与对比表示中极为关键：普通欧氏距离 $\|x-\mu\|_2$ 盲目地认为所有方向同等重要，而马氏距离先用 $\Sigma^{-1/2}$ 把数据“白化”成各向同性球体，再量欧氏距离——等价于\textbf{在每个主轴方向上以“该方向的标准差 $\sqrt{\lambda_i}$”为尺子去衡量偏离程度}！因此，沿着数据天然拉长的长轴偏离 $2$ 个单位可能很常见，但沿着数据紧密压缩的短轴偏离 $2$ 个单位却是极罕见的异常点。

\subsection{三种经典协方差结构与参数量权衡}

在真实高维 AI 模型（如高斯混合模型 GMM、VAE 编码器、扩散模型）中，根据精度与算力的权衡，协方差矩阵 $\Sigma$ 通常取三种层级（图 \ref{fig:ch2-mvg-contours}）：
\begin{enumerate}
  \item \textbf{各向同性协方差（Isotropic）} $\Sigma = \sigma^2 I_d$：仅 $1$ 个方差参数，等高线为\textbf{正圆/超球}，各维独立且方差相同；
  \item \textbf{对角协方差（Diagonal）} $\Sigma = \diag(\sigma_1^2, \dots, \sigma_d^2)$：仅 $d$ 个参数，等高线为\textbf{与坐标轴平行的椭圆}，各维相互独立但尺度不同（VAE 隐空间最常用的参数化方式）；
  \item \textbf{全协方差（Full Covariance）} $\Sigma \succ 0$：含 $\frac{d(d+1)}{2}$ 个独立参数，等高线为\textbf{任意倾斜的椭圆}，能完整刻画维度之间的线性相关性，但在高维（如 $d = 10{,}000$）时存储与求逆代价高昂。
\end{enumerate}

\begin{figure}[htbp]
  \centering
  \begin{tikzpicture}[font=\footnotesize\sffamily]
    % -------- Panel 1: Isotropic --------
    \begin{scope}[xshift=0cm]
      \node[font=\small\sffamily\bfseries, color=tealDark] at (0, 2.76) {(a) 各向同性 $\Sigma = \sigma^2 I$};
      \draw[-{Stealth[length=1.8mm]}, slate!70] (-2.05, 0) -- (2.15, 0) node[right, font=\scriptsize] {$x_1$};
      \draw[-{Stealth[length=1.8mm]}, slate!70] (0, -1.95) -- (0, 2.05) node[above, font=\scriptsize] {$x_2$};

      \fill[tealLight, opacity=0.45] (0,0) circle (1.50cm);
      \fill[tealLight, opacity=0.70] (0,0) circle (1.00cm);
      \draw[teal!50, thick] (0,0) circle (1.50cm);
      \draw[teal, thick] (0,0) circle (1.00cm);
      \draw[tealDark, very thick] (0,0) circle (0.50cm);
      \fill[ink] (0,0) circle (1.4pt);

      \draw[{Stealth[length=1.6mm]}-{Stealth[length=1.6mm]}, ochre, thick] (0,0) -- (0.707, 0.707) node[midway, above left, xshift=0pt, yshift=2pt, fill=white, inner sep=0.6pt, font=\scriptsize, text=ochre] {$\sigma$};
      \node[below, font=\scriptsize, color=slate, align=center] at (0, -2.10) {正圆等高线，各向等方差\\参数量：$1$};
    \end{scope}

    % -------- Panel 2: Diagonal --------
    \begin{scope}[xshift=5.3cm]
      \node[font=\small\sffamily\bfseries, color=navy] at (0, 2.76) {(b) 对角阵 $\Sigma = \diag(\sigma_1^2, \sigma_2^2)$};
      \draw[-{Stealth[length=1.8mm]}, slate!70] (-2.15, 0) -- (2.25, 0) node[right, font=\scriptsize] {$x_1$};
      \draw[-{Stealth[length=1.8mm]}, slate!70] (0, -1.95) -- (0, 2.05) node[above, font=\scriptsize] {$x_2$};

      \fill[navyLight, opacity=0.55] (0,0) ellipse (1.80cm and 0.75cm);
      \draw[navy!50, thick] (0,0) ellipse (1.80cm and 0.75cm);
      \draw[navy, thick] (0,0) ellipse (1.20cm and 0.50cm);
      \draw[navy, very thick] (0,0) ellipse (0.60cm and 0.25cm);
      \fill[ink] (0,0) circle (1.4pt);

      \draw[-{Stealth[length=1.6mm]}, ochre, thick] (0,0) -- (1.20, 0) node[below, fill=white, inner sep=0.7pt, pos=0.68, font=\scriptsize, text=ochre] {$\sigma_1$};
      \draw[-{Stealth[length=1.6mm]}, ochre, thick] (0,0) -- (0, 0.50) node[left, fill=white, inner sep=0.7pt, pos=0.65, font=\scriptsize, text=ochre] {$\sigma_2$};
      \node[below, font=\scriptsize, color=slate, align=center] at (0, -2.10) {轴对齐椭圆，$X_1 \indep X_2$\\参数量：$d$};
    \end{scope}

    % -------- Panel 3: Full Covariance --------
    \begin{scope}[xshift=10.75cm]
      \node[font=\small\sffamily\bfseries, color=ochre] at (0, 2.76) {(c) 全协方差 $\Sigma = Q\Lambda Q^T$};
      \draw[-{Stealth[length=1.8mm]}, slate!70] (-2.15, 0) -- (2.25, 0) node[right, font=\scriptsize] {$x_1$};
      \draw[-{Stealth[length=1.8mm]}, slate!70] (0, -1.95) -- (0, 2.05) node[above, font=\scriptsize] {$x_2$};

      % Tilted ellipse at 25.1 deg (matching [[4, 1.8],[1.8, 1.0]] aspect ratio ~ 5.5 -> softened slightly for clarity)
      \begin{scope}[rotate=27]
        \fill[ochreLight, opacity=0.65] (0,0) ellipse (1.85cm and 0.62cm);
        \draw[ochre!55, thick] (0,0) ellipse (1.85cm and 0.62cm);
        \draw[ochre, thick] (0,0) ellipse (1.24cm and 0.42cm);
        \draw[ochre, very thick] (0,0) ellipse (0.62cm and 0.21cm);
        % Eigenvectors
        \draw[-{Stealth[length=2mm]}, tealDark, very thick] (0,0) -- (1.24, 0) node[right, xshift=1pt, yshift=3pt, font=\scriptsize\bfseries, color=tealDark] {$\sqrt{\lambda_1}\,q_1$};
        \draw[-{Stealth[length=2mm]}, navy, very thick] (0,0) -- (0, 0.42) node[above left, xshift=4pt, yshift=1pt, font=\scriptsize\bfseries, color=navy] {$\sqrt{\lambda_2}\,q_2$};
      \end{scope}
      \fill[ink] (0,0) circle (1.4pt);
      \node[below, font=\scriptsize, color=slate, align=center] at (0, -2.10) {倾斜椭圆，主轴沿特征向量\\参数量：$d(d+1)/2$};
    \end{scope}
  \end{tikzpicture}
  \caption{三种经典协方差结构对应的二维高斯等密度轮廓线（$d_M(x,\mu) = 0.5, 1.0, 1.5$）。全协方差情形中，谱分解 $\Sigma = Q\Lambda Q^T$ 直接指明了倾斜椭圆的正交主轴方向 $q_1, q_2$ 与半轴尺度 $\sqrt{\lambda_1}, \sqrt{\lambda_2}$。}
  \label{fig:ch2-mvg-contours}
\end{figure}

\begin{failBox}{奇异协方差矩阵陷阱与“拒绝显式求逆”原则}
在真实高维场景中，若两列特征线性相关（如同时放进“身高（厘米）”与“身高（米）”），或者样本量少于特征维度（$N < d$），经验协方差矩阵 $\hat{\Sigma}$ 的最小特征值必为 $0$（$\det\hat{\Sigma} = 0$，仅半正定而非正定）。此时 $\hat{\Sigma}^{-1}$ 不存在，高斯密度函数在全空间 $\R^d$ 上退化为一个集中在低维超平面上的奇异分布。

\textbf{两大工程铁律}：（1）当协方差矩阵可能病态或奇异时，添加微量对角正则项 $\Sigma_\epsilon = \Sigma + \epsilon I_d$（称为 \textbf{jitter / 岭平滑}，常见 $\epsilon = 10^{-6}$）确保严格正定；（2）在计算马氏距离 $(x-\mu)^T\Sigma^{-1}(x-\mu)$ 与对数行列式 $\log\det\Sigma$ 时，\textbf{永远不要显式调用 \texttt{np.linalg.inv(Sigma)}}！正确做法是先做 Cholesky 分解 $\Sigma = L L^T$，解下三角方程组 $L v = x - \mu$，立得：
\[
  d_M^2(x, \mu) = \|v\|_2^2, \qquad
  \log\det\Sigma = 2\sum_{i=1}^d \log L_{ii}.
\]
此法既快一倍、数值精度高出数个量级，还能彻底避免 $\det\Sigma$ 在高维下溢为 $0$！
\end{failBox}

% =========================================================================
\section{仿射封闭性、重参数化技巧与条件/边际高斯分布}
\label{sec:ch2-affine-cond}

多元高斯之所以成为概率图模型、状态估计与深度生成模型的“万金油”，不仅因为它的密度公式优美，更因为它在\textbf{仿射变换、边际化、条件化}这三大基本运算下\textbf{全部保持封闭——结果依然是精确的多元高斯}！

\subsection{仿射变换封闭性、Cholesky 采样与 VAE 重参数化}

\begin{thmBox}
\begin{theorem}[多元高斯的仿射变换封闭性]
\label{thm:ch2-affine}
设 $X \sim \mathcal{N}(\mu, \Sigma)$（$X \in \R^d$），$A \in \R^{m\times d}$ 为任意确定性矩阵，$b \in \R^m$ 为任意确定性向量，则仿射变换后的随机向量 $Y := A X + b \in \R^m$ 依然服从多元高斯分布：
\begin{equation}
  Y = A X + b \;\sim\; \mathcal{N}\bigl(A\mu + b,\; A\Sigma A^T\bigr). \label{eq:ch2-affine}
\end{equation}
（当 $A\Sigma A^T$ 仅半正定时，该分布理解为支撑在其列空间上的退化多元高斯分布。）
\end{theorem}
\end{thmBox}

\begin{proof}[核心验证（均值、协方差与可逆情形密度）]
由命题 \ref{prop:ch1-cov-psd}，无需任何额外假设已严格证得 $\E[Y] = A\mu + b$ 与 $\Cov(Y) = A\Sigma A^T$。下面对 $m=d$ 且 $A$ 可逆的情形直接验证密度形式：由反变换 $x = A^{-1}(y - b)$，雅可比行列式为 $|\det(A^{-1})| = |\det A|^{-1}$，代入式 \eqref{eq:ch2-mvg-pdf}：
\begin{align*}
  p_Y(y) &= p_X\bigl(A^{-1}(y - b)\bigr)\cdot |\det A|^{-1} \\
  &= \frac{1}{(2\pi)^{d/2}\bigl(\det(A\Sigma A^T)\bigr)^{1/2}}
    \exp\!\left(-\frac{1}{2}(y - A\mu - b)^T (A\Sigma A^T)^{-1}(y - A\mu - b)\right),
\end{align*}
其中用到了 $(\det\Sigma)^{1/2}|\det A| = (\det(A\Sigma A^T))^{1/2}$ 与 $(A^{-1})^T\Sigma^{-1}A^{-1} = (A\Sigma A^T)^{-1}$。一般非方阵或降秩矩阵 $A$ 的完整证明借助特征函数 $\varphi_X(t) = \E[e^{it^T X}] = \exp(it^T\mu - \frac{1}{2}t^T\Sigma t)$，代入 $\varphi_Y(s) = e^{is^T b}\varphi_X(A^T s)$ 一行即得。
\end{proof}

定理 \ref{thm:ch2-affine} 直接送给我们人工智能中最常用的两项实战绝技：

\textbf{绝技一：Cholesky 采样（如何从任意 $\mathcal{N}(\mu, \Sigma)$ 生成样本）}。计算机最擅长生成独立标准正态随机数 $Z \sim \mathcal{N}(0, I_d)$。给定任意目标正定协方差阵 $\Sigma \succ 0$，对其做 Cholesky 分解 $\Sigma = L L^T$（$L$ 为对角元全正的下三角矩阵），然后令
\begin{equation}
  X = \mu + L Z. \label{eq:ch2-cholesky-sample}
\end{equation}
由定理 \ref{thm:ch2-affine}，立即有 $\E[X] = \mu + L\cdot 0 = \mu$，且 $\Cov(X) = L I_d L^T = L L^T = \Sigma$！这正是 \texttt{numpy} 底层生成相关多元高斯样本的标准算法。

\begin{insightBox}{绝技二：第 5 卷 VAE 的重参数化技巧（Reparameterization Trick）}
在变分自编码器（VAE，第 5 卷）中，编码器网络输出隐变量分布的参数 $\mu_\phi(x) \in \R^d$ 与标准差向量 $\sigma_\phi(x) \in \R_+^d$，随后需要采样 $z \sim \mathcal{N}\bigl(\mu_\phi(x),\; \diag(\sigma_\phi^2(x))\bigr)$ 送入解码器算损失。如果直接把“采样”塞进计算图，反向传播走到随机采样节点就会断链——你无法对随机数发生器求导！

定理 \ref{thm:ch2-affine} 提供了天才般的绕道方案：把随机性抽离成与网络参数 $\phi$ 无关的标准噪声 $\varepsilon \sim \mathcal{N}(0, I_d)$，再通过确定性的仿射变换构造隐变量：
\[
  z = \mu_\phi(x) + \sigma_\phi(x) \had \varepsilon, \qquad \varepsilon \sim \mathcal{N}(0, I_d).
\]
由于每个分量满足 $\Var(\sigma_{\phi,i}\varepsilon_i) = \sigma_{\phi,i}^2$，$z$ 拥有与直接采样\textbf{完全相同的概率分布}；但现在 $z$ 成了参数 $(\mu_\phi, \sigma_\phi)$ 的光滑可微函数（$\frac{\partial z}{\partial \mu_\phi} = I$，$\frac{\partial z}{\partial \sigma_\phi} = \diag(\varepsilon)$），梯度便畅通无阻地穿过采样层回传至编码器！
\end{insightBox}

\subsection{多元高斯的边际分布与条件分布}

在预测与贝叶斯推断中，我们最常面对这样的问题：已知一个联合高斯向量 $X$ 的一部分分量 $X_2 = x_2$（已观测特征或历史状态），如何精确推断未观测部分 $X_1$（待预测目标或隐藏状态）的概率分布？把 $X \in \R^d$ 按维度切分为两块 $X = \begin{pmatrix} X_1 \\ X_2 \end{pmatrix}$（$X_1 \in \R^{d_1},\; X_2 \in \R^{d_2},\; d_1+d_2=d$），并将均值与协方差矩阵同步分块：
\[
  \begin{pmatrix} X_1 \\ X_2 \end{pmatrix}
  \;\sim\;
  \mathcal{N}\!\left(
    \begin{pmatrix} \mu_1 \\ \mu_2 \end{pmatrix},\;
    \begin{pmatrix} \Sigma_{11} & \Sigma_{12} \\ \Sigma_{21} & \Sigma_{22} \end{pmatrix}
  \right), \qquad \Sigma_{21} = \Sigma_{12}^T,\; \Sigma \succ 0.
\]

\begin{thmBox}
\begin{theorem}[多元高斯的边际分布与条件分布]
\label{thm:ch2-cond-mvg}
在上述分块多元高斯模型下：
\begin{enumerate}
  \item \textbf{边际分布（Marginal Distribution）}：直接截取对应分块即可，
  \begin{equation}
    X_1 \sim \mathcal{N}(\mu_1,\; \Sigma_{11}), \qquad X_2 \sim \mathcal{N}(\mu_2,\; \Sigma_{22}). \label{eq:ch2-marginal}
  \end{equation}
  \item \textbf{条件分布（Conditional Distribution）}：给定观测 $X_2 = x_2$ 后，$X_1$ 仍服从精确的多元高斯分布：
  \begin{equation}
    X_1 \mid (X_2 = x_2) \;\sim\; \mathcal{N}\bigl(\mu_{1\mid 2},\; \Sigma_{1\mid 2}\bigr), \label{eq:ch2-cond}
  \end{equation}
  其中条件均值与条件协方差矩阵分别为：
  \begin{align}
    \mu_{1\mid 2} &= \mu_1 + \Sigma_{12}\Sigma_{22}^{-1}(x_2 - \mu_2), \label{eq:ch2-cond-mean} \\
    \Sigma_{1\mid 2} &= \Sigma_{11} - \Sigma_{12}\Sigma_{22}^{-1}\Sigma_{21}. \label{eq:ch2-cond-cov}
  \end{align}
\end{enumerate}
\end{theorem}
\end{thmBox}

边际分布 \eqref{eq:ch2-marginal} 的证明仅需一行：令 $A = \begin{pmatrix} I_{d_1} & 0 \end{pmatrix}$，于是 $X_1 = A X$，直接套用定理 \ref{thm:ch2-affine} 即得 $A\mu = \mu_1$ 与 $A\Sigma A^T = \Sigma_{11}$！

在深入证明条件公式 \eqref{eq:ch2-cond-mean}--\eqref{eq:ch2-cond-cov} 之前，请先静心品味这两条公式揭示的三条深刻直觉（它们将反复出现在卡尔曼滤波、高斯过程与线性回归中）：
\begin{itemize}
  \item \textbf{条件均值是观测的线性修正}：未看到 $x_2$ 前，对 $X_1$ 的最佳猜测是先验均值 $\mu_1$；看到 $x_2$ 偏离预期 $\mu_2$ 的“新息（innovation）” $x_2 - \mu_2$ 后，乘以\textbf{增益矩阵} $K := \Sigma_{12}\Sigma_{22}^{-1}$ 对猜测做线性修正。若 $X_1$ 与 $X_2$ 不相关（$\Sigma_{12} = 0$），修正项归零，说明在联合高斯下“不相关 $\implies$ 独立”！
  \item \textbf{条件方差与具体观测值 $x_2$ 无关}：式 \eqref{eq:ch2-cond-cov} 右侧根本不含 $x_2$！无论 $X_2$ 测出的是 $3$ 还是 $30$，只要知道它落在某个精确值上，残留的不确定度 $\Sigma_{1\mid 2}$ 都是同一个常数矩阵。
  \item \textbf{观测永不增加不确定性}：由于 $\Sigma_{22} \succ 0$，减去的项 $\Sigma_{12}\Sigma_{22}^{-1}\Sigma_{21} = (\Sigma_{12}\Sigma_{22}^{-1/2})(\Sigma_{12}\Sigma_{22}^{-1/2})^T \succeq 0$，故必有 $\Sigma_{1\mid 2} \preceq \Sigma_{11}$！只要 $X_1$ 与 $X_2$ 存在相关性（$\Sigma_{12} \neq 0$），观测 $X_2$ 就严格压缩 $X_1$ 的方差。
\end{itemize}

\begin{example}[二维相关高斯的精确条件切片与 50,000 样本实测校验]
\label{ex:ch2-cond-2d}
考虑二维联合高斯分布 $X = (X_1, X_2)^T \sim \mathcal{N}(\mu, \Sigma)$，其中：
\[
  \mu = \begin{pmatrix} 1.0 \\ 2.0 \end{pmatrix}, \qquad
  \Sigma = \begin{pmatrix} 4.0 & 1.8 \\ 1.8 & 1.0 \end{pmatrix}.
\]
先检查正定性：$\Sigma_{11} = 4.0 > 0$ 且 $\det\Sigma = 4.0\times 1.0 - 1.8^2 = 0.76 > 0$，合法正定！两维相关系数高达 $\rho = \frac{1.8}{\sqrt{4.0}\cdot\sqrt{1.0}} = 0.90$。在没有任何观测前，$X_1$ 的边际分布是 $\mathcal{N}(1.0, 4.0)$（标准差 $\sigma_1 = 2.0$，相当宽泛）。

现在假设我们精确观测到 $X_2 = 3.0$（比均值 $\mu_2 = 2.0$ 高了 $1$ 个单位）。代入式 \eqref{eq:ch2-cond-mean}--\eqref{eq:ch2-cond-cov}：
\begin{align*}
  \mu_{1\mid 2} &= 1.0 + \frac{1.8}{1.0}\,(3.0 - 2.0) = \mathbf{2.80}, \\
  \Var(X_1 \mid X_2 = 3.0) &= 4.0 - \frac{1.8^2}{1.0} = 4.0 - 3.24 = \mathbf{0.76}\quad (\sigma_{1\mid 2} \approx 0.872).
\end{align*}
仅仅观测了 $X_2 = 3.0$，对 $X_1$ 的预期中心就从 $1.0$ 精准拉升到 $2.80$，而方差则从 $4.00$ 暴跌至 $0.76$——\textbf{不确定性缩减了整整 $81\%$}（恰好等于 $\rho^2 = 0.90^2 = 0.81$）！

用配套脚本 \texttt{code/vol2\_experiments.py} 通过 Cholesky 因子 $L = \begin{psmallmatrix} 2.0 & 0 \\ 0.9 & \sqrt{0.19} \end{psmallmatrix}$ 独立采样 $50{,}000$ 个二维样本，筛选落在薄条带 $|X_2 - 3.0| < 0.05$ 内的 $n = 1{,}214$ 个样本统计 $X_1$ 的经验均值与方差，实测结果为：
\[
  \hat{\mu}_{1\mid 2} = \mathbf{2.775}\;(\text{理论值 } 2.800), \qquad
  \widehat{\Var}_{1\mid 2} = \mathbf{0.778}\;(\text{理论值 } 0.760).
\]
理论公式与大样本几何切片完美吻合（图 \ref{fig:ch2-cond-slice}）。
\end{example}

\begin{figure}[htbp]
  \centering
  \begin{tikzpicture}[font=\footnotesize\sffamily]
    % -------- Left Panel: 2D Contour + Horizontal Slice X2 = 3 --------
    \begin{scope}[xshift=0cm]
      \node[font=\small\sffamily\bfseries, color=tealDark] at (1.5, 3.95) {(a) 二维联合高斯等高线与切片 $X_2 = 3$};

      % Coordinate mapping: draw in scaled cm: x_draw = 0.72 * x1, y_draw = 0.82 * x2
      % Axes: x1 in [-3.5, 5.8] -> [-2.52, 4.18], x2 in [-0.2, 4.2] -> [-0.16, 3.44]
      \draw[-{Stealth[length=2mm]}, slate] (-2.2, 0) -- (4.35, 0) node[right, font=\scriptsize] {$X_1$};
      \draw[-{Stealth[length=2mm]}, slate] (0, -0.3) -- (0, 3.35) node[above, font=\scriptsize] {$X_2$};

      % Ticks
      \foreach \x in {-2, 1, 3, 5} {
        \draw[slate] (0.72*\x, 0.06) -- (0.72*\x, -0.06) node[below, font=\scriptsize] {$\x$};
      }
      \foreach \y in {1, 2, 3} {
        \draw[slate] (0.06, 0.82*\y) -- (-0.06, 0.82*\y) node[left, font=\scriptsize] {$\y$};
      }

      % Mean (1, 2) -> (0.72, 1.64).
      % True Sigma = [[4, 1.8],[1.8, 1]]: eigvals = 4.843, 0.157 -> sqrt = 2.201, 0.396, angle 25.1 deg.
      % Under (0.72, 0.82) scale, plot parametric ellipses directly:
      % x(t) = mu + c * L * (cos t, sin t)^T with L = [[2, 0],[0.9, 0.4359]]
      \draw[teal!45, thick, fill=tealLight, fill opacity=0.55, domain=0:360, samples=90, variable=\t]
        plot ({0.72*(1 + 2.0*(2.0*cos(\t)))}, {0.82*(2 + 2.0*(0.9*cos(\t) + 0.4359*sin(\t)))});
      \draw[teal, thick, domain=0:360, samples=90, variable=\t]
        plot ({0.72*(1 + 1.3*(2.0*cos(\t)))}, {0.82*(2 + 1.3*(0.9*cos(\t) + 0.4359*sin(\t)))});
      \draw[tealDark, very thick, domain=0:360, samples=90, variable=\t]
        plot ({0.72*(1 + 0.65*(2.0*cos(\t)))}, {0.82*(2 + 0.65*(0.9*cos(\t) + 0.4359*sin(\t)))});

      % Mean dot
      \fill[tealDark] (0.72, 1.64) circle (2pt);
      \node[below right, font=\scriptsize, color=tealDark] at (0.37, 0.65) {$\mu=(1,2)^T$};

      % Horizontal slice line X2 = 3.0 -> y_draw = 2.46
      \draw[ochre, very thick, dashed] (-2.0, 2.46) -- (4.2, 2.46);
      \node[above left, font=\scriptsize\bfseries, color=ochre] at (4.25, 3.36) {观测切片 $X_2 = 3.0$};

      % Conditional mean point (2.8, 3.0) -> (2.016, 2.46)
      \fill[ochre] (2.016, 2.46) circle (2.3pt);
      \draw[ochre, dotted, thick] (2.016, 2.46) -- (2.016, 0);
      \node[below, font=\scriptsize\bfseries, color=ochre] at (2.016, -0.32) {$\mu_{1\mid 2}=2.8$};
    \end{scope}

    % -------- Right Panel: Marginal vs Conditional 1D Gaussian PDFs --------
    \begin{scope}[xshift=6.8cm, yshift=-0.1cm]
      \begin{axis}[
        width=7.1cm,
        height=4.85cm,
        title={\small\sffamily\bfseries\color{navy}(b) 边际分布 vs 条件切片分布},
        title style={yshift=-2pt},
        xlabel={$x_1$},
        ylabel={密度 $p(x_1)$},
        xmin=-4, xmax=6.5,
        ymin=0, ymax=0.73,
        xtick={-3, -1, 1, 2.8, 5},
        xticklabels={$-3$, $-1$, $1$, $2.8$, $5$},
        grid=major,
        grid style={dashed, gray!25},
        legend pos=north west,
        legend style={font=\scriptsize, fill=white, draw=codeborder, rounded corners=1pt},
        tick label style={font=\scriptsize},
        label style={font=\scriptsize\sffamily}
      ]
        % Marginal N(1, 4): sigma=2
        \addplot[color=teal, thick, dashed, domain=-4:6.5, samples=100]
          {exp(-0.5*((x-1.0)/2.0)^2) / (sqrt(2*3.14159)*2.0)};
        \addlegendentry{边际 $X_1\sim\mathcal{N}(1.0,\,4.00)$}

        % Conditional N(2.8, 0.76): sigma=0.8718
        \addplot[color=ochre, very thick, domain=-4:6.5, samples=120]
          {exp(-0.5*((x-2.8)/0.8718)^2) / (sqrt(2*3.14159)*0.8718)};
        \addlegendentry{条件 $X_1\mid(X_2=3)\sim\mathcal{N}(2.8,\,0.76)$}

        % Vertical guides
        \draw[teal, dotted] (axis cs:1.0, 0) -- (axis cs:1.0, 0.20);
        \draw[ochre, dotted] (axis cs:2.8, 0) -- (axis cs:2.8, 0.46);
      \end{axis}
    \end{scope}
  \end{tikzpicture}
  \caption{例 \ref{ex:ch2-cond-2d} 的几何可视化。（a）沿 $X_2=3.0$ 对倾斜二维高斯椭球作水平切片，切面中心精准右移至 $\mu_{1\mid 2}=2.8$；（b）对比边际密度（虚线，方差 $4.00$）与条件切片密度（实线，方差 $0.76$），直观展示观测相关变量如何同时修正均值并大幅压缩不确定性。}
  \label{fig:ch2-cond-slice}
\end{figure}

\begin{extBox}{条件高斯的优雅证明：分块矩阵去相关与 Schur 补}
很多教材通过繁琐的分块矩阵求逆硬拆二次型来证明定理 \ref{thm:ch2-cond-mvg}，但借助\textbf{线性去相关 + 仿射封闭性}，只需五步就能优雅证完，并自然引出线性代数中的\textbf{Schur 补（Schur Complement）}！

\textbf{第 1 步：构造与 $X_2$ 不相关的残差向量 $U$}。设 $K := \Sigma_{12}\Sigma_{22}^{-1} \in \R^{d_1\times d_2}$，构造块下三角可逆线性变换：
\[
  \begin{pmatrix} U \\ X_2 \end{pmatrix}
  := \begin{pmatrix} I_{d_1} & -K \\ 0 & I_{d_2} \end{pmatrix} \begin{pmatrix} X_1 \\ X_2 \end{pmatrix}
  = \begin{pmatrix} X_1 - K X_2 \\ X_2 \end{pmatrix}.
\]
\textbf{第 2 步：由定理 \ref{thm:ch2-affine} 算联合协方差矩阵}。令变换阵为 $M = \begin{psmallmatrix} I & -K \\ 0 & I \end{psmallmatrix}$，直接做分块矩阵乘法：
\begin{align*}
  M \Sigma M^T
  &= \begin{pmatrix} I & -K \\ 0 & I \end{pmatrix}
     \begin{pmatrix} \Sigma_{11} & \Sigma_{12} \\ \Sigma_{21} & \Sigma_{22} \end{pmatrix}
     \begin{pmatrix} I & 0 \\ -K^T & I \end{pmatrix} \\
  &= \begin{pmatrix} \Sigma_{11} - K\Sigma_{21} & \Sigma_{12} - K\Sigma_{22} \\ \Sigma_{21} & \Sigma_{22} \end{pmatrix}
     \begin{pmatrix} I & 0 \\ -K^T & I \end{pmatrix}
   = \begin{pmatrix} \Sigma_{11} - \Sigma_{12}\Sigma_{22}^{-1}\Sigma_{21} & 0 \\ 0 & \Sigma_{22} \end{pmatrix}.
\end{align*}
右上角 $\Sigma_{12} - K\Sigma_{22} = \Sigma_{12} - \Sigma_{12}\Sigma_{22}^{-1}\Sigma_{22} = 0$ 彻底消去！左上角留下的矩阵 $\Sigma / \Sigma_{22} := \Sigma_{11} - \Sigma_{12}\Sigma_{22}^{-1}\Sigma_{21}$ 正是分块矩阵 $\Sigma$ 关于 $\Sigma_{22}$ 的 \textbf{Schur 补}。

\textbf{第 3 步：联合高斯 + 零协方差 $\implies$ 严格独立}。由于 $(U, X_2)$ 是 $X$ 的仿射变换，它们依然服从联合多元高斯分布，且协方差阵为块对角阵，于是其联合密度完美分解为 $p(u, x_2) = p(u)\,p(x_2)$，即 \textbf{$U \indep X_2$}！同时，$U \sim \mathcal{N}\bigl(\mu_1 - K\mu_2,\; \Sigma / \Sigma_{22}\bigr)$。

\textbf{第 4 步：还原 $X_1$ 并固定 $X_2 = x_2$}。反写 $X_1 = U + K X_2$。在给定 $X_2 = x_2$ 的条件下，$K x_2$ 变成确定性常向量，而 $U$ 与 $X_2$ 独立故分布丝毫不变！因此
\begin{align*}
  X_1 \mid (X_2 = x_2) = U + K x_2
  &\;\sim\; \mathcal{N}\bigl(\mu_1 - K\mu_2 + K x_2,\; \Sigma / \Sigma_{22}\bigr) \\
  &\;=\; \mathcal{N}\bigl(\mu_1 + \Sigma_{12}\Sigma_{22}^{-1}(x_2 - \mu_2),\; \Sigma_{11} - \Sigma_{12}\Sigma_{22}^{-1}\Sigma_{21}\bigr).
\end{align*}
顺带一提，对 $M\Sigma M^T = \diag(\Sigma / \Sigma_{22}, \Sigma_{22})$ 两边取行列式，因 $\det M = 1$，立刻附赠一条重要恒等式：$\det\Sigma = (\det\Sigma_{22})\det(\Sigma_{11} - \Sigma_{12}\Sigma_{22}^{-1}\Sigma_{21})$，证明了只要 $\Sigma \succ 0$，条件协方差 $\Sigma_{1\mid 2}$ 必然严格正定！
\end{extBox}

\subsection{纯 NumPy 实战：Cholesky 采样、无求逆马氏距离与条件切片检验}

代码 \ref{lst:ch2-mvg-code} 完整实现了本章核心算法：用 Cholesky 生成例 \ref{ex:ch2-cond-2d} 的 $50{,}000$ 个二维高斯样本、用 \texttt{np.linalg.solve}（绝不调用 \texttt{inv}）批量计算马氏距离并检验 $\E[d_M^2]=d=2$，以及验证条件切片 $X_1\mid(X_2\approx 3.0)$ 的均值与方差。

\begin{lstlisting}[caption={纯 NumPy 多元高斯 Cholesky 采样、无求逆马氏距离与条件切片实测}, label={lst:ch2-mvg-code}]
import numpy as np

rng = np.random.default_rng(42)
mu = np.array([1.0, 2.0])
Sigma = np.array([[4.0, 1.8],
                  [1.8, 1.0]])

# 1. Cholesky 采样（对齐 vol2_experiments.py）
L = np.linalg.cholesky(Sigma)
Z = rng.normal(size=(50000, 2))
X = mu + Z @ L.T
print("Empirical mean:", np.round(X.mean(axis=0), 3))
print("Empirical cov:\n", np.round(np.cov(X, rowvar=False), 3))

# 2. 无显式求逆的批量马氏距离：解 L @ V = (X - mu)^T
V = np.linalg.solve(L, (X - mu).T)          # shape (2, 50000)
d_maha_sq = np.sum(V ** 2, axis=0)          # 平方马氏距离
log_det = 2.0 * np.sum(np.log(np.diag(L)))  # 稳定的 log det(Sigma)
print(f"Mean d_M^2 = {d_maha_sq.mean():.3f} (exact d=2.000), log|Sigma| = {log_det:.4f}")

# 3. 条件分布切片 X1 | (X2 ~ 3.0) 验证
# 理论精确值: mu_{1|2} = 1 + (1.8/1.0)*(3-2) = 2.800, var_{1|2} = 4 - 1.8^2/1 = 0.760
mask = np.abs(X[:, 1] - 3.0) < 0.05
print(f"Slice n={mask.sum()}: emp_mean={X[mask, 0].mean():.3f} (exact 2.800), "
      f"emp_var={X[mask, 0].var():.3f} (exact 0.760)")
\end{lstlisting}

% =========================================================================
\section{本章习题}
\label{sec:ch2-exercises}

\exTier{A. 概念辨析与反例查错}{合上书独立作答，精准锁定多元高斯与核心分布的边界条件}

\begin{enumerate}[label=\textbf{2.\arabic*}, leftmargin=2.4em]
  \item \textbf{（核心分布与多元高斯概念辨析）} 判断以下四个命题正误；对错误命题给出反例或反驳理由，对正确命题简述依据：
  \begin{enumerate}[label=(\alph*), leftmargin=1.8em]
    \item 若随机变量 $X_1 \sim \mathcal{N}(0, 1)$ 且 $X_2 \sim \mathcal{N}(0, 1)$，则二维随机向量 $(X_1, X_2)^T$ 必服从二维联合高斯分布；
    \item 对任意速率参数 $\lambda > 0$，Poisson 分布 $\mathrm{Pois}(\lambda)$ 的均值与方差恒相等；
    \item 若随机向量 $(X_1, X_2)^T$ 服从二维联合高斯分布且 $\Cov(X_1, X_2) = 0$，则 $X_1$ 与 $X_2$ 必相互独立；
    \item 在分块多元高斯条件分布 $X_1 \mid (X_2 = x_2)$ 中，条件协方差矩阵 $\Sigma_{1\mid 2}$ 随观测值 $x_2$ 偏离 $\mu_2$ 越远而增大。
  \end{enumerate}

  \item \textbf{（非法协方差阵与欧氏 vs 马氏距离反直觉查错）}
  \begin{enumerate}[label=(\arabic*), leftmargin=1.8em]
    \item 某同学试图用 $\Sigma = \begin{pmatrix} 1.0 & 1.2 \\ 1.2 & 1.0 \end{pmatrix}$ 作为二维高斯分布的协方差矩阵，请用两种方法（特征值 / 相关系数）迅速指出其错误所在；
    \item 给定对角高斯 $\mathcal{N}\bigl(0, \diag(16, 1)\bigr)$，比较两点 $a = (4, 0)^T$ 与 $b = (0, 2)^T$ 到原点的欧氏距离、马氏距离及概率密度大小，解释为何“欧氏距离更近的点 $b$ 反而落在了概率密度更低的尾部区域”。
  \end{enumerate}
\end{enumerate}

\exTier{B. 小例手算与数值直觉}{亲手打通离散分布、Beta 先验与二维高斯条件推断}

\begin{enumerate}[label=\textbf{2.\arabic*}, leftmargin=2.4em, resume]
  \item \textbf{（二项评测概率与 Beta 先验手算）}
  \begin{enumerate}[label=(\arabic*), leftmargin=1.8em]
    \item 设某二分类器在单条独立测试样本上的准确率为 $p = 0.8$，现抽取 $n = 5$ 条独立测试样本，求答对题数 $S_5$ 的期望、方差，以及“至少答对 $4$ 题”的精确概率 $\Prob(S_5 \ge 4)$；
    \item 考虑先验分布 $X \sim \mathrm{Beta}(3, 2)$：手算归一化常数 $B(3, 2)$、密度函数显式表达式、均值 $\E[X]$、方差 $\Var(X)$ 与众数（密度极大值点 $x^*$），并解释该分布对应多少次“先验伪计数”。
  \end{enumerate}

  \item \textbf{（2D 多元高斯椭球主轴与条件分布精确手算）} 给定二维多元高斯分布
  \[
    X = \begin{pmatrix} X_1 \\ X_2 \end{pmatrix}
    \;\sim\;
    \mathcal{N}\!\left(\begin{pmatrix} 0 \\ 1 \end{pmatrix},\; \begin{pmatrix} 5 & 2 \\ 2 & 2 \end{pmatrix}\right).
  \]
  \begin{enumerate}[label=(\arabic*), leftmargin=1.8em]
    \item 手算协方差矩阵 $\Sigma$ 的全部特征值与单位正交特征向量，指出等密度椭圆 $d_M^2(x, \mu) = 1$ 的长短半轴长度；
    \item 求点 $x = (2, 3)^T$ 到均值 $\mu$ 的平方马氏距离 $d_M^2(x, \mu)$；
    \item 若观测到 $X_2 = 3$，求条件分布 $X_1 \mid (X_2 = 3)$ 的精确均值与方差，并计算观测 $X_2$ 带来的方差缩减比例。
  \end{enumerate}
\end{enumerate}

\exTier{C. 实用推导与性质证明}{高维高斯理论与深度学习高频用到的核心结论}

\begin{enumerate}[label=\textbf{2.\arabic*}, leftmargin=2.4em, resume]
  \item \textbf{（独立高斯之和与样本均值的 $1/\sqrt{n}$ 标准误）}
  \begin{enumerate}[label=(\arabic*), leftmargin=1.8em]
    \item 设 $X_1 \sim \mathcal{N}(\mu_1, \Sigma_1)$ 与 $X_2 \sim \mathcal{N}(\mu_2, \Sigma_2)$ 是 $\R^d$ 中相互独立的两个随机向量。先写出拼接向量 $\begin{pmatrix} X_1 \\ X_2 \end{pmatrix} \in \R^{2d}$ 的联合高斯分布，再仅用定理 \ref{thm:ch2-affine}（选取恰当矩阵 $A$）一步证明：$X_1 + X_2 \sim \mathcal{N}(\mu_1 + \mu_2,\; \Sigma_1 + \Sigma_2)$；
    \item 推广至 $n$ 个独立同分布一维高斯样本 $X_1, \dots, X_n \stackrel{\text{i.i.d.}}{\sim} \mathcal{N}(\mu, \sigma^2)$，求出样本均值 $\bar{X}_n = \frac{1}{n}\sum_{i=1}^n X_i$ 的精确分布，并指出为何样本量扩大 $4$ 倍才能把估计标准误减半。
  \end{enumerate}
  \exHint{第 (1) 问因 $X_1 \indep X_2$，拼接向量的联合协方差矩阵是块对角阵 $\diag(\Sigma_1, \Sigma_2)$；取 $A = \begin{pmatrix} I_d & I_d \end{pmatrix}$ 即可。}

  \item \textbf{（白化等价性与高维高斯薄球壳现象）} 设 $X \sim \mathcal{N}(\mu, \Sigma)$（$X \in \R^d,\; \Sigma \succ 0$），定义白化随机向量 $Z := \Sigma^{-1/2}(X - \mu)$。
  \begin{enumerate}[label=(\arabic*), leftmargin=1.8em]
    \item 证明 $Z \sim \mathcal{N}(0, I_d)$，且平方马氏距离满足 $d_M^2(X, \mu) = \|Z\|_2^2 = \sum_{i=1}^d Z_i^2$；
    \item 利用标准正态四阶矩 $\E[Z_i^4] = 3$，证明 $\E[d_M^2(X, \mu)] = d$ 且 $\Var\bigl(d_M^2(X, \mu)\bigr) = 2d$；
    \item 结合第 (2) 问解释高维空间的反直觉现象：虽然高斯密度在均值 $\mu$ 处最高，但当维度 $d = 1{,}000$ 时，绝大多数随机样本并不在 $\mu$ 附近，而是聚集在马氏半径约 $\sqrt{d} \approx 31.6$ 的薄球壳上！
  \end{enumerate}
  \exHint{对独立的 $Z_i \sim \mathcal{N}(0,1)$，$\Var(Z_i^2) = \E[Z_i^4] - (\E[Z_i^2])^2 = 3 - 1^2 = 2$，再对 $i=1,\dots,d$ 用方差可加性。}
\end{enumerate}

\exTier{D. 编程实战与可自检实验}{用纯 NumPy 验证高维高斯采样、条件推断与异常检测}

\begin{enumerate}[label=\textbf{2.\arabic*}, leftmargin=2.4em, resume]
  \item \textbf{（3D 高斯 Cholesky 采样、白化还原与多元条件切片验证）} 设 3 维多元高斯参数为：
  \[
    \mu = (1.0, -1.0, 2.0)^T, \qquad
    \Sigma = \begin{pmatrix} 4.0 & 1.2 & -0.8 \\ 1.2 & 2.0 & 0.6 \\ -0.8 & 0.6 & 1.5 \end{pmatrix}.
  \]
  使用 \texttt{np.random.default\_rng(2026)} 通过 Cholesky 分解采样 $N = 200{,}000$ 个三维样本，完成以下三项验证：
  \begin{enumerate}[label=(\arabic*), leftmargin=1.8em]
    \item 验证经验均值、经验协方差与理论 $\mu, \Sigma$ 的最大绝对误差；
    \item 通过 \texttt{np.linalg.solve(L, (X - mu).T)} 将样本白化还原为 $W$，计算 $W$ 的经验协方差与平方马氏距离样本均值/方差（对比习题 2.6 的理论值 $d=3$ 与 $2d=6$）；
    \item 用定理 \ref{thm:ch2-cond-mvg} 精确算出 $X_1 \mid (X_2 = 0.0,\; X_3 = 2.5)$ 的理论条件均值与方差，并与经验薄切片 $\{|X_2 - 0.0|<0.15 \text{ 且 } |X_3 - 2.5|<0.15\}$ 的统计量对照。
  \end{enumerate}
  \par\smallskip\noindent{\footnotesize\sffamily\color{tealDark}\textbf{【自检标准】}\,经验协方差与白化协方差的逐元素最大绝对误差均须小于 $0.03$；平方马氏距离均值须落在 $[2.95, 3.05]$、方差落在 $[5.85, 6.15]$；理论条件均值应为 $1.4242$、理论条件方差为 $2.2606$，切片经验均值与方差的绝对偏差均须小于 $0.15$。}

  \item \textbf{（欧氏距离 vs 马氏距离的隐蔽异常检测实战）} 构造二维零均值相关高斯 $\mathcal{N}(0, \Sigma)$，其中 $\Sigma = \begin{psmallmatrix} 4.0 & 1.8 \\ 1.8 & 1.0 \end{psmallmatrix}$：
  \begin{enumerate}[label=(\arabic*), leftmargin=1.8em]
    \item 用 \texttt{np.random.default\_rng(2026)} 生成 $2{,}000$ 个正常样本；
    \item 求出 $\Sigma$ 最小特征值对应的单位短轴特征向量 $u_{\min}$，沿 $\pm 1.8\,u_{\min}$ 附近（加标准差 $0.05$ 的微小高斯抖动）注入 $20$ 个“欧氏距离很小、但违背变量相关性”的隐蔽异常点；
    \item \textbf{禁止调用 \texttt{np.linalg.inv}}，仅用 \texttt{np.linalg.cholesky} 与 \texttt{np.linalg.solve} 实现向量化马氏距离；分别按欧氏距离和马氏距离选出得分最高的 $25$ 个候选点，统计两种方法对 $20$ 个隐蔽异常点的召回数量（Recall@25）。
  \end{enumerate}
  \par\smallskip\noindent{\footnotesize\sffamily\color{tealDark}\textbf{【自检标准】}\,由于异常点欧氏半径仅约 $1.8$（小于正常长轴标准差 $\sqrt{\lambda_1}\approx 2.20$），欧氏距离的 Recall@25 应接近 $0/20$（$\le 2/20$）；而异常点马氏距离高达 $\approx 4.54$，马氏距离的 Recall@25 必须达到 $20/20$（$100\%$）。}
\end{enumerate}

\chapter[极限定理、浓度不等式与蒙特卡洛]{极限定理、浓度不等式与蒙特卡洛方法}
\label{ch:ch3}

\chapEpigraph{即便是最迟钝的人，凭某种自然本能也能明白：观察次数越多，偏离真实比例的危险就越小。}{雅各布·伯努利（Jacob Bernoulli），《猜度术》}

\begin{chapIntro}{本章定位：从理想期望通往有限样本与可计算 AI 的三座桥梁}
第 1--2 章教会了我们如何用随机变量、条件概率与多元高斯去描述“不确定性”。但在真实人工智能系统中，我们几乎永远碰不到“上帝视角”的真实期望 $\E[L(f(X), Y)]$——手头只有有限的 $n$ 条测试样本，或者一个维度高达成千上万、根本算不出解析积分的高维分布。

本章正是为了回答三个直击 AI 工程命门的基础问题：
\begin{enumerate}
  \item \textbf{凸性与变分下界的起点}：为什么在概率建模里，\textbf{Jensen 不等式}仅凭“先映射再平均 $\ge$ 先平均再映射”一句话，就能一手催生出第 4 章的 KL 散度非负性与第 4--5 卷的 EM/VAE 证据下界（ELBO）？
  \item \textbf{有限样本的可靠性}：只用 $n$ 条测试集算出的经验准确率 $\widehat{\mathrm{Acc}}_n$，凭什么能代表真实泛化准确率？到底需要多少测试样本，才能以 $95\%$ 的把握把误差锁死在 $\pm 3\%$ 以内？答案藏在\textbf{浓度不等式三阶梯（Markov $\to$ Chebyshev $\to$ Hoeffding）}里。
  \item \textbf{高维期望的可计算性}：当期望积分维度 $d = 1000$、传统网格积分彻底遭遇维数灾难时，\textbf{中心极限定理（CLT）、蒙特卡洛采样与重要性采样}如何用随机性以 $O(1/\sqrt{M})$ 的维数无关速度破局？
\end{enumerate}
\end{chapIntro}

% =========================================================================
\section{Jensen 不等式：整个概率机器学习的凸性引擎}
\label{sec:ch3-jensen}

在展开大数定律和概率边界之前，我们先装配一把贯穿整套 AI 数学底座的“瑞士军刀”——\textbf{Jensen 不等式}。表面看它只是微积分里“凸函数图形位于切线上方”的一句概率翻译，实际上它几乎统治了信息论、统计学习与变分推断里所有最关键的不等式推导。

\begin{recallBox}{一元凸函数的等价刻画（承接第 1 卷第 5 章）}
设函数 $g: I \to \R$ 定义在区间 $I \subseteq \R$ 上。若对任意 $x_1, x_2 \in I$ 与任意权重 $\lambda \in [0, 1]$，均有
\[
  g\bigl(\lambda x_1 + (1-\lambda) x_2\bigr) \le \lambda g(x_1) + (1-\lambda) g(x_2),
\]
则称 $g$ 为\textbf{凸函数}（几何上：任意两点间的\textbf{割线位于函数曲线上方}）；若不等号对 $x_1 \neq x_2, \lambda \in (0,1)$ 严格成立，则称 $g$ 为\textbf{严格凸函数}。若 $-g$ 为凸函数，则称 $g$ 为\textbf{凹函数}。

对于二阶可微函数，$g$ 为凸 $\iff g''(x) \ge 0$；且对任意内部点 $\mu \in I$，必存在过 $(\mu, g(\mu))$ 的\textbf{支撑切线}（可微时斜率 $m = g'(\mu)$，不可微时为次梯度 $m$），使得曲线处处位于切线上方：
\begin{equation}
  g(x) \ge g(\mu) + m(x - \mu), \qquad \forall x \in I. \label{eq:ch3-supporting-line}
\end{equation}
典型例子：$x^2,\; e^x,\; -\ln x,\; |x|,\; x\ln x$ 为凸函数；$\ln x,\; \sqrt{x}$ 为凹函数。
\end{recallBox}

\subsection{从两点掷币看懂“先映射再平均”与“先平均再映射”}

先看一个无需任何符号负担的小例子。设随机变量 $X$ 以相等概率取值 $1$ 或 $3$：$\Prob(X=1) = \Prob(X=3) = \tfrac{1}{2}$，显然均值 $\E[X] = \tfrac{1+3}{2} = 2$。现在分别考察一个凸函数和一个凹函数：
\begin{itemize}
  \item 取凸函数 $g(x) = x^2$：“先平方再求平均”得 $\E[X^2] = \tfrac{1^2 + 3^2}{2} = 5$；“先求平均再平方”得 $(\E[X])^2 = 2^2 = 4$。前者更大：$5 > 4$。
  \item 取凹函数 $h(x) = \ln x$：“先取对数再求平均”得 $\E[\ln X] = \tfrac{\ln 1 + \ln 3}{2} \approx 0.5493$；“先求平均再取对数”得 $\ln(\E[X]) = \ln 2 \approx 0.6931$。后者更大：$0.5493 < 0.6931$。
\end{itemize}

把两点随机变量 $(x_1, x_2)$ 的概率写成 $\lambda = \Prob(X = x_1)$、$1-\lambda = \Prob(X = x_2)$，你会立刻发现：凸函数的定义式 $\lambda g(x_1) + (1-\lambda) g(x_2) \ge g(\lambda x_1 + (1-\lambda) x_2)$ 根本就是\textbf{两点分布上的 Jensen 不等式} $\E[g(X)] \ge g(\E[X])$！图 \ref{fig:ch3-jensen} 把这一几何机制画得一清二楚：割线上的点高度是 $\E[g(X)]$，曲线上的点高度是 $g(\E[X])$，两者之间的垂直差值就是非负的 \textbf{Jensen 间隙（Jensen gap）}；而支撑切线则从下方托住整条曲线，提供了把两点推广到\textbf{任意离散或连续分布}的钥匙。

\begin{figure}[htbp]
  \centering
  \begin{tikzpicture}[scale=1.02]
    % Axes
    \draw[-{Stealth[length=2.3mm]}, slate, thick] (-0.3, 0) -- (6.8, 0) node[right, font=\small\sffamily, ink] {$x$};
    \draw[-{Stealth[length=2.3mm]}, slate, thick] (0, -0.3) -- (0, 4.7) node[above, font=\small\sffamily, ink] {$y$};

    % Convex curve g(x) = 0.24*(x - 1.4)^2 + 0.55
    \draw[teal, line width=1.6pt, domain=0.25:5.75, samples=80]
      plot (\x, {0.24*(\x - 1.4)*(\x - 1.4) + 0.55});
    \node[tealDark, font=\small\sffamily\bfseries, anchor=west] at (5.35, 4.35) {凸函数 $y = g(x)$};

    % Points x1 = 1.0, x2 = 5.2, weight p1 = 0.45, p2 = 0.55 => mu = 3.31
    % g(1.0) = 0.24*0.16 + 0.55 = 0.5884
    % g(5.2) = 0.24*14.44 + 0.55 = 4.0156
    % mu = 3.31 -> g(mu) = 0.24*(1.91)^2 + 0.55 = 1.4255
    % E[g(X)] = 0.45*0.5884 + 0.55*4.0156 = 2.4734
    % Tangent at mu=3.31: slope m = 0.48*(3.31 - 1.4) = 0.9168
    % y_tan(x) = 1.4255 + 0.9168*(x - 3.31)
    \coordinate (P1) at (1.0, 0.5884);
    \coordinate (P2) at (5.2, 4.0156);
    \coordinate (Pmu) at (3.31, 1.4255);
    \coordinate (Pchord) at (3.31, 2.4734);

    % Supporting tangent line at mu
    \draw[navy, line width=1.1pt, densely dashed, domain=1.55:5.65]
      plot (\x, {1.4255 + 0.9168*(\x - 3.31)});
    \node[navy, font=\footnotesize\sffamily, anchor=west] at (5.70, 3.38) {支撑切线 $g(\mu)+m(x-\mu)$};

    % Chord connecting P1 and P2
    \draw[ochre, line width=1.25pt] (P1) -- (P2);
    \node[ochre, font=\footnotesize\sffamily, rotate=38, anchor=south] at (2.22, 1.72) {割线（两点凸组合）};

    % Vertical projection lines
    \draw[slate, dotted, thick] (1.0, 0) -- (P1);
    \draw[slate, dotted, thick] (5.2, 0) -- (P2);
    \draw[slate, densely dotted, thick] (3.31, 0) -- (Pchord);

    % Horizontal guides to y-axis
    \draw[teal!70, densely dotted] (0, 1.4255) -- (Pmu);
    \draw[ochre!80, densely dotted] (0, 2.4734) -- (Pchord);

    % Jensen gap brace
    \draw[decorate, decoration={brace, amplitude=4.5pt, mirror}, ochre, thick]
      (3.43, 1.44) -- (3.43, 2.45)
      node[midway, right=6pt, xshift=21mm, yshift=-5mm, font=\footnotesize\sffamily, ochre, align=left]
      {Jensen 间隙\\$\E[g(X)] - g(\E[X]) \ge 0$};

    \draw[ochre, thin] (3.69,1.96) -- (4.32,1.42) -- (5.52,1.42);
    % Dots
    \fill[teal] (P1) circle (2.4pt);
    \fill[teal] (P2) circle (2.4pt);
    \fill[tealDark] (Pmu) circle (2.8pt);
    \fill[ochre] (Pchord) circle (2.8pt);

    % Axis ticks & labels
    \node[below=2pt, font=\footnotesize, ink] at (1.0, 0) {$x_1$};
    \node[below=2pt, font=\footnotesize\bfseries, tealDark] at (3.31, 0) {$\mu = \E[X]$};
    \node[below=2pt, font=\footnotesize, ink] at (5.2, 0) {$x_2$};

    \node[left=2pt, font=\footnotesize, tealDark] at (0, 1.4255) {$g(\E[X])$};
    \node[left=2pt, font=\footnotesize, ochre] at (0, 2.4734) {$\E[g(X)]$};
  \end{tikzpicture}
  \caption{Jensen 不等式的几何本质：凸函数 $g(x)$ 处处位于支撑切线上方且位于任意割线下方，因此“先映射再平均”的高度 $\E[g(X)]$ 恒不低于“先平均再映射”的高度 $g(\E[X])$。}
  \label{fig:ch3-jensen}
\end{figure}

\subsection{Jensen 不等式与支撑切线完整证明}

\begin{thmBox}
\begin{theorem}[Jensen 不等式]
\label{thm:ch3-jensen}
设 $g: I \to \R$ 是区间 $I$ 上的凸函数，随机变量 $X$ 取值于 $I$，且期望 $\E[X]$ 与 $\E[g(X)]$ 均有限。则必有
\begin{equation}
  \E[g(X)] \ge g(\E[X]). \label{eq:ch3-jensen}
\end{equation}
若 $g$ 为凹函数，则不等号反向：$\E[g(X)] \le g(\E[X])$。进一步地，当 $g$ 为\textbf{严格凸}（或严格凹）函数时，等号成立当且仅当 $X$ 几乎必然为常数，即 $\Prob(X = \E[X]) = 1$。
\end{theorem}
\end{thmBox}

\begin{proof}
记 $\mu := \E[X]$。若 $X$ 几乎必然取 $I$ 的某个端点常数，则 $\E[g(X)] = g(\mu)$ 平凡成立。否则 $\mu$ 必落在 $I$ 的内部，由凸函数的支撑切线性质 \eqref{eq:ch3-supporting-line}，存在常数 $m \in \R$（当 $g$ 在 $\mu$ 处可微时即取 $m = g'(\mu)$），使得对一切 $x \in I$ 有
\[
  g(x) \ge g(\mu) + m(x - \mu).
\]
将上式中的确定变量 $x$ 替换为随机变量 $X$，得到逐样本点成立的随机变量不等式：
\[
  g(X) \ge g(\mu) + m(X - \mu).
\]
对不等式两边同时取期望，由第 1 章期望的单调性与线性性立即得到
\[
  \E[g(X)] \ge \E\bigl[g(\mu) + m(X - \mu)\bigr] = g(\mu) + m\bigl(\E[X] - \mu\bigr) = g(\mu) = g(\E[X]).
\]
若 $g$ 为凹函数，对凸函数 $-g$ 应用 \eqref{eq:ch3-jensen} 得 $\E[-g(X)] \ge -g(\E[X])$，两边同乘 $-1$ 即得 $\E[g(X)] \le g(\E[X])$。

最后证明严格凸时的等号条件：令非负随机变量 $Z := g(X) - g(\mu) - m(X - \mu) \ge 0$。若 $\E[g(X)] = g(\mu)$，则 $\E[Z] = 0$。非负随机变量期望为零蕴含 $\Prob(Z = 0) = 1$。但 $g$ 严格凸意味着支撑切线仅在切点 $x = \mu$ 处与曲线相交，即对一切 $x \neq \mu$ 严格有 $g(x) > g(\mu) + m(x - \mu)$（$Z > 0$）。因此 $\Prob(Z = 0) = 1$ 当且仅当 $\Prob(X = \mu) = 1$。
\end{proof}

回味一下这个证明：一共只有四行核心代数，却对离散分布、连续分布、混合分布乃至任意维随机向量（只需将 $m(x-\mu)$ 换成内积 $\nabla g(\mu)^T(x-\mu)$）\textbf{无条件通吃}！这就是寻找正确几何视角（支撑切线）的威力。

\begin{corollary}[Jensen 不等式的三大高频推论]
\label{cor:ch3-jensen-cor}
由定理 \ref{thm:ch3-jensen} 直接代入常用凸/凹函数，立得：
\begin{enumerate}
  \item \textbf{方差非负性}（取凸函数 $g(x)=x^2$）：$\E[X^2] \ge (\E[X])^2$，即 $\Var(X) = \E[X^2] - (\E[X])^2 \ge 0$；
  \item \textbf{对数期望不等式}（取严格凹函数 $g(x)=\ln x$，$X > 0$）：
  \begin{equation}
    \E[\ln X] \le \ln \E[X], \label{eq:ch3-log-jensen}
  \end{equation}
  等号当且仅当 $X$ 几乎必然为正数常量时成立；
  \item \textbf{倒数期望不等式}（取严格凸函数 $g(x)=1/x$，$X > 0$）：$\E[1/X] \ge 1/\E[X]$。
\end{enumerate}
\end{corollary}

\begin{insightBox}{一根 Jensen 线串起信息论、EM 算法与 VAE}
不要低估式 \eqref{eq:ch3-log-jensen} 的 $\E[\ln X] \le \ln \E[X]$——它是整个概率机器学习里“出场率第一”的推导引擎：
\begin{itemize}
  \item \textbf{本卷第 4 章（Gibbs 不等式）}：把 $X = Q/P$ 代入 $\E_P[-\ln X] \ge -\ln \E_P[X]$，3 行即证 KL 散度非负性 $\DKL(P\|Q) \ge 0$；
  \item \textbf{第 4 卷（EM 算法）与第 5 卷（变分自编码器 VAE）}：面对含隐变量 $z$ 的难算对数似然 $\ln p_\theta(x) = \ln \int p_\theta(x,z)\,dz$，引入任意变分分布 $q(z\mid x)$ 做一次恒等改写再套用 \eqref{eq:ch3-log-jensen}：
  \[
    \ln p_\theta(x)
    = \ln \E_{z\sim q}\!\left[\frac{p_\theta(x,z)}{q(z\mid x)}\right]
    \ge \E_{z\sim q}\!\left[\ln \frac{p_\theta(x,z)}{q(z\mid x)}\right]
    =: \mathrm{ELBO}(q, \theta).
  \]
  这就是大名鼎鼎的\textbf{证据下界（Evidence Lower Bound, ELBO）}的出生证明！
\end{itemize}
\end{insightBox}

\begin{failBox}{把非线性运算“直接穿进期望”是新手头号暗坑}
工程代码与公式推导中，极易在无意间写出 $\E[1/X] = 1/\E[X]$、$\E[\sqrt{X}] = \sqrt{\E[X]}$ 或 $\E[\ln \hat{p}] = \ln \E[\hat{p}]$ 之类的错误等式。例如，某模型在两条样本上预测正确类别的概率分别是 $0.01$ 与 $0.99$：
\begin{itemize}
  \item 先平均概率再取负对数：$-\ln\bigl(\tfrac{0.01+0.99}{2}\bigr) = -\ln 0.5 \approx 0.693$；
  \item 真实的平均交叉熵（先取负对数再平均）：$\tfrac{-\ln 0.01 - \ln 0.99}{2} \approx 2.308$！
\end{itemize}
两者相差超过 $3$ 倍！由于 $-\ln x$ 是严格凸函数，真实交叉熵损失永远严格惩罚极端自信的错误预测，绝不能把样本概率先做平均再代进对数。
\end{failBox}

% =========================================================================
\section{浓度不等式三阶梯：Markov、Chebyshev 与 Hoeffding}
\label{sec:ch3-concentration}

设独立同分布（i.i.d.）样本 $X_1, X_2, \dots, X_n$ 具有真实均值 $\mu = \E[X_1]$，我们计算经验样本均值 $\bar{X}_n = \frac{1}{n}\sum_{i=1}^n X_i$。在 AI 评测中，$X_i = \mathbf{1}\{\hat{y}_i = y_i\}$ 表示第 $i$ 条测试样本是否预测正确，那么 $\bar{X}_n$ 就是测试集准确率，而 $\mu$ 是模型在真实数据分布上的泛化准确率。

我们最想控制的是\textbf{大偏差尾概率}：
\[
  \Prob\bigl(|\bar{X}_n - \mu| \ge \varepsilon\bigr) \le \text{?},
\]
即“经验均值偏离真实期望至少 $\varepsilon$ 的坏运气概率，最多能有多大？”下面我们沿着信息量递增的三级阶梯，一步步把尾概率上界从粗糙压到指数级锐利。

\subsection{第一阶梯：Markov 不等式（只需非负与均值）}

先看一个生活直觉：如果某推理服务所有请求的\textbf{平均延迟}是 $20\,\text{ms}$（延迟显然非负），那么延迟 $\ge 100\,\text{ms}$ 的慢请求占比绝对不可能超过 $20\%$——否则光这些慢请求就已经把全体平均值拉超过 $0.20 \times 100 = 20\,\text{ms}$ 了！把这段朴素算术翻译成概率语言，就是 Markov 不等式。

\begin{thmBox}
\begin{theorem}[Markov 不等式]
\label{thm:ch3-markov}
设非负随机变量 $X \ge 0$ 的期望 $\E[X] < \infty$，则对任意常数 $a > 0$，有
\begin{equation}
  \Prob(X \ge a) \le \frac{\E[X]}{a}. \label{eq:ch3-markov}
\end{equation}
\end{theorem}
\end{thmBox}

\begin{proof}
构造指示随机变量 $\mathbf{1}_{\{X \ge a\}}$（当 $X \ge a$ 时取 $1$，否则取 $0$）。由于 $X \ge 0$，在任意样本点上都逐点满足
\[
  a \cdot \mathbf{1}_{\{X \ge a\}} \le X.
\]
（当 $X \ge a$ 时左边为 $a \le X$；当 $0 \le X < a$ 时左边为 $0 \le X$。）对两边取期望，利用 $\E[\mathbf{1}_{\{X \ge a\}}] = \Prob(X \ge a)$，即得
\[
  a \cdot \Prob(X \ge a) \le \E[X] \implies \Prob(X \ge a) \le \frac{\E[X]}{a}.
\]
\end{proof}

Markov 不等式的优势在于条件极弱（不要求方差存在、不要求任何分布族假设），且上界本身是紧的（取两点分布 $\Prob(X=a)=p, \Prob(X=0)=1-p$ 即可取等号）。但正因为只用了“一阶矩（均值）”，它的衰减速率仅有 $O(1/a)$，面对样本均值集中问题远远不够尖锐。注意\textbf{非负前提 $X \ge 0$ 不可省略}：若允许 $X$ 取负值，即使 $\E[X] = 0$ 也完全可以有 $\Prob(X \ge 1) = \tfrac{1}{2} > 0$。

\subsection{第二阶梯：Chebyshev 不等式与弱大数定律（WLLN）}

如果我们额外知道随机变量的二阶矩（方差 $\sigma^2 = \Var(X)$），就能立刻把 Markov 不等式套在非负随机变量 $Y = (X - \mu)^2$ 上，得到第二级阶梯。

\begin{thmBox}
\begin{theorem}[Chebyshev 不等式]
\label{thm:ch3-chebyshev}
设随机变量 $X$ 具有有限均值 $\mu = \E[X]$ 与有限方差 $\sigma^2 = \Var(X)$，则对任意 $\varepsilon > 0$，有
\begin{equation}
  \Prob\bigl(|X - \mu| \ge \varepsilon\bigr) \le \frac{\Var(X)}{\varepsilon^2}. \label{eq:ch3-chebyshev}
\end{equation}
\end{theorem}
\end{thmBox}

\begin{proof}
事件 $\{|X - \mu| \ge \varepsilon\}$ 与 $\{(X - \mu)^2 \ge \varepsilon^2\}$ 完全等价。由于 $(X - \mu)^2 \ge 0$，直接对其应用定理 \ref{thm:ch3-markov}（取 $a = \varepsilon^2$），立得
\[
  \Prob\bigl(|X - \mu| \ge \varepsilon\bigr)
  = \Prob\bigl((X - \mu)^2 \ge \varepsilon^2\bigr)
  \le \frac{\E[(X - \mu)^2]}{\varepsilon^2}
  = \frac{\Var(X)}{\varepsilon^2}.
\]
\end{proof}

把 Chebyshev 不等式立刻应用到 i.i.d. 样本均值 $\bar{X}_n = \frac{1}{n}\sum_{i=1}^n X_i$ 上，我们三行就能收获概率论的基石之一——\textbf{弱大数定律}！

\begin{thmBox}
\begin{theorem}[弱大数定律, Weak Law of Large Numbers, WLLN]
\label{thm:ch3-wlln}
设 $X_1, X_2, \dots, X_n$ 独立同分布，具有有限均值 $\mu$ 与有限方差 $\sigma^2$。则样本均值 $\bar{X}_n = \frac{1}{n}\sum_{i=1}^n X_i$ 依概率收敛于 $\mu$（记作 $\bar{X}_n \xrightarrow{p} \mu$）：对任意 $\varepsilon > 0$，
\begin{equation}
  \Prob\bigl(|\bar{X}_n - \mu| \ge \varepsilon\bigr) \le \frac{\sigma^2}{n \varepsilon^2} \xrightarrow{n \to \infty} 0. \label{eq:ch3-wlln}
\end{equation}
\end{theorem}
\end{thmBox}

\begin{proof}
由第 1 章期望线性性与独立变量方差可加性，有
\[
  \E[\bar{X}_n] = \frac{1}{n}\sum_{i=1}^n \E[X_i] = \mu,
  \qquad
  \Var(\bar{X}_n) = \frac{1}{n^2}\sum_{i=1}^n \Var(X_i) = \frac{\sigma^2}{n}.
\]
将 $\bar{X}_n$ 代入 Chebyshev 不等式 \eqref{eq:ch3-chebyshev} 即得 $\Prob(|\bar{X}_n - \mu| \ge \varepsilon) \le \frac{\sigma^2}{n\varepsilon^2}$。令 $n \to \infty$，右端趋于 $0$，证毕。
\end{proof}

\textbf{停下来算一笔现实账}：WLLN 虽然保证了“样本越多越准”，但 Chebyshev 给出的衰减速率只有多项式级 $O(1/n)$。假设 $X_i \sim \mathrm{Bernoulli}(p)$（方差 $p(1-p) \le \tfrac{1}{4}$），如果我们希望误差 $\varepsilon = 0.03$ 的失败概率不超过 $\delta = 0.05$（$95\%$ 置信度），Chebyshev 要求
\[
  \frac{0.25}{n \times 0.03^2} \le 0.05 \implies n \ge 5556.
\]
更糟的是，若要求工业高可靠场景的 $99.9\%$ 置信度（$\delta = 0.001$），Chebyshev 因为对 $\delta$ 呈 $1/\delta$ 依赖，竟然要求高达 $n \ge 277{,}778$ 条测试样本！为什么 Chebyshev 这么保守？因为它只利用了二阶矩，不得不为“极端长尾分布”买单。而在 AI 中，准确率指示器 $X_i \in [0, 1]$、截断损失 $L_i \in [0, B]$ 都是\textbf{严格有界}的！能否用有界性换取\textbf{指数级衰减}？

\subsection{第三阶梯：Hoeffding 不等式（有界独立变量的指数级集中）}

把平方函数 $(X-\mu)^2$ 换成指数函数 $e^{s(\bar{X}_n - \mu)}$（其中待定参数 $s > 0$），再套用 Markov 不等式，就能触发所谓的 \textbf{Chernoff 指数技巧}——指数函数会把独立变量的\textbf{和}瞬间变成期望的\textbf{乘积}！

\begin{thmBox}
\begin{theorem}[Hoeffding 不等式]
\label{thm:ch3-hoeffding}
设 $X_1, X_2, \dots, X_n$ 是相互独立的随机变量，且每个变量几乎必然落在有界闭区间内：$a_i \le X_i \le b_i$（$i=1,\dots,n$）。记样本均值 $\bar{X}_n = \frac{1}{n}\sum_{i=1}^n X_i$ 与其期望 $\mu = \E[\bar{X}_n]$。则对任意 $\varepsilon > 0$，单侧与双侧尾概率分别满足：
\begin{align}
  \Prob\bigl(\bar{X}_n - \mu \ge \varepsilon\bigr) &\le \exp\!\left(-\frac{2 n^2 \varepsilon^2}{\sum_{i=1}^n (b_i - a_i)^2}\right), \label{eq:ch3-hoeffding-1}\\
  \Prob\bigl(|\bar{X}_n - \mu| \ge \varepsilon\bigr) &\le 2\exp\!\left(-\frac{2 n^2 \varepsilon^2}{\sum_{i=1}^n (b_i - a_i)^2}\right). \label{eq:ch3-hoeffding-2}
\end{align}
特别地，当所有 $X_i \in [0, 1]$（如二元准确率、归一化损失）时，上式化简为极易心算的黄金公式：
\begin{equation}
  \Prob\bigl(|\bar{X}_n - \mu| \ge \varepsilon\bigr) \le 2 e^{-2 n \varepsilon^2}. \label{eq:ch3-hoeffding-01}
\end{equation}
\end{theorem}
\end{thmBox}

\begin{proof}[主线证明（Chernoff 技巧 + Hoeffding 引理）]
记 $Z_i := X_i - \E[X_i]$，则 $\E[Z_i] = 0$，且 $Z_i$ 落在长度为 $b_i - a_i$ 的区间内。对任意待优化的常数 $s > 0$，利用指数函数严格单调递增与 Markov 不等式（定理 \ref{thm:ch3-markov}）：
\[
  \Prob(\bar{X}_n - \mu \ge \varepsilon)
  = \Prob\!\left(\sum_{i=1}^n Z_i \ge n\varepsilon\right)
  = \Prob\!\left(e^{s \sum_{i=1}^n Z_i} \ge e^{s n \varepsilon}\right)
  \le e^{-s n \varepsilon}\, \E\!\left[\prod_{i=1}^n e^{s Z_i}\right].
\]
由 $Z_1, \dots, Z_n$ 相互独立，乘积的期望等于期望的乘积：
\[
  \Prob(\bar{X}_n - \mu \ge \varepsilon)
  \le e^{-s n \varepsilon} \prod_{i=1}^n \E\bigl[e^{s Z_i}\bigr].
\]
引用下方选读框中的 \textbf{Hoeffding 引理}：对任意零均值且取值于长度 $b_i - a_i$ 区间内的随机变量 $Z_i$，均有 $\E[e^{s Z_i}] \le \exp\!\bigl(\frac{s^2 (b_i - a_i)^2}{8}\bigr)$。代入得
\[
  \Prob(\bar{X}_n - \mu \ge \varepsilon)
  \le \exp\!\left(-s n \varepsilon + \frac{s^2}{8}\sum_{i=1}^n (b_i - a_i)^2\right).
\]
指数括号里是关于 $s$ 的开口向上二次函数！求导令零即得最优参数 $s^* = \frac{4 n \varepsilon}{\sum_{i=1}^n (b_i - a_i)^2} > 0$，代回立得单侧界 \eqref{eq:ch3-hoeffding-1}。将 $X_i$ 替换为 $-X_i \in [-b_i, -a_i]$ 得到完全相同的下尾界 $\Prob(\bar{X}_n - \mu \le -\varepsilon) \le \exp\bigl(-\frac{2n^2\varepsilon^2}{\sum (b_i-a_i)^2}\bigr)$，由并集界 $\Prob(A \cup B) \le \Prob(A) + \Prob(B)$ 将两式相加即得双侧界 \eqref{eq:ch3-hoeffding-2}。
\end{proof}

\begin{extBox}{Hoeffding 引理的完整零跳步证明（对数矩母函数 + Taylor 展开）}
\textbf{引理}：设随机变量 $Z$ 满足 $\E[Z] = 0$ 且几乎必然落在区间 $[c, d]$ 中（其中 $d - c = b - a$），则对任意 $s \in \R$ 有
\[
  \E\bigl[e^{s Z}\bigr] \le \exp\!\left(\frac{s^2 (b-a)^2}{8}\right).
\]
\textbf{证明}：定义对数矩母函数 $\psi(s) := \ln \E[e^{s Z}]$。易见 $\psi(0) = \ln 1 = 0$。对 $s$ 求一阶与二阶导数：
\[
  \psi'(s) = \frac{\E[Z e^{sZ}]}{\E[e^{sZ}]},
  \qquad
  \psi''(s) = \frac{\E[Z^2 e^{sZ}]}{\E[e^{sZ}]} - \left(\frac{\E[Z e^{sZ}]}{\E[e^{sZ}]}\right)^2.
\]
代入 $s = 0$ 得 $\psi'(0) = \E[Z] = 0$。奇妙的是，若定义一个新的“指数倾斜”概率分布 $Q_s$（密度或质量正比于 $e^{sz} p(z)$），则上式恰好就是随机变量 $Z$ 在新分布 $Q_s$ 下的方差：
\[
  \psi''(s) = \Var_{Q_s}(Z).
\]
而对\textbf{任何}取值在 $[c, d]$ 内的随机变量 $Y$，以区间中点 $m_0 = \frac{c+d}{2}$ 为参照恒有 $|Y - m_0| \le \frac{d-c}{2} = \frac{b-a}{2}$，从而
\[
  \Var_{Q_s}(Z) = \E_{Q_s}\!\bigl[(Z - \E_{Q_s}[Z])^2\bigr] \le \E_{Q_s}\!\bigl[(Z - m_0)^2\bigr] \le \frac{(b-a)^2}{4}.
\]
（第一步利用了“均值使均方误差最小”，即 $\E[(Z-c_0)^2] = \Var(Z) + (\E[Z]-c_0)^2 \ge \Var(Z)$。）最后对 $\psi(s)$ 在 $0$ 处做带拉格朗日余项的二阶 Taylor 展开：存在 $\xi \in [0, s]$ 使得
\[
  \psi(s) = \psi(0) + \psi'(0)s + \frac{1}{2}\psi''(\xi) s^2 \le 0 + 0 + \frac{1}{2}\cdot\frac{(b-a)^2}{4} s^2 = \frac{s^2 (b-a)^2}{8}.
\]
两边取指数即得 $\E[e^{sZ}] \le \exp\!\bigl(\frac{s^2(b-a)^2}{8}\bigr)$。证毕！
\end{extBox}

\subsection{AI 实战算账与本机实测：测试集到底需要多少样本？}

现在我们用 Hoeffding 黄金公式 \eqref{eq:ch3-hoeffding-01} 亲手做一次 AI 评测样本量设计。

\begin{example}[分类模型测试集样本量计算]
\label{ex:ch3-test-size}
某二分类模型在独立测试集 $\{(x_i, y_i)\}_{i=1}^n$ 上评测，记指示器 $X_i = \mathbf{1}\{\hat{y}_i = y_i\} \in [0, 1]$，经验准确率 $\widehat{\mathrm{Acc}}_n = \bar{X}_n$，真实泛化准确率 $\mathrm{Acc}^* = \E[X_i]$。我们希望以至少 $1 - \delta$ 的概率保证 $|\widehat{\mathrm{Acc}}_n - \mathrm{Acc}^*| < \varepsilon$。由 \eqref{eq:ch3-hoeffding-01}，令
\[
  2 e^{-2 n \varepsilon^2} \le \delta
  \iff 2 n \varepsilon^2 \ge \ln\!\left(\frac{2}{\delta}\right)
  \iff n \ge \frac{1}{2\varepsilon^2}\ln\!\left(\frac{2}{\delta}\right).
\]
代入典型需求 $\varepsilon = 0.03$（误差不超过 $3$ 个百分点）、$\delta = 0.05$（$95\%$ 置信度）：
\[
  n \ge \frac{\ln 40}{2 \times 0.03^2} = \frac{3.6889}{0.0018} \approx 2049.4 \implies n \ge 2050.
\]
如果把置信度一口气拉高到 $99.9\%$（$\delta = 0.001$），由于对数 $\ln(2/\delta)$ 增长极慢，Hoeffding 只需要
\[
  n \ge \frac{\ln 2000}{0.0018} \approx 4222.7 \implies n \ge 4223,
\]
相比刚才 Chebyshev 索要的 $277{,}778$ 条样本，直接节省了超过 \textbf{65 倍} 的标注成本！
\end{example}

表 \ref{tab:ch3-conc-compare} 系统对比了浓度不等式三阶梯的假设、形式与样本量代价；图 \ref{fig:ch3-hoeffding} 则直接读取本卷配套实验脚本 \texttt{code/vol2\_experiments.py} 生成的实测数据 \texttt{data/hoeffding\_tail.csv}（$X_i \stackrel{\text{i.i.d.}}{\sim} \mathrm{Bernoulli}(0.5)$，偏差阈值 $\varepsilon = 0.1$，$20{,}000$ 次重复蒙特卡洛试验），在半对数坐标下直观展示 Chebyshev 与 Hoeffding 的鸿沟。

\begin{table}[htbp]
  \centering
  \caption{浓度不等式三阶梯核心对比（针对 $X_i \in [0,1]$ 的独立样本均值 $\bar{X}_n$）}
  \label{tab:ch3-conc-compare}
  \small
  \begin{tabularx}{\linewidth}{@{}p{22mm}p{30mm}p{31mm}p{27mm}X@{}}
    \toprule
    \textbf{阶梯} & \textbf{所需假设} & \textbf{尾概率上界} & \textbf{样本量 $n(\varepsilon,\delta)$} & \textbf{典型用途与局限} \\
    \midrule
    Markov & 非负 + 有限均值 & $\E[\bar{X}_n] / a$ & —— & 仅控单侧绝对大小，最通用也最松 \\
    Chebyshev & 有限方差 $\sigma^2$ & $\sigma^2 / (n \varepsilon^2)$ & $O\bigl(\tfrac{1}{\varepsilon^2 \delta}\bigr)$ & 3 行证 WLLN；对 $\delta$ 呈 $1/\delta$ 依赖过慢 \\
    Hoeffding & 独立 + 有界 $[a_i,b_i]$ & $2\exp(-2n\varepsilon^2)$ & $O\bigl(\tfrac{1}{\varepsilon^2}\ln\tfrac{2}{\delta}\bigr)$ & 指数级衰减；测试集规模与泛化界主力 \\
    \bottomrule
  \end{tabularx}
\end{table}

\begin{figure}[htbp]
  \centering
  \begin{tikzpicture}
    \begin{axis}[
      width=0.88\linewidth,
      height=6.6cm,
      ymode=log,
      xlabel={样本量 $n$},
      ylabel={尾概率 $\Prob(|\bar{X}_n - 0.5| \ge 0.1)$（对数刻度）},
      xmin=0, xmax=315,
      ymin=2e-4, ymax=2.2,
      grid=both,
      major grid style={slate!22},
      minor grid style={slate!10},
      legend pos=south west,
      legend style={font=\footnotesize\sffamily, fill=white, fill opacity=0.9, draw=codeborder},
      label style={font=\small\sffamily, ink},
      tick label style={font=\footnotesize, slate}
    ]
      % Horizontal reference line at probability = 1
      \addplot[slate!65, dotted, thick, domain=0:315, forget plot] {1.0};
      \node[slate, font=\scriptsize\sffamily, anchor=south east] at (axis cs:308, 1.05) {平凡概率上界 $1$};

      % Chebyshev bound
      \addplot[navy, thick, dashed, mark=square*, mark size=1.8pt]
        table[x=n, y=chebyshev, col sep=comma] {data/hoeffding_tail.csv};
      \addlegendentry{Chebyshev 上界 $\frac{0.25}{n\varepsilon^2}$（多项式 $O(1/n)$）}

      % Hoeffding bound
      \addplot[teal, very thick, mark=*, mark size=1.9pt]
        table[x=n, y=hoeffding, col sep=comma] {data/hoeffding_tail.csv};
      \addlegendentry{Hoeffding 上界 $2e^{-2n\varepsilon^2}$（指数衰减）}

      % Empirical tail
      \addplot[ochre, very thick, mark=triangle*, mark size=2.3pt]
        table[x=n, y=emp_tail, col sep=comma] {data/hoeffding_tail.csv};
      \addlegendentry{本机 20{,}000 次蒙特卡洛实测尾概率}

      % Callout annotations at n=300
      \node[navy, font=\scriptsize\sffamily, anchor=south east] at (axis cs:296, 0.105) {$0.0833$};
      \node[tealDark, font=\scriptsize\sffamily, anchor=south east] at (axis cs:296, 0.011) {$0.00496$};
      \node[ochre, font=\scriptsize\sffamily, anchor=south east] at (axis cs:296, 0.00088) {$0.00040$};
    \end{axis}
  \end{tikzpicture}
  \caption{本机实测（\texttt{data/hoeffding\_tail.csv}）：对 $\mathrm{Bernoulli}(0.5)$ 与 $\varepsilon = 0.1$，Chebyshev 上界在半对数图上越走越平缓（$n=300$ 时仍高达 $0.0833$），而 Hoeffding 上界呈笔直向下的指数衰减（$n=300$ 时降至 $0.00496$，比 Chebyshev 紧 $16.8$ 倍），真实经验尾概率则已低至 $0.00040$。注：$n < 35$ 时未经截断的 $2e^{-2n\varepsilon^2} > 1$，此时自然以 $1$ 截断。}
  \label{fig:ch3-hoeffding}
\end{figure}

\begin{insightBox}{一步通向第 4 卷 PAC 泛化理论（Hoeffding + 并集界）}
如果手里不是只有 $1$ 个固定模型，而是要在包含 $M = |\mathcal{H}|$ 个候选模型的假设空间 $\mathcal{H}$ 里挑训练误差最小的模型，怎么防止“挑花了眼（过拟合）”？只需对每个模型 $h \in \mathcal{H}$ 套用 Hoeffding 不等式，再用 1 行并集界把 $M$ 个坏事件的概率相加：
\[
  \Prob\!\left(\exists h \in \mathcal{H}:\; |\widehat{R}_n(h) - R(h)| \ge \varepsilon\right) \le 2 M e^{-2 n \varepsilon^2} \le \delta
  \implies
  |\widehat{R}_n(h) - R(h)| \le \sqrt{\frac{\ln(2M / \delta)}{2n}}.
\]
你看，整个统计学习理论（PAC 泛化界）的骨架就是\textbf{今天学的 Hoeffding 不等式 + 并集界}：模型复杂度只以 $\sqrt{\ln M}$ 的温和速度增长，这就是第 4 卷机器学习泛化理论的真正起点！
\end{insightBox}

% =========================================================================
\section{中心极限定理（CLT）、蒙特卡洛估计与重要性采样}
\label{sec:ch3-clt-mc}

Hoeffding 不等式强在“对任意有限样本 $n$ 严格成立、不假设任何分布形状”，但它为了守住最坏情况，有时会偏保守（比如图 \ref{fig:ch3-hoeffding} 中 $n=300$ 时 Hoeffding 为 $0.00496$，而实测只有 $0.00040$）。如果我们既想知道样本均值“有多集中”，又想知道它的\textbf{精确渐近分布形状}，就需要请出概率论的另一座丰碑——中心极限定理。

\subsection{中心极限定理（CLT）与渐近置信区间}

\begin{thmBox}
\begin{theorem}[中心极限定理, Central Limit Theorem, CLT]
\label{thm:ch3-clt}
设 $X_1, X_2, \dots, X_n$ 独立同分布，具有有限均值 $\mu = \E[X_1]$ 与有限正方差 $0 < \sigma^2 = \Var(X_1) < \infty$。记样本均值 $\bar{X}_n = \frac{1}{n}\sum_{i=1}^n X_i$，则当 $n \to \infty$ 时，标准化后的随机变量依分布收敛于标准正态分布：
\begin{equation}
  Z_n := \frac{\sqrt{n}(\bar{X}_n - \mu)}{\sigma} \xrightarrow{d} \mathcal{N}(0, 1). \label{eq:ch3-clt}
\end{equation}
等价地，对任意实数 $z \in \R$，$\lim_{n\to\infty} \Prob(Z_n \le z) = \Phi(z) = \int_{-\infty}^z \frac{1}{\sqrt{2\pi}}e^{-t^2/2}\,dt$。
\end{theorem}
\end{thmBox}

CLT 令人着迷的魔力在于\textbf{普适性}：无论单条样本 $X_i$ 服从离散的 Bernoulli 分布、严重右偏的指数分布，还是某个未知形态的神经网络单样本交叉熵损失，只要 $X_i$ 独立同分布且方差有限，把大量样本求平均后，$\bar{X}_n$ 的涨落分布就会自动被“磨平”成钟形正态分布 $\mathcal{N}\bigl(\mu,\, \sigma^2/n\bigr)$！

由此立即可构造工程里最常用的\textbf{大样本 $95\%$ 渐近置信区间}：由于标准正态分布满足 $\Prob(|Z| \le 1.96) \approx 0.95$，用样本标准差 $\hat{\sigma}$ 替代未知 $\sigma$，即得
\begin{equation}
  \bar{X}_n \pm 1.96\,\frac{\hat{\sigma}}{\sqrt{n}}. \label{eq:ch3-clt-ci}
\end{equation}
拿例 \ref{ex:ch3-test-size} 的 $n = 2050$ 做个对照：若测得模型准确率 $\widehat{\mathrm{Acc}} = 0.85$，样本标准差 $\hat{\sigma} = \sqrt{0.85 \times 0.15} \approx 0.3571$，则 CLT 给出的 $95\%$ 置信区间半宽仅为
\[
  1.96 \times \frac{0.3571}{\sqrt{2050}} \approx 0.0155,
\]
差不多只有 Hoeffding 最坏情况界（$0.0300$）的一半！原因一目了然：Hoeffding 用的是区间宽度 $b-a=1$ 对应的最坏方差上界 $0.25$ 且承诺对\textbf{每一个有限 $n$} 绝对成立；而 CLT 既利用了实际更小的方差 $0.85\times 0.15 = 0.1275$，又用上了正态分布的精确尾面积（代价是依赖大样本近似，样本太小时可能失准；第 5 章我们将用 \textbf{Bootstrap} 在无需正态近似时直接重采样估计置信区间）。

\begin{failBox}{CLT 不是免死金牌：方差无穷与样本相关时的彻底失效}
套用 $\bar{X}_n \pm 1.96\hat{\sigma}/\sqrt{n}$ 前，务必核对两条红线：
\begin{enumerate}
  \item \textbf{重尾导致方差无穷}：例如标准\textbf{柯西分布（Cauchy）} $p(x) = \frac{1}{\pi(1+x^2)}$，它的期望与方差都不存在。可以证明，$n$ 个独立柯西变量的样本均值 $\bar{X}_n$ 与单个 $X_1$ 服从\textbf{完全相同的柯西分布}——无论取一万还是十亿个样本，均值都永远不会集中，也不会变成高斯！
  \item \textbf{样本不独立}：在时间序列预测、股票收益或同一用户多次出现的推荐日志里，样本间存在强正相关性，真实方差 $\Var(\bar{X}_n)$ 可能远大于 $\sigma^2/n$，盲用 CLT 公式会严重低估不确定性（给出虚假的窄置信区间）。
\end{enumerate}
\end{failBox}

\subsection{蒙特卡洛估计：用随机性打破维数灾难}

在现代 AI 里，大量核心计算都可以写成某个高维期望积分：
\[
  I = \E_{x \sim p}[f(x)] = \int_{\R^d} f(x) p(x)\,dx.
\]
例如：计算复杂后验分布的均值、评估策略网络在随机环境中的期望回报、估计生成模型的归一化常数。如果按传统微积分做\textbf{确定性网格数值积分}（每条坐标轴切 $k$ 个格子），总计算点数高达 $M = k^d$，换算成关于预算 $M$ 的误差阶为 $O(k^{-1}) = O(M^{-1/d})$。一旦维度 $d = 50$，哪怕每轴只切可怜的 $k = 2$ 个点，也需要 $2^{50} \approx 1.1 \times 10^{15}$ 次计算——这就是令人绝望的\textbf{维数灾难（Curse of Dimensionality）}。

\textbf{蒙特卡洛方法（Monte Carlo Estimation）}反其道而行之：放弃规则网格，直接从分布 $p(x)$ 中独立抽取 $M$ 个随机样本 $x_1, \dots, x_M \stackrel{\text{i.i.d.}}{\sim} p$，用样本均值去估计期望。

\begin{defBox}
\begin{definition}[蒙特卡洛估计量]
\label{def:ch3-mc}
给定目标分布 $p$ 与可积函数 $f$，独立抽取 $x_1, \dots, x_M \stackrel{\text{i.i.d.}}{\sim} p$，定义\textbf{蒙特卡洛估计量}为
\begin{equation}
  \hat{I}_M := \frac{1}{M}\sum_{m=1}^M f(x_m). \label{eq:ch3-mc}
\end{equation}
\end{definition}
\end{defBox}

\begin{proposition}[蒙特卡洛估计的无偏性与维数无关收敛速度]
\label{prop:ch3-mc}
设 $I = \E_{x\sim p}[f(x)]$ 与 $\sigma_f^2 = \Var_{x\sim p}(f(x))$ 均有限，则蒙特卡洛估计量 $\hat{I}_M$ 满足：
\begin{enumerate}
  \item \textbf{严格无偏}：$\E[\hat{I}_M] = \frac{1}{M}\sum_{m=1}^M \E[f(x_m)] = I$；
  \item \textbf{标准误与维度 $d$ 完全无关}：
  \begin{equation}
    \Var(\hat{I}_M) = \frac{\sigma_f^2}{M},
    \qquad
    \mathrm{SE}(\hat{I}_M) := \sqrt{\Var(\hat{I}_M)} = \frac{\sigma_f}{\sqrt{M}} = O\!\left(M^{-1/2}\right). \label{eq:ch3-mc-se}
  \end{equation}
  \item \textbf{相合性与渐近正态}：由 WLLN（定理 \ref{thm:ch3-wlln}）与 CLT（定理 \ref{thm:ch3-clt}），$\hat{I}_M \xrightarrow{p} I$ 且大样本下近似服从 $\mathcal{N}(I, \sigma_f^2/M)$。
\end{enumerate}
\end{proposition}

仔细品味 \eqref{eq:ch3-mc-se}：无论变量 $x$ 是 1 维、100 维还是 100 万维，蒙特卡洛的均方根误差阶永远是 $O(1/\sqrt{M})$，彻底摆脱了对维度 $d$ 的指数诅咒！不过它也有一把双刃剑：想把误差缩小到原来的 $1/10$，采样数 $M$ 必须暴涨 $100$ 倍。这就引出了降低方差、解决难采样问题的核心武器——\textbf{重要性采样}。

\subsection{重要性采样（Importance Sampling）与 PPO 比率前瞻}

设想两个常见困境：(a) 目标分布 $p(x)$ 过于复杂，无法直接从中采样（比如复杂策略或后验分布）；(b) 即使能从 $p(x)$ 采样，我们关心的区域 $\{f(x) \neq 0\}$ 在 $p$ 下是\textbf{极小概率的罕见事件}！

举个震撼的例子：用标准正态 $Z \sim \mathcal{N}(0, 1)$ 估计罕见尾概率 $I = \Prob(Z > 4) \approx 3.167 \times 10^{-5}$。若直接用朴素蒙特卡洛抽 $M = 10{,}000$ 个样本，期望命中次数只有 $10000 \times 3.167\times 10^{-5} \approx 0.32$ 次——大概率一次都抽不到，直接输出 $\hat{I}_{\mathrm{MC}} = 0$（相对误差 $100\%$）！

解决思路极其巧妙：找一个容易采样、且在“关键区域”落点密集的\textbf{提议分布（Proposal Distribution）} $q(x)$，从 $q$ 中采样，再乘上一个修正权重把分布“搬回” $p$。

\begin{thmBox}
\begin{proposition}[重要性采样无偏恒等式]
\label{prop:ch3-is}
设提议分布 $q(x)$ 满足\textbf{支集覆盖条件}：对所有 $x$，只要 $p(x)f(x) \neq 0$，必有 $q(x) > 0$。定义\textbf{重要性权重（Importance Weight）} $w(x) := \frac{p(x)}{q(x)}$，则
\begin{equation}
  \E_{x \sim p}[f(x)]
  = \int f(x) p(x)\,dx
  = \int f(x)\,\frac{p(x)}{q(x)}\,q(x)\,dx
  = \E_{x \sim q}\bigl[f(x)\, w(x)\bigr]. \label{eq:ch3-is}
\end{equation}
相应地，抽取 $x_1, \dots, x_M \stackrel{\text{i.i.d.}}{\sim} q$ 构造的\textbf{重要性采样估计量}
\[
  \hat{I}_{\mathrm{IS}} := \frac{1}{M}\sum_{m=1}^M f(x_m)\,\frac{p(x_m)}{q(x_m)}
\]
是 $I = \E_p[f(x)]$ 的严格无偏估计量。
\end{proposition}
\end{thmBox}

回到刚才的罕见尾概率 $\Prob(Z > 4)$！我们把提议分布中心直接平移到悬崖边缘，取 $q = \mathcal{N}(4, 1)$。此时约有一半样本（约 $5{,}000$ 个）精准落在 $Z > 4$ 区域内，且每个样本的权重解析极简：
\[
  w(z) = \frac{\frac{1}{\sqrt{2\pi}}e^{-z^2/2}}{\frac{1}{\sqrt{2\pi}}e^{-(z-4)^2/2}} = \exp(-4z + 8).
\]
在下方实测代码中，同样只用 $M = 10{,}000$ 个样本，朴素 MC 颗粒无收（估计为 $0$），而重要性采样一步给出 $3.004 \times 10^{-5} \pm 0.126 \times 10^{-5}$，与精确值 $3.167 \times 10^{-5}$ 高度吻合！

\begin{failBox}{重要性采样的头号杀手：提议尾巴过瘦与权重塌缩}
虽然只要满足支集覆盖条件，$\hat{I}_{\mathrm{IS}}$ 就永远\textbf{无偏}，但它的方差
\[
  \Var_q\bigl(f(X)w(X)\bigr) = \E_q\!\left[f(X)^2\frac{p(X)^2}{q(X)^2}\right] - I^2
\]
可能彻底爆炸！最经典的翻车场景是\textbf{提议分布 $q$ 的尾部比目标分布 $p$ 衰减得更快}：例如 $p = \mathcal{N}(0, 2^2)$ 而 $q = \mathcal{N}(0, 1)$，权重 $w(x) \propto \exp(\tfrac{3}{8}x^2)$ 在尾部呈超指数爆炸，直接导致 $\E_q[w(X)^2] = +\infty$（方差无穷大，详见习题 3.2）！

在高维问题中，即使方差有限，如果 $q$ 与 $p$ 稍有错位，也极易出现“$9999$ 个样本权重接近 $0$，仅 $1$ 个幸运样本独占几乎全部权重”的\textbf{权重塌缩（Weight Collapse）}现象。工程中常用\textbf{有效样本量（Effective Sample Size, ESS）}来实时体检：
\[
  \mathrm{ESS} := \frac{\left(\sum_{m=1}^M w_m\right)^2}{\sum_{m=1}^M w_m^2} \in [1, M].
\]
所有权重相等时 $\mathrm{ESS} = M$（完美利用）；若某单一样本权重独大，则 $\mathrm{ESS} \approx 1$，提示提议分布 $q$ 必须重新设计。
\end{failBox}

\begin{insightBox}{直通第 6 卷强化学习王牌算法 PPO 的截断重要性比率}
在深度强化学习与大模型 RLHF 对齐（第 6 卷）中，采样环境交互轨迹极其昂贵，因此我们只能用\textbf{旧策略 $\pi_{\text{old}}(a\mid s)$} 采集一批轨迹，却想用这批数据连续多步更新\textbf{新策略 $\pi_\theta(a\mid s)$}。这正是标准的重要性采样场景——用比率
\[
  r_t(\theta) := \frac{\pi_\theta(a_t \mid s_t)}{\pi_{\text{old}}(a_t \mid s_t)}
\]
对优势函数 $A_t$ 重加权，把旧策略下的期望搬运到新策略上！一旦 $\pi_\theta$ 走得离 $\pi_{\text{old}}$ 太远，比率 $r_t(\theta)$ 就会发生上面警告过的\textbf{方差爆炸}。OpenAI 提出的 \textbf{PPO（Proximal Policy Optimization）} 的神来之笔，就是把 $r_t(\theta)$ 强行截断在 $[1-\epsilon, 1+\epsilon]$ 区间内（典型 $\epsilon = 0.2$），用一行 \texttt{clip} 稳住了大模型训练的半壁江山！
\end{insightBox}

\subsection{纯 NumPy 实战：蒙特卡洛积分、CLT 收敛与稀尾重要性采样}

下面给出一段零外部依赖、本机可直接运行的 NumPy 验证脚本，一次打通本节三大实验：(1) 蒙特卡洛投针法估 $\pi$ 并核验 $O(1/\sqrt{M})$ 标度律；(2) 用严重右偏的指数分布 $\mathrm{Exp}(1)$ 验证中心极限定理的 $95\%$ 覆盖率；(3) 对比朴素 MC 与平移正态重要性采样在罕见尾事件 $\Prob(Z > 4)$ 上的天壤之别。

\begin{lstlisting}[caption={蒙特卡洛估 $\pi$、指数分布 CLT 验证与 $\Prob(Z>4)$ 重要性采样（纯 NumPy）},label={lst:ch3-mc-is}]
import math
import numpy as np

rng = np.random.default_rng(2026)

# 1. Monte Carlo estimation of pi: verify O(1 / sqrt(M)) standard error
for M in [1_000, 10_000, 100_000]:
    xy = rng.uniform(-1.0, 1.0, size=(M, 2))
    inside = (xy[:, 0]**2 + xy[:, 1]**2 <= 1.0).astype(float)
    pi_hat = 4.0 * inside.mean()
    se = 4.0 * inside.std(ddof=1) / np.sqrt(M)
    print(f"MC pi (M={M:6d}): {pi_hat:.5f} +/- {1.96*se:.5f}")
# Verified output:
# MC pi (M=  1000): 3.20000 +/- 0.09922
# MC pi (M= 10000): 3.15480 +/- 0.03201  (SE shrinks ~3.1x per 10x M)
# MC pi (M=100000): 3.14792 +/- 0.01015

# 2. CLT check on skewed Exponential(1) (mean=1, std=1), n=50, 20000 trials
n_clt, trials = 50, 20_000
means_exp = rng.exponential(scale=1.0, size=(trials, n_clt)).mean(axis=1)
z_scores = np.sqrt(n_clt) * (means_exp - 1.0) / 1.0
cov95 = np.mean(np.abs(z_scores) <= 1.96)
print(f"CLT Exp(1): mean={z_scores.mean():.4f}, std={z_scores.std():.4f}, P(|Z|<=1.96)={cov95:.4f}")
# Verified output: mean=-0.0061, std=0.9987, P(|Z|<=1.96)=0.9513

# 3. Rare event P(Z > 4) for Z ~ N(0, 1): Naive MC vs Importance Sampling q = N(4, 1)
true_p = 0.5 * math.erfc(4.0 / math.sqrt(2.0))   # 3.167124e-05
M_rare = 10_000
p_mc = np.mean(rng.normal(0.0, 1.0, size=M_rare) > 4.0)

z_is = rng.normal(4.0, 1.0, size=M_rare)
w = np.exp(-4.0 * z_is + 8.0)                    # p(z) / q(z)
vals_is = (z_is > 4.0).astype(float) * w
p_is = vals_is.mean()
se_is = vals_is.std(ddof=1) / np.sqrt(M_rare)
print(f"True={true_p:.6e} | Naive MC={p_mc:.6e} | IS={p_is:.6e} +/- {1.96*se_is:.6e}")
# Verified output: True=3.167124e-05 | Naive MC=0.000000e+00 | IS=3.004483e-05 +/- 1.264907e-06
\end{lstlisting}

% =========================================================================
\section{本章习题}
\label{sec:ch3-exercises}

\exTier{A. 概念辨析与反例查错}{严守不等式前提条件，识破方差爆炸与 CLT 失效陷阱}

\begin{enumerate}[label=\textbf{3.\arabic*}, leftmargin=2.4em]
  \item \textbf{（极限、浓度不等式与采样核心概念辨析）} 判断以下四个命题正误；对错误命题给出反例，对正确命题简述理由：
  \begin{enumerate}[label=(\alph*), leftmargin=1.8em]
    \item 对任意实值随机变量 $X$（不必非负）和任意 $a > 0$，只要 $\E[X] > 0$，Markov 不等式 $\Prob(X \ge a) \le \E[X]/a$ 就必然成立；
    \item 若随机变量 $X > 0$ 且不几乎必然为常数，则严格有 $\E[\ln X] < \ln \E[X]$ 且 $\E[1/X] > 1/\E[X]$；
    \item 只要 $X_1,\dots,X_n$ 相互独立同分布，样本均值 $\bar{X}_n$ 在 $n \to \infty$ 时就必然渐近服从正态分布；
    \item 在重要性采样中，若提议分布 $q(x)$ 在目标区域 $\{x: p(x)f(x) \neq 0\}$ 的某个正概率子集上取值为 $0$，那么即使采样数 $M \to \infty$，重要性采样估计量一般也会收敛到错误的值（存在系统性偏差）。
  \end{enumerate}

  \item \textbf{（重要性采样“无偏但方差无穷”反例）} 设目标分布为 $p = \mathcal{N}(0, 2^2)$，提议分布误选为尾部更瘦的标准正态 $q = \mathcal{N}(0, 1)$，权重为 $w(x) = p(x)/q(x)$。
  \begin{enumerate}[label=(\arabic*), leftmargin=1.8em]
    \item 写出 $w(x)$ 的显式表达式，直接积分验证 $\E_{X\sim q}[w(X)] = 1$（说明估计常数函数 $f(x)=1$ 时严格无偏）；
    \item 计算二阶矩 $\E_{X\sim q}[w(X)^2] = \int_{-\infty}^{+\infty} \frac{p(x)^2}{q(x)}\,dx$，证明该积分发散到 $+\infty$，从而 $\Var_q(w(X)) = +\infty$；
    \item 据此总结选择提议分布 $q$ 时的尾部设计原则：为什么工程上要求 $q$ 的尾部“至少与 $p$ 一样厚”？
  \end{enumerate}
\end{enumerate}

\exTier{B. 小例手算与数值直觉}{亲手对比三阶梯上界松紧，算清真实测试集样本量}

\begin{enumerate}[label=\textbf{3.\arabic*}, leftmargin=2.4em, resume]
  \item \textbf{（三阶梯浓度不等式正面交锋）} 抛掷 $n = 100$ 次均匀硬币，$X_i \stackrel{\text{i.i.d.}}{\sim} \mathrm{Bernoulli}(0.5)$，记正面比例 $\bar{X}_{100} = \frac{1}{100}\sum_{i=1}^{100} X_i$（均值 $\mu = 0.5$，方差 $\Var(\bar{X}_{100}) = 0.25/100 = 0.0025$）。请分别手算以下四个上界（保留四位小数）并比较松紧程度：
  \begin{enumerate}[label=(\arabic*), leftmargin=1.8em]
    \item 用 \textbf{Markov 不等式} 给出单侧概率 $\Prob(\bar{X}_{100} \ge 0.65)$ 的上界；
    \item 用 \textbf{Chebyshev 不等式} 给出双侧概率 $\Prob(|\bar{X}_{100} - 0.5| \ge 0.15)$ 的上界；
    \item 用 \textbf{Hoeffding 不等式} 分别给出单侧 $\Prob(\bar{X}_{100} \ge 0.65)$ 与双侧 $\Prob(|\bar{X}_{100} - 0.5| \ge 0.15)$ 的上界。
  \end{enumerate}

  \item \textbf{（AI 评测测试集规模设计：Hoeffding vs CLT）} 某视觉分类模型准备上线，团队希望测试集经验准确率 $\widehat{\mathrm{Acc}}_n$ 与真实准确率 $\mathrm{Acc}^*$ 的偏差不超过 $\varepsilon = 0.02$（$2$ 个百分点）。
  \begin{enumerate}[label=(\arabic*), leftmargin=1.8em]
    \item 使用 Hoeffding 不等式，分别计算 $95\%$ 置信度（$\delta = 0.05$）和 $99\%$ 置信度（$\delta = 0.01$）所需的最小测试集样本量 $n$；
    \item 若已知模型真实准确率大约在 $\mathrm{Acc}^* \approx 0.90$ 附近，改用 CLT 渐近 $95\%$ 置信条件 $1.96\sqrt{p(1-p)/n}$ $\le 0.02$，重新估算所需最小样本量 $n_{\mathrm{CLT}}$；
    \item 从“方差信息”与“有限样本保证 vs 渐近近似”两个角度，解释为何 $n_{\mathrm{CLT}}$ 远小于 Hoeffding 样本量。
  \end{enumerate}
\end{enumerate}

\exTier{C. 实用推导与性质证明}{用 Jensen 拿下 Softplus 与 AM--GM，推导零方差最优提议分布}

\begin{enumerate}[label=\textbf{3.\arabic*}, leftmargin=2.4em, resume]
  \item \textbf{（Softplus 凸性与加权算术-几何平均不等式）} 结合 Jensen 不等式（定理 \ref{thm:ch3-jensen}）完成以下两项证明：
  \begin{enumerate}[label=(\arabic*), leftmargin=1.8em]
    \item 证明神经网络与逻辑回归中高频出现的 Softplus 函数 $f(x) = \ln(1 + e^x)$ 在 $\R$ 上严格凸，从而对任意具有有限期望的随机变量 $X$ 恒有 $\E[\ln(1 + e^X)] \ge \ln(1 + e^{\E[X]})$；
    \item 对任意正实数 $a_1, \dots, a_k > 0$ 与非负权重 $w_1, \dots, w_k \ge 0$（满足 $\sum_{i=1}^k w_i = 1$），构造一个离散随机变量 $X$ 并对严格凹函数 $\ln x$ 应用 Jensen 不等式，一步证明\textbf{加权算术-几何平均不等式（AM--GM）}：
    \[
      \prod_{i=1}^k a_i^{w_i} \le \sum_{i=1}^k w_i a_i,
    \]
    并说明等号成立的充要条件。
  \end{enumerate}
  \exHint{第 (1) 问计算 $f'(x) = \sigma(x) = \frac{1}{1+e^{-x}}$ 与 $f''(x) = \sigma(x)(1-\sigma(x)) > 0$；第 (2) 问令 $\Prob(X = a_i) = w_i$。}

  \item \textbf{（重要性采样的理论最优提议分布与“零方差悖论”）} 设 $f(x) > 0$ 恒正，我们希望用重要性采样估计 $I = \E_{x\sim p}[f(x)] > 0$。
  \begin{enumerate}[label=(\arabic*), leftmargin=1.8em]
    \item 证明单样本重要性估计量 $Y = f(X)\frac{p(X)}{q(X)}$（$X \sim q$）的方差为 $\Var_q(Y) = \E_q\!\bigl[\frac{f(X)^2 p(X)^2}{q(X)^2}\bigr] - I^2$；
    \item 对任意合法提议分布 $q$，用 Jensen 不等式证明 $\E_q\!\bigl[\bigl(f(X)\frac{p(X)}{q(X)}\bigr)^2\bigr] \ge I^2$，并直接验证当且仅当选取
    \[
      q^*(x) = \frac{f(x) p(x)}{I} = \frac{f(x) p(x)}{\int f(t) p(t)\,dt}
    \]
    时，方差达到绝对下界 $\Var_{q^*}(Y) = 0$（只需 $M = 1$ 个样本就能零误差命中真值！）；
    \item 解释为什么 $q^*(x)$ 无法在真实计算中直接拿来采样，它又给工程设计提议分布 $q$ 带来了怎样的指导原则？
  \end{enumerate}
  \exHint{第 (2) 问中直接将 $q^*(x)$ 代入 $f(x)p(x)/q^*(x)$ 会发现结果恒等于常数 $I$；对一般 $q$ 直接应用 $\E_q[Y^2] \ge (\E_q[Y])^2 = I^2$。}
\end{enumerate}

\exTier{D. 编程实战与可自检实验}{用 NumPy 复现浓度衰减曲线、CLT 形态演化与权重塌缩诊断}

\begin{enumerate}[label=\textbf{3.\arabic*}, leftmargin=2.4em, resume]
  \item \textbf{（浓度不等式实测对比与右偏分布 CLT 演化实验）} 编写纯 NumPy 脚本（固定随机种子 \texttt{2026}）完成以下两项实验：
  \begin{enumerate}[label=(\arabic*), leftmargin=1.8em]
    \item 对 $n \in \{20, 60, 100, 200, 300\}$，各做 $20{,}000$ 次独立重复试验计算 $\mathrm{Bernoulli}(0.5)$ 样本均值 $\bar{X}_n$，打印每个 $n$ 下的实测尾概率 $\Prob(|\bar{X}_n - 0.5| \ge 0.1)$、Chebyshev 上界与 Hoeffding 上界，验证实测值恒不超过两者且 Hoeffding 最终显著紧于 Chebyshev；
    \item 对严重右偏的指数分布 $X_i \sim \mathrm{Exp}(1)$，分别取 $n \in \{1, 5, 30, 100\}$，各做 $20{,}000$ 次试验计算标准化统计量 $Z_n = \sqrt{n}(\bar{X}_n - 1)$，统计并输出每个 $n$ 下的样本偏度 $\frac{1}{T}\sum Z_n^3$ 与落入 $[-1.96, 1.96]$ 的经验覆盖率。
  \end{enumerate}
  {\rightskip=0pt plus 1em\relax \par\smallskip\noindent{\footnotesize\sffamily\color{tealDark}\textbf{【自检标准】}\,第 (1) 问在 $n=300$ 时，Chebyshev 上界须为 $0.08333$、Hoeffding 上界为 $0.00496$、实测尾概率应落在 $[0.0001, 0.0010]$ 之间；第 (2) 问随 $n$ 从 $1$ 增至 $100$，偏度应从约 $2.0$ 单调降至约 $0.20$（理论值 $2/\sqrt{n}$），$[-1.96, 1.96]$ 覆盖率应从约 $0.948$ 稳定收敛至 $0.950 \pm 0.005$。}\par}

  \item \textbf{（罕见尾事件重要性采样与权重塌缩 ESS 诊断）} 编写纯 NumPy 脚本完成以下两项实验：
  \begin{enumerate}[label=(\arabic*), leftmargin=1.8em]
    \item 以精确值 $p_{\mathrm{true}} = \Prob_{Z\sim\mathcal{N}(0,1)}(Z > 4) \approx 3.1671 \times 10^{-5}$ 为基准，在 $M = 20{,}000$ 采样预算下，重复 $50$ 次独立试验对比朴素蒙特卡洛与平移正态提议 $q = \mathcal{N}(4, 1)$ 重要性采样的均值、标准差与平均相对误差；
    \item 目标估计 $I = \E_{X\sim\mathcal{N}(0, 2^2)}[X^2] = 4.0$。取 $M = 10{,}000$ 重复 $50$ 次试验，分别使用“瘦尾提议” $q_1 = \mathcal{N}(0, 1)$ 与“厚尾提议” $q_2 = \mathcal{N}(0, 3^2)$ 计算重要性采样估计值、$50$ 次重复的经验标准差，以及平均相对有效样本量 $\mathrm{ESS}/M$。
  \end{enumerate}
  \par\smallskip\noindent{\footnotesize\sffamily\color{tealDark}\textbf{【自检标准】}\,第 (1) 问中，朴素 MC 约有一半轮次输出 $0$（标准差与真值同量级、超过 $3\times 10^{-5}$），而 $q=\mathcal{N}(4,1)$ 重要性采样在 $50$ 轮中的均值应在 $(3.167 \pm 0.05)\times 10^{-5}$ 内且单轮标准差 $< 7\times 10^{-7}$（平均相对误差 $< 2\%$）；第 (2) 问中，厚尾提议 $q_2 = \mathcal{N}(0, 3^2)$ 须稳定估计 $4.00 \pm 0.05$、经验标准差 $< 0.05$ 且 $\mathrm{ESS}/M > 0.75$，而瘦尾提议 $q_1 = \mathcal{N}(0, 1)$ 会出现剧烈波动（经验标准差 $> 0.5$、$\mathrm{ESS}/M < 0.25$），直观印证习题 3.2 的无穷方差陷阱。}
\end{enumerate}

\chapter[信息论：熵、交叉熵与 KL 散度]{信息论核心：熵、交叉熵、\\KL 散度与互信息}
\label{ch:ch4}

\chapEpigraph{通信的基本问题，是在一点精确或近似地重现另一点所选择的消息。}{克劳德·香农（Claude E. Shannon），《通信的数学理论》（1948）}

\begin{chapIntro}{本章定位：现代人工智能损失函数的通用语法}
1948 年，Claude Shannon 为了解决电话线上的信号压缩与噪声传输问题，创立了\textbf{信息论（Information Theory）}。谁也没有想到，七十余年后，这套原本属于通信工程的数学语言，竟成了整个人工智能领域的\textbf{通用目标函数语法}：
\begin{itemize}
  \item 图像分类与大语言模型训练的\textbf{交叉熵损失（Cross-Entropy）}，衡量真实标签与预测分布的平均编码代价；
  \item 大模型评测的\textbf{困惑度（Perplexity）}，本质是交叉熵的指数化“等效词表分支数”；
  \item 第 4 卷\textbf{决策树的信息增益}，正是本章的\textbf{互信息（Mutual Information）}；
  \item 第 5 卷\textbf{变分自编码器（VAE）}与第 6 卷\textbf{大模型对齐（RLHF/PPO/DPO）}，核心正则项都是\textbf{KL 散度（KL Divergence）}，而正向 KL 与反向 KL 的不对称性，直接决定了模型是“广泛覆盖”还是“精准锁峰”。
\end{itemize}
本章将承接第 3 章的 Jensen 不等式，从零构建“自信息 $\to$ 熵 $\to$ 互信息 $\to$ KL 散度 $\to$ 交叉熵与困惑度”的完整严密链条，让你再也不必死记任何一条损失函数公式。
\end{chapIntro}

% =========================================================================
\section{自信息、Shannon 熵与最大熵直觉}
\label{sec:ch4-entropy}

\subsection{自信息（Surprise）：为什么非用 $-\log p(x)$ 不可？}

先问一个直觉问题：如何量化“得知一个随机事件发生后，我们获得了多少信息（或者感到多意外）”？对比三条新闻：
\begin{enumerate}
  \item “明天太阳从东边升起”（$p \approx 1$）：毫无意外，几乎不含新信息；
  \item “均匀硬币掷出正面”（$p = 0.5$）：消除了二选一的不确定性，记作 $1\;\text{bit}$ 信息；
  \item “你买的单注彩票中了一等奖”（$p = 10^{-6}$）：极其罕见，包含巨大的惊讶度与信息量。
\end{enumerate}

若希望信息量函数 $I(p)$ 同时满足三条朴素公理：
\begin{enumerate}
  \item \textbf{必然事件无信息}：$I(1) = 0$；
  \item \textbf{越罕见信息越多}：$I(p)$ 关于 $p \in (0, 1]$ 严格单调递减；
  \item \textbf{独立事件信息可加}：若独立事件 $A, B$ 同时发生（概率 $p_A p_B$），总信息量应等于两者之和，即 $I(p_A p_B) = I(p_A) + I(p_B)$。
\end{enumerate}
微积分中满足乘积变加法且连续单调的函数，在差一个正常数倍（对数底数）的意义下\textbf{唯一}就是负对数函数！因此定义事件 $x$ 的\textbf{自信息（Self-Information）}为
\begin{equation}
  I(x) := -\log P(x). \label{eq:ch4-self-info}
\end{equation}
\textbf{对数底数与单位约定}：
\begin{itemize}
  \item 取以 $2$ 为底 $\log_2$ 时，单位为 \textbf{bit（比特）}：每次二分查询最多提供 $1\;\text{bit}$ 信息，适合信息编码、决策树手算与直觉理解；
  \item 取自然对数 $\ln$ 时，单位为 \textbf{nat（奈特）}：$1\;\text{nat} = \frac{1}{\ln 2} \approx 1.4427\;\text{bits}$。由于 $\frac{d}{dp}\ln p = \frac{1}{p}$ 不带多余的 $\ln 2$ 因子，现代深度学习框架（PyTorch/JAX）与理论推导默认使用 $\ln$。本章中若公式写 $\log$ 而未标底数，代表对 $\log_2$ 与 $\ln$ 同时成立（只要全式底数统一）。
\end{itemize}

\subsection{Shannon 熵与二元熵函数 $H_2(p)$}

自信息 $I(X) = -\log P(X)$ 度量的是\textbf{单个结果}的惊讶度；对整个分布取期望，就得到了衡量整个随机变量“平均不确定性”的 \textbf{Shannon 熵}。

\begin{defBox}
\begin{definition}[Shannon 熵, Shannon Entropy]
\label{def:ch4-entropy}
设离散随机变量 $X$ 取值于有限集合 $\mathcal{X}$，概率分布为 $P$。定义 $X$（或分布 $P$）的 \textbf{Shannon 熵} 为自信息的期望：
\begin{equation}
  H(X) = H(P) := \E_{X \sim P}\bigl[-\log P(X)\bigr] = -\sum_{x \in \mathcal{X}} P(x)\log P(x). \label{eq:ch4-entropy}
\end{equation}
当某个取值满足 $P(x) = 0$ 时，依据极限 $\lim_{t \to 0^+} t\log t = 0$，约定 $0\log 0 := 0$（不可能发生的事件不贡献不确定性）。
\end{definition}
\end{defBox}

\textbf{最直观的解剖标本：二元熵函数 $H_2(p)$}。对两点分布 $X \sim \mathrm{Bernoulli}(p)$，其以比特为单位的熵称为\textbf{二元熵}：
\[
  H_2(p) := -p\log_2 p - (1-p)\log_2(1-p), \qquad p \in [0, 1].
\]
代入几个代表点立刻见真章：
\begin{itemize}
  \item $p = 0$ 或 $p = 1$（必出反面或必出正面，毫无悬念）：$H_2(0) = H_2(1) = 0\;\text{bit}$；
  \item $p = 0.1$（高度偏向一侧）：$H_2(0.1) = -0.1\log_2 0.1 - 0.9\log_2 0.9 \approx 0.4690\;\text{bit}$；
  \item $p = 0.5$（完全对半开、最难猜）：$H_2(0.5) = -0.5(-1) - 0.5(-1) = 1.0000\;\text{bit}$。
\end{itemize}
图 \ref{fig:ch4-binary-entropy} 绘制了 $H_2(p)$ 的严格上凸（凹函数）拱形曲线：两端陡峭趋零，中央平缓封顶。

\begin{figure}[htbp]
  \centering
  \begin{tikzpicture}
    \begin{axis}[
      width=0.76\linewidth,
      height=5.4cm,
      xlabel={正类概率 $p = \Prob(X = 1)$},
      ylabel={二元熵 $H_2(p)$（单位：bit）},
      xmin=-0.03, xmax=1.03,
      ymin=-0.04, ymax=1.18,
      xtick={0, 0.1, 0.25, 0.5, 0.75, 0.9, 1.0},
      ytick={0, 0.25, 0.5, 0.75, 1.0},
      grid=major,
      major grid style={slate!18},
      label style={font=\small\sffamily, ink},
      tick label style={font=\footnotesize, slate}
    ]
      % Shaded area under curve
      \addplot[fill=tealLight, draw=none, domain=0.001:0.999, samples=160]
        {-x*ln(x)/ln(2) - (1-x)*ln(1-x)/ln(2)} \closedcycle;

      % Main binary entropy curve
      \addplot[teal, line width=1.6pt, domain=0.001:0.999, samples=160]
        {-x*ln(x)/ln(2) - (1-x)*ln(1-x)/ln(2)};

      % Vertical dashed line at p=0.5
      \draw[ochre, dashed, thick] (axis cs:0.5, 0) -- (axis cs:0.5, 1.0);
      \draw[slate, dotted, thick] (axis cs:0, 1.0) -- (axis cs:0.5, 1.0);

      % Highlighted points
      \fill[ochre] (axis cs:0.5, 1.0) circle (2.6pt);
      \fill[tealDark] (axis cs:0, 0) circle (2.3pt);
      \fill[tealDark] (axis cs:1, 0) circle (2.3pt);
      \fill[navy] (axis cs:0.1, 0.4690) circle (2.3pt);

      % Annotations (well spaced, zero overlap)
      \node[ochre, font=\footnotesize\sffamily\bfseries, anchor=south]
        at (axis cs:0.5, 1.03) {最大不确定性：$(0.5,\; 1.0\;\text{bit})$};
      \node[navy, font=\scriptsize\sffamily, anchor=west]
        at (axis cs:0.13, 0.43) {$(0.1,\; 0.469)$};
      \node[tealDark, font=\scriptsize\sffamily, anchor=south west]
        at (axis cs:0.01, 0.04) {确定性 $H_2=0$};
      \node[tealDark, font=\scriptsize\sffamily, anchor=south east]
        at (axis cs:0.99, 0.04) {确定性 $H_2=0$};
    \end{axis}
  \end{tikzpicture}
  \caption{二元熵函数 $H_2(p) = -p\log_2 p - (1-p)\log_2(1-p)$：在 $p = 0.5$ 处唯一达到最大值 $1\;\text{bit}$，向两端确定性状态 $p \to 0$ 或 $p \to 1$ 对称衰减至 $0$。}
  \label{fig:ch4-binary-entropy}
\end{figure}

\subsection{熵的上下界与最大熵直觉（由第 3 章 Jensen 一步证出）}

推广到 $K$ 个离散类别时，“最确定的分布熵为 $0$，最均匀的分布熵最大”依然严格成立，而且用第 3 章的 Jensen 不等式三行就能证毕。

\begin{thmBox}
\begin{theorem}[离散熵的上下界与最大熵原理]
\label{thm:ch4-entropy-bounds}
设随机变量 $X \sim P$ 取值于有限集合 $\mathcal{X}$（共有 $K = |\mathcal{X}|$ 个可能状态），则
\begin{equation}
  0 \le H(P) \le \log K. \label{eq:ch4-entropy-bounds}
\end{equation}
其中：左端等号成立当且仅当 $P$ 为\textbf{确定性分布}（存在唯一 $x_0$ 使 $P(x_0)=1$，即 one-hot 分布）；右端等号成立当且仅当 $P$ 为\textbf{均匀分布}（对所有 $x \in \mathcal{X}$ 有 $P(x) = 1/K$）。
\end{theorem}
\end{thmBox}

\begin{proof}
\textbf{下界}：对每个 $x \in \mathcal{X}$ 有 $0 \le P(x) \le 1$，故 $-P(x)\log P(x) \ge 0$，求和即得 $H(P) \ge 0$。等号成立当且仅当每项均为 $0$，即对每个 $x$ 要么 $P(x)=0$ 要么 $P(x)=1$；结合 $\sum_x P(x)=1$，必存在唯一 $x_0$ 使 $P(x_0)=1$。

\textbf{上界}：记正概率支集 $\mathcal{S} := \{x \in \mathcal{X} : P(x) > 0\}$，显然 $|\mathcal{S}| \le K$。将熵写成关于正随机变量 $1/P(X)$ 的期望，并对严格凹函数 $\log(\cdot)$ 应用第 3 章定理 \ref{thm:ch3-jensen}（Jensen 不等式）：
\begin{align*}
  H(P)
  &= \sum_{x \in \mathcal{S}} P(x)\log\frac{1}{P(x)}
   = \E_{P}\!\left[\log\frac{1}{P(X)}\right]
   \le \log\!\left(\E_P\!\left[\frac{1}{P(X)}\right]\right) \\
  &= \log\!\left(\sum_{x \in \mathcal{S}} P(x)\cdot \frac{1}{P(x)}\right)
   = \log |\mathcal{S}|
   \le \log K.
\end{align*}
由于 $\log$ 严格凹，第一个等号成立当且仅当 $1/P(X)$ 在 $\mathcal{S}$ 上几乎必然为常数（即在 $\mathcal{S}$ 上均匀分布）；第二个等号成立当且仅当 $|\mathcal{S}| = K$。两者合起来恰好等价于 $P(x) = 1/K$（$\forall x \in \mathcal{X}$）。
\end{proof}

这一结论揭示了人工智能中极富哲学意味的\textbf{最大熵原理（Principle of Maximum Entropy）}：在没有任何额外先验知识时，最诚实、不乱做假设的概率分布就是熵最大的均匀分布；而在第 6 章的 Softmax 回归里，如果对模型预测额外加上熵正则（鼓励适度保留不确定性，如标签平滑 Label Smoothing），就能有效防止神经网络变得过度自信而过拟合。

\begin{extBox}{连续随机变量的微分熵（Differential Entropy）与一个重要提醒}
对连续概率密度 $p(x)$，类比定义\textbf{微分熵} $h(X) := -\int p(x)\ln p(x)\,dx$。例如一维高斯 $\mathcal{N}(\mu, \sigma^2)$ 的微分熵为 $h(\mathcal{N}(\mu,\sigma^2)) = \frac{1}{2}\ln(2\pi e \sigma^2)$（方差越大越不确定，且在给定方差的所有连续分布中，\textbf{高斯分布微分熵最大}！）。但务必注意：与离散熵 $H(X) \ge 0$ 不同，\textbf{微分熵可以为负数}！例如 $X \sim \mathrm{Uniform}(0, 0.5)$ 的密度恒为 $p(x)=2$，其微分熵为 $h(X) = -\int_0^{0.5} 2\ln 2\,dx = -\ln 2 < 0$；而且做尺度缩放 $Y = cX$ 会令 $h(cX) = h(X) + \ln|c|$。幸运的是，下文将要定义的\textbf{连续 KL 散度与连续互信息}是两个密度之比的对数期望，缩放因子恰好完全抵消，因此永远非负且坐标无关！
\end{extBox}

% =========================================================================
\section{联合熵、条件熵与互信息（决策树信息增益的真身）}
\label{sec:ch4-mi}

现实 AI 任务从来不是孤立地看单个变量 $Y$，而是\textbf{利用观测特征 $X$ 去消除目标标签 $Y$ 的不确定性}。下面把熵的概念自然推广到两个变量 $(X, Y)$。

\subsection{联合熵、条件熵与熵的链式法则}

\begin{defBox}
\begin{definition}[联合熵与条件熵]
\label{def:ch4-joint-cond-entropy}
设离散随机变量对 $(X, Y)$ 具有联合分布 $P(x, y)$，边缘分布 $P(x), P(y)$ 与条件分布 $P(y\mid x)$。定义：
\begin{enumerate}
  \item \textbf{联合熵（Joint Entropy）}：
  \[
    H(X, Y) := -\sum_{x, y} P(x, y)\log P(x, y);
  \]
  \item \textbf{条件熵（Conditional Entropy）}：给定具体观测 $X = x$ 后 $Y$ 的剩余熵为 $H(Y \mid X = x) := -\sum_y P(y\mid x)\log P(y\mid x)$；再对 $X$ 取平均，即得条件熵：
  \begin{equation}
    H(Y \mid X) := \sum_{x} P(x)\, H(Y \mid X = x) = -\sum_{x, y} P(x, y)\log P(y \mid x). \label{eq:ch4-cond-entropy}
  \end{equation}
\end{enumerate}
\end{definition}
\end{defBox}

把第 1 章概率乘法公式 $P(x, y) = P(x)P(y \mid x)$ 取负对数再求期望，乘法瞬间拆成加法——这就是优美无比的\textbf{熵的链式法则}。

\begin{thmBox}
\begin{theorem}[熵的链式法则, Chain Rule for Entropy]
\label{thm:ch4-chain-rule}
对任意离散随机变量对 $(X, Y)$，恒有
\begin{equation}
  H(X, Y) = H(X) + H(Y \mid X) = H(Y) + H(X \mid Y). \label{eq:ch4-chain-rule}
\end{equation}
推广到 $n$ 个变量 $(X_1, \dots, X_n)$，有 $H(X_1, \dots, X_n) = \sum_{t=1}^n H(X_t \mid X_1, \dots, X_{t-1})$。
\end{theorem}
\end{thmBox}

\begin{proof}
由 $P(x, y) = P(x)P(y\mid x)$（对 $P(x,y)>0$ 的项）直接展开：
\begin{align*}
  H(X, Y)
  &= -\sum_{x,y} P(x,y)\bigl[\log P(x) + \log P(y\mid x)\bigr] \\
  &= -\sum_x \left(\sum_y P(x,y)\right)\log P(x) - \sum_{x,y} P(x,y)\log P(y\mid x) \\
  &= H(X) + H(Y \mid X).
\end{align*}
交换 $X$ 与 $Y$ 的角色即得第二式；对 $n$ 个变量归纳即得一般链式法则。
\end{proof}

\subsection{互信息（Mutual Information）与信息韦恩图}

知道了 $X$ 之后，$Y$ 的不确定性从 $H(Y)$ 下降到了 $H(Y\mid X)$。这个\textbf{不确定性的平均减少量}，就是特征 $X$ 所携带的关于标签 $Y$ 的信息量——\textbf{互信息}！

\begin{defBox}
\begin{definition}[互信息, Mutual Information]
\label{def:ch4-mi}
随机变量 $X$ 与 $Y$ 之间的\textbf{互信息}定义为
\begin{align}
  I(X; Y)
  &:= H(Y) - H(Y \mid X) \nonumber\\
  &= H(X) - H(X \mid Y)
   = H(X) + H(Y) - H(X, Y) \nonumber\\
  &= \sum_{x, y} P(x, y)\log\frac{P(x, y)}{P(x)P(y)}. \label{eq:ch4-mi}
\end{align}
\end{definition}
\end{defBox}

注意式 \eqref{eq:ch4-mi} 展现出的完美对称性：$I(X; Y) = I(Y; X)$（$X$ 告诉我们多少关于 $Y$ 的信息，恰好等于 $Y$ 告诉我们多少关于 $X$ 的信息）。图 \ref{fig:ch4-venn} 用经典的\textbf{信息韦恩图（I-Diagram）}把这五个量之间的关系一眼定格——只要记住这张图，链式法则与互信息的全部等价式永远不会背混。

\begin{figure}[htbp]
  \centering
  \begin{tikzpicture}[scale=1.0]
    % Left circle H(X) and Right circle H(Y)
    \begin{scope}
      \fill[tealLight, opacity=0.88] (-1.35, 0) ellipse (2.35cm and 1.48cm);
      \fill[navyLight, opacity=0.82] ( 1.35, 0) ellipse (2.35cm and 1.48cm);
    \end{scope}

    % Highlight intersection I(X;Y)
    \begin{scope}
      \clip (-1.35, 0) ellipse (2.35cm and 1.48cm);
      \fill[ochreLight, opacity=0.92] ( 1.35, 0) ellipse (2.35cm and 1.48cm);
    \end{scope}

    % Outlines
    \draw[teal, line width=1.4pt] (-1.35, 0) ellipse (2.35cm and 1.48cm);
    \draw[navy, line width=1.4pt] ( 1.35, 0) ellipse (2.35cm and 1.48cm);

    % Region labels
    \node[tealDark, font=\small\sffamily\bfseries, align=center] at (-2.25, 0.12) {$H(X \mid Y)$};
    \node[teal, font=\scriptsize\sffamily, align=center] at (-2.25, -0.35) {仅 $X$ 独有的不确定性};

    \node[ochre, font=\small\sffamily\bfseries, align=center] at (0, 0.15) {$I(X; Y)$};
    \node[ochre, font=\scriptsize\sffamily, align=center] at (0, -0.32) {共享信息\\（信息增益）};

    \node[navy, font=\small\sffamily\bfseries, align=center] at (2.25, 0.12) {$H(Y \mid X)$};
    \node[navy, font=\scriptsize\sffamily, align=center] at (2.25, -0.35) {已知 $X$ 后 $Y$ 剩余熵};

    % Top braces / labels for H(X) and H(Y)
    \draw[decorate, decoration={brace, amplitude=5pt}, teal, thick]
      (-3.68, 1.66) -- (0.95, 1.66)
      node[midway, above=5pt, font=\small\sffamily\bfseries, tealDark] {$H(X) = H(X\mid Y) + I(X;Y)$};

    \draw[decorate, decoration={brace, amplitude=5pt}, navy, thick]
      (-0.95, -1.66) -- (3.68, -1.66);
    \node[navy, font=\small\sffamily\bfseries] at (1.36, -2.18) {$H(Y) = H(Y\mid X) + I(X;Y)$};

    % Bottom total brace for H(X,Y)
    \draw[decorate, decoration={brace, amplitude=6pt, mirror}, slate, thick]
      (-3.70, -2.55) -- (3.70, -2.55)
      node[midway, below=6pt, font=\small\sffamily\bfseries, ink]
      {联合熵 $H(X, Y) = H(X\mid Y) + I(X; Y) + H(Y\mid X)$};
  \end{tikzpicture}
  \caption{信息论韦恩图（I-Diagram）：两个随机变量的熵、条件熵、联合熵与互信息满足直观的“集合加减”几何关系。}
  \label{fig:ch4-venn}
\end{figure}

\begin{thmBox}
\begin{theorem}[互信息非负性与“条件化平均不增熵”]
\label{thm:ch4-mi-nonneg}
对任意随机变量 $X, Y$，恒有
\begin{equation}
  I(X; Y) \ge 0 \qquad \Longleftrightarrow \qquad H(Y \mid X) \le H(Y), \label{eq:ch4-mi-nonneg}
\end{equation}
且等号成立当且仅当 $X$ 与 $Y$ 相互独立（$X \indep Y$）。
\end{theorem}
\end{thmBox}

\begin{proof}
记正概率支集 $\mathcal{S} = \{(x, y) : P(x, y) > 0\}$。在 $\mathcal{S}$ 上必有 $P(x) > 0$ 且 $P(y) > 0$。以自然对数计算（换底仅差正因子 $\frac{1}{\ln 2}$），对严格凸函数 $-\ln(\cdot)$ 应用第 3 章定理 \ref{thm:ch3-jensen}（Jensen 不等式）：
\begin{align*}
  I(X; Y)
  &= \sum_{(x,y)\in\mathcal{S}} P(x,y)\left(-\ln\frac{P(x)P(y)}{P(x,y)}\right)
   \ge -\ln\!\left(\sum_{(x,y)\in\mathcal{S}} P(x,y)\cdot\frac{P(x)P(y)}{P(x,y)}\right) \\
  &= -\ln\!\left(\sum_{(x,y)\in\mathcal{S}} P(x)P(y)\right)
   \ge -\ln 1 = 0.
\end{align*}
等号成立当且仅当两处不等号同时取等：(i) $\frac{P(x)P(y)}{P(x,y)} = c$ 在 $\mathcal{S}$ 上为常数；(ii) $\sum_{\mathcal{S}} P(x)P(y) = 1$（即所有 $P(x)P(y)>0$ 的点都在 $\mathcal{S}$ 内）。将 $P(x)P(y) = c P(x,y)$ 两边求和得 $1 = c \cdot 1 \implies c = 1$，故 $P(x,y) = P(x)P(y)$ 处处成立，即 $X \indep Y$。
\end{proof}

\begin{example}[垃圾邮件关键词过滤的信息增益手算]
\label{ex:ch4-spam-mi}
设 $Y \in \{0, 1\}$ 表示邮件是否为垃圾邮件，$X \in \{0, 1\}$ 表示标题是否包含关键词“限时免费”，两者联合分布为：$P(0,0) = 0.40,\; P(0,1) = 0.10,\; P(1,0) = 0.10,\; P(1,1) = 0.40$。
\begin{itemize}
  \item 边缘分布：$P(X=1) = 0.50,\; P(Y=1) = 0.50$，因此 $H(X) = H(Y) = H_2(0.5) = 1.0000\;\text{bit}$；
  \item 给定 $X=1$ 时，$P(Y=1\mid X=1) = \frac{0.40}{0.50} = 0.80$，故 $H(Y\mid X=1) = H_2(0.8) \approx 0.7219\;\text{bit}$；同理 $H(Y\mid X=0) = H_2(0.2) \approx 0.7219\;\text{bit}$；
  \item 条件熵 $H(Y\mid X) = 0.5 \times 0.7219 + 0.5 \times 0.7219 = 0.7219\;\text{bit}$；
  \item 互信息（信息增益）：$I(X; Y) = H(Y) - H(Y\mid X) = 1.0000 - 0.7219 = 0.2781\;\text{bit} > 0$。
\end{itemize}
仅观察这一个关键词，就把判断垃圾邮件所需的平均不确定性减少了 $27.8\%$！
\end{example}

\begin{failBox}{反直觉陷阱：“条件化平均不增熵”绝不意味着“每次观测都不增熵”！}
定理 \ref{thm:ch4-mi-nonneg} 证明的是对 $X$ \textbf{取期望后}有 $H(Y\mid X) \le H(Y)$，但对某一个\textbf{具体取值} $X = x_0$，完全可能出现 $H(Y \mid X = x_0) > H(Y)$！

\textbf{极简反例}：设 $P(X=0, Y=0) = 0.8,\; P(X=1, Y=0) = 0.1,\; P(X=1, Y=1) = 0.1$（其余为 $0$）。
\begin{itemize}
  \item 不看 $X$ 时，$P(Y=1) = 0.1$，不确定性较小：$H(Y) = H_2(0.1) \approx 0.4690\;\text{bit}$；
  \item 一旦观察到 $X = 1$，$P(Y=1 \mid X=1) = 0.5$，瞬间变成五五开，不确定性\textbf{暴涨}：$H(Y \mid X=1) = H_2(0.5) = 1.0000\;\text{bit} > H(Y)$！
  \item 但取期望后，$H(Y \mid X) = 0.8 \cdot H(Y\mid X=0) + 0.2 \cdot H(Y\mid X=1) = 0.8\times 0 + 0.2\times 1 = 0.2000\;\text{bit} < 0.4690\;\text{bit}$，定理依然稳如泰山！
\end{itemize}
直觉就像医学初筛：没做检查前你知道某罕见病大概率没得（熵很小），一旦初筛亮黄灯（$X=1$），得与不得瞬间变成五五开（熵反而变大）；但平均而言，做初筛总能减少不确定性。
\end{failBox}

\begin{insightBox}{直连第 4 卷决策树信息增益与现代对比学习}
在第 4 卷经典机器学习中，ID3/C4.5 决策树在每个节点挑选“最优分裂特征 $A$”所用的\textbf{信息增益（Information Gain）}
\[
  \mathrm{IG}(Y, A) := H(Y) - H(Y \mid A) = I(A; Y),
\]
分毫不差就是样本分布上的\textbf{互信息}！此外，互信息超越了传统皮尔逊相关系数“只能捕捉线性相关”的局限，能度量任意非线性依赖关系，因此也成为现代自监督表征学习（如 InfoNCE 对比学习最大化互信息下界）的核心理论框架。
\end{insightBox}

% =========================================================================
\section{KL 散度（相对熵）与正/反向 KL 的本质差异}
\label{sec:ch4-kl}

机器学习的核心任务，本质上是用参数化模型分布 $Q_\theta$ 去逼近未知的真实数据分布 $P$。那么，该如何定量衡量“两个概率分布 $P$ 和 $Q$ 到底差了多远”？答案是整个现代 AI 出现频率最高的“分布差距标尺”——\textbf{KL 散度（Kullback--Leibler Divergence）}。

\subsection{定义与 Gibbs 不等式（完整严谨证明）}

\begin{defBox}
\begin{definition}[KL 散度 / 相对熵, Kullback--Leibler Divergence]
\label{def:ch4-kl}
设 $P$ 与 $Q$ 是定义在同一空间上的两个概率分布。离散情形与连续情形下，从 $Q$ 到 $P$ 的 \textbf{KL 散度} 分别定义为：
\begin{equation}
  \DKL(P \,\|\, Q) := \sum_{x \in \mathcal{X}} P(x)\log\frac{P(x)}{Q(x)},
  \qquad
  \DKL(p \,\|\, q) := \int p(x)\log\frac{p(x)}{q(x)}\,dx. \label{eq:ch4-kl}
\end{equation}
按连续性约定 $0\log\frac{0}{q} := 0$；若存在某点使 $P(x) > 0$ 而 $Q(x) = 0$，则约定 $P(x)\log\frac{P(x)}{0} := +\infty$。因此 $\DKL(P\|Q)$ 有限的必要条件是\textbf{支集覆盖}：$\mathrm{supp}(P) \subseteq \mathrm{supp}(Q)$。
\end{definition}
\end{defBox}

\textbf{编码意义}：若真实事件服从分布 $P$，最优编码平均长度是 $H(P) = \E_P[-\log P(X)]$；如果你误用了针对分布 $Q$ 设计的编码方案（给结果 $x$ 分配长度 $-\log Q(x)$），平均每个符号多浪费的额外比特/奈特数，恰好就是 $\E_P[-\log Q(X) - (-\log P(X))] = \DKL(P \| Q)$！

既然是“多浪费的代价”，直觉告诉我们它绝不可能为负，且只有当 $Q = P$ 时才能做到零浪费。下面我们用第 3 章的 Jensen 不等式给出严格证明——这就是大名鼎鼎的 \textbf{Gibbs 不等式}。

\begin{thmBox}
\begin{theorem}[Gibbs 不等式：KL 散度非负性]
\label{thm:ch4-gibbs}
对任意两个概率分布 $P$ 与 $Q$，恒有
\begin{equation}
  \DKL(P \,\|\, Q) \ge 0, \label{eq:ch4-gibbs}
\end{equation}
且等号成立当且仅当 $P = Q$（连续情形下为几乎处处 $p(x) = q(x)$）。
\end{theorem}
\end{thmBox}

\begin{proof}
（以下按自然对数 $\ln$ 证明，换用 $\log_2$ 仅差正因子 $\frac{1}{\ln 2}$。）记 $P$ 的正支集 $\mathcal{S} := \{x : P(x) > 0\}$。若存在 $x \in \mathcal{S}$ 满足 $Q(x) = 0$，按定义 $\DKL(P\|Q) = +\infty > 0$，结论平凡成立。现设对一切 $x \in \mathcal{S}$ 均有 $Q(x) > 0$。将 KL 散度改写为关于正随机变量 $Z = \frac{Q(X)}{P(X)}$ 的期望，再对严格凸函数 $g(z) = -\ln z$ 应用第 3 章定理 \ref{thm:ch3-jensen}（Jensen 不等式）：
\begin{align*}
  \DKL(P \,\|\, Q)
  &= \sum_{x \in \mathcal{S}} P(x)\left(-\ln\frac{Q(x)}{P(x)}\right)
   = \E_{X \sim P}\!\left[-\ln\frac{Q(X)}{P(X)}\right] \\
  &\ge -\ln\!\left(\E_{X \sim P}\!\left[\frac{Q(X)}{P(X)}\right]\right)
   = -\ln\!\left(\sum_{x \in \mathcal{S}} P(x)\cdot\frac{Q(x)}{P(x)}\right) \\
  &= -\ln\!\left(\sum_{x \in \mathcal{S}} Q(x)\right)
   \ge -\ln 1 = 0.
\end{align*}
最后一个不等号利用了 $\sum_{x \in \mathcal{S}} Q(x) \le \sum_{x \in \mathcal{X}} Q(x) = 1$ 与 $-\ln(\cdot)$ 单调递减。

\textbf{等号充要条件}：$\DKL(P\|Q) = 0$ 当且仅当上述两处不等号同时取等：(1) 由 $-\ln$ 严格凸，Jensen 取等当且仅当在 $\mathcal{S}$ 上 $\frac{Q(x)}{P(x)} \equiv c$（某个正常数）；(2) 第二处取等当且仅当 $\sum_{x \in \mathcal{S}} Q(x) = 1$（即 $Q$ 在 $\mathcal{S}$ 之外处处为 $0$）。将 $Q(x) = c P(x)$ 在 $\mathcal{S}$ 上求和得 $1 = \sum_{x\in\mathcal{S}} Q(x) = c\sum_{x\in\mathcal{S}} P(x) = c$，于是 $c = 1$，从而对所有 $x \in \mathcal{X}$ 有 $Q(x) = P(x)$。（连续情形将求和改为积分，证明逐字相同。）
\end{proof}

有了 KL 散度，前两节的核心结论立刻实现了大一统：
\begin{itemize}
  \item \textbf{互信息就是联合分布与独立乘积分布的 KL 散度}：$I(X; Y) = \DKL\bigl(P_{X,Y} \,\|\, P_X P_Y\bigr) \ge 0$；
  \item \textbf{最大熵原理的 KL 版本}：设 $U_K$ 是 $K$ 点均匀分布，则 $\DKL(P \| U_K) = \log K - H(P) \ge 0 \implies H(P) \le \log K$！
\end{itemize}

\subsection{一维高斯 KL 散度解析闭式（第 5 卷 VAE 的核心零件）}

在深度生成模型（尤其是第 5 卷变分自编码器 VAE）中，编码器输出高斯后验 $q_\phi(z\mid x) = \mathcal{N}(\mu, \sigma^2)$，正则项需要把每个隐维度拉向标准高斯先验 $\mathcal{N}(0, 1)$，因此下面的解析闭式公式必须信手拈来。

\begin{proposition}[两个一维高斯分布之间的 KL 散度闭式]
\label{prop:ch4-gaussian-kl}
设 $P = \mathcal{N}(\mu_1, \sigma_1^2)$ 与 $Q = \mathcal{N}(\mu_2, \sigma_2^2)$（$\sigma_1, \sigma_2 > 0$），则以自然对数计，
\begin{equation}
  \DKL\!\bigl(\mathcal{N}(\mu_1, \sigma_1^2) \,\|\, \mathcal{N}(\mu_2, \sigma_2^2)\bigr)
  = \ln\frac{\sigma_2}{\sigma_1} + \frac{\sigma_1^2 + (\mu_1 - \mu_2)^2}{2\sigma_2^2} - \frac{1}{2}. \label{eq:ch4-gaussian-kl}
\end{equation}
特别地，当 $Q = \mathcal{N}(0, 1)$ 时（VAE 损失函数标准形式）：
\begin{equation}
  \DKL\!\bigl(\mathcal{N}(\mu, \sigma^2) \,\|\, \mathcal{N}(0, 1)\bigr)
  = \frac{1}{2}\bigl(\mu^2 + \sigma^2 - 1 - \ln \sigma^2\bigr). \label{eq:ch4-vae-kl}
\end{equation}
\end{proposition}

\begin{proof}
直接写出对数密度比：
\[
  \ln\frac{p(x)}{q(x)}
  = \ln\frac{\sigma_2}{\sigma_1} - \frac{(x - \mu_1)^2}{2\sigma_1^2} + \frac{(x - \mu_2)^2}{2\sigma_2^2}.
\]
对 $X \sim P = \mathcal{N}(\mu_1, \sigma_1^2)$ 取期望，利用 $\E_P[(X - \mu_1)^2] = \sigma_1^2$ 以及
\[
  \E_P\bigl[(X - \mu_2)^2\bigr]
  = \E_P\bigl[((X - \mu_1) + (\mu_1 - \mu_2))^2\bigr]
  = \sigma_1^2 + (\mu_1 - \mu_2)^2,
\]
代入即得
\[
  \DKL(P \| Q)
  = \ln\frac{\sigma_2}{\sigma_1} - \frac{\sigma_1^2}{2\sigma_1^2} + \frac{\sigma_1^2 + (\mu_1 - \mu_2)^2}{2\sigma_2^2}
  = \ln\frac{\sigma_2}{\sigma_1} + \frac{\sigma_1^2 + (\mu_1 - \mu_2)^2}{2\sigma_2^2} - \frac{1}{2}.
\]
令 $\mu_1 = \mu,\; \sigma_1 = \sigma,\; \mu_2 = 0,\; \sigma_2 = 1$，立刻化简为 \eqref{eq:ch4-vae-kl}。
\end{proof}

\subsection{非对称性与正/反向 KL 的本质分野：Mode-Covering vs Mode-Seeking}

仔细看一眼 \eqref{eq:ch4-gaussian-kl} 就会发现：交换 $(\mu_1,\sigma_1)$ 与 $(\mu_2,\sigma_2)$ 后结果并不相等！事实上，KL 散度既\textbf{不满足对称性}（一般 $\DKL(P\|Q) \neq \DKL(Q\|P)$），也\textbf{不满足三角不等式}，因此它在数学上不是“距离度量（Metric）”，而称为“散度（Divergence）”。

\textbf{千万别以为不对称是个缺陷——它恰恰是现代 AI 算法设计里最具威力的一道分水岭！}

设想一个极有代表性的建模冲突：真实数据分布 $P(x)$ 是一个\textbf{双峰高斯混合}：
\[
  P(x) = 0.5\,\mathcal{N}(-2,\, 0.6^2) + 0.5\,\mathcal{N}(2,\, 0.6^2),
\]
而我们的近似模型 $Q_\theta(x) = \mathcal{N}(\mu, \sigma^2)$ 只有\textbf{单个高斯}的表达能力（欠参数化，不可能完美拟合两个峰）。现在分别按两种方向最小化 KL 散度，结果截然相反！

\begin{enumerate}
  \item \textbf{正向 KL（Forward KL / 覆盖模式，Mode-Covering）}：求解 $\min_Q \DKL(P \| Q) = \E_{x\sim P}\!\bigl[\ln\frac{P(x)}{Q(x)}\bigr]$。
  观察被积函数：只要某个区域 $P(x) > 0$，一旦你敢把 $Q(x)$ 设成接近 $0$，惩罚项 $P(x)\ln\frac{P(x)}{Q(x)}$ 就会瞬间冲向 $+\infty$！因此正向 KL 强迫 $Q$ 必须“\textbf{雨露均沾}（Zero-Avoiding）”，把身躯拉得足够胖以同时罩住 $P$ 的左右两座山峰。事实上（见习题 4.5），单高斯极小化正向 KL 的精确最优解正是\textbf{矩匹配（Moment Matching）}：
  \[
    \mu_{\mathrm{fwd}} = \E_P[X] = 0,
    \qquad
    \sigma_{\mathrm{fwd}}^2 = \Var_P(X) = 0.6^2 + 2^2 = 4.36 \implies \sigma_{\mathrm{fwd}} \approx 2.09.
  \]
  代价是：$q_{\mathrm{fwd}}$ 把最高密度堆在了真实分布几乎没数据的“低谷” $x = 0$ 处！

  \item \textbf{反向 KL（Reverse KL / 寻峰模式，Mode-Seeking）}：求解 $\min_Q \DKL(Q \| P) = \E_{x\sim Q}\!\bigl[\ln\frac{Q(x)}{P(x)}\bigr]$。
  此刻外层期望换成了 $Q(x)$！凡是 $P(x) \approx 0$ 的无人区（比如中央低谷 $x = 0$），只要 $Q(x) > 0$ 就会遭遇巨额惩罚（Zero-Forcing）；相反，在 $Q(x) \approx 0$ 的地方，即使 $P(x)$ 很高也完全不扣分！因此反向 KL 宁可彻底放弃另一座山峰，也要死死锁住其中一座高峰：
  \[
    q_{\mathrm{rev}} = \mathcal{N}(2.0,\; 0.60^2) \qquad (\text{或对称地 }\mathcal{N}(-2.0,\; 0.60^2)).
  \]
  它生成的样本个个逼真（绝不落在低谷），代价则是丢失了另一半数据模态（\textbf{模式坍缩, Mode Collapse}）！
\end{enumerate}

图 \ref{fig:ch4-kl-bimodal} 直接读取本卷实验脚本生成的真实曲线数据 \texttt{data/kl\_bimodal.csv}，表 \ref{tab:ch4-kl-bimodal} 列出两组解在正/反向 KL 下的精确数值对比——请把这幅图像刻进你的 AI 直觉库里！

\begin{figure}[htbp]
  \centering
  \begin{tikzpicture}
    \begin{axis}[
      width=0.88\linewidth,
      height=6.2cm,
      xlabel={$x$},
      ylabel={概率密度},
      xmin=-5.8, xmax=5.8,
      ymin=0, ymax=0.76,
      grid=major,
      major grid style={slate!18},
      legend pos=north west,
      legend style={font=\footnotesize\sffamily, fill=white, fill opacity=0.92, draw=codeborder},
      label style={font=\small\sffamily, ink},
      tick label style={font=\footnotesize, slate}
    ]
      % Shaded target P(x)
      \addplot[fill=slate!14, draw=none, forget plot]
        table[x=x, y=p, col sep=comma] {data/kl_bimodal.csv} \closedcycle;

      % Target bimodal P(x)
      \addplot[ink, very thick]
        table[x=x, y=p, col sep=comma] {data/kl_bimodal.csv};
      \addlegendentry{真实双峰 $P = \tfrac{1}{2}\mathcal{N}(-2, 0.6^2) + \tfrac{1}{2}\mathcal{N}(2, 0.6^2)$}

      % Forward KL (mode-covering)
      \addplot[teal, very thick, dashed]
        table[x=x, y=qfwd, col sep=comma] {data/kl_bimodal.csv};
      \addlegendentry{正向 KL 最优 $q_{\mathrm{fwd}} = \mathcal{N}(0, 2.09^2)$（覆盖双峰）}

      % Reverse KL (mode-seeking)
      \addplot[ochre, very thick]
        table[x=x, y=qrev, col sep=comma] {data/kl_bimodal.csv};
      \addlegendentry{反向 KL 最优 $q_{\mathrm{rev}} = \mathcal{N}(2, 0.60^2)$（单峰锁定）}

      % Clean annotations
      \node[tealDark, font=\scriptsize\sffamily, align=center] at (axis cs:0.0, 0.255) {低谷处误堆密度\\（Mode-Covering）};
      \node[ochre, font=\scriptsize\sffamily, anchor=west, align=left] at (axis cs:2.55, 0.57) {高纯度锁定右峰\\忽略左峰（Mode-Seeking）};
    \end{axis}
  \end{tikzpicture}
  \caption{本机实测（\texttt{data/kl\_bimodal.csv}）：用单高斯近似双峰分布 $P(x)$ 时，正向 KL（$\min \DKL(P\|Q)$）给出宽平的矩匹配分布 $q_{\mathrm{fwd}}$，而反向 KL（$\min \DKL(Q\|P)$）锐利锁定单个模态 $q_{\mathrm{rev}}$。}
  \label{fig:ch4-kl-bimodal}
\end{figure}

\begin{table}[htbp]
  \centering
  \caption{双峰高斯近似实验中正向 KL 与反向 KL 数值对照（来自 \texttt{vol2\_experiments.py} 实测）}
  \label{tab:ch4-kl-bimodal}
  \small
  \begin{tabularx}{\linewidth}{@{}p{34mm}p{24mm}p{24mm}X@{}}
    \toprule
    \textbf{近似单高斯模型 $Q$} & $\DKL(P \,\|\, Q)$ & $\DKL(Q \,\|\, P)$ & \textbf{行为特征与对应 AI 范式} \\
    \midrule
    $q_{\mathrm{fwd}} = \mathcal{N}(0,\, 2.09^2)$ & \textbf{0.555}（胜出） & 1.175 & \textbf{Mode-Covering}：覆盖全部模态，对应 MLE 与自回归预训练 \\
    $q_{\mathrm{rev}} = \mathcal{N}(2,\, 0.60^2)$ & 10.306（巨额罚分） & \textbf{0.692}（胜出） & \textbf{Mode-Seeking}：聚焦高密度峰，对应变分推断与 RLHF/DPO \\
    \bottomrule
  \end{tabularx}
\end{table}

\begin{insightBox}{贯穿第 5--6 卷：为什么预训练爱正向 KL，而对齐爱反向 KL？}
\begin{itemize}
  \item \textbf{极大似然 / 大模型预训练（正向 KL）}：下节立刻证明，MLE 等价于最小化 $\DKL(P_{\mathrm{data}}\|Q_\theta)$。因此预训练语言模型倾向于\textbf{覆盖一切可能的人类表达}（多样性强，但容易在多模态答案之间“和稀泥”或产生幻觉）；
  \item \textbf{变分推断（VAE 编码器）与 RLHF/DPO 对齐（反向 KL）}：在强化学习微调中，目标函数加入反向 KL 惩罚 $-\beta\,\DKL(\pi_\theta \,\|\, \pi_{\mathrm{ref}})$。由于反向 KL 具有\textbf{锁峰特性}，微调后的大模型回复会变得高度干净、集中、高胜率，但过度强化也会导致回答千篇一律（多样性下降 / 模式坍缩）——这就是今天这幅双峰图在千亿参数大模型上的真实投影！
\end{itemize}
\end{insightBox}

% =========================================================================
\section{交叉熵、对数似然等价性与语言模型困惑度}
\label{sec:ch4-ce-ppl}

现在我们把全章收口到日常写代码时最常调用的一行函数——\textbf{交叉熵损失（Cross-Entropy Loss）}，并一次性打通它与 KL 散度、极大似然估计（MLE）以及语言模型困惑度（Perplexity）之间的严格等价关系。

\subsection{交叉熵定义与核心分解恒等式}

\begin{defBox}
\begin{definition}[交叉熵, Cross-Entropy]
\label{def:ch4-ce}
设真实分布为 $P$、模型预测分布为 $Q$，定义从 $P$ 到 $Q$ 的\textbf{交叉熵}为用 $Q$ 的编码长度对 $P$ 取期望：
\begin{equation}
  H(P, Q) := \E_{X \sim P}\bigl[-\log Q(X)\bigr] = -\sum_{x \in \mathcal{X}} P(x)\log Q(x). \label{eq:ch4-ce}
\end{equation}
\end{definition}
\end{defBox}

把对数里的 $-\log Q(x)$ 拆成 $-\log P(x) + \log\frac{P(x)}{Q(x)}$，立刻得到信息论中最核心的分解恒等式：

\begin{proposition}[交叉熵 = 真实熵 + KL 散度]
\label{prop:ch4-ce-kl}
对任意概率分布 $P, Q$，恒有
\begin{equation}
  H(P, Q) = H(P) + \DKL(P \,\|\, Q) \ge H(P), \label{eq:ch4-ce-kl}
\end{equation}
且等号成立当且仅当 $Q = P$。
\end{proposition}

\begin{proof}
直接展开并结合定理 \ref{thm:ch4-gibbs}（Gibbs 不等式）：
\[
  H(P, Q)
  = -\sum_x P(x)\log P(x) + \sum_x P(x)\log\frac{P(x)}{Q(x)}
  = H(P) + \DKL(P \,\|\, Q)
  \ge H(P).
\]
等号成立当且仅当 $\DKL(P\|Q) = 0 \iff Q = P$。
\end{proof}

这条恒等式的物理图像清澈见底：模型能达到的最低交叉熵极限不是 $0$，而是数据本身不可消除的内禀不确定性 $H(P)$；交叉熵高出 $H(P)$ 的那一部分，分毫不差就是模型偏离真实分布的惩罚项 $\DKL(P\|Q)$！

\subsection{统一大定理：最小交叉熵 $\equiv$ 最小正向 KL $\equiv$ 最大对数似然}

在监督学习与语言建模中，真实数据分布 $P_{\mathrm{data}}$ 是固定不变的，我们只优化模型参数 $\theta$。由此立刻得到贯通第 4--6 章的\textbf{统一大定理}：

\begin{thmBox}
\begin{theorem}[最小交叉熵、最小正向 KL 与最大对数似然的三重等价]
\label{thm:ch4-ce-kl-mle}
给定独立同分布观测样本 $x_1, \dots, x_N$，定义其经验分布 $P_{\mathrm{data}}(x) := \frac{1}{N}\sum_{i=1}^N \mathbf{1}\{x_i = x\}$，并设参数化概率模型为 $Q_\theta$。则以下三个优化问题具有\textbf{完全相同的最优参数集合}：
\begin{equation}
  \argmin_\theta H(P_{\mathrm{data}}, Q_\theta)
  \;=\;
  \argmin_\theta \DKL(P_{\mathrm{data}} \,\|\, Q_\theta)
  \;=\;
  \argmax_\theta \frac{1}{N}\sum_{i=1}^N \log Q_\theta(x_i). \label{eq:ch4-unifying}
\end{equation}
\end{theorem}
\end{thmBox}

\begin{proof}
由命题 \ref{prop:ch4-ce-kl}，有 $\DKL(P_{\mathrm{data}} \| Q_\theta) = H(P_{\mathrm{data}}, Q_\theta) - H(P_{\mathrm{data}})$。由于 $H(P_{\mathrm{data}})$ 与参数 $\theta$ 无关，仅为一个固定常数，因此最小化 $\DKL(P_{\mathrm{data}} \| Q_\theta)$ 与最小化 $H(P_{\mathrm{data}}, Q_\theta)$ 完全等价。另一方面，将经验分布代入交叉熵定义：
\[
  H(P_{\mathrm{data}}, Q_\theta)
  = -\sum_{x} P_{\mathrm{data}}(x)\log Q_\theta(x)
  = -\frac{1}{N}\sum_{i=1}^N \log Q_\theta(x_i).
\]
最小化负的平均对数似然，显然等价于最大化平均对数似然 $\frac{1}{N}\sum_{i=1}^N \log Q_\theta(x_i)$。证毕。
\end{proof}

\textbf{立即连线第 6 章分类损失}：在多分类任务中，单个样本的真实标签通常表示为一个 one-hot 分布 $P = e_c$（正确类别 $c$ 处概率为 $1$，其余为 $0$）。此时 $H(P) = 0$，交叉熵瞬间化简为极其熟悉的\textbf{负对数似然损失}：
\[
  H(e_c, Q_\theta) = -\ln Q_\theta(c).
\]
对二分类标签 $y \in \{0, 1\}$ 与预测概率 $\hat{p} = Q_\theta(y=1\mid x)$，就是第 6 章将亲手实现的\textbf{二元交叉熵（BCE）}：
\[
  \ell_{\mathrm{BCE}}(y, \hat{p}) = -\bigl[y\ln \hat{p} + (1 - y)\ln(1 - \hat{p})\bigr].
\]

\subsection{语言模型困惑度（Perplexity, PPL）}

训练大语言模型时，盯着一条下降的交叉熵曲线（比如从 $3.2$ 降到 $2.3\;\text{nats}$）往往不够直观。工程与论文中常把平均每 token 交叉熵取指数，转化为\textbf{困惑度（Perplexity）}。

\begin{defBox}
\begin{definition}[语言模型困惑度, Perplexity]
\label{def:ch4-ppl}
给定测试文本序列 $w_1, w_2, \dots, w_T$，语言模型 $Q_\theta$ 逐词预测条件概率 $Q_\theta(w_t \mid w_{<t})$。记平均每词交叉熵（以 nat 为单位）为
\[
  \bar{H}_{\mathrm{nat}} := -\frac{1}{T}\sum_{t=1}^T \ln Q_\theta(w_t \mid w_{<t}),
\]
则模型在该文本上的\textbf{困惑度}定义为
\begin{equation}
  \mathrm{PPL} := \exp(\bar{H}_{\mathrm{nat}}) = 2^{\bar{H}_{\mathrm{bit}}} = \left(\prod_{t=1}^T Q_\theta(w_t \mid w_{<t})\right)^{-1/T}. \label{eq:ch4-ppl}
\end{equation}
\end{definition}
\end{defBox}

\textbf{绝妙的直觉：等效词表分支数}。注意 $\mathrm{PPL}$ 与对数底数无关（用 $\ln$ 配 $\exp$，或用 $\log_2$ 配 $2^{(\cdot)}$，得到的都是几何平均概率的倒数！）。
\begin{itemize}
  \item \textbf{词表大小 $V = 50{,}000$ 的随机盲猜模型}：每步均匀猜 $Q = 1/50000$，平均交叉熵 $\bar{H}_{\mathrm{nat}} = \ln 50000 \approx 10.82\;\text{nats}$，于是 $\mathrm{PPL} = e^{10.82} = 50{,}000$；
  \item \textbf{强语言模型 $\bar{H}_{\mathrm{nat}} = 2.30\;\text{nats}$}：$\mathrm{PPL} = e^{2.30} \approx 10.0$——意思是每预测下一个词时，模型的不确定性\textbf{等效于在 $10$ 个候选词里均匀盲猜}！
  \item \textbf{完美预言机}：每个词概率都给到 $1$，则 $\bar{H} = 0$，$\mathrm{PPL} = e^0 = 1$（理论最低下界）。
\end{itemize}

\begin{failBox}{信息论编程三大数值雷区：$0\log 0$、$\log 0$ 爆炸与底数错配}
亲手写熵、KL 与困惑度代码时，以下三个坑几乎人人踩过一次：
\begin{enumerate}
  \item \textbf{直接写 \texttt{p * np.log(p)} 触发 \texttt{NaN}}：在 IEEE 754 浮点数中，当 \texttt{p == 0.0} 时，\texttt{np.log(0.0)} 返回 \texttt{-inf}，而 \texttt{0.0 * (-inf)} 返回 \texttt{nan}！正确做法是先用布尔掩码 \texttt{mask = p > 0} 过滤掉零概率项。
  \item \textbf{预测概率触零导致损失爆炸}：一旦某条样本真实类别对应的预测概率 $Q(x)$ 因下溢变成 \texttt{0.0}，\texttt{-np.log(0.0)} 立刻输出 \texttt{inf} 并摧毁梯度。工程中要么对概率做微小截断（如 \texttt{np.clip(q, 1e-15, 1.0)}），要么像第 6 章那样直接在 logit 空间使用 \texttt{logaddexp} / \texttt{log\_softmax} 避免先算指数再取对数。
  \item \textbf{底数与指数错配}：PyTorch 的 \texttt{CrossEntropyLoss} 默认用自然对数 \texttt{ln}（单位 nat），算困惑度必须用 \texttt{math.exp(loss)}；若误写成 \texttt{2 ** loss}，会严重低估真实困惑度！
\end{enumerate}
\end{failBox}

\subsection{纯 NumPy 实战：信息论工具箱与双峰正/反向 KL 实测}

下方代码实现一套零依赖、安全处理零概率的信息论基础工具箱，并当场验证分解恒等式 $H(P,Q) = H(P) + \DKL(P\|Q)$、例 \ref{ex:ch4-spam-mi} 的互信息以及图 \ref{fig:ch4-kl-bimodal} 中的双峰正/反向 KL 数值对比。

\begin{lstlisting}[caption={安全版信息论工具箱与双峰正/反向 KL 实测（纯 NumPy）},label={lst:ch4-info-tools}]
import numpy as np

def entropy(p, base=np.e):
    p = np.asarray(p, dtype=float)
    mask = p > 0
    return float(-np.sum(p[mask] * np.log(p[mask])) / np.log(base))

def kl_div(p, q, base=np.e):
    p, q = np.asarray(p, dtype=float), np.asarray(q, dtype=float)
    mask = p > 0
    if np.any(q[mask] <= 0):
        return float("inf")
    return float(np.sum(p[mask] * np.log(p[mask] / q[mask])) / np.log(base))

def cross_entropy(p, q, base=np.e):
    p, q = np.asarray(p, dtype=float), np.asarray(q, dtype=float)
    mask = p > 0
    if np.any(q[mask] <= 0):
        return float("inf")
    return float(-np.sum(p[mask] * np.log(q[mask])) / np.log(base))

def mutual_info(joint_xy, base=np.e):
    joint = np.asarray(joint_xy, dtype=float)
    px = joint.sum(axis=1, keepdims=True)
    py = joint.sum(axis=0, keepdims=True)
    return kl_div(joint.ravel(), (px @ py).ravel(), base=base)

# 1. Verify H(P, Q) = H(P) + KL(P || Q) with zero-entry safety
P = np.array([0.5, 0.25, 0.25, 0.0])
Q = np.array([0.25, 0.25, 0.25, 0.25])
H_p, D_pq, CE_pq = entropy(P, 2), kl_div(P, Q, 2), cross_entropy(P, Q, 2)
print(f"H(P)={H_p:.4f}, KL={D_pq:.4f}, CE={CE_pq:.4f}, gap={abs(CE_pq-(H_p+D_pq)):.1e}")
# Verified output: H(P)=1.5000, KL=0.5000, CE=2.0000, gap=0.0e+00

# 2. Verify Example 4.2 mutual information
P_xy = np.array([[0.40, 0.10], [0.10, 0.40]])
print(f"I(X; Y) = {mutual_info(P_xy, base=2):.4f} bits")  # 0.2781 bits

# 3. Bimodal Forward KL vs Reverse KL numerical check
x = np.linspace(-6.0, 6.0, 601)
dx = x[1] - x[0]
norm_pdf = lambda z, m, s: np.exp(-0.5 * ((z - m) / s)**2) / (np.sqrt(2*np.pi) * s)
p = 0.5 * norm_pdf(x, -2.0, 0.6) + 0.5 * norm_pdf(x, 2.0, 0.6)
p /= np.sum(p) * dx

for tag, q in [("q_fwd(0, 2.09)", norm_pdf(x, 0.0, np.sqrt(4.36))),
               ("q_rev(2, 0.60)", norm_pdf(x, 2.0, 0.6))]:
    kl_pq = float(np.sum(p * np.log((p + 1e-15) / (q + 1e-15))) * dx)
    kl_qp = float(np.sum(q * np.log((q + 1e-15) / (p + 1e-15))) * dx)
    print(f"{tag}: KL(P||Q)={kl_pq:.3f}, KL(Q||P)={kl_qp:.3f}")
# Verified output:
# q_fwd(0, 2.09): KL(P||Q)=0.555, KL(Q||P)=1.175
# q_rev(2, 0.60): KL(P||Q)=10.306, KL(Q||P)=0.692
\end{lstlisting}

% =========================================================================
\section{本章习题}
\label{sec:ch4-exercises}

\exTier{A. 概念辨析与反例查错}{看清散度非对称、微分熵可负与逐点条件熵的微妙边界}

\begin{enumerate}[label=\textbf{4.\arabic*}, leftmargin=2.4em]
  \item \textbf{（信息论核心概念辨析）} 判断以下四个命题正误；对错误命题给出反例，对正确命题简述理由：
  \begin{enumerate}[label=(\alph*), leftmargin=1.8em]
    \item 因为 KL 散度满足 $\DKL(P\|Q) \ge 0$ 且 $\DKL(P\|Q)=0 \iff P=Q$，所以 $\DKL$ 是概率分布空间上的合法距离度量（满足对称性与三角不等式）；
    \item 任何离散随机变量的 Shannon 熵必然非负，但连续随机变量的微分熵 $h(X) = -\int p(x)\ln p(x)\,dx$ 可以为负数；
    \item 因为条件熵满足 $H(Y \mid X) \le H(Y)$，所以只要观测到任意具体特征值 $X = x_0$，条件熵 $H(Y \mid X = x_0)$ 就必然不超过边缘熵 $H(Y)$；
    \item 对同一条测试文本和同一个语言模型，用 $\log_2$ 算平均交叉熵再取 $2^{(\cdot)}$ 与用 $\ln$ 算平均交叉熵再取 $\exp(\cdot)$，得到的困惑度 $\mathrm{PPL}$ 数值完全一致。
  \end{enumerate}

  \item \textbf{（KL 散度非对称与支集爆炸反例）} 在二元状态空间 $\mathcal{X} = \{0, 1\}$ 上完成两项反例构造与验算（均用自然对数 $\ln$）：
  \begin{enumerate}[label=(\arabic*), leftmargin=1.8em]
    \item 取 $P = (0.9, 0.1)$ 与 $Q = (0.5, 0.5)$，手算 $\DKL(P\|Q)$ 与 $\DKL(Q\|P)$（保留四位小数），直观说明两者为何不等；
    \item 构造一对分布 $P, Q$，使得 $\DKL(P\|Q)$ 为有限正数，但反向散度 $\DKL(Q\|P) = +\infty$，并用“正向 KL 零回避 / 反向 KL 零强制”解释其物理含义。
  \end{enumerate}
\end{enumerate}

\exTier{B. 小例手算与数值直觉}{亲手打通二维联合熵、互信息、语言模型困惑度与高斯 KL}

\begin{enumerate}[label=\textbf{4.\arabic*}, leftmargin=2.4em, resume]
  \item \textbf{（$2\times 2$ 联合分布全套信息量手算）} 设二元随机变量 $(X, Y)$ 的联合分布为：
  \[
    P(X=0, Y=0) = \frac{1}{2}, \quad P(X=0, Y=1) = 0, \quad P(X=1, Y=0) = \frac{1}{4}, \quad P(X=1, Y=1) = \frac{1}{4}.
  \]
  请以 $\log_2$（单位：bit）为单位，逐步手算以下六个量（写出精确表达式并保留四位小数）：
  \[
    H(X),\quad H(Y),\quad H(X, Y),\quad H(Y \mid X),\quad H(X \mid Y),\quad I(X; Y),
  \]
  并直接核验链式法则 $H(X,Y) = H(X) + H(Y\mid X) = H(Y) + H(X\mid Y)$ 与互信息对称性 $I(X;Y) = H(Y) - H(Y\mid X) = H(X) - H(X\mid Y)$。

  \item \textbf{（语言模型困惑度与一维高斯 KL 手算）} 完成以下两项独立手算：
  \begin{enumerate}[label=(\arabic*), leftmargin=1.8em]
    \item 某语言模型在一条含 $4$ 个词的测试句上，对真实下一个词给出的预测条件概率依次为 $(0.5,\allowbreak\; 0.25,\allowbreak\; 0.125,\allowbreak\; 0.5)$。手算平均每词交叉熵 $\bar{H}_{\mathrm{bit}}$（单位 bit）、$\bar{H}_{\mathrm{nat}}$（单位 nat）以及困惑度 $\mathrm{PPL}$；
    \item 利用命题 \ref{prop:ch4-gaussian-kl} 的高斯 KL 闭式公式，分别手算 $\DKL(\mathcal{N}(1, 2^2)\,\|\,\mathcal{N}(0, 1))$ 与 $\DKL(\mathcal{N}(0, 1)\,\|\,\mathcal{N}(1, 2^2))$（单位 nat，保留四位小数），体会把宽分布逼近窄分布与把窄分布逼近宽分布的不对称代价。
  \end{enumerate}
\end{enumerate}

\exTier{C. 实用推导与性质证明}{证明正向 KL 矩匹配最优性，掌握对称化 Jensen--Shannon 散度}

\begin{enumerate}[label=\textbf{4.\arabic*}, leftmargin=2.4em, resume]
  \item \textbf{（单高斯正向 KL 极小化的精确矩匹配定理）} 设任意目标分布 $P$ 具有有限均值 $\mu_P = \E_P[X]$ 与有限正方差 $\sigma_P^2 = \Var_P(X)$。考虑用单高斯分布 $Q_{\mu,\sigma^2} = \mathcal{N}(\mu, \sigma^2)$ 去最小化正向 KL 散度：
  \[
    (\mu^*, (\sigma^*)^2) = \argmin_{\mu \in \R,\; \sigma^2 > 0} \DKL\bigl(P \,\|\, \mathcal{N}(\mu, \sigma^2)\bigr).
  \]
  \begin{enumerate}[label=(\arabic*), leftmargin=1.8em]
    \item 将目标函数展开并去掉与 $(\mu, \sigma^2)$ 无关的项，证明其等价于最小化
    \[
      J(\mu, \sigma^2) = \frac{1}{2}\ln(2\pi \sigma^2) + \frac{\sigma_P^2 + (\mu_P - \mu)^2}{2\sigma^2};
    \]
    \item 对 $\mu$ 与 $\sigma^2$ 分别求偏导并令其为零，严格证明唯一全局最优解就是\textbf{矩匹配解} $\mu^* = \mu_P,\; (\sigma^*)^2 = \sigma_P^2$；
    \item 利用此结论，三行算出双峰混合 $P = 0.5\mathcal{N}(-2, 0.6^2) + 0.5\mathcal{N}(2, 0.6^2)$ 的正向 KL 最优参数 $\mu_{\mathrm{fwd}} = 0$ 与 $\sigma_{\mathrm{fwd}} = \sqrt{4.36} \approx 2.088$。
  \end{enumerate}
  \exHint{第 (1) 问用到恒等式 $\E_P[(X - \mu)^2] = \Var_P(X) + (\mu_P - \mu)^2$；第 (3) 问用全期望与全方差公式：$\E_P[X^2] = 0.5(0.6^2 + (-2)^2) + 0.5(0.6^2 + 2^2) = 4.36$。}

  \item \textbf{（对称且有界的 Jensen--Shannon 散度）} 为克服 KL 散度不对称且可能发散到 $+\infty$ 的缺陷，定义 \textbf{Jensen--Shannon 散度（JSD）}：记混合分布 $M := \frac{1}{2}(P + Q)$，
  \[
    \mathrm{JSD}(P \,\|\, Q) := \frac{1}{2}\DKL(P \,\|\, M) + \frac{1}{2}\DKL(Q \,\|\, M).
  \]
  \begin{enumerate}[label=(\arabic*), leftmargin=1.8em]
    \item 证明 $\mathrm{JSD}(P\|Q) = \mathrm{JSD}(Q\|P) \ge 0$，且 $\mathrm{JSD}(P\|Q) = 0 \iff P = Q$；
    \item 构造辅助二元随机变量 $Z \sim \mathrm{Bernoulli}(0.5)$，规定条件分布 $X \mid (Z=0) \sim P$、$X \mid (Z=1) \sim Q$。证明 $\mathrm{JSD}(P\|Q)$ 恰好等于互信息 $I(X; Z)$；
    \item 结合 (2) 的结果，两行证明对\textbf{任意}两个分布 $P, Q$（哪怕支集完全不相交！），恒有上界 $\mathrm{JSD}(P\|Q) \le \ln 2$（以 $\log_2$ 计则为 $1\;\text{bit}$）。
  \end{enumerate}
  \exHint{第 (2) 问直接计算 $H(X) - H(X\mid Z) = H(M) - \frac{1}{2}H(P) - \frac{1}{2}H(Q)$ 并与 $\mathrm{JSD}$ 展开式核对；第 (3) 问利用 $I(X; Z) = H(Z) - H(Z\mid X) \le H(Z) = \ln 2$。}
\end{enumerate}

\exTier{D. 编程实战与可自检实验}{用 NumPy 实现决策树信息增益排位与双峰 KL 网格寻优}

\begin{enumerate}[label=\textbf{4.\arabic*}, leftmargin=2.4em, resume]
  \item \textbf{（安全信息论工具箱与决策树根节点信息增益选特征）} 编写纯 NumPy 脚本完成以下两项任务：
  \begin{enumerate}[label=(\arabic*), leftmargin=1.8em]
    \item 实现安全处理零概率的 \texttt{entropy}、\texttt{kl\_div}、\texttt{cross\_entropy} 与 \texttt{mutual\_info} 四个函数；随机生成 $100$ 对长度 $K = 8$ 的离散分布（含零元素，确保支集覆盖）验证 $|H(P, Q) - (H(P) + \DKL(P\|Q))| < 10^{-12}$ 且 $\DKL(P\|Q) \ge -10^{-15}$；
    \item 给定一个含 $8$ 条样本的二分类小数据集（标签与三个候选二值特征如下）：
    \begin{align*}
      y &= [1, 1, 1, 1, 0, 0, 0, 0], & F_1 &= [1, 1, 1, 0, 0, 0, 0, 1], \\
      F_2 &= [1, 1, 0, 0, 1, 1, 0, 0], & F_3 &= [1, 1, 1, 1, 0, 0, 1, 0].
    \end{align*}
    计算每个特征与标签 $y$ 的信息增益 $\mathrm{IG}(y, F_j) = I(F_j; y)$（单位 bit），选出最优决策树根节点分裂特征。
  \end{enumerate}
  \par\smallskip\noindent{\footnotesize\sffamily\color{tealDark}\textbf{【自检标准】}\,第 (1) 问所有 $100$ 对分布的恒等式最大绝对误差须 $< 10^{-12}$ 且不能出现任何 \texttt{NaN}；第 (2) 问的经验信息增益须精确命中：$\mathrm{IG}(y, F_1) = 0.1887\;\text{bit}$、$\mathrm{IG}(y, F_2) = 0.0000\;\text{bit}$（与标签完全独立）、$\mathrm{IG}(y, F_3) = 0.5488\;\text{bit}$，最优分裂特征为 $F_3$。}

  \item \textbf{（双峰分布正向 vs 反向 KL 数值寻优复现）} 在一维网格 $x \in [-6, 6]$（$601$ 个等距节点）上构造真实双峰分布 $P(x) = 0.5\,\mathcal{N}(-2, 0.6^2) + 0.5\,\mathcal{N}(2, 0.6^2)$（并在网格上归一化使 $\sum p(x)\Delta x = 1$）。编写纯 NumPy 网格搜索脚本，在 $\mu \in [-2.5, 2.5]$（步长 $0.05$）与 $\sigma \in [0.3, 2.8]$（步长 $0.025$）组成的 $101 \times 101$ 参数网格上：
  \begin{enumerate}[label=(\arabic*), leftmargin=1.8em]
    \item 寻找使正向散度 $\DKL(P \| \mathcal{N}(\mu, \sigma^2))$ 最小的最优参数 $(\mu_{\mathrm{fwd}}, \sigma_{\mathrm{fwd}})$；
    \item 寻找使反向散度 $\DKL(\mathcal{N}(\mu, \sigma^2) \| P)$ 最小的最优参数 $(\mu_{\mathrm{rev}}, \sigma_{\mathrm{rev}})$；
    \item 分别报告两组最优解下的 $\DKL(P\|Q)$ 与 $\DKL(Q\|P)$ 数值，核验其与表 \ref{tab:ch4-kl-bimodal} 的吻合程度。
  \end{enumerate}
  \par\smallskip\noindent{\footnotesize\sffamily\color{tealDark}\textbf{【自检标准】}\,正向 KL 搜索结果须满足 $\mu_{\mathrm{fwd}} \in [-0.05, 0.05]$、$\sigma_{\mathrm{fwd}} \in [2.05, 2.13]$，且对应 $\DKL(P\|Q) \approx 0.555 \pm 0.005$、$\DKL(Q\|P) \approx 1.175 \pm 0.015$；反向 KL 搜索结果须满足 $|\mu_{\mathrm{rev}}| \in [1.95, 2.05]$、$\sigma_{\mathrm{rev}} \in [0.57, 0.63]$，且对应 $\DKL(Q\|P) \approx 0.692 \pm 0.005$、$\DKL(P\|Q) > 9.5$（精确峰处约 $10.306$）。}
\end{enumerate}

\chapter[统计估计：MLE、MAP 与 Bootstrap]{统计估计方法论：MLE、MAP、共轭先验与 Bootstrap}
\label{ch:ch5-estimation-bootstrap}

\chapEpigraph{所有模型都是错的，但其中一些是有用的。}{乔治·博克斯（George E. P. Box），1976}

\begin{chapIntro}{本章定位：从“已知分布算概率”到“手握数据反推模型”}
前四章我们都在做同一件事——\textbf{给定概率分布的参数，推导概率、期望、浓度界与信息散度}。但在真实人工智能任务里，方向恰好反过来：自然界不会把硬币正面概率 $p$、高斯噪声方差 $\sigma^2$ 或分类器权重 $w$ 印在说明书上，它只给我们一份有限的数据集 $\mathcal{D} = \{x_1, \dots, x_n\}$。本章建立从数据反推参数与量化不确定性的完整统计方法论：
\begin{itemize}
  \item 用\textbf{偏差-方差分解}回答“什么样的估计量才算好”，并彻底讲透为何样本方差要除以 $n-1$；
  \item 用\textbf{极大似然估计（MLE）}统一常见分布参数估计与第 4 章的最小 KL 散度，同时用“连抛 3 次全正面”反例揭露其小样本过拟合陷阱；
  \item 用\textbf{共轭先验与最大后验估计（MAP）}把领域知识转化为伪计数与正则化，为第 6 章导出岭回归与 Lasso 铺好数学桥梁；
  \item 用\textbf{非参数 Bootstrap 重采样}为准确率、$F_1$、AUC 等任意 AI 指标构造无需正态假设的置信区间。
\end{itemize}
\end{chapIntro}

% ====================================================================
\section{估计量的偏差、方差与均方误差分解}
\label{sec:ch5-bias-var-mse}

\begin{recallBox}{第 2--3 章核心结论速忆}
对独立同分布（i.i.d.）随机变量 $X_1, \dots, X_n$，若 $\E[X_i] = \mu$、$\Var(X_i) = \sigma^2 < \infty$，则样本均值 $\bar{X} = \frac{1}{n}\sum_{i=1}^n X_i$ 满足
\[
  \E[\bar{X}] = \mu, \qquad \Var(\bar{X}) = \frac{\sigma^2}{n}.
\]
同时恒有方差展开式 $\Var(Z) = \E[Z^2] - (\E[Z])^2$ 与期望线性性 $\E[aZ + b] = a\E[Z] + b$。
\end{recallBox}

先看一个贴近工程的场景：你想测量某个广告推荐位真实的点击率 $\theta \in [0,1]$。由于不可能把广告展示给全世界所有人，你抽样观测了 $n=100$ 次曝光记录 $X_1, \dots, X_{100} \stackrel{\text{i.i.d.}}{\sim} \mathrm{Bernoulli}(\theta)$，然后用某种规则（例如点击比例 $\hat{\theta}_n = \bar{X}$）给出一个数值猜测。

这里必须严格区分三个初学极易混淆的对象：
\begin{enumerate}
  \item \textbf{真实参数（True Parameter）$\theta$}：在频率派视角下，它是一个未知但固定的常数；
  \item \textbf{点估计量（Point Estimator）$\hat{\theta}_n = g(X_1, \dots, X_n)$}：它是随机样本的函数，因此在抽样前本身是一个\textbf{随机变量}，拥有自己的期望与方差；
  \item \textbf{估计值（Estimate）}：把具体采集到的数据 $x_1, \dots, x_n$ 代入后算出的那个具体实数，例如 $0.23$。
\end{enumerate}

既然 $\hat{\theta}_n$ 是一个随机变量，如果今天重新采样 100 次曝光，估计值可能变成 $0.19$；明天再采一次，可能是 $0.27$。那么我们如何科学地评判两个不同估计量 $\hat{\theta}_n^{(1)}$ 和 $\hat{\theta}_n^{(2)}$ 谁更“准”？

\begin{defBox}
\begin{definition}[偏差、方差、均方误差与相合性]
\label{def:ch5-estimator-metrics}
设 $\hat{\theta}_n = g(X_1, \dots, X_n)$ 是未知标量参数 $\theta \in \R$ 的点估计量。定义：
\begin{itemize}
  \item \textbf{偏差（Bias）}：衡量估计量的平均瞄准位置与真值的系统性偏移，
  \[
    \Bias(\hat{\theta}_n) \triangleq \E[\hat{\theta}_n] - \theta.
  \]
  若对参数空间中所有 $\theta$ 均有 $\Bias(\hat{\theta}_n) = 0$（即 $\E[\hat{\theta}_n] = \theta$），则称 $\hat{\theta}_n$ 为\textbf{无偏估计量（Unbiased Estimator）}。
  \item \textbf{方差（Variance）}：衡量不同次随机抽样导致估计值自身抖动的幅度，
  \[
    \Var(\hat{\theta}_n) \triangleq \E\!\left[\big(\hat{\theta}_n - \E[\hat{\theta}_n]\big)^2\right].
  \]
  \item \textbf{均方误差（Mean Squared Error, MSE）}：衡量估计量到真实参数 $\theta$ 的平均平方距离，
  \[
    \mathrm{MSE}(\hat{\theta}_n) \triangleq \E\!\left[(\hat{\theta}_n - \theta)^2\right].
  \]
  \item \textbf{相合性（Consistency）}：若对任意 $\epsilon > 0$ 有 $\lim_{n\to\infty}\Prob\big(\abs{\hat{\theta}_n - \theta} > \epsilon\big) = 0$（记作 $\hat{\theta}_n \xrightarrow{P} \theta$），则称 $\hat{\theta}_n$ 是相合估计量。由第 3 章 Chebyshev 不等式可知，只要 $\lim_{n\to\infty}\mathrm{MSE}(\hat{\theta}_n) = 0$，估计量必然相合。
\end{itemize}
\end{definition}
\end{defBox}

如果把估计参数比作打靶：真值 $\theta$ 是靶心，每次抽样算出的估计值是一发子弹。$\Bias(\hat{\theta}_n)$ 是弹着点中心偏离靶心的距离，$\Var(\hat{\theta}_n)$ 是弹着点自身的散开程度，而 $\mathrm{MSE}(\hat{\theta}_n)$ 则是所有子弹到靶心平均平方距离的总账。下面这条极其优美的四行定理揭示了三者之间的精确关系。

\begin{thmBox}
\begin{theorem}[估计量的均方误差偏差-方差分解]
\label{thm:ch5-mse-decomp}
对任意具有有限二阶矩的估计量 $\hat{\theta}$ 与真实参数 $\theta$，恒有
\begin{equation}
  \mathrm{MSE}(\hat{\theta}) = \Bias(\hat{\theta})^2 + \Var(\hat{\theta}).
  \label{eq:ch5-mse-decomp}
\end{equation}
\end{theorem}
\end{thmBox}

\begin{proof}
记 $m = \E[\hat{\theta}]$（它是一个确定性常数）。在 $(\hat{\theta} - \theta)^2$ 中添项减项 $m$ 并展开：
\begin{align*}
  \mathrm{MSE}(\hat{\theta})
  &= \E\!\left[\big((\hat{\theta} - m) + (m - \theta)\big)^2\right] \\
  &= \E\!\left[(\hat{\theta} - m)^2\right] + 2(m - \theta)\,\E[\hat{\theta} - m] + (m - \theta)^2 \\
  &= \Var(\hat{\theta}) + 2(m - \theta)\cdot 0 + \Bias(\hat{\theta})^2 \\
  &= \Bias(\hat{\theta})^2 + \Var(\hat{\theta}).
\end{align*}
其中第二步利用了 $m - \theta$ 为常数可提出期望号，第三步利用了 $\E[\hat{\theta} - m] = \E[\hat{\theta}] - \E[\hat{\theta}] = 0$。
\end{proof}

\begin{insightBox}{直通第 4 卷监督学习的核心法则：偏差-方差权衡}
式~\eqref{eq:ch5-mse-decomp} 是整个机器学习理论中最著名的\textbf{偏差-方差权衡（Bias-Variance Tradeoff）}在参数估计层的直接前身！很多人直觉上认为“无偏估计量一定是最好的”，但式~\eqref{eq:ch5-mse-decomp} 明确指出：总误差 $\mathrm{MSE}$ 是\textbf{偏差平方与方差之和}。如果一个估计量稍微引入一点点可控偏差（例如将参数向零收缩），却能把方差砍掉一大半，它的总体 $\mathrm{MSE}$ 反而会\textbf{严格低于}无偏估计量！我们在第 5.3 节的贝叶斯 MAP 估计与第 6 章的 $L_2$ 岭回归中，都会反复见证“以微小偏差换取方差暴跌”的魔力。
\end{insightBox}

\subsection{经典推导：为什么样本方差要除以 $n-1$？}
\label{subsec:ch5-bessel}

几乎每位理工科学生都背过“无偏样本方差要除以 $n-1$（Bessel 校正）”，但到底少了哪 $1$ 个自由度？直接除以 $n$ 又会偏小多少？让我们先看一个 $n=2$ 的极简特例：设 $X_1, X_2 \stackrel{\text{i.i.d.}}{\sim} \mathcal{N}(0, 1)$，真方差 $\sigma^2 = 1$。此时 $\bar{X} = \frac{X_1+X_2}{2}$，于是
\[
  \frac{1}{2}\sum_{i=1}^2 (X_i - \bar{X})^2
  = \frac{1}{2}\!\left[\left(\frac{X_1-X_2}{2}\right)^2 + \left(\frac{X_2-X_1}{2}\right)^2\right]
  = \frac{1}{4}(X_1 - X_2)^2.
\]
因为 $X_1 - X_2 \sim \mathcal{N}(0, 2)$，其二阶矩为 $2$，所以除以 $n=2$ 得到的期望竟只有 $\frac{1}{4}\times 2 = 0.5$——系统性地\textbf{低估了整整一半}！直觉上，样本均值 $\bar{X}$ 恰好位于样本 $X_1, \dots, X_n$ 的“正中心”，它天然比未知的真实中心 $\mu$ 更贴近这批样本，从而人为压小了偏差平方和。下面的命题精确量化了这个低估量。

\begin{thmBox}
\begin{proposition}[样本方差的无偏校正与 MLE 偏差]
\label{prop:ch5-sample-var}
设 $X_1, \dots, X_n$（$n \ge 2$）i.i.d. 且均值为 $\mu$、方差为 $\sigma^2 < \infty$。定义直接除以 $n$ 的经验方差与除以 $n-1$ 的样本方差分别为
\[
  \hat{\sigma}_n^2 \triangleq \frac{1}{n}\sum_{i=1}^n (X_i - \bar{X})^2,
  \qquad
  S^2 \triangleq \frac{1}{n-1}\sum_{i=1}^n (X_i - \bar{X})^2.
\]
则恒有：
\[
  \E[\hat{\sigma}_n^2] = \frac{n-1}{n}\sigma^2,
  \qquad
  \Bias(\hat{\sigma}_n^2) = -\frac{\sigma^2}{n},
  \qquad
  \E[S^2] = \sigma^2 \;\;(\text{无偏}).
\]
\end{proposition}
\end{thmBox}

\begin{proof}
对离差平方和做“以 $\bar{X}$ 为中转”的恒等分解。首先将 $X_i - \mu = (X_i - \bar{X}) + (\bar{X} - \mu)$ 两边平方并对 $i=1,\dots,n$ 求和：
\begin{align*}
  \sum_{i=1}^n (X_i - \mu)^2
  &= \sum_{i=1}^n (X_i - \bar{X})^2 + 2(\bar{X} - \mu)\underbrace{\sum_{i=1}^n (X_i - \bar{X})}_{=\,0} + \sum_{i=1}^n (\bar{X} - \mu)^2 \\
  &= \sum_{i=1}^n (X_i - \bar{X})^2 + n(\bar{X} - \mu)^2.
\end{align*}
移项得核心恒等式：
\[
  \sum_{i=1}^n (X_i - \bar{X})^2 = \sum_{i=1}^n (X_i - \mu)^2 - n(\bar{X} - \mu)^2.
\]
两边同时取期望，利用 $\E[(X_i - \mu)^2] = \Var(X_i) = \sigma^2$ 与 $\E[(\bar{X} - \mu)^2] = \Var(\bar{X}) = \sigma^2/n$：
\[
  \E\!\left[\sum_{i=1}^n (X_i - \bar{X})^2\right]
  = \sum_{i=1}^n \sigma^2 - n\cdot\frac{\sigma^2}{n}
  = (n - 1)\sigma^2.
\]
因此 $\E[\hat{\sigma}_n^2] = \frac{n-1}{n}\sigma^2$，其偏差为 $\Bias(\hat{\sigma}_n^2) = \frac{n-1}{n}\sigma^2 - \sigma^2 = -\frac{\sigma^2}{n}$；而 $S^2 = \frac{1}{n-1}\sum_{i=1}^n (X_i - \bar{X})^2$ 满足 $\E[S^2] = \sigma^2$，精确无偏。
\end{proof}

\begin{remark}[几何视角与“无偏未必 MSE 最小”]
{\rightskip=0pt plus 1em\relax 回想第 1 卷第 2 章的正交投影：把样本向量 $x = (X_1,\dots,X_n)^T$（$\in \R^n$）减去均值 $\bar{X}\mathbf{1}$，本质上是把 $x - \mu\mathbf{1}$ 正交投影到与全一向量 $\mathbf{1}$ 垂直的 $(n-1)$ 维超平面上，投影正好丢掉了沿 $\mathbf{1}$ 方向的 $1$ 个自由度（能量恰好等于 $\sigma^2$）！\par}

更有趣的是，如果数据服从高斯分布 $\mathcal{N}(\mu, \sigma^2)$，可以证明 $\Var\!\left(\sum_{i=1}^n (X_i-\bar{X})^2\right) = 2(n-1)\sigma^4$。于是对形如 $c\sum_{i=1}^n (X_i-\bar{X})^2$ 的缩放估计量，应用定理~\ref{thm:ch5-mse-decomp} 得
\[
  \mathrm{MSE}(c) = \sigma^4\!\left[\big(c(n-1) - 1\big)^2 + 2c^2(n-1)\right].
\]
对 $c$ 求导令其为零，竟解出使 $\mathrm{MSE}$ 最小的系数是 $c^* = \frac{1}{n+1}$！以 $n=5, \sigma^2=4$ 为例（见习题 5.7 实测验证）：无偏估计量 $S^2$（除以 $4$）的 $\mathrm{MSE} = 8.000$，除以 $5$ 的 $\hat{\sigma}_n^2$ 的 $\mathrm{MSE} = 5.760$，而有偏更甚的除以 $6$ 估计量的 $\mathrm{MSE}$ 反倒最低，仅 $5.333$！这再次提醒我们：在 AI 建模里，切忌迷信“无偏”，要时刻盯紧偏差与方差的合账。
\end{remark}

% ====================================================================
\section{极大似然估计（MLE）及其小样本过拟合陷阱}
\label{sec:ch5-mle}

现在我们提出一个系统性的通用方法来构造估计量：假设观测数据 $\mathcal{D} = \{x_1, \dots, x_n\}$ 独立同分布地采样自参数化分布族 $\{p(x; \theta)\}_{\theta \in \Theta}$，怎么挑出最合理的 $\theta$？

先看一个直觉小例：抛一枚质地不均的硬币 $10$ 次，观测到 $7$ 次正面、$3$ 次反面。在候选参数 $p \in \{0.1,\allowbreak 0.5,\allowbreak 0.7,\allowbreak 0.9\}$ 中，哪一个最“合理”？让我们直接算出每个 $p$ 下出现这一观测序列的概率 $p^7(1-p)^3$：
\begin{itemize}
  \item 若 $p = 0.1$：$0.1^7 \times 0.9^3 \approx 7.29 \times 10^{-8}$（几乎不可能发生）；
  \item 若 $p = 0.5$：$0.5^{10} \approx 9.77 \times 10^{-4}$；
  \item 若 $p = 0.7$：$0.7^7 \times 0.3^3 \approx 2.22 \times 10^{-3}$（最大！）；
  \item 若 $p = 0.9$：$0.9^7 \times 0.1^3 \approx 4.78 \times 10^{-4}$。
\end{itemize}
把“选出让已经发生的观测数据出现概率（或密度）最大的那个参数”这一朴素原则形式化，就得到了现代统计与机器学习最核心的参数估计准则——\textbf{极大似然估计}。

\begin{defBox}
\begin{definition}[似然函数、对数似然与极大似然估计（MLE）]
\label{def:ch5-mle}
设样本 $\mathcal{D} = \{x_1, \dots, x_n\}$ 独立同分布于概率质量函数（离散）或概率密度函数（连续）$p(x; \theta)$。固定观测数据 $\mathcal{D}$，将联合概率视为参数 $\theta$ 的函数，称为\textbf{似然函数（Likelihood Function）}：
\[
  L(\theta; \mathcal{D}) \triangleq \prod_{i=1}^n p(x_i; \theta).
\]
取自然对数得到\textbf{对数似然函数（Log-Likelihood）}：
\begin{equation}
  \ell(\theta) \triangleq \ln L(\theta; \mathcal{D}) = \sum_{i=1}^n \ln p(x_i; \theta).
  \label{eq:ch5-loglik}
\end{equation}
\textbf{极大似然估计（Maximum Likelihood Estimation, MLE）}定义为最大化似然（等价于最大化对数似然）的参数：
\[
  \hat{\theta}_{\text{MLE}} \triangleq \argmax_{\theta \in \Theta} L(\theta; \mathcal{D}) = \argmax_{\theta \in \Theta} \ell(\theta).
\]
\end{definition}
\end{defBox}

为什么在理论推导和代码实现中几乎永远使用\textbf{对数似然} $\ell(\theta)$ 而非原始连乘 $L(\theta)$？原因有二：（1）对数函数 $\ln(\cdot)$ 严格单调递增，不改变极大值点位置；（2）把连乘变成求和后，不仅求导极为方便，更避免了数值下溢——哪怕每个样本概率高达 $0.5$，连乘 $n=10000$ 个也会得到 $0.5^{10000} \approx 10^{-3010}$，远低于 IEEE 754 双精度浮点数最小正数 $2.2\times 10^{-308}$ 而直接下溢成 \texttt{0.0}！

\subsection{MLE 与第 4 章最小 KL 散度的等价性}
\label{subsec:ch5-mle-kl}

在第 4 章我们建立了信息论中的交叉熵与 KL 散度；现在把第 4 章和本章接通，会看到一个贯通整个深度学习的黄金恒等式。

\begin{thmBox}
\begin{proposition}[极大似然估计 $\equiv$ 最小化经验 KL 散度]
\label{prop:ch5-mle-kl}
设 $\hat{P}_{\text{data}}$ 表示由样本 $\{x_1, \dots, x_n\}$ 定义的经验分布（每个观测点权重 $1/n$），$P_\theta$ 表示模型分布 $p(\cdot; \theta)$。则最大化平均对数似然完全等价于最小化经验交叉熵，也等价于最小化从数据分布到模型分布的\textbf{正向 KL 散度}：
\begin{equation}
  \argmax_\theta \frac{1}{n}\sum_{i=1}^n \ln p(x_i; \theta)
  = \argmin_\theta H(\hat{P}_{\text{data}}, P_\theta)
  = \argmin_\theta \DKL(\hat{P}_{\text{data}} \| P_\theta).
  \label{eq:ch5-mle-kl}
\end{equation}
\end{proposition}
\end{thmBox}

\begin{proof}
由第 4 章交叉熵与 KL 散度定义（使用自然对数底，单位 nat）：
\[
  \DKL(\hat{P}_{\text{data}} \| P_\theta)
  = \underbrace{-\frac{1}{n}\sum_{i=1}^n \ln p(x_i; \theta)}_{=\, H(\hat{P}_{\text{data}}, P_\theta)} - H(\hat{P}_{\text{data}}).
\]
由于经验分布的熵 $H(\hat{P}_{\text{data}})$ 与待优化参数 $\theta$ 无关，最小化 $\DKL(\hat{P}_{\text{data}} \| P_\theta)$ 恰好等价于最小化负平均对数似然 $-\frac{1}{n}\ell(\theta)$，即最大化对数似然 $\ell(\theta)$。
\end{proof}

这条等价性解释了为什么大语言模型做“下一个词预测”（本质是自回归 MLE）优化的正是交叉熵损失，也解释了为何第 4 章强调\textbf{正向 KL 散度（mode-covering）}——用 MLE 拟合模型时，只要数据里出现过的样本 $x_i$，模型绝不敢给它赋予零概率（否则 $-\ln 0 = +\infty$）。

\subsection{三个核心分布的 MLE 完整推导}
\label{subsec:ch5-mle-examples}

\begin{example}[Bernoulli 分布的 MLE]
\label{ex:ch5-mle-bern}
设 $x_1, \dots, x_n \in \{0, 1\}$ 独立同分布于 $\mathrm{Bernoulli}(p)$（$p \in (0,1)$），记正面总次数 $k = \sum_{i=1}^n x_i$。似然与对数似然为
\[
  L(p) = \prod_{i=1}^n p^{x_i}(1-p)^{1-x_i} = p^k (1-p)^{n-k},
  \qquad
  \ell(p) = k\ln p + (n - k)\ln(1 - p).
\]
对 $p$ 求一阶导并令其为零：
\[
  \ell'(p) = \frac{k}{p} - \frac{n - k}{1 - p} = 0
  \;\Longrightarrow\;
  k(1 - p) = (n - k)p
  \;\Longrightarrow\;
  \hat{p}_{\text{MLE}} = \frac{k}{n}.
\]
检验二阶导数 $\ell''(p) = -\frac{k}{p^2} - \frac{n-k}{(1-p)^2} < 0$（对 $0 < k < n$ 严格凹），故 $\hat{p}_{\text{MLE}} = k/n$ 确为唯一全局极大值点；而当边界情况 $k=0$ 或 $k=n$ 时，直接观察 $L(p) = (1-p)^n$ 或 $L(p) = p^n$ 在 $[0,1]$ 上的单调性，可知 $\hat{p}_{\text{MLE}} = 0$ 或 $1$，统一写作 $\hat{p}_{\text{MLE}} = k/n$。
\end{example}

\begin{example}[一维高斯分布 $\mathcal{N}(\mu, \sigma^2)$ 的 MLE]
\label{ex:ch5-mle-1dgauss}
设 $x_1, \dots, x_n \stackrel{\text{i.i.d.}}{\sim} \mathcal{N}(\mu, \sigma^2)$，令 $v = \sigma^2 > 0$。写出对数似然：
\[
  \ell(\mu, v)
  = \sum_{i=1}^n \left[-\frac{1}{2}\ln(2\pi) - \frac{1}{2}\ln v - \frac{(x_i - \mu)^2}{2v}\right]
  = -\frac{n}{2}\ln(2\pi) - \frac{n}{2}\ln v - \frac{1}{2v}\sum_{i=1}^n (x_i - \mu)^2.
\]
分别对 $\mu$ 和 $v$ 求偏导并令其为零：
\[
  \frac{\partial \ell}{\partial \mu} = \frac{1}{v}\sum_{i=1}^n (x_i - \mu) = 0
  \;\Longrightarrow\;
  \hat{\mu}_{\text{MLE}} = \frac{1}{n}\sum_{i=1}^n x_i = \bar{x}.
\]
将 $\mu = \bar{x}$ 代入对 $v$ 的偏导方程：
\[
  \frac{\partial \ell}{\partial v} = -\frac{n}{2v} + \frac{1}{2v^2}\sum_{i=1}^n (x_i - \bar{x})^2 = 0
  \;\Longrightarrow\;
  \hat{\sigma}^2_{\text{MLE}} = \hat{v}_{\text{MLE}} = \frac{1}{n}\sum_{i=1}^n (x_i - \bar{x})^2.
\]
请注意：高斯方差的极大似然估计自动给出的是\textbf{除以 $n$} 的 $\hat{\sigma}_n^2$！结合命题~\ref{prop:ch5-sample-var} 立即得到一个重要结论：\textbf{MLE 并不保证无偏}（$\hat{\sigma}^2_{\text{MLE}}$ 有偏差 $-\sigma^2/n$），但随着 $n\to\infty$，偏差迅速趋近于零，因此它是渐近无偏且相合的。
\end{example}

\begin{example}[多元高斯分布 $\mathcal{N}(\mu, \Sigma)$ 的 MLE]
\label{ex:ch5-mle-mvg}
设 $d$ 维向量样本 $x_1, \dots, x_n \stackrel{\text{i.i.d.}}{\sim} \mathcal{N}(\mu, \Sigma)$，其中 $\Sigma \succ 0$ 为对称正定协方差矩阵（第 2 章）。直接将一维结果推广到矩阵形式，其极大似然估计同样简洁优美：
\begin{equation}
  \hat{\mu}_{\text{MLE}} = \bar{x} = \frac{1}{n}\sum_{i=1}^n x_i,
  \qquad
  \widehat{\Sigma}_{\text{MLE}} = \frac{1}{n}\sum_{i=1}^n (x_i - \bar{x})(x_i - \bar{x})^T.
  \label{eq:ch5-mvg-mle}
\end{equation}
注意当样本数 $n \le d$ 时，$\widehat{\Sigma}_{\text{MLE}}$ 的秩至多为 $n-1 < d$，必然是\textbf{奇异（不可逆）矩阵}——这正是在高维小样本下使用高斯判别分析（GDA）或 GMM 时必须对协方差矩阵加对角正则项 $\widehat{\Sigma} + \epsilon I$ 的根本原因！
\end{example}

\begin{extBox}{多元高斯 MLE 的矩阵微分推导（衔接第 1 卷第 5 章迹技巧）}
引入对称正定\textbf{精度矩阵} $\Lambda \triangleq \Sigma^{-1} \succ 0$，利用 $\ln\det(\Sigma^{-1}) = -\ln\det\Sigma$ 与标量的迹恒等式 $a^T \Lambda a = \tr(\Lambda a a^T)$，多元高斯对数似然写为
\[
  \ell(\mu, \Lambda)
  = -\frac{nd}{2}\ln(2\pi) + \frac{n}{2}\ln\det\Lambda - \frac{1}{2}\sum_{i=1}^n (x_i - \mu)^T \Lambda (x_i - \mu).
\]
先对 $\mu$ 求梯度：$\nabla_\mu \ell = \Lambda\sum_{i=1}^n (x_i - \mu) = 0$，因 $\Lambda \succ 0$ 可逆，立得 $\hat{\mu}_{\text{MLE}} = \bar{x}$。

再记散度矩阵 $S_0 \triangleq \sum_{i=1}^n (x_i - \bar{x})(x_i - \bar{x})^T$，利用迹的线性性把二次型之和改写为 $-\frac{1}{2}\tr(\Lambda S_0)$。结合第 1 卷第 5 章的矩阵求导法则 $\nabla_\Lambda \tr(\Lambda S_0) = S_0^T = S_0$ 与雅可比行列式求导公式 $\nabla_\Lambda \ln\det\Lambda = \Lambda^{-T} = \Lambda^{-1}$（对对称正定矩阵 $\Lambda$）：
\[
  \nabla_\Lambda \ell(\bar{x}, \Lambda) = \frac{n}{2}\Lambda^{-1} - \frac{1}{2}S_0 = 0
  \;\Longrightarrow\;
  \widehat{\Sigma}_{\text{MLE}} = \widehat{\Lambda}_{\text{MLE}}^{-1} = \frac{1}{n}S_0 = \frac{1}{n}\sum_{i=1}^n (x_i - \bar{x})(x_i - \bar{x})^T.
\]
一气呵成！用精度矩阵 $\Lambda = \Sigma^{-1}$ 换元避免了对 $\Sigma^{-1}$ 做复杂的矩阵求逆微分，是概率图与高斯建模中极其常用的化简技巧。
\end{extBox}

\begin{failBox}{小样本过拟合反例：连抛 3 次全正面，反面概率真是 0？}
假设你拿到一枚陌生硬币，只抛了 $n=3$ 次，结果碰巧全是正面（$k=3$）。按例~\ref{ex:ch5-mle-bern} 的 MLE 公式，你会毫不犹豫地输出
\[
  \hat{p}_{\text{MLE}} = \frac{3}{3} = 1.000,
  \qquad
  \widehat{\Prob}(\text{反面}) = 1 - \hat{p}_{\text{MLE}} = 0.000.
\]
这个看似“完全忠于数据”的估计在 AI 部署中是一场\textbf{灾难}：
\begin{enumerate}
  \item \textbf{零概率崩溃（Zero-Probability Disaster）}：只要测试集里第 1 次出现反面，模型给出的对数似然就是 $\ln(0) = -\infty$，交叉熵损失与语言模型困惑度（Perplexity）瞬间炸成无穷大！
  \item \textbf{真实 NLP / 朴素贝叶斯事故}：在垃圾邮件过滤（第 1 章朴素贝叶斯）或 $n$-gram 语言模型中，词表高达数万词，训练集哪怕再大也总会有低频词从未出现过（$k=0$）。若用裸 MLE 把未登录词概率估成 $0$，整个句子的连乘联合概率会立刻被这一个词一票否决归零！
\end{enumerate}
根因非常清晰：MLE 是一种\textbf{唯数据论}的估计器，当样本量 $n$ 很小时，它会把抽样中的偶然波动（噪声）当成绝对真理来死记硬背——这就是统计学习中\textbf{过拟合（Overfitting）}最原始的形态。
\end{failBox}

% ====================================================================
\section{贝叶斯估计、共轭先验与最大后验（MAP）}
\label{sec:ch5-bayes-map}

怎么拯救“连抛 3 次全正面就断定反面永不出现”的 MLE？只需要让模型保留一点人类常识：在抛硬币之前，我们心里本来就有一个模糊的\textbf{先验信念（Prior）}——正常硬币正反面大概率接近 $0.5$，不太可能极端到 $0$ 或 $1$；除非你给我看连续 50 次正面，否则单凭 3 次正面不足以让我相信反面概率为零！

这正是\textbf{贝叶斯参数估计}的思想飞跃：不再把参数 $\theta$ 当成冰冷的未知常数，而是把它也看作一个随机变量，用先验分布 $p(\theta)$ 表达见数据前的信念；看到数据 $\mathcal{D}$ 后，用第 1 章的 Bayes 定理把先验更新为\textbf{后验分布（Posterior）}：
\begin{equation}
  p(\theta \mid \mathcal{D})
  = \frac{p(\mathcal{D} \mid \theta)\,p(\theta)}{p(\mathcal{D})}
  = \frac{p(\mathcal{D} \mid \theta)\,p(\theta)}{\int p(\mathcal{D} \mid \theta')\,p(\theta')\,\mathrm{d}\theta'}
  \propto p(\mathcal{D} \mid \theta)\,p(\theta).
  \label{eq:ch5-bayes-post}
\end{equation}
这里分母 $p(\mathcal{D})$ 称为\textbf{边缘似然（Evidence）}，它只依赖数据 $\mathcal{D}$ 而与待估参数 $\theta$ 无关，仅起到把后验密度积分为 $1$ 的归一化作用。

拿到完整后验分布 $p(\theta \mid \mathcal{D})$ 后，如果工程上只需要输出一个具体的点估计值，有两种核心做法：

\begin{defBox}
\begin{definition}[最大后验估计（MAP）与全贝叶斯后验均值]
\label{def:ch5-map-pm}
给定先验分布 $p(\theta)$ 与数据似然 $p(\mathcal{D} \mid \theta)$：
\begin{enumerate}
  \item \textbf{最大后验估计（Maximum A Posteriori, MAP）}取后验分布的\textbf{众数（峰值）}：
  \begin{equation}
    \hat{\theta}_{\text{MAP}}
    \triangleq \argmax_{\theta} p(\theta \mid \mathcal{D})
    = \argmax_{\theta}\Big[\underbrace{\ln p(\mathcal{D} \mid \theta)}_{\text{对数似然（拟合数据）}} + \underbrace{\ln p(\theta)}_{\text{对数先验（正则惩罚）}}\Big].
    \label{eq:ch5-map}
  \end{equation}
  因为分母 $p(\mathcal{D})$ 不含 $\theta$，求解 MAP 根本不需要计算高维归一化积分！
  \item \textbf{全贝叶斯后验均值（Posterior Mean / MMSE 估计）}取后验分布的\textbf{期望}：
  \[
    \hat{\theta}_{\text{Bayes}} \triangleq \E[\theta \mid \mathcal{D}] = \int \theta\, p(\theta \mid \mathcal{D})\,\mathrm{d}\theta.
  \]
  可以证明它使后验均方误差 $\E[(\theta - \hat{\theta})^2 \mid \mathcal{D}]$ 达到最小，且在预测新样本 $x_{\text{new}}$ 时，全贝叶斯做法甚至会对整个后验积分求边缘预测分布 $p(x_{\text{new}}\mid \mathcal{D}) = \int p(x_{\text{new}}\mid \theta)\,p(\theta\mid \mathcal{D})\,\mathrm{d}\theta$。
\end{enumerate}
\end{definition}
\end{defBox}

表~\ref{tab:ch5-mle-map-bayes} 简明对比了 MLE、MAP 与全贝叶斯三种范式在工程上的定位差异：当参数维度高达数百万（如神经网络）时，全贝叶斯高维积分几乎无法精确求解，而 MAP 只需在 MLE 的损失函数后面加上一项 $-\ln p(\theta)$，就能以几乎零额外计算成本获得强大的正则化保护！

\begin{table}[htbp]
\centering
\caption{MLE、MAP 与全贝叶斯估计的核心对比}
\label{tab:ch5-mle-map-bayes}
\footnotesize
\begin{tabularx}{\linewidth}{@{}p{21mm}p{26mm}p{24mm}p{24mm}X@{}}
\toprule
\textbf{估计范式} & \textbf{优化/计算目标} & \textbf{是否用先验} & \textbf{输出形式} & \textbf{核心优势与 AI 典型场景} \\
\midrule
\textbf{MLE} & $\max_\theta \ln p(\mathcal{D}\mid\theta)$ & 否（等价均匀先验） & 单一点估计 & 计算最简、大样本一致；小样本易过拟合（无正则回归、裸词频） \\
\textbf{MAP} & $\max_\theta [\ln p(\mathcal{D}\mid\theta)+\ln p(\theta)]$ & 是 & 后验峰值（点估计） & 免高维积分、直接对应正则化损失（岭回归、Lasso、权重衰减） \\
\textbf{全贝叶斯} & 计算 $p(\theta\mid\mathcal{D})$ 并积分 & 是 & 完整后验分布/均值 & 自动量化参数不确定性；共轭时有闭式解，高维需采样/变分近似 \\
\bottomrule
\end{tabularx}
\end{table}

\subsection{Beta-Bernoulli 共轭先验：伪计数直觉与实测对比}
\label{subsec:ch5-beta-bern}

有没有办法让全贝叶斯后验分布的推导\textbf{不用算任何积分}，还能一眼看出解析形式？答案就是\textbf{共轭先验（Conjugate Prior）}：如果先验分布 $p(\theta)$ 与似然 $p(\mathcal{D}\mid\theta)$ 相乘后，得到的后验分布 $p(\theta\mid\mathcal{D})$ 仍然属于与先验\textbf{相同的分布族}（仅仅更新了超参数），我们就称该先验与该似然是共轭的。

回忆第 2 章的 $\mathrm{Beta}(\alpha, \beta)$ 分布（定义在 $p \in [0,1]$ 上，超参数 $\alpha > 0, \beta > 0$）：
\[
  p(p; \alpha, \beta) = \frac{1}{B(\alpha, \beta)}\, p^{\alpha - 1}(1 - p)^{\beta - 1},
  \qquad
  \E[p] = \frac{\alpha}{\alpha + \beta},
  \qquad
  \mathrm{Mode}(p) = \frac{\alpha - 1}{\alpha + \beta - 2}\;\;(\alpha, \beta > 1).
\]
把 $\mathrm{Beta}(\alpha, \beta)$ 乘上 Bernoulli 似然，奇迹立刻发生！

\begin{thmBox}
\begin{theorem}[Beta-Bernoulli 共轭更新定理]
\label{thm:ch5-beta-bern}
设硬币正面概率 $p$ 服从先验 $\mathrm{Beta}(\alpha, \beta)$（$\alpha, \beta > 0$）。独立抛掷 $n$ 次观测到 $k$ 次正面与 $n-k$ 次反面，则后验分布精确为
\begin{equation}
  p \mid (n, k) \;\sim\; \mathrm{Beta}(\alpha + k,\; \beta + n - k).
  \label{eq:ch5-beta-post}
\end{equation}
由此直接读出全贝叶斯后验均值与 MAP 估计（当 $\alpha + k > 1$ 且 $\beta + n - k > 1$ 时）：
\begin{equation}
  \hat{p}_{\text{Bayes}} = \frac{k + \alpha}{n + \alpha + \beta},
  \qquad
  \hat{p}_{\text{MAP}} = \frac{k + \alpha - 1}{n + \alpha + \beta - 2}.
  \label{eq:ch5-beta-est}
\end{equation}
\end{theorem}
\end{thmBox}

\begin{proof}
依据式~\eqref{eq:ch5-bayes-post}，略去所有不含 $p$ 的归一化常数：
\[
  p(p \mid n, k)
  \propto p(n, k \mid p)\, p(p)
  \propto \big[p^k (1 - p)^{n-k}\big] \cdot \big[p^{\alpha - 1}(1 - p)^{\beta - 1}\big]
  = p^{(\alpha + k) - 1} (1 - p)^{(\beta + n - k) - 1}.
\]
对比 Beta 密度核函数的指数形式，由于概率密度函数在 $[0,1]$ 上积分为 $1$ 的归一化常数唯一，故后验分布必然严格等于 $\mathrm{Beta}(\alpha + k,\, \beta + n - k)$。代入 Beta 分布的均值与众数公式即得式~\eqref{eq:ch5-beta-est}。
\end{proof}

式~\eqref{eq:ch5-beta-est} 赋予了 Beta 先验极为通透的\textbf{伪计数（Pseudo-Count）直觉}：
\begin{itemize}
  \item 在 MAP 眼里，设定 $\mathrm{Beta}(\alpha, \beta)$ 先验，完全等同于\textbf{在真实数据之外凭空预先塞进了 $\alpha - 1$ 次虚拟正面与 $\beta - 1$ 次虚拟反面}！特别地，当取对称平滑先验 $\mathrm{Beta}(2, 2)$ 时，$\alpha - 1 = \beta - 1 = 1$，MAP 公式变成 $\frac{k + 1}{n + 2}$——这正是 NLP 和朴素贝叶斯中鼎鼎大名的\textbf{拉普拉斯加一平滑（Laplace / Add-1 Smoothing）}！
  \item 在后验均值眼里，即便取毫无偏向的均匀先验 $\mathrm{Beta}(1, 1)$，MAP 退化为裸 MLE $k/n$，但全贝叶斯后验均值依然给出 $\frac{k+1}{n+2}$（历史著名的\textbf{拉普拉斯继承法则}，例如太阳连续升起 $n$ 天后明天继续升起的预测概率），天然杜绝零概率灾难！
\end{itemize}

让我们把刚才“连抛 $n=3$ 次全是正面（$k=3$）”的棘手反例代入 $\mathrm{Beta}(2, 2)$ 先验（这与配套脚本 \texttt{code/vol2\_experiments.py} 的 \texttt{exp5\_mle\_map\_bootstrap} 实测完全一致）：
\begin{itemize}
  \item \textbf{裸 MLE}：$\hat{p}_{\text{MLE}} = \frac{3}{3} = \mathbf{1.0000}$（断定反面永不出现）；
  \item \textbf{Beta(2,2) 后验分布}：$\mathrm{Beta}(2+3,\, 2+0) = \mathrm{Beta}(5, 2)$，密度函数为 $30\,p^4(1-p)$；
  \item \textbf{MAP 估计}：$\hat{p}_{\text{MAP}} = \frac{3 + 2 - 1}{3 + 2 + 2 - 2} = \frac{4}{5} = \mathbf{0.8000}$；
  \item \textbf{全贝叶斯后验均值}：$\hat{p}_{\text{Bayes}} = \frac{3 + 2}{3 + 2 + 2} = \frac{5}{7} \approx \mathbf{0.7143}$。
\end{itemize}
图~\ref{fig:ch5-beta-update} 直观绘制了这次贝叶斯更新中先验、归一化似然与后验三条曲线的几何拉锯：后验 $\mathrm{Beta}(5,2)$ 正好是先验 $\mathrm{Beta}(2,2)$ 与似然 $p^3$ 逐点相乘的折中产物——它在 $p=1$ 处被先验强制拉回 $0$，从而既尊重了“连出 3 正面”的偏高证据，又保留了反面出现的合理余地。

\begin{figure}[htbp]
\centering
\begin{tikzpicture}
\begin{axis}[
  width=0.84\linewidth,
  height=6.2cm,
  xmin=0, xmax=1.03,
  ymin=0, ymax=4.35,
  xlabel={硬币正面概率参数 $p$},
  ylabel={概率密度（似然已归一化以便同图对比）},
  xlabel style={font=\small\sffamily, yshift=-1mm},
  ylabel style={font=\small\sffamily},
  tick label style={font=\footnotesize},
  grid=major,
  grid style={dashed, gray!25},
  legend pos=north west,
  legend style={font=\footnotesize\sffamily, fill=white, fill opacity=0.92, draw=tealLine},
  axis lines=left,
  enlargelimits=false,
  clip=false
]
  % Prior Beta(2,2): 6*x*(1-x)
  \addplot[navy, very thick, dashed, domain=0:1, samples=120] {6*x*(1-x)};
  \addlegendentry{先验 $\mathrm{Beta}(2,2)$（均值 $0.50$）}

  % Normalized likelihood for k=3, n=3: 4*x^3 (Beta(4,1))
  \addplot[ochre, very thick, dotted, domain=0:1, samples=120] {4*x^3};
  \addlegendentry{归一化似然 $4p^3$（$n=3, k=3$）}

  % Posterior Beta(5,2): 30*x^4*(1-x)
  \addplot[teal, ultra thick, domain=0:1, samples=140] {30*x^4*(1-x)};
  \addlegendentry{后验 $\mathrm{Beta}(5,2)$}

  % Soft fill under posterior
  \addplot[teal, fill=tealLight, fill opacity=0.45, draw=none, domain=0:1, samples=100] {30*x^4*(1-x)} \closedcycle;

  % Vertical reference lines and non-overlapping annotations
  % Posterior Mean = 5/7 = 0.7143, density = 30*(5/7)^4*(2/7) = 2.230
  \draw[tealDark, thin, densely dashed] (axis cs:0.7143, 0) -- (axis cs:0.7143, 2.230);
  \fill[tealDark] (axis cs:0.7143, 2.230) circle (2.2pt);
  \node[font=\footnotesize\sffamily\bfseries, color=tealDark, align=center, anchor=east]
    at (axis cs:0.68, 2.55) {后验均值\\$5/7 \approx 0.714$};
  \draw[-{Stealth[length=1.8mm]}, thin, tealDark] (axis cs:0.685, 2.40) -- (axis cs:0.710, 2.26);

  % MAP = 0.80, density = 30*(0.8)^4*(0.2) = 2.4576
  \draw[teal, thin, densely dashed] (axis cs:0.80, 0) -- (axis cs:0.80, 2.4576);
  \fill[teal] (axis cs:0.80, 2.4576) circle (2.4pt);
  \node[font=\footnotesize\sffamily\bfseries, color=teal, align=center, anchor=south]
    at (axis cs:0.76, 3.05) {MAP 峰值 $= 0.800$};
  \draw[-{Stealth[length=1.8mm]}, thin, teal] (axis cs:0.78, 3.02) -- (axis cs:0.798, 2.52);

  % MLE = 1.00, normalized likelihood peak = 4.0
  \fill[ochre] (axis cs:1.0, 4.0) circle (2.4pt);
  \node[font=\footnotesize\sffamily\bfseries, color=ochre, align=right, anchor=east]
    at (axis cs:0.93, 3.92) {MLE $= 1.000$\\（过拟合崩溃点）};
  \draw[-{Stealth[length=1.8mm]}, thin, ochre] (axis cs:0.935, 3.96) -- (axis cs:0.988, 3.99);
\end{axis}
\end{tikzpicture}
\caption{“连抛 3 次全正面”的 Beta-Bernoulli 共轭更新：$\mathrm{Beta}(2,2)$ 先验把极端的 MLE 峰值 $1.000$ 拉回至 $\mathrm{MAP}=0.800$ 与后验均值 $0.7143$，彻底消除零概率灾难。}
\label{fig:ch5-beta-update}
\end{figure}

读者也许会担心：“先验是我主观选的，万一我先验猜错了怎么办？”把后验均值写成一个加权平均就能彻底打消顾虑：
\[
  \hat{p}_{\text{Bayes}}
  = \frac{k + \alpha}{n + \alpha + \beta}
  = \underbrace{\frac{n}{n + \alpha + \beta}}_{\text{数据权重}\to 1}\,\hat{p}_{\text{MLE}}
  \;+\;
  \underbrace{\frac{\alpha + \beta}{n + \alpha + \beta}}_{\text{先验权重}\to 0}\,\frac{\alpha}{\alpha + \beta}.
\]
当样本量 $n$ 很小时（$n=3$），先验权重高达 $4/7$，果断保护模型免受噪声绑架；一旦样本量增大到 $n=1000$，先验权重骤降为 $4/1004 < 0.4\%$，\textbf{海量数据会彻底淹没任何合理的有限先验}，MAP 与后验均值双双收敛到 MLE 与真值 $p$！

\subsection{高斯-高斯共轭：为第 6 章岭回归铺路}
\label{subsec:ch5-gauss-gauss}

除了离散硬币，连续变量中最核心的共轭组合是\textbf{高斯先验 + 高斯似然}。设观测噪声方差 $\sigma^2$ 已知，我们想估计未知均值 $\mu$：给定零均值高斯先验 $\mu \sim \mathcal{N}(0, \tau^2)$ 与样本 $x_1, \dots, x_n \stackrel{\text{i.i.d.}}{\sim} \mathcal{N}(\mu, \sigma^2)$。写出后验对数密度（略去与 $\mu$ 无关的常数）：
\begin{align*}
  \ln p(\mu \mid x_1, \dots, x_n)
  &= -\frac{1}{2\sigma^2}\sum_{i=1}^n (x_i - \mu)^2 - \frac{\mu^2}{2\tau^2} + \text{const} \\
  &= -\frac{1}{2}\left(\frac{n}{\sigma^2} + \frac{1}{\tau^2}\right)\mu^2 + \left(\frac{n\bar{x}}{\sigma^2}\right)\mu + \text{const}.
\end{align*}
对 $\mu$ 配完全平方，立即看出后验分布仍然是高斯分布 $\mu \mid \mathcal{D} \sim \mathcal{N}(\mu_n, \tau_n^2)$，且满足极为优雅的\textbf{精度可加律}：
\begin{equation}
  \frac{1}{\tau_n^2} = \frac{1}{\tau^2} + \frac{n}{\sigma^2},
  \qquad
  \hat{\mu}_{\text{MAP}} = \hat{\mu}_{\text{Bayes}} = \mu_n
  = \frac{\frac{n}{\sigma^2}\bar{x}}{\frac{1}{\tau^2} + \frac{n}{\sigma^2}}
  = \frac{n}{n + \sigma^2/\tau^2}\,\bar{x}.
  \label{eq:ch5-gauss-conj}
\end{equation}
由于高斯分布左右对称，后验众数（MAP）与后验均值完全重合！仔细端详式~\eqref{eq:ch5-gauss-conj} 中的收缩因子 $\frac{n}{n + \lambda}$（其中 $\lambda \triangleq \sigma^2/\tau^2$）：
\begin{itemize}
  \item 它自动把样本均值 $\bar{x}$ 向零收缩；
  \item 当观测噪声很大（$\sigma^2$ 大）或先验非常笃定 $\mu$ 靠近零（$\tau^2$ 小）时，$\lambda = \sigma^2/\tau^2$ 变大，收缩力度更强；
  \item 当样本量 $n \to \infty$ 时，收缩因子趋于 $1$，恢复无偏均值 $\bar{x}$。
\end{itemize}

\begin{insightBox}{直通第 6 章第 6.1 节：高斯先验 MAP $\equiv$ $L_2$ 岭回归！}
请记住式~\eqref{eq:ch5-gauss-conj} 里刚刚诞生的比值 $\lambda = \sigma^2 / \tau^2$！在第 6 章一开篇，我们把标量均值 $\mu$ 升级为高维回归权重向量 $w \sim \mathcal{N}(0, \tau^2 I)$，对数先验 $-\ln p(w) = \frac{1}{2\tau^2}\norm{w}_2^2 + \text{const}$ 会原封不动地长成第 1 卷已经用过的\textbf{岭回归 $L_2$ 正则项 $\frac{\lambda}{2}\norm{w}_2^2$}，而神秘的超参数 $\lambda$ 恰好就是噪声方差与先验信号方差之比 $\sigma^2/\tau^2$！
\end{insightBox}

% ====================================================================
\section{非参数重采样：Bootstrap 置信区间}
\label{sec:ch5-bootstrap}

假设你刚刚用第 6 章的逻辑回归在 $n=180$ 条验证样本上跑出了 $153$ 条预测正确，即验证集准确率 $\widehat{\text{Acc}} = 153/180 = \mathbf{85.00\%}$，$F_1 = 0.8439$，ROC-AUC $= 0.9160$。紧接着导师或审稿人追问你一个致命问题：
\begin{quote}
  \itshape “如果换一批同样来自真实世界的 180 条测试样本，你的准确率会波动多大？你说你的模型比基线 84.0\% 高 1 个百分点，这到底是真实进步，还是随机抽样的运气？”
\end{quote}

这就需要给点估计配上一个\textbf{置信区间（Confidence Interval, CI）}。在传统数理统计课上，你只能对样本均值这类极简单的统计量套用正态近似或 $t$ 分布查表；可面对 $F_1$ 分数、ROC-AUC、中位数推理延迟、甚至复杂深度学习评测指标（如 BLEU、mAP）时，它们根本没有闭式的有限样本抽样分布！

1979 年统计学家 Bradley Efron 提出了一个堪称“用算力换公式”的天才构想——\textbf{自助法（Bootstrap）}，其核心哲学称为\textbf{代入原则（Plug-in Principle）}：
\begin{itemize}
  \item \textbf{理想世界（无法实现）}：从未知的真实总体分布 $P$ 中反复独立抽取 $B=5000$ 份大小为 $n$ 的测试集，算出 $5000$ 个指标值，直接看它们的波动范围；
  \item \textbf{Bootstrap 世界（可计算）}：既然手里唯一的样本集 $\mathcal{D} = \{z_1, \dots, z_n\}$ 所构成的\textbf{经验分布 $\hat{P}_n$}（每个观测样本权重 $1/n$）是对未知总体 $P$ 最好的非参数近似，那就改为\textbf{从 $\mathcal{D}$ 中有放回地随机抽取 $n$ 个样本}来模拟一次新抽样，重复 $B$ 次！
\end{itemize}

\begin{remark}[为什么必须“有放回”？顺便预埋第 4 卷袋外误差]
如果无放回抽取 $n$ 个样本，每次抽出的都是原数据集的完全拷贝，算出的指标纹丝不动（方差恒为 0）。而\textbf{有放回}抽取 $n$ 次时，某条特定样本一次都没被抽中的概率是 $(1 - 1/n)^n \xrightarrow{n\to\infty} e^{-1} \approx 36.8\%$，因此每份 Bootstrap 重采样集平均包含约 $63.2\%$ 的原始不重复样本（部分样本出现 $2$ 次或多次）。这个神奇的 $36.8\%$ 将在第 4 卷随机森林的\textbf{袋外验证（Out-of-Bag, OOB）}中再次登场！
\end{remark}

\begin{defBox}
\begin{definition}[百分位 Bootstrap 置信区间算法（Percentile Bootstrap CI）]
\label{def:ch5-bootstrap}
给定含 $n$ 条独立样本的数据集 $\mathcal{D} = \{z_1, \dots, z_n\}$、任意待评统计量 $T(\mathcal{D})$（如准确率、$F_1$、AUC 等）、重采样次数 $B$（工程上常取 $B = 2000 \sim 10000$）与置信水平 $1 - \alpha$（如 $95\%$ 对应 $\alpha = 0.05$）：
\begin{enumerate}
  \item \textbf{有放回重采样}：对 $b = 1, \dots, B$，独立地从索引 $\{1, \dots, n\}$ 中有放回随机抽取 $n$ 个索引，构造第 $b$ 份自助样本集 $\mathcal{D}^{*(b)}$；
  \item \textbf{重算统计量}：在每份自助样本上计算 $\hat{\theta}^{*(b)} = T(\mathcal{D}^{*(b)})$，得到经验抽样分布 $\{\hat{\theta}^{*(1)}, \dots, \hat{\theta}^{*(B)}\}$；
  \item \textbf{估计标准误}：自助标准误 $\widehat{\mathrm{SE}}_{\text{boot}} = \mathrm{std}\big(\hat{\theta}^{*(1)}, \dots, \hat{\theta}^{*(B)}\big)$；
  \item \textbf{截取分位数}：将 $B$ 个自助估计值从小到大排序，取 $\frac{\alpha}{2}$ 与 $1 - \frac{\alpha}{2}$ 分位数（$95\%$ CI 即 $2.5\%$ 与 $97.5\%$ 分位数），输出区间 $\big[q_{\alpha/2},\; q_{1-\alpha/2}\big]$。
\end{enumerate}
\end{definition}
\end{defBox}

\begin{figure}[htbp]
\centering
\resizebox{0.97\linewidth}{!}{%
\begin{tikzpicture}[
  node distance=5.5mm and 7.5mm,
  box/.style={
    rectangle, rounded corners=1.8mm, draw=teal, thick, fill=tealLight,
    minimum height=14mm, align=center, font=\sffamily\footnotesize, text=ink, inner sep=5pt
  },
  rep/.style={
    rectangle, rounded corners=1.4mm, draw=navy, thick, fill=navyLight,
    text width=32mm, minimum height=7.2mm, align=center, font=\sffamily\scriptsize, text=ink, inner sep=3pt
  },
  res/.style={
    rectangle, rounded corners=1.8mm, draw=ochre, very thick, fill=ochreLight,
    text width=44mm, minimum height=17mm, align=center, font=\sffamily\footnotesize, text=ink, inner sep=5pt
  },
  arr/.style={-{Stealth[length=2.1mm]}, thick, teal}
]
  \node[box, text width=31mm] (orig) {\textbf{原始验证集 $\mathcal{D}$}\\$n=180$ 条样本\\命中 $153$ 条 ($85.00\%$)};

  \node[rep, right=12mm of orig, yshift=11mm] (b1) {重采样 $\mathcal{D}^{*(1)}$ ($n=180$)\\$\Rightarrow \hat{\theta}^{*(1)} = 82.78\%$};
  \node[rep, right=12mm of orig, yshift=0mm]  (b2) {重采样 $\mathcal{D}^{*(2)}$ ($n=180$)\\$\Rightarrow \hat{\theta}^{*(2)} = 86.11\%$};
  \node[rep, right=12mm of orig, yshift=-11mm](bB) {重采样 $\mathcal{D}^{*(5000)}$ ($n=180$)\\$\Rightarrow \hat{\theta}^{*(5000)} = 84.44\%$};
  \node[font=\footnotesize\sffamily, text=slate] at ($(b2.south)!0.5!(bB.north)$) {$\vdots$};

  \node[res, right=12mm of b2] (ci) {\textbf{\color{ochre}百分位 95\% 置信区间}\\汇总 $B=5000$ 个 $\hat{\theta}^{*(b)}$\\截取 $[2.5\%,\, 97.5\%]$ 分位：\\\textbf{$[80.00\%,\; 90.00\%]$}\\{\scriptsize $\widehat{\mathrm{SE}}_{\text{boot}} = 2.63\%$}};

  \draw[arr] (orig.east) -- ++(4mm,0) |- (b1.west);
  \draw[arr] (orig.east) -- (b2.west);
  \draw[arr] (orig.east) -- ++(4mm,0) |- (bB.west);
  \node[font=\scriptsize\sffamily, text=tealDark, above=1.5mm] at ($(orig.east)+(4mm,11mm)$) {有放回抽样};

  \draw[arr] (b1.east) -| ($(ci.west)+(-4mm,0)$) -- (ci.west);
  \draw[arr] (b2.east) -- (ci.west);
  \draw[arr] (bB.east) -| ($(ci.west)+(-4mm,0)$) -- (ci.west);
\end{tikzpicture}}
\caption{非参数百分位 Bootstrap 流程：以原始验证集经验分布为代理，并行重采样 $B=5000$ 次，直接读出 $153/180$ 验证准确率的 95\% 置信区间 $[80.00\%, 90.00\%]$。}
\label{fig:ch5-bootstrap-flow}
\end{figure}

运行 \texttt{code/vol2\_experiments.py} 中的 \texttt{exp5\_mle\_map\_bootstrap}（$n=180$ 条验证预测，153 对 27 错，$B=5000$，随机种子 \texttt{2026}），实测结果如下：
\[
  \widehat{\text{Acc}} = 85.00\%, \qquad
  \text{Bootstrap 95\% CI} = \mathbf{[80.00\%,\; 90.00\%]}, \qquad
  \widehat{\mathrm{SE}}_{\text{boot}} = \mathbf{2.63\%}.
\]
（作为交叉验证，二项分布理论标准误为 $\sqrt{0.85 \times 0.15 / 180} = 2.66\%$，两者高度吻合！）

\begin{failBox}{科研与工程避坑：CI 重叠莫吹牛，且 Bootstrap 并非万能}
\begin{enumerate}
  \item \textbf{别把 1\% 的随机涨落当成 SOTA 突破}：上面的实测告诉我们，在 $n=180$ 的验证集上，准确率的 95\% 置信区间宽达整整 $10$ 个百分点（$\pm 5\%$）！如果模型 A 测出 $85\%$、模型 B 测出 $86\%$，两者的置信区间大面积重叠——\textbf{绝不能直接宣称“模型 B 显著优于模型 A”}！正确做法是增大测试集样本量，或在同一批测试样本上做\textbf{配对 Bootstrap（Paired Bootstrap）}检验准确率差值 $\Delta = \text{Acc}_B - \text{Acc}_A$ 的 95\% CI 是否严格不含 $0$（详见习题 5.8 及第 7 卷实验统计检验专题）。
  \item \textbf{Bootstrap 的适用边界}：Bootstrap 依赖“经验分布 $\hat{P}_n$ 足以代表总体 $P$”。当样本量极小（如 $n=3$ 连出 3 正面，所有重采样全都是 $\{1,1,1\}$，CI 坍缩成毫无意义的 $[100\%, 100\%]$！）或统计量属于极端边界量（如样本最大值 $\max_i X_i$）以及样本存在强时间序列相关性时，朴素百分位 Bootstrap 会失效。
\end{enumerate}
\end{failBox}

\subsection{纯 NumPy 实现：Beta-Bernoulli 更新与向量化 Bootstrap}
\label{subsec:ch5-code}

代码~\ref{lst:ch5-beta-bootstrap} 给出极简、无依赖且全向量化的 NumPy 实现：既可复现本章的 Beta-Bernoulli 对比与 $153/180$ 准确率 95\% 置信区间，也可无缝替换为任意自定义指标函数（例如第 6 章的 $F_1$ 分数）。

\begin{lstlisting}[caption={纯 NumPy 实现 Beta-Bernoulli 共轭更新与向量化百分位 Bootstrap 置信区间}, label={lst:ch5-beta-bootstrap}]
import numpy as np

def beta_bernoulli_summary(k: int, n: int, alpha: float = 2.0, beta: float = 2.0):
    """计算 Bernoulli 观测下的 MLE、Beta 共轭 MAP 与后验均值。"""
    mle = k / n
    a_post, b_post = alpha + k, beta + (n - k)
    post_mean = a_post / (a_post + b_post)
    map_est = (a_post - 1.0) / (a_post + b_post - 2.0) if (a_post > 1 and b_post > 1) else mle
    return {"MLE": mle, "MAP": map_est, "PostMean": post_mean, "post_params": (a_post, b_post)}

def bootstrap_ci_accuracy(hits: np.ndarray, B: int = 5000, alpha: float = 0.05, seed: int = 2026):
    """向量化百分位 Bootstrap：计算二值命中序列 (0/1) 准确率的 (1-alpha) 置信区间。"""
    rng = np.random.default_rng(seed)
    n = len(hits)
    # 一次性抽取 (B, n) 矩阵：比 Python for 循环快约 30 倍
    boot_samples = rng.choice(hits, size=(B, n), replace=True)
    boot_accs = boot_samples.mean(axis=1) * 100.0
    lo, hi = np.percentile(boot_accs, [100.0 * alpha / 2.0, 100.0 * (1.0 - alpha / 2.0)])
    return float(hits.mean() * 100.0), float(lo), float(hi), float(boot_accs.std())

def bootstrap_ci_general(y_true: np.ndarray, y_pred: np.ndarray, metric_fn, B: int = 2000, seed: int = 2026):
    """通用指标（如 F1、Precision、Recall）的百分位 Bootstrap 95% 置信区间。"""
    rng = np.random.default_rng(seed)
    n = len(y_true)
    vals = np.empty(B, dtype=np.float64)
    for b in range(B):
        idx = rng.integers(0, n, size=n)
        vals[b] = metric_fn(y_true[idx], y_pred[idx])
    lo, hi = np.percentile(vals, [2.5, 97.5])
    return float(metric_fn(y_true, y_pred)), float(lo), float(hi), float(vals.std())

if __name__ == "__main__":
    res = beta_bernoulli_summary(k=3, n=3, alpha=2.0, beta=2.0)
    print(f"3/3 heads: MLE={res['MLE']:.4f}, MAP={res['MAP']:.4f}, PostMean={res['PostMean']:.4f}")
    # 实测输出: MLE=1.0000, MAP=0.8000, PostMean=0.7143

    hits = np.array([1] * 153 + [0] * 27)
    acc, lo, hi, se = bootstrap_ci_accuracy(hits, B=5000, seed=2026)
    print(f"Acc={acc:.2f}%, 95% CI=[{lo:.2f}%, {hi:.2f}%], boot_std={se:.2f}%")
    # 实测输出: Acc=85.00%, 95% CI=[80.00%, 90.00%], boot_std=2.63%
\end{lstlisting}

% ====================================================================
\section{本章习题}
\label{sec:ch5-exercises}

\exTier{A. 概念辨析与反例查错}{击破统计估计与置信区间中的直觉陷阱}
\begin{enumerate}[label=\textbf{5.\arabic*}]
  \item \textbf{判断正误并纠错}：请判断以下三句话的真假，并分别用本章定理或具体反例给出严密理由：
  (a) “在所有估计量中，无偏估计量在均方误差（MSE）意义下一定优于有偏估计量。”
  (b) “极大似然估计（MLE）永远是无偏估计量。”
  (c) “在 $n=180$ 条验证样本上，模型 B 的准确率 $86.0\%$ 高于模型 A 的 $85.0\%$，因此可以断定模型 B 在真实数据分布上显著优于模型 A。”

  \item \textbf{Bootstrap 失效反例诊断}：
  (a) 在 $n=3$ 次抛硬币全为正面的观测序列 $\{1, 1, 1\}$ 上运行百分位 Bootstrap（$B=5000$），得到的 95\% 置信区间是什么？为什么这与直观常识严重违背？对比 $\mathrm{Beta}(2,2)$ 先验是如何化解该问题的；
  (b) 若 $X_1, \dots, X_n \stackrel{\text{i.i.d.}}{\sim} \mathrm{Uniform}(0, \theta)$，请定性说明为何用朴素 Bootstrap 估计样本最大值 $T = \max_i X_i$ 的置信区间同样会失效（提示：重采样样本的最大值永远不可能超过原样本最大值，且恰好等于原最大值的概率是多少？）。
\end{enumerate}

\exTier{B. 具体小例手算}{亲手推演常见分布的 MLE 与共轭顺序更新}
\begin{enumerate}[label=\textbf{5.\arabic*}, resume]
  \item \textbf{指数分布与均匀分布的 MLE 手算}：
  (a) 设 $X_1, \dots, X_n \stackrel{\text{i.i.d.}}{\sim} \mathrm{Exp}(\lambda)$（密度 $p(x; \lambda) = \lambda e^{-\lambda x},\; x \ge 0,\; \lambda > 0$）。写出对数似然并求 $\hat{\lambda}_{\text{MLE}}$；当 $n=4$、观测为 $\{1, 2, 3, 2\}$ 时算出具体值；
  (b) 设 $X_1, \dots, X_n \stackrel{\text{i.i.d.}}{\sim} \mathrm{Uniform}(0, \theta)$（密度 $p(x;\theta) = \frac{1}{\theta}\mathbf{1}_{[0,\theta]}(x)$）。求 $\hat{\theta}_{\text{MLE}}$，说明为何此处\textbf{不能}简单令 $\ell'(\theta)=0$ 求解；并已知 $M_n = \max_i X_i$ 的密度为 $f_{M_n}(t) = \frac{n t^{n-1}}{\theta^n}$（$0\le t\le\theta$），计算 $n=3$ 时 $\hat{\theta}_{\text{MLE}}$ 的期望与偏差。

  \item \textbf{Beta-Bernoulli 共轭的顺序更新性质}：某新上线推荐位采用先验 $p \sim \mathrm{Beta}(2, 8)$（先验均值 $0.20$）。第一批曝光 $10$ 次获得 $4$ 次点击；第二批继续曝光 $90$ 次获得 $26$ 次点击。
  (a) 计算第一批数据后的后验分布、MLE、MAP 与全贝叶斯后验均值；
  (b) 验证“以第一批后验作为新先验再吸收第二批数据”与“直接以原始先验 $\mathrm{Beta}(2,8)$ 一次性吸收全部 $100$ 次曝光（$30$ 次点击）”得到\textbf{完全相同}的后验分布，并对比 $n=10$ 与 $n=100$ 时 MLE 与 MAP 的差距变化。
\end{enumerate}

\exTier{C. AI 核心推演与实用证明}{夯实收缩估计量与共轭推导功底}
\begin{enumerate}[label=\textbf{5.\arabic*}, resume]
  \item \textbf{收缩估计量以微偏换低 MSE（正则化前奏）}：设 $X_1, \dots, X_n \stackrel{\text{i.i.d.}}{\sim} \mathcal{N}(\mu, \sigma^2)$（$\mu \neq 0$），考虑向零收缩的线性估计量 $\hat{\mu}_\alpha = \alpha \bar{X}$（常数 $\alpha \in (0, 1]$）。
  (a) 计算 $\Bias(\hat{\mu}_\alpha)$、$\Var(\hat{\mu}_\alpha)$，并用定理~\ref{thm:ch5-mse-decomp} 写出 $\mathrm{MSE}(\hat{\mu}_\alpha)$；
  (b) 求使 $\mathrm{MSE}(\hat{\mu}_\alpha)$ 最小的最优收缩因子 $\alpha^*$，证明 $\alpha^* < 1$ 且 $\mathrm{MSE}(\hat{\mu}_{\alpha^*}) < \mathrm{MSE}(\bar{X}) = \sigma^2/n$；结合式~\eqref{eq:ch5-gauss-conj} 说明这正是高斯先验 MAP 的收缩机理。
  \exHint{$\mathrm{MSE}(\hat{\mu}_\alpha) = (\alpha - 1)^2\mu^2 + \alpha^2\frac{\sigma^2}{n}$，这是关于 $\alpha$ 的开口向上二次函数，求导立得 $\alpha^* = \frac{n\mu^2}{n\mu^2 + \sigma^2} < 1$。}

  \item \textbf{泊松-伽马共轭与多元高斯均值 MAP 推导}：
  (a) 设到达率 $\lambda > 0$ 服从先验 $\mathrm{Gamma}(\alpha, \beta)$（密度 $p(\lambda) \propto \lambda^{\alpha - 1} e^{-\beta\lambda}$），观测到独立泊松计数 $x_1, \dots, x_n \stackrel{\text{i.i.d.}}{\sim} \mathrm{Poisson}(\lambda)$（$p(x;\lambda) = e^{-\lambda}\lambda^x / x!$）。证明后验分布仍为 Gamma 分布并写出其超参数、后验均值与 MAP（$\alpha + \sum x_i > 1$）；
  (b) 设 $d$ 维向量 $x_1, \dots, x_n \stackrel{\text{i.i.d.}}{\sim} \mathcal{N}(\mu, \sigma^2 I_d)$，均值先验 $\mu \sim \mathcal{N}(0, \tau^2 I_d)$。写出对数后验并对 $\mu$ 求梯度，证明 $\hat{\mu}_{\text{MAP}} = \frac{n}{n + \sigma^2/\tau^2}\bar{x}$。
  \exHint{(a) 只保留含 $\lambda$ 的因子相乘即可直接辨认 Gamma 核函数；(b) $\ln p(\mu\mid\mathcal{D}) = -\frac{1}{2\sigma^2}\sum_{i=1}^n \norm{x_i-\mu}_2^2 - \frac{1}{2\tau^2}\norm{\mu}_2^2 + \text{const}$。}
\end{enumerate}

\exTier{D. 可自检编程实验}{动手验证方差估计量 MSE 与配对 Bootstrap 检验}
\begin{enumerate}[label=\textbf{5.\arabic*}, resume]
  \item \textbf{蒙特卡洛实测三种方差估计量的偏差与 MSE}：固定随机种子 \texttt{2026}，重复 $M=20000$ 次从 $\mathcal{N}(0, \sigma^2=4.0)$ 中独立抽取 $n=5$ 个样本，每次计算 $\hat{\sigma}_n^2 = \frac{1}{n}\sum (X_i-\bar{X})^2$、无偏 $S^2 = \frac{1}{n-1}\sum (X_i-\bar{X})^2$ 与收缩版 $\tilde{\sigma}^2 = \frac{1}{n+1}\sum (X_i-\bar{X})^2$，统计三者的经验均值、偏差与经验 MSE，并与理论值对比。\\
  \textbf{自检标准}：断言 $|\text{mean}(S^2) - 4.0| < 0.08$、$|\text{mean}(\hat{\sigma}_n^2) - 3.2| < 0.08$，且经验 MSE 严格满足 $\mathrm{MSE}(\tilde{\sigma}^2) < \mathrm{MSE}(\hat{\sigma}_n^2) < \mathrm{MSE}(S^2)$（精确理论值分别为 $5.333 < 5.760 < 8.000$）。
  \exHint{利用 NumPy 一次生成形状为 \texttt{(20000, 5)} 的矩阵并沿 \texttt{axis=1} 向量化计算离差平方和，全程不到 0.05 秒。}

\par\pagebreak[4]
  \item \textbf{实现配对 Bootstrap 比较两个分类器}：在 $n=180$ 条验证样本上比较模型 A 与模型 B，考虑以下两种配对命中情形：
  \begin{enumerate}[label=(\alph*)]
    \item \textbf{微弱差异}：151 条同对、2 条仅 A 对、4 条仅 B 对、23 条同错（A 准确率 $85.00\%$，B 准确率 $86.11\%$，差值 $+1.11\%$）；
    \item \textbf{显著差异}：145 条同对、2 条仅 A 对、14 条仅 B 对、19 条同错（A 准确率 $81.67\%$，B 准确率 $88.33\%$，差值 $+6.67\%$）。
  \end{enumerate}
  固定随机种子 \texttt{2026}，对两种情形各运行 $B=10000$ 次\textbf{配对 Bootstrap}（每次抽取同一组长度 $180$ 的有放回索引同时作用于 A 与 B），计算差值 $\Delta = \text{Acc}_B - \text{Acc}_A$ 的 95\% 置信区间及 $\Prob^*(\Delta > 0)$。\\
  \textbf{自检标准}：情形 (a) 的 95\% CI 必须包含 $0$（约 $[-1.67\%, +3.89\%]$），不能宣称显著胜出；情形 (b) 的 95\% CI 下界严格大于 $0$（约 $[+2.78\%, +11.11\%]$）且 $\Prob^*(\Delta > 0) > 0.99$。
  \exHint{配对 Bootstrap 的关键是用\textbf{同一份索引矩阵} \texttt{idx = rng.integers(0, n, size=(B, n))} 同时索引 \texttt{hits\_a[idx]} 与 \texttt{hits\_b[idx]}，从而消除样本难易度共变噪声、大幅缩窄差值方差！}
\end{enumerate}

\chapter[纯 NumPy 回归/分类端到端闭环]{底座大合流实战：纯 NumPy 线性回归与逻辑/Softmax 回归全闭环}
\label{ch:ch6-logreg-pipeline}

\chapEpigraph{凡我不能亲手创造之物，我便尚未真正理解。}{理查德·费曼（Richard Feynman），加州理工学院黑板遗言（1988）}

\begin{chapIntro}{本章定位：第一阶段（第 1--2 卷）数理与工程底座的终极合流}
走到这里，你已经完成了本科人工智能数理底座的全部核心拼图：第 1 卷装备了 \textbf{NumPy 向量化、正交投影最小二乘、实对称谱定理/PSD、矩阵微分与梯度下降/有限差分校验}；第 2 卷前五章装备了 \textbf{Bayes 推断、高斯/Bernoulli/Categorical 分布、交叉熵/KL 散度、MLE/MAP 与 Bootstrap 置信区间}。本章不引入任何深度学习框架（不用 PyTorch，不用 scikit-learn），将两卷工具彻底合流：
\begin{itemize}
  \item 从概率建模出发，一举统一\textbf{最小二乘、$L_2$ 岭回归与 $L_1$ Lasso}；
  \item 用严格反例与导数证明揭示\textbf{线性回归}与 \textbf{Sigmoid + MSE} 在二分类上的双重失效；
  \item 完整推导\textbf{逻辑回归与 Softmax 回归}的交叉熵梯度、全局严格凸性、权重发散定理与 \textbf{Log-Sum-Exp} 数值稳定写法；
  \item 亲手实测\textbf{纯噪声数据泄漏灾难}、构建不平衡分类评测体系，并交付一套纯 NumPy 端到端训练、扫参与 Bootstrap 置信区间闭环！
\end{itemize}
\end{chapIntro}

% ====================================================================
\section{线性回归的概率统一：MLE $\equiv$ 最小二乘，MAP $\equiv$ 岭回归 / Lasso}
\label{sec:ch6-linreg-prob}

在第 1 卷第 2 章中，我们从几何正交投影出发，把线性回归写成最小化平方残差和 $\frac{1}{2}\sum_{i=1}^n (y_i - w^T x_i - b)^2$，又为了挽救病态矩阵 $X^T X$ 而“人为”加上了岭回归惩罚项 $\frac{\lambda}{2}\norm{w}_2^2$。当时留下的悬念——\textbf{为什么偏偏是平方误差、为什么 $L_2$ 惩罚对应权重缩小、$\lambda$ 的物理含义究竟是什么}——现在用第 5 章的 MLE 与 MAP 工具只需十行推导即可彻底揭晓！

考虑如下\textbf{生成式概率回归模型}：给定 $n$ 条独立观测 $\{(x_i, y_i)\}_{i=1}^n$（$x_i \in \R^d,\; y_i \in \R$），假设目标值由线性信号叠加独立同分布零均值高斯噪声产生：
\begin{equation}
  y_i = w^T x_i + b + \varepsilon_i,
  \qquad
  \varepsilon_i \stackrel{\text{i.i.d.}}{\sim} \mathcal{N}(0, \sigma^2)
  \;\;\Longleftrightarrow\;\;
  y_i \mid x_i;\, w, b \;\sim\; \mathcal{N}(w^T x_i + b,\; \sigma^2),
  \label{eq:ch6-gauss-noise}
\end{equation}
其中噪声方差 $\sigma^2 > 0$ 为固定常数。

\begin{thmBox}
\begin{theorem}[线性回归的三重概率等价性：OLS、Ridge 与 Lasso]
\label{thm:ch6-linreg-prob}
在高斯噪声模型~\eqref{eq:ch6-gauss-noise} 下（偏置 $b$ 取无信息平坦先验，不加惩罚）：
\begin{enumerate}
  \item \textbf{高斯 MLE $\equiv$ 普通最小二乘（OLS）}：最大化条件对数似然 $\sum_{i=1}^n \ln p(y_i \mid x_i; w, b)$ 完全等价于
  \[
    (\hat{w}_{\text{MLE}}, \hat{b}_{\text{MLE}})
    = \argmin_{w, b}\; \frac{1}{2}\sum_{i=1}^n \big(y_i - w^T x_i - b\big)^2.
  \]
  \item \textbf{各向同性高斯先验 MAP $\equiv$ $L_2$ 岭回归（Ridge）}：若对权重赋予先验 $w \sim \mathcal{N}(0, \tau^2 I_d)$（$\tau > 0$），则最大后验估计严格等价于
  \begin{equation}
    (\hat{w}_{\text{MAP}}, \hat{b}_{\text{MAP}})
    = \argmin_{w, b}\; \frac{1}{2}\sum_{i=1}^n \big(y_i - w^T x_i - b\big)^2 + \frac{\lambda}{2}\norm{w}_2^2,
    \qquad
    \lambda = \frac{\sigma^2}{\tau^2}.
    \label{eq:ch6-map-ridge}
  \end{equation}
  \item \textbf{独立拉普拉斯先验 MAP $\equiv$ $L_1$ Lasso 回归}：若各权重分量独立服从零均值拉普拉斯先验 $p(w_j) = \frac{1}{2b_0}\exp\!\left(-\frac{|w_j|}{b_0}\right)$（尺度参数 $b_0 > 0$），则最大后验估计严格等价于
  \begin{equation}
    (\hat{w}_{\text{MAP}}, \hat{b}_{\text{MAP}})
    = \argmin_{w, b}\; \frac{1}{2}\sum_{i=1}^n \big(y_i - w^T x_i - b\big)^2 + \lambda_1 \norm{w}_1,
    \qquad
    \lambda_1 = \frac{\sigma^2}{b_0}.
    \label{eq:ch6-map-lasso}
  \end{equation}
\end{enumerate}
\end{theorem}
\end{thmBox}

\begin{proof}
由条件独立性，写出高斯对数似然：
\[
  \ell(w, b)
  = \sum_{i=1}^n \ln \left[\frac{1}{\sqrt{2\pi\sigma^2}}\exp\!\left(-\frac{(y_i - w^T x_i - b)^2}{2\sigma^2}\right)\right]
  = -\frac{n}{2}\ln(2\pi\sigma^2) - \frac{1}{2\sigma^2}\sum_{i=1}^n (y_i - w^T x_i - b)^2.
\]
由于第一项为常数且 $\sigma^2 > 0$，最大化 $\ell(w,b)$ 等价于最小化 $-\sigma^2\ell(w,b)$（去掉常数），立得结论 1（注意最优 $(\hat{w}_{\text{MLE}}, \hat{b}_{\text{MLE}})$ 竟不依赖于未知噪声方差 $\sigma^2$！）。

对于结论 2，高斯先验 $p(w) = (2\pi\tau^2)^{-d/2}\exp\!\left(-\frac{1}{2\tau^2}\norm{w}_2^2\right)$ 给出对数后验：
\[
  \ln p(w, b \mid \mathcal{D})
  = -\frac{1}{2\sigma^2}\sum_{i=1}^n (y_i - w^T x_i - b)^2 - \frac{1}{2\tau^2}\norm{w}_2^2 + \text{const}.
\]
两边同乘 $-\sigma^2$（正数不改变极值点方向）并最大化转最小化，直接得到式~\eqref{eq:ch6-map-ridge}，且正则化强度精确等于 $\lambda = \sigma^2 / \tau^2$。

对于结论 3，独立拉普拉斯先验给出 $\ln p(w) = \sum_{j=1}^d \ln p(w_j) = -d\ln(2b_0) - \frac{1}{b_0}\sum_{j=1}^d |w_j| = -\frac{1}{b_0}\norm{w}_1 + \text{const}$。同样对 $\ell(w,b) + \ln p(w)$ 同乘 $-\sigma^2$ 转为最小化，即得式~\eqref{eq:ch6-map-lasso} 且 $\lambda_1 = \sigma^2 / b_0$。
\end{proof}

\begin{insightBox}{读懂 $\lambda = \sigma^2 / \tau^2$ 与“为何不惩罚偏置 $b$”}
\begin{enumerate}
  \item \textbf{信噪比倒数定律}：$\lambda = \sigma^2/\tau^2$ 正是\textbf{观测噪声方差与先验信号方差之比}！数据越嘈杂（$\sigma^2$ 大），或我们越相信真实权重应当很小（$\tau^2$ 小），$\lambda$ 就越大，模型越不敢大起大落地去拟合噪声；反之当噪声极小时，$\lambda \to 0$，退化回无偏 OLS。
  \item \textbf{为何拉普拉斯先验偏爱稀疏解（$w_j = 0$）}：相比高斯密度在 $0$ 处的圆滑抛物线顶，拉普拉斯密度 $\propto e^{-|w_j|/b_0}$ 在 $w_j = 0$ 处有一个\textbf{尖锐的不可微峰尖}，把更多概率质量压向精确的零点，从而在第 4 卷的 Lasso 特征选择中能自动将无关噪声特征的权重清零！
  \item \textbf{为何绝不惩罚偏置 $b$}：如果把所有房价标签统一加上 $100$ 万元，真正应该平移的只是截距 $b$，特征权重 $w$ 丝毫不该改变。若对 $b$ 也加 $\frac{\lambda}{2}b^2$ 惩罚，就会把预测整体强行拉向 $0$，破坏标签平移不变性。因此标准工程实现中，\textbf{正则化永远只作用于权重 $w$，不作用于偏置 $b$}！
\end{enumerate}
\end{insightBox}

\begin{table}[htbp]
\centering
\caption{线性模型中噪声分布、权重先验与损失/正则化的对应词典}
\label{tab:ch6-prob-dict}
\footnotesize
\begin{tabularx}{\linewidth}{@{}p{27mm}p{27mm}p{35mm}X@{}}
\toprule
\textbf{似然 / 噪声假设} & \textbf{权重先验 $p(w)$} & \textbf{等价优化目标} & \textbf{模型名称与行为特征} \\
\midrule
高斯 $\mathcal{N}(w^Tx+b, \sigma^2)$ & 无（均匀先验） & $\frac{1}{2}\norm{Xw+b\mathbf{1}-y}_2^2$ & 普通最小二乘（OLS），高维小样本易过拟合 \\
高斯 $\mathcal{N}(w^Tx+b, \sigma^2)$ & 高斯 $\mathcal{N}(0, \tau^2 I)$ & $\frac{1}{2}\norm{Xw+b\mathbf{1}-y}_2^2 + \frac{\sigma^2}{2\tau^2}\norm{w}_2^2$ & $L_2$ 岭回归（Ridge），平滑压缩权重、改善条件数 \\
高斯 $\mathcal{N}(w^Tx+b, \sigma^2)$ & 拉普拉斯 $\mathrm{Lap}(0, b_0)$ & $\frac{1}{2}\norm{Xw+b\mathbf{1}-y}_2^2 + \frac{\sigma^2}{b_0}\norm{w}_1$ & $L_1$ Lasso 回归，诱导稀疏权重、自动特征筛选 \\
Bernoulli($\sigma(w^Tx+b)$) & 高斯 $\mathcal{N}(0, \tau^2 I)$ & $\mathrm{BCE}(w,b) + \frac{\lambda}{2}\norm{w}_2^2$ & $L_2$ 逻辑回归（本章第 6.3 节主角！） \\
\bottomrule
\end{tabularx}
\end{table}

% ====================================================================
\section{二分类困境：为何线性回归与 Sigmoid + MSE 都会失效}
\label{sec:ch6-binary-traps}

既然线性回归已经有了如此优雅的闭式解与梯度下降，面对\textbf{二分类问题}（标签 $y_i \in \{0, 1\}$，例如垃圾邮件识别、肿瘤良恶性判断），初学者的第一反应通常是以下两种偷懒方案：
\begin{itemize}
  \item \textbf{方案一}：直接把 $0/1$ 当成实数跑线性回归 $\hat{y} = w^T x + b$，以 $0.5$ 为阈值判别正负类；
  \item \textbf{方案二}：套一个 Sigmoid 函数 $\sigma(z) = \frac{1}{1+e^{-z}} \in (0, 1)$ 把输出压到 $(0,1)$ 内，再沿用熟悉的均方误差（MSE）$\frac{1}{2}(\sigma(z) - y)^2$ 做梯度下降。
\end{itemize}
让我们用两个硬核反例彻底击碎这两个直觉陷阱——它们不仅解释了为何必须改用交叉熵，更为第 5 卷深度神经网络的梯度消失问题打下伏笔！

\subsection{陷阱一：远方“铁定正确”的正样本把线性回归决策边界拖歪！}
\label{subsec:ch6-linreg-fail}

线性回归用于分类的第一个表面缺陷是输出 $\hat{y} = w^T x + b \in (-\infty, +\infty)$，根本无法解释为概率。但更致命的是\textbf{平方损失会重罚“太正确”的样本}！看一个可以手算验算的 1 维反例：

\begin{example}[加一条毫无争议的远方正样本，线性回归竟然翻车！]
\label{ex:ch6-ols-drag}
考虑 6 个完全线性可分的一维分类点：负类（$y=0$）位于 $x \in \{-2, -1, 0\}$，正类（$y=1$）位于 $x \in \{1, 2, 3\}$。
\begin{enumerate}
  \item \textbf{加入远点前（6 个点）}：由最小二乘法精确解得回归直线
  \[
    \hat{y}_0(x) = \frac{9}{35}x + \frac{13}{35} \approx 0.2571x + 0.3714.
  \]
  令 $\hat{y}_0(x) = 0.5$，决策边界正好落在正负类正中央 $x_0^* = \mathbf{0.5000}$，分类准确率 $100\%$！
  \item \textbf{仅增加 1 个极易分辨的远方正样本 $(30, 1)$（共 7 个点）}：如果按原直线，在 $x=30$ 处的预测值高达 $\hat{y}_0(30) \approx 8.09$，远远超过标签 $1$，产生巨大的无谓平方残差 $(8.09 - 1)^2 \approx 50.2$！为了拼命压平右端直线以讨好 $(30, 1)$，最小二乘被迫把整条直线放倒为
  \[
    \hat{y}_1(x) = \frac{15}{668}x + \frac{311}{668} \approx 0.0225x + 0.4656.
  \]
  此时 $0.5$ 阈值决策边界被狠狠向右拖拽到
  \[
    x_1^* = \frac{0.5 - 311/668}{15/668} = \frac{23}{15} \approx \mathbf{1.5333} > 1!
  \]
  结果原本清晰无比的正类样本 $(1, 1)$ 的预测值跌至 $\hat{y}_1(1) = \frac{163}{334} \approx 0.4880 < 0.5$，直接被\textbf{错判为负类}！
\end{enumerate}
\end{example}

\subsection{陷阱二：Sigmoid + MSE 的非凸性与梯度消失严格证明}
\label{subsec:ch6-sigmse-fail}

既然线性回归是因为输出无界而被远点拖歪，那么加上值域在 $(0,1)$ 之间的 Sigmoid 函数 $\sigma(z) = \frac{1}{1+e^{-z}}$（满足导数恒等式 $\sigma'(z) = \sigma(z)(1-\sigma(z))$），然后继续用 MSE 损失行不行？绝对不行！

\begin{failBox}{Sigmoid + MSE 的两大死穴：非凸陷阱 + 错得越离谱梯度越趋近于 0！}
固定一个正类样本 $y = 1$，令 $p \triangleq \sigma(z) \in (0, 1)$，考察单样本 Sigmoid + MSE 损失 $\ell_{\text{MSE}}(z) = \frac{1}{2}(p - 1)^2$。利用链式法则与 $\frac{\mathrm{d}p}{\mathrm{d}z} = p(1 - p)$，逐阶求导：
\begin{align}
  \ell'_{\text{MSE}}(z)
  &= (p - 1)\cdot \frac{\mathrm{d}p}{\mathrm{d}z}
  = -p(1 - p)^2, \label{eq:ch6-sigmse-d1} \\
  \ell''_{\text{MSE}}(z)
  &= \frac{\mathrm{d}}{\mathrm{d}p}\!\big[-p(1 - p)^2\big]\cdot p(1 - p)
  = -\big[(1 - p)^2 - 2p(1 - p)\big]\cdot p(1 - p)
  = p(1 - p)^2(3p - 1). \label{eq:ch6-sigmse-d2}
\end{align}
式~\eqref{eq:ch6-sigmse-d1}--\eqref{eq:ch6-sigmse-d2} 立刻暴露出两个致命病症（图~\ref{fig:ch6-mse-vs-bce} 可视化了完整对比）：
\begin{enumerate}
  \item \textbf{破坏凸性（Non-Convexity）}：由式~\eqref{eq:ch6-sigmse-d2}，只要 $p < \frac{1}{3}$（即 $z < -\ln 2 \approx -0.693$），就有 $\ell''_{\text{MSE}}(z) < 0$，损失函数\textbf{严格凹}！例如在 $z = -2$（$p \approx 0.1192$）处，二阶导数
  \[
    \ell''_{\text{MSE}}(-2) \approx \mathbf{-0.0594} < 0.
  \]
  非凸损失意味着梯度下降极易陷入平坦高原或糟糕的局部极小值！
  \item \textbf{自信犯错时的梯度消失（Gradient Vanishing）}：当真实标签是 $y=1$，但模型却极度自信地判错，输出 $z = -6$（预测概率仅 $p \approx 0.002473$）时，我们最需要一个巨大的梯度把它拉回正轨；然而式~\eqref{eq:ch6-sigmse-d1} 中额外乘上的饱和因子 $p = \sigma(-6)$ 把梯度活活压死到了
  \[
    \abs{\ell'_{\text{MSE}}(-6)} \approx \mathbf{0.00246} \qquad \text{vs.} \qquad \abs{\ell'_{\text{BCE}}(-6)} = 1 - p \approx \mathbf{0.99753}!
  \]
  在最该猛烈纠错的时刻，MSE 的梯度竟然比交叉熵（BCE）微弱了整整 \textbf{405 倍}！
\end{enumerate}
\end{failBox}

\begin{figure}[htbp]
\centering
\begin{tikzpicture}
  % Left subplot: Loss curves for y=1
  \begin{axis}[
    name=leftplot,
    width=0.47\linewidth,
    height=5.6cm,
    xmin=-6.1, xmax=4.1,
    ymin=-0.05, ymax=6.3,
    xlabel={线性得分 $z = w^T x + b$（真实标签 $y=1$）},
    ylabel={单样本损失值 $\ell(z)$},
    xlabel style={font=\footnotesize\sffamily},
    ylabel style={font=\footnotesize\sffamily},
    tick label style={font=\scriptsize},
    grid=major,
    grid style={dashed, gray!25},
    legend pos=north east,
    legend style={font=\scriptsize\sffamily, fill=white, fill opacity=0.92, draw=tealLine},
    axis lines=left
  ]
    % BCE: ln(1 + exp(-z))
    \addplot[teal, ultra thick, domain=-6:4, samples=120] {ln(1 + exp(-x))};
    \addlegendentry{交叉熵 $\ell_{\text{BCE}}(z)$（严格凸）}

    % MSE: 0.5*(1/(1+exp(-z)) - 1)^2
    \addplot[ochre, very thick, dashed, domain=-6:4, samples=120] {0.5*(1/(1+exp(-x)) - 1)^2};
    \addlegendentry{Sigmoid+MSE $\ell_{\text{MSE}}(z)$}

    % Highlight non-convex region z < -ln(2)
    \draw[ochre!70, dotted, thick] (axis cs:-0.693, 0) -- (axis cs:-0.693, 2.3);
    \node[font=\scriptsize\sffamily, color=ochre, align=center] at (axis cs:-3.3, 1.35) {非凸平坦区\\$z < -\ln 2\approx -0.693$\\($\ell''_{\text{MSE}}(-2)\approx -0.059$)};
    \draw[-{Stealth[length=1.6mm]}, thin, ochre] (axis cs:-2.5, 0.85) -- (axis cs:-2.0, 0.44);
  \end{axis}

  % Right subplot: Gradient magnitude |ell'(z)|
  \begin{axis}[
    at={(leftplot.east)}, anchor=west, xshift=13mm,
    width=0.47\linewidth,
    height=5.6cm,
    xmin=-6.1, xmax=4.1,
    ymin=-0.03, ymax=1.16,
    xlabel={线性得分 $z = w^T x + b$（真实标签 $y=1$）},
    ylabel={梯度幅值 $\abs{\ell'(z)}$},
    xlabel style={font=\footnotesize\sffamily},
    ylabel style={font=\footnotesize\sffamily},
    tick label style={font=\scriptsize},
    grid=major,
    grid style={dashed, gray!25},
    legend pos=north east,
    legend style={font=\scriptsize\sffamily, fill=white, fill opacity=0.92, draw=tealLine},
    axis lines=left
  ]
    % |ell'_BCE| = 1 - sigma(z) = 1/(1 + exp(z))
    \addplot[teal, ultra thick, domain=-6:4, samples=120] {1/(1 + exp(x))};
    \addlegendentry{$\abs{\ell'_{\text{BCE}}} = 1 - p$}

    % |ell'_MSE| = p*(1-p)^2
    \addplot[ochre, very thick, dashed, domain=-6:4, samples=120] {(1/(1+exp(-x)))*(1/(1+exp(x)))^2};
    \addlegendentry{$\abs{\ell'_{\text{MSE}}} = p(1-p)^2$}

    % Annotate z=-6 gradient gap
    \fill[teal] (axis cs:-6, 0.9975) circle (2pt);
    \fill[ochre] (axis cs:-6, 0.0025) circle (2pt);
    \draw[{Stealth[length=1.5mm]}-{Stealth[length=1.5mm]}, thick, navy] (axis cs:-5.85, 0.04) -- (axis cs:-5.85, 0.96);
    \node[font=\scriptsize\sffamily\bfseries, color=navy, align=left, anchor=west] at (axis cs:-5.65, 0.56) {在 $z=-6$ 处：\\BCE $= 0.9975$\\MSE $= 0.00246$\\（相差 \textbf{405 倍}！）};
  \end{axis}
\end{tikzpicture}
\caption{当 $y=1$ 时，交叉熵损失 $\ell_{\text{BCE}}(z)$ 与 Sigmoid+MSE 损失 $\ell_{\text{MSE}}(z)$ 的损失曲线（左）与梯度幅值（右）对比：交叉熵全局严格凸且在严重错判时保持满额纠错梯度，而 Sigmoid+MSE 在 $z < -\ln 2$ 处失去凸性且梯度崩溃消失。}
\label{fig:ch6-mse-vs-bce}
\end{figure}

% ====================================================================
\section{逻辑回归与 Softmax 回归：交叉熵、凸性证明与 Log-Sum-Exp}
\label{sec:ch6-logreg-softmax}

既然拍脑袋凑损失会撞上非凸与梯度消失，那就回到最根本的第一性原理——\textbf{按第 2、4、5 章的概率与信息论体系自然推导二分类与多分类模型}！

\subsection{对数几率、二元交叉熵（BCE）的双重推导与梯度抵消奇迹}
\label{subsec:ch6-bce}

二分类中条件分布 $y \mid x$ 天然服从 $\mathrm{Bernoulli}(p)$。为了把实数线性打分 $z = w^T x + b \in (-\infty, +\infty)$ 无缝映射到概率 $p \in (0, 1)$，我们假设\textbf{对数几率（Log-Odds / Logit）}是特征的线性函数：
\[
  \ln\frac{p}{1 - p} = w^T x + b = z
  \;\;\Longleftrightarrow\;\;
  \frac{p}{1 - p} = e^z
  \;\;\Longleftrightarrow\;\;
  p = \sigma(z) \triangleq \frac{1}{1 + e^{-z}}.
\]
这就是\textbf{逻辑回归（Logistic Regression）}（名字里虽然带“回归”，但它是地道的概率二分类模型！）。它满足两条极其好用的恒等式：$\sigma(-z) = 1 - \sigma(z)$ 与 $\sigma'(z) = \sigma(z)(1 - \sigma(z))$。

给定训练集 $\{(x_i, y_i)\}_{i=1}^N$（$y_i \in \{0, 1\}$），记 $z_i = w^T x_i + b$、$\hat{p}_i = \sigma(z_i)$。我们从\textbf{两条殊途同归的路径}导出它的目标函数：
\begin{enumerate}
  \item \textbf{路径一（第 5 章 Bernoulli MLE + 高斯先验 MAP）}：单个样本的 Bernoulli 似然为 $\hat{p}_i^{y_i}(1 - \hat{p}_i)^{1 - y_i}$。对平均负对数似然加上权重高斯先验 $-\ln p(w)$（缩放后记为 $\frac{\lambda}{2}\norm{w}_2^2$，$\lambda \ge 0$，偏置 $b$ 不加惩罚），得
  \begin{equation}
    \mathcal{L}(w, b)
    = -\frac{1}{N}\sum_{i=1}^N \Big[y_i\ln \hat{p}_i + (1 - y_i)\ln(1 - \hat{p}_i)\Big] + \frac{\lambda}{2}\norm{w}_2^2.
    \label{eq:ch6-bce}
  \end{equation}
  利用 $\ln \hat{p}_i = -\ln(1+e^{-z_i})$ 与 $\ln(1-\hat{p}_i) = -\ln(1+e^{z_i})$ 以及 $\ln(1+e^{-z_i}) = \ln(1+e^{z_i}) - z_i$，可将其化简为不含对数除法的\textbf{直接 Logit 形式}：
  \begin{equation}
    \mathcal{L}(w, b)
    = \frac{1}{N}\sum_{i=1}^N \Big[\ln\big(1 + e^{z_i}\big) - y_i z_i\Big] + \frac{\lambda}{2}\norm{w}_2^2.
    \label{eq:ch6-bce-logit}
  \end{equation}
  \item \textbf{路径二（第 4 章交叉熵与 KL 散度）}：对每条样本，真实标签 $y_i \in \{0,1\}$ 对应退化分布 $P_i = \mathrm{Bernoulli}(y_i)$（其熵 $H(P_i) = 0$），模型预测为 $Q_i = \mathrm{Bernoulli}(\hat{p}_i)$，则第 $i$ 项损失恰好等于
  \[
    \ell_i = -\big[y_i\ln\hat{p}_i + (1 - y_i)\ln(1 - \hat{p}_i)\big] = H(P_i, Q_i) = \DKL(P_i \| Q_i)!
  \]
  因此最小化\textbf{二元交叉熵（Binary Cross-Entropy, BCE）}就是让模型预测概率分布与真实标签分布之间的平均 KL 散度最小！
\end{enumerate}

现在计算单样本损失 $\ell_i(z_i) = \ln(1 + e^{z_i}) - y_i z_i$ 对 $z_i$ 的导数——见证\textbf{梯度抵消奇迹}：
\[
  \frac{\partial \ell_i}{\partial z_i}
  = \frac{e^{z_i}}{1 + e^{z_i}} - y_i
  = \frac{1}{1 + e^{-z_i}} - y_i
  = \hat{p}_i - y_i!
\]
（如果你从原式~\eqref{eq:ch6-bce} 求导，$\frac{\partial\ell_i}{\partial \hat{p}_i} = \frac{\hat{p}_i - y_i}{\hat{p}_i(1-\hat{p}_i)}$ 的分母 $\hat{p}_i(1-\hat{p}_i)$ 恰好与 Sigmoid 导数 $\sigma'(z_i) = \hat{p}_i(1-\hat{p}_i)$ 完美约分！）那个在 MSE 里把梯度扼杀 405 倍的饱和因子彻底消失了！由第 1 卷第 5 章链式法则，矩阵化批量梯度简洁得令人惊叹：
\begin{equation}
  \nabla_w \mathcal{L}(w, b) = \frac{1}{N} X^T (\hat{p} - y) + \lambda w,
  \qquad
  \frac{\partial \mathcal{L}(w, b)}{\partial b} = \frac{1}{N}\sum_{i=1}^N (\hat{p}_i - y_i),
  \label{eq:ch6-bce-grad}
\end{equation}
其中 $X \in \R^{N\times d}$ 为特征矩阵，$\hat{p} = \sigma(Xw + b\mathbf{1}) \in \R^N$，$y \in \{0,1\}^N$。你会发现它和线性回归的梯度 $\frac{1}{N}X^T(\hat{y} - y) + \lambda w$ 具有\textbf{完全相同的“误差乘特征”结构}！

\subsection{全局严格凸性与线性可分权重发散定理}
\label{subsec:ch6-convex-diverge}

与 Sigmoid + MSE 坑洼遍地的非凸地形截然不同，逻辑回归的交叉熵损失具有整个优化理论梦寐以求的\textbf{全局严格凸性}。

\begin{thmBox}
\begin{theorem}[带 $L_2$ 正则化逻辑回归的全局严格凸性]
\label{thm:ch6-logreg-convex}
对任意样本矩阵 $X \in \R^{N\times d}$ 与标签 $y \in \{0,1\}^N$，二元交叉熵目标函数 $\mathcal{L}(w, b)$ 关于权重 $w$ 的 Hessian 矩阵为
\begin{equation}
  H_w(w, b) \triangleq \nabla_w^2 \mathcal{L}(w, b) = \frac{1}{N} X^T D X + \lambda I_d,
  \qquad
  D = \diag\big(\hat{p}_1(1-\hat{p}_1), \dots, \hat{p}_N(1-\hat{p}_N)\big).
  \label{eq:ch6-logreg-hessian}
\end{equation}
当 $\lambda = 0$ 时，$H_w \succeq 0$（全局凸）；当 $\lambda > 0$ 时，$H_w \succ 0$ 处处成立，进而联合 Hessian $\nabla_{(w,b)}^2 \mathcal{L} \succ 0$，故 $\mathcal{L}(w, b)$ 是\textbf{全局严格凸函数}，存在\textbf{唯一全局极小值点}，绝无任何局部极小值陷阱！
\end{theorem}
\end{thmBox}

\begin{proof}
对式~\eqref{eq:ch6-bce-grad} 的第 $j$ 个分量再次对 $w_k$ 求偏导，注意到 $\frac{\partial \hat{p}_i}{\partial w_k} = \hat{p}_i(1-\hat{p}_i) x_{ik}$：
\[
  \frac{\partial^2 \mathcal{L}}{\partial w_j \partial w_k}
  = \frac{1}{N}\sum_{i=1}^N x_{ij}\, \hat{p}_i(1-\hat{p}_i)\, x_{ik} + \lambda\,\delta_{jk},
\]
写成矩阵形式恰好是式~\eqref{eq:ch6-logreg-hessian}。为检验正定性（第 1 卷第 3 章二次型判据），任取非零向量 $v \in \R^d$：
\[
  v^T H_w v
  = \frac{1}{N}\sum_{i=1}^N \underbrace{\hat{p}_i(1-\hat{p}_i)}_{>\,0}(x_i^T v)^2 + \lambda\norm{v}_2^2
  \ge \lambda\norm{v}_2^2.
\]
因为对任意有限 $z_i \in \R$ 均有 $\hat{p}_i \in (0,1)$，故每个权重因子 $d_i = \hat{p}_i(1-\hat{p}_i) > 0$。当 $\lambda = 0$ 时 $v^T H_w v \ge 0$，即 $H_w \succeq 0$；当 $\lambda > 0$ 时 $v^T H_w v \ge \lambda\norm{v}_2^2 > 0$，即 $H_w \succ 0$。进一步对增广向量 $(v, c) \in \R^{d+1} \setminus \{0\}$，令 $\tilde{x}_i = (x_i, 1)^T$，联合二次型为 $\frac{1}{N}\sum_{i=1}^N d_i (x_i^T v + c)^2 + \lambda\norm{v}_2^2$：若 $v \neq 0$ 则第二项严格为正；若 $v = 0$ 且 $c \neq 0$ 则第一项等于 $\frac{c^2}{N}\sum_i d_i > 0$。故联合 Hessian 处处严格正定，全局严格凸性得证。
\end{proof}

不过请千万留神：当 $\lambda = 0$（无正则化）且数据\textbf{完全线性可分}时，虽然损失函数仍是凸的，但\textbf{有限的最优权重根本不存在}！

\begin{thmBox}
\begin{proposition}[线性可分数据上无正则逻辑回归的权重发散定理]
\label{prop:ch6-weight-diverge}
设 $\lambda = 0$，且训练集线性可分：存在 $(w_0, b_0)$ 与间隔 $\gamma > 0$ 使得对所有样本 $i = 1, \dots, N$ 均有
\[
  s_i (w_0^T x_i + b_0) \ge \gamma > 0, \qquad \text{其中 } s_i \triangleq 2y_i - 1 \in \{-1, +1\}.
\]
则沿射线 $(t w_0, t b_0)$（$t > 0$），随着 $t \to +\infty$，无正则损失严格单调趋近于 $0$：
\[
  \lim_{t \to +\infty} \mathcal{L}(t w_0, t b_0) = 0,
\]
而对任意有限 $(w, b)$ 恒有 $\mathcal{L}(w, b) > 0$。因此极小值在有限参数空间中\textbf{不可达}，梯度下降会驱使权重范数 $\norm{w}_2 \to +\infty$！
\end{proposition}
\end{thmBox}

\begin{proof}
利用 $s_i = 2y_i - 1 \in \{-1, +1\}$ 容易验证单样本 BCE 损失可统一写作 $\ell_i(z_i) = \ln\!\big(1 + e^{-s_i z_i}\big) > 0$（当 $y_i=1$ 时为 $\ln(1+e^{-z_i})$，当 $y_i=0$ 时为 $\ln(1+e^{z_i})$）。将 $(tw_0, tb_0)$ 代入得
\[
  0 < \mathcal{L}(t w_0, t b_0)
  = \frac{1}{N}\sum_{i=1}^N \ln\!\Big(1 + e^{-t\,s_i(w_0^T x_i + b_0)}\Big)
  \le \frac{1}{N}\sum_{i=1}^N \ln\!\big(1 + e^{-t\gamma}\big)
  = \ln\!\big(1 + e^{-t\gamma}\big).
\]
令 $t \to +\infty$，因 $\gamma > 0$ 故 $e^{-t\gamma} \to 0$，右边 $\ln(1 + e^{-t\gamma}) \to \ln 1 = 0$。由夹逼准则即得 $\lim_{t\to+\infty}\mathcal{L}(t w_0, t b_0) = 0$，而任何有限 $(w, b)$ 的损失皆严格大于 $0$，因此最优解必在无穷远处发散！
\end{proof}

这解释了我们在第 6.5 节实测中将会亲眼看到的惊人现象：在 30 维高维噪声数据上，如果不加 $L_2$ 正则化（$\lambda = 0$），梯度下降不断把权重 $\norm{w}_2$ 推爆到 $5.825$ 以上，把训练集硬生生拟合到极端自信，验证集交叉熵损失却一路飙升！而只要加上哪怕很小的 $\lambda > 0$（高斯先验 MAP），二次惩罚项 $\frac{\lambda}{2}t^2\norm{w_0}_2^2 \to +\infty$ 瞬间锁死权重边界，恢复唯一有限最优解！

\subsection{多分类 Softmax 回归与 Log-Sum-Exp 数值稳定化}
\label{subsec:ch6-softmax-lse}

当类别数从 $2$ 类推广到 $K \ge 3$ 类（如手写数字 $0$--$9$、ImageNet 千类分类、大模型词表预测）时，样本 $x_i \in \R^d$ 经权重矩阵 $W \in \R^{K\times d}$ 与偏置 $b \in \R^K$ 算出 $K$ 维打分向量（logits）$z_i = W x_i + b \in \R^K$，再用第 2 章的 \textbf{Softmax 函数}转为合法 Categorical 概率分布：
\[
  \hat{p}_{ik} = \mathrm{softmax}(z_i)_k \triangleq \frac{e^{z_{ik}}}{\sum_{c=1}^K e^{z_{ic}}},
  \qquad k = 1, \dots, K.
\]
记真实类别 $c_i \in \{1,\dots,K\}$ 的 one-hot 向量为 $y_i \in \{0,1\}^K$（仅第 $c_i$ 位为 $1$）。由 Categorical 分布极大似然（同样等价于最小化真实 one-hot 分布与预测分布的交叉熵 / KL 散度），单样本\textbf{多分类交叉熵损失}为
\begin{equation}
  \ell_i(z_i, y_i)
  = -\sum_{k=1}^K y_{ik}\ln \hat{p}_{ik}
  = -\ln \hat{p}_{i, c_i}
  = \ln\!\left(\sum_{c=1}^K e^{z_{ic}}\right) - z_{i, c_i}.
  \label{eq:ch6-softmax-loss}
\end{equation}
直接对式~\eqref{eq:ch6-softmax-loss} 关于任意分量 $z_{ik}$ 求偏导（详细展开见习题 6.5）：
\[
  \frac{\partial \ell_i}{\partial z_{ik}}
  = \frac{e^{z_{ik}}}{\sum_{c=1}^K e^{z_{ic}}} - y_{ik}
  = \hat{p}_{ik} - y_{ik}.
\]
仍然是完美简洁的“\textbf{预测概率减去真实标签}”！对整批样本 $X \in \R^{N\times d}$、one-hot 标签矩阵 $Y \in \{0,1\}^{N\times K}$、概率矩阵 $\hat{P} \in \R^{N\times K}$，加上 $L_2$ 正则项 $\frac{\lambda}{2}\norm{W}_F^2$ 后的矩阵梯度为
\[
  \nabla_W \mathcal{L} = \frac{1}{N}(\hat{P} - Y)^T X + \lambda W,
  \qquad
  \nabla_b \mathcal{L} = \frac{1}{N}\sum_{i=1}^N (\hat{p}_i - y_i).
\]

\begin{failBox}{浮点溢出红线与 Log-Sum-Exp 数值稳定恒等式}
如果你天真地按课本公式直接写 \texttt{np.exp(z) / np.sum(np.exp(z))} 或 \texttt{-np.log(p)}，训练几步就可能爆出满屏 \texttt{nan}：
\begin{itemize}
  \item 在 IEEE 754 \texttt{float64} 中，$e^{710} > 1.797 \times 10^{308}$，直接上溢为 \texttt{inf}，进而 \texttt{inf / inf = nan}；
  \item 当 $z = -800$ 时 $e^{z}$ 下溢为 \texttt{0.0}，随后 \texttt{np.log(0.0) = -inf}，损失立刻炸毁！
\end{itemize}
\textbf{工业界标准解法——Log-Sum-Exp（LSE）平移不变技巧}：记 $m = \max_{1\le c\le K} z_c$，在指数内部统一减去最大值 $m$：
\begin{equation}
  \operatorname{LSE}(z) \triangleq \ln\sum_{c=1}^K e^{z_c} = m + \ln\sum_{c=1}^K e^{z_c - m},
  \qquad
  \mathrm{softmax}(z)_k = \frac{e^{z_k - m}}{\sum_{c=1}^K e^{z_c - m}}.
  \label{eq:ch6-lse}
\end{equation}
因为每一项 $z_c - m \le 0$，故 $e^{z_c - m} \in (0, 1]$ 绝不会上溢；同时至少有一项 $z_{c^*} - m = 0$ 使 $e^0 = 1$，故求和结果必 $\ge 1$，绝不会下溢成 $0$！对于二元交叉熵式~\eqref{eq:ch6-bce-logit}，NumPy 已经内置了专门的两数 Log-Sum-Exp 原语 \texttt{np.logaddexp(0.0, z)} 精确计算 $\ln(1 + e^z)$——哪怕传入 $z = +1000$ 或 $z = -1000$，也能给出分毫不差的有限浮点数！
\end{failBox}

% ====================================================================
\section{实验工程红线：数据泄漏实测与分类评测体系}
\label{sec:ch6-leakage-metrics}

数学公式无懈可击，并不等于模型实验就可靠。无数初学者在竞赛、论文甚至工业项目里栽在两大工程暗礁上：\textbf{数据泄漏（Data Leakage）}与\textbf{不平衡准确率陷阱（Imbalanced Accuracy Trap）}。

\subsection{红线一：纯噪声也能跑出 80\% 测试准确率？数据泄漏实测！}
\label{subsec:ch6-leakage}

机器学习评估的最高铁律只有一句话：\textbf{验证集/测试集在最终评估之前必须完全“隐身”}！一切依赖数据统计量的预处理——无论是 Z-score 标准化（均值 $\mu$、标准差 $\sigma$）、缺失值填充、PCA 降维，还是特征筛选——都\textbf{只能仅在训练集上 \texttt{fit}}，然后把训练集学到的参数死板地 \texttt{transform} 到验证/测试集上。

如果你觉得“先对整个数据集挑一挑最相关的特征再划分训练测试集，能有多大影响？”请看配套脚本 \texttt{code/vol2\_experiments.py} 中 \texttt{exp6\_leakage\_and\_logreg} 的震撼实测：
\begin{itemize}
  \item \textbf{实验设定}：生成 $N=60$ 条样本、$D=500$ 维\textbf{纯随机标准正态白噪声特征} $X_{ij} \sim \mathcal{N}(0,1)$，以及与 $X$ 完全独立的随机 $0/1$ 标签 $y_i \sim \mathrm{Bernoulli}(0.5)$。由于特征与标签毫无任何真实关系，理论上\textbf{全宇宙最强的贝叶斯最优分类器准确率也只能是 $50.00\%$}！
  \item \textbf{对比两种流水线（重复 50 次独立实验，40 训练 / 20 测试）}：
  \begin{enumerate}
    \item \textbf{泄漏流水线（错误）}：先在全部 60 条样本上计算每个特征与 $y$ 的相关性，挑出相关性最高的 10 个特征，再切分前 40 条训练、后 20 条测试；
    \item \textbf{干净流水线（正确）}：先切分 40 训练 / 20 测试，严格只在 40 条训练集上挑前 10 相关特征，再在 20 条测试集上评估。
  \end{enumerate}
\end{itemize}

\begin{failBox}{纯噪声数据泄漏实测结果：凭空造出 $+30.4\%$ 虚假幻觉！}
\begin{center}
\footnotesize
\begin{tabular}{@{}lccc@{}}
\toprule
\textbf{实验流水线（$N=60, D=500$ 纯随机白噪声，50 次平均）} & \textbf{实测测试准确率} & \textbf{理论真实准确率} & \textbf{虚假虚高幅度} \\
\midrule
泄漏版：先全量挑 Top-10 特征，再分 40/20 & \textbf{\color{ochre}80.40\%} & 50.00\% & \textbf{\color{ochre}+30.40\%} \\
干净版：先分 40/20，仅在 40 条训练集挑特征 & \textbf{\color{teal}49.50\%} & 50.00\% & $-0.50\%$（诚实！） \\
\bottomrule
\end{tabular}
\end{center}
\textbf{原理解剖}：在 $500$ 维纯噪声中，仅凭巧合就必然存在十几个在全部 60 条样本上与标签高度相关的“伪黄金特征”——它们之所以在后 20 条测试集上表现好，仅仅是因为你在挑特征时已经偷看过后 20 条测试集的标签！一旦上线面对第 61 条真正的新样本，这个号称 80.40\% 准确率的模型会瞬间跌回 50\% 瞎猜！
\end{failBox}

\subsection{红线二：混淆矩阵、96\% 负类准确率陷阱与精确评测体系}
\label{subsec:ch6-metrics}

哪怕杜绝了数据泄漏，如果选错了评估指标，依然会被表象骗得团团转。假设你开发一套信用卡欺诈检测（或罕见癌筛查）系统，真实数据中 \textbf{96\% 是正常负类（$y=0$）}，仅有 \textbf{4\% 是欺诈正类（$y=1$）}。此时一个什么都不学的“偷懒模型”只需写一行 \texttt{return 0}（永远预测正常），就能轻松拿到 \textbf{96.0\% 的准确率（Accuracy）}——听起来极其惊艳，可它抓住了几笔欺诈？\textbf{0 笔！}

要真正看清二分类器在正负两类上的功过得失，必须先拆解 \textbf{$2\times 2$ 混淆矩阵（Confusion Matrix）}（表~\ref{tab:ch6-cm}）：

\begin{table}[htbp]
\centering
\caption{二分类 $2\times 2$ 混淆矩阵（TP / FP / FN / TN）}
\label{tab:ch6-cm}
\footnotesize
\begin{tabular}{@{}lcc@{}}
\toprule
 & \textbf{预测为正类（$\hat{y} = 1$）} & \textbf{预测为负类（$\hat{y} = 0$）} \\
\midrule
\textbf{真实为正类（$y = 1$）} & \textbf{TP}（真正例，抓对了） & \textbf{FN}（假负例，漏报了） \\
\textbf{真实为负类（$y = 0$）} & \textbf{FP}（假正例，误报了） & \textbf{TN}（真负例，放对了） \\
\bottomrule
\end{tabular}
\end{table}

基于混淆矩阵的四个计数，现代分类工程构建了四把核心尺子：
\begin{enumerate}
  \item \textbf{精确率（Precision，“报正的里有多少真阳”）}：
  \[
    \mathrm{Precision} = \frac{\mathrm{TP}}{\mathrm{TP} + \mathrm{FP}}.
  \]
  垃圾邮件过滤、推荐拦截等“怕误杀好人”的场景极度看重 Precision。
  \item \textbf{召回率（Recall / TPR，“真阳里抓住了多少”）}：
  \[
    \mathrm{Recall} = \mathrm{TPR} = \frac{\mathrm{TP}}{\mathrm{TP} + \mathrm{FN}}.
  \]
  恶性肿瘤初筛、金融反欺诈、自动驾驶障碍物检测等“绝不能漏报”的场景极度看重 Recall。
  \item \textbf{$F_1$ 分数（Precision 与 Recall 的调和平均）}：
  \[
    F_1 = \frac{2\cdot \mathrm{Precision}\cdot \mathrm{Recall}}{\mathrm{Precision} + \mathrm{Recall}} = \frac{2\mathrm{TP}}{2\mathrm{TP} + \mathrm{FP} + \mathrm{FN}}.
  \]
  为何用调和平均而非算术平均？因为调和平均\textbf{严惩偏科}——在刚才 96\% 负类的偷懒模型上，$\mathrm{Recall} = 0$，算术平均或许还能糊弄出一个 $0.5$，而 $F_1$ 直接给出诚实的 $0.000$！
  \item \textbf{阈值无关曲线：ROC-AUC 与 PR-AUC}：
  逻辑回归输出的是连续概率 $\hat{p} \in (0, 1)$，改变判别阈值 $\tau \in [0, 1]$（$\hat{p} \ge \tau \Rightarrow \hat{y}=1$）会得到一整条权衡曲线：
  \begin{itemize}
    \item \textbf{ROC 曲线与 ROC-AUC}：以假正率 $\mathrm{FPR} = \frac{\mathrm{FP}}{\mathrm{FP}+\mathrm{TN}}$ 为横轴、真正率 $\mathrm{TPR} = \mathrm{Recall}$ 为纵轴绘制曲线，其曲线下面积即 ROC-AUC（随机瞎猜为 $0.5$，完美排序为 $1.0$）。它有一个极其优美的概率含义：\textbf{随机抽一条正样本和一条负样本，模型给正样本打分高于负样本的概率}！
    \item \textbf{PR 曲线与 PR-AUC}：以 Recall 为横轴、Precision 为纵轴。在正负极度失衡（如 $1:99$）时，庞大的 TN 会把 $\mathrm{FPR} = \frac{\mathrm{FP}}{\mathrm{FP}+\mathrm{TN}}$ 稀释得很小，使 ROC-AUC 看起来虚高；而 Precision 不含 TN，因此\textbf{极度不平衡任务中必须优先盯紧 PR-AUC 与 $F_1$}！
  \end{itemize}
\end{enumerate}

% ====================================================================
\section{端到端纯 NumPy 实战项目、实测 $L_2$ 扫参与置信区间}
\label{sec:ch6-capstone}

现在我们把两卷全部知识拧成一股绳，在真实运行的纯 NumPy 工程闭环中完成阶段一的毕业实战（完全对应 \texttt{code/vol2\_experiments.py} 中的 \texttt{exp6\_leakage\_and\_logreg}）：
\begin{itemize}
  \item \textbf{任务设定（高维多尺度噪声干扰分类）}：共 $N=300$ 条样本（正负类各 $150$ 条），特征维度 $D=30$，其中仅前 $2$ 维是携带类别信息的二维相关高斯信号，其余 $28$ 维全是尺度在 $[0.05, 80.0]$ 之间剧烈跨越的纯随机噪声！
  \item \textbf{无泄漏切分与标准化}：随机打乱后划分为 $120$ 条训练集与 $180$ 条验证集，构建 \texttt{StandardScaler} 严格仅在 $120$ 条训练集上拟合 $(\mu, \sigma)$ 再变换验证集。
  \item \textbf{梯度校验 + 全批量梯度下降训练 + $L_2$ 正则化扫参}：训练前先用第 1 卷第 6 章的中心有限差分 \texttt{check\_gradient} 校验解析梯度（实测相对误差 $\mathbf{2.20 \times 10^{-11}}$，彻底杜绝求导 Bug！），随后以学习率 $\eta = 0.25$ 训练 $800$ 轮，系统扫描 $\lambda \in \{0.00, 0.01, 0.10, 0.50, 2.00, 5.00\}$。
\end{itemize}

\subsection{模块化纯 NumPy 完整代码实现}
\label{subsec:ch6-code}

代码~\ref{lst:ch6-pipeline} 展示了高内聚、零框架依赖的面向对象实现：包含防泄漏标准化器、Log-Sum-Exp 稳定 $L_2$ 逻辑回归类、完整分类评测函数，以及第 1 卷有限差分梯度校验钩子。

\begin{lstlisting}[caption={纯 NumPy 端到端防泄漏逻辑回归流水线与有限差分梯度校验}, label={lst:ch6-pipeline}]
import numpy as np

class StandardScaler:
    """防泄漏 Z-score 标准化器：严格仅在训练集上 fit，再 transform 验证/测试集。"""
    def fit(self, X: np.ndarray):
        self.mu_ = X.mean(axis=0)
        self.sd_ = np.where(X.std(axis=0) > 1e-12, X.std(axis=0), 1.0)
        return self

    def transform(self, X: np.ndarray) -> np.ndarray:
        return (X - self.mu_) / self.sd_

class LogisticRegressionL2:
    """带 L2 正则化（高斯先验 MAP）与数值稳定 logaddexp 的纯 NumPy 逻辑回归。"""
    def __init__(self, lam: float = 0.10, lr: float = 0.25, epochs: int = 800):
        self.lam, self.lr, self.epochs = lam, lr, epochs

    @staticmethod
    def sigmoid(z: np.ndarray) -> np.ndarray:
        return 1.0 / (1.0 + np.exp(-np.clip(z, -500.0, 500.0)))

    def loss_and_grad(self, X: np.ndarray, y: np.ndarray, w: np.ndarray, b: float):
        z = X @ w + b
        # 数值稳定 BCE: mean(ln(1 + exp(z)) - y*z) + 0.5 * lam * ||w||^2 (偏置 b 不加惩罚)
        loss = float(np.mean(np.logaddexp(0.0, z) - y * z) + 0.5 * self.lam * np.dot(w, w))
        err = self.sigmoid(z) - y
        gw = (X.T @ err) / len(y) + self.lam * w
        gb = float(np.mean(err))
        return loss, gw, gb

    def fit(self, X: np.ndarray, y: np.ndarray):
        self.w = np.zeros(X.shape[1], dtype=np.float64)
        self.b = 0.0
        for _ in range(self.epochs):
            _, gw, gb = self.loss_and_grad(X, y, self.w, self.b)
            self.w -= self.lr * gw
            self.b -= self.lr * gb
        return self

    def predict_proba(self, X: np.ndarray) -> np.ndarray:
        return self.sigmoid(X @ self.w + self.b)

    def predict(self, X: np.ndarray, threshold: float = 0.5) -> np.ndarray:
        return (self.predict_proba(X) >= threshold).astype(np.float64)

def classification_metrics(y_true: np.ndarray, y_pred: np.ndarray):
    tp = int(np.sum((y_pred == 1) & (y_true == 1)))
    fp = int(np.sum((y_pred == 1) & (y_true == 0)))
    fn = int(np.sum((y_pred == 0) & (y_true == 1)))
    tn = int(np.sum((y_pred == 0) & (y_true == 0)))
    acc = (tp + tn) / len(y_true) * 100.0
    prec = tp / max(tp + fp, 1)
    rec = tp / max(tp + fn, 1)
    f1 = 2 * tp / max(2 * tp + fp + fn, 1)
    return {"Acc": acc, "P": prec, "R": rec, "F1": f1, "CM": (tp, fp, fn, tn)}

def verify_logreg_gradient(lam: float = 0.10, eps: float = 1e-5) -> float:
    """衔接第 1 卷第 6 章中心有限差分梯度校验：正式训练前必跑单元测试！"""
    rng = np.random.default_rng(7)
    Xm, ym = rng.normal(size=(20, 6)), rng.integers(0, 2, size=20).astype(np.float64)
    theta = rng.normal(scale=0.5, size=7)
    model = LogisticRegressionL2(lam=lam)
    _, gw, gb = model.loss_and_grad(Xm, ym, theta[:-1], float(theta[-1]))
    g_ana = np.append(gw, gb)
    g_num = np.zeros_like(theta)
    for i in range(len(theta)):
        tp, tm = theta.copy(), theta.copy()
        tp[i] += eps; tm[i] -= eps
        lp = model.loss_and_grad(Xm, ym, tp[:-1], float(tp[-1]))[0]
        lm = model.loss_and_grad(Xm, ym, tm[:-1], float(tm[-1]))[0]
        g_num[i] = (lp - lm) / (2.0 * eps)
    rel_err = float(np.linalg.norm(g_ana - g_num) / max(np.linalg.norm(g_ana) + np.linalg.norm(g_num), 1e-12))
    assert rel_err < 1e-9, f"Gradient check failed: {rel_err:.2e}"
    return rel_err  # 实测 rel_err = 2.20e-11
\end{lstlisting}

\subsection{实测 $L_2$ 扫参结果、U 形过拟合曲线与 Bootstrap 闭环}
\label{subsec:ch6-results}

在 $120$ 训练 / $180$ 验证、$30$ 维多尺度噪声数据集上运行 \texttt{code/vol2\_experiments.py}，完整扫参实测结果汇于表~\ref{tab:ch6-l2-sweep}，而无正则（$\lambda=0.00$）与最优正则（$\lambda=0.10$）的训练/验证交叉熵演化轨迹由真实生成的 \texttt{data/logreg\_curve.csv} 绘制于图~\ref{fig:ch6-logreg-curve}。

\begin{table}[htbp]
\centering
\caption{30 维多尺度噪声分类任务（120 训练 / 180 验证）的 $L_2$ 正则化扫参实测总表}
\label{tab:ch6-l2-sweep}
\footnotesize
\begin{tabular}{@{}ccccccc@{}}
\toprule
\textbf{$\lambda$（MAP 强度）} & \textbf{$\norm{w}_2$} & \textbf{训练准确率} & \textbf{验证准确率} & \textbf{验证 BCE 损失} & \textbf{混淆矩阵 (TP, FP, FN, TN)} & \textbf{验证 $F_1$} \\
\midrule
$0.00$（裸 MLE） & $5.825$ & $95.83\%$ & $78.89\%$ & $0.7219$ & $(66,\; 13,\; 25,\; 76)$ & $0.7765$ \\
$0.01$ & $3.226$ & $93.33\%$ & $81.67\%$ & $0.4576$ & $(69,\; 11,\; 22,\; 78)$ & $0.8070$ \\
\textbf{\color{teal}0.10（最优）} & \textbf{\color{teal}1.309} & \textbf{\color{teal}91.67\%} & \textbf{\color{teal}85.00\%} & \textbf{\color{teal}0.4060} & \textbf{\color{teal}(73, 9, 18, 80)} & \textbf{\color{teal}0.8439} \\
$0.50$ & $0.568$ & $89.17\%$ & $84.44\%$ & $0.5108$ & $(74,\; 11,\; 17,\; 78)$ & $0.8409$ \\
$2.00$ & $0.208$ & $86.67\%$ & $83.33\%$ & $0.6144$ & $(71,\; 10,\; 20,\; 79)$ & $0.8256$ \\
$5.00$（过强欠拟合） & $0.093$ & $85.83\%$ & $81.67\%$ & $0.6563$ & $(66,\; 8,\; 25,\; 81)$ & $0.8000$ \\
\bottomrule
\end{tabular}
\end{table}

\begin{figure}[htbp]
\centering
\begin{tikzpicture}
\begin{axis}[
  width=0.86\linewidth,
  height=6.5cm,
  xmin=0, xmax=790,
  ymin=0.10, ymax=0.91,
  xlabel={梯度下降迭代轮数（Epoch）},
  ylabel={未加惩罚的平均二元交叉熵（BCE Loss）},
  xlabel style={font=\small\sffamily},
  ylabel style={font=\small\sffamily},
  tick label style={font=\footnotesize},
  grid=major,
  grid style={dashed, gray!25},
  legend pos=north east,
  legend style={font=\footnotesize\sffamily, fill=white, fill opacity=0.92, draw=tealLine, legend columns=2},
  axis lines=left
]
  % lambda=0.00 Train & Val (real CSV!)
  \addplot[ochre, thick, dashed, mark=none] table[x=epoch, y=tr0, col sep=comma] {data/logreg_curve.csv};
  \addlegendentry{$\lambda=0.00$ 训练集}

  \addplot[ochre, ultra thick, mark=none] table[x=epoch, y=va0, col sep=comma] {data/logreg_curve.csv};
  \addlegendentry{$\lambda=0.00$ 验证集（U 形过拟合！）}

  % lambda=0.10 Train & Val (real CSV!)
  \addplot[teal, thick, dashed, mark=none] table[x=epoch, y=tr01, col sep=comma] {data/logreg_curve.csv};
  \addlegendentry{$\lambda=0.10$ 训练集}

  \addplot[teal, ultra thick, mark=none] table[x=epoch, y=va01, col sep=comma] {data/logreg_curve.csv};
  \addlegendentry{$\lambda=0.10$ 验证集（平稳最优）}

  % Annotate U-shape trough and overfitting divergence
  \fill[ochre] (axis cs:40, 0.38219) circle (2.2pt);
  \node[font=\scriptsize\sffamily\bfseries, color=ochre, align=center] at (axis cs:170, 0.56) {约 40--60 轮触底\\随后记忆 28 维噪声\\验证损失反升至 $0.715$！};
  \draw[-{Stealth[length=1.7mm]}, thin, ochre] (axis cs:95, 0.50) -- (axis cs:46, 0.395);

  \node[font=\scriptsize\sffamily\bfseries, color=tealDark, align=center] at (axis cs:597, 0.500) {$\lambda=0.10$（MAP 高斯先验）\\验证 BCE 稳稳锁定 $0.4060$\\验证准确率 $85.00\%,\; F_1=0.8439$};
  \draw[-{Stealth[length=1.7mm]}, thin, tealDark] (axis cs:597, 0.445) -- (axis cs:597, 0.412);

  \node[font=\scriptsize\sffamily, color=ochre!90!black] at (axis cs:620, 0.205) {$\lambda=0$ 训练损失一路跌至 $0.149$（死记硬背）};
\end{axis}
\end{tikzpicture}
\caption{实测训练与验证交叉熵曲线（数据源：\texttt{data/logreg\_curve.csv}）：裸 MLE（$\lambda=0.00$）在训练 60 轮后呈现教科书级 U 形过拟合，验证 BCE 恶化至 $0.7219$；加入高斯先验 MAP（$\lambda=0.10$）后，验证损失平稳收敛在 $0.4060$。}
\label{fig:ch6-logreg-curve}
\end{figure}

请停下来细细品味表~\ref{tab:ch6-l2-sweep} 与图~\ref{fig:ch6-logreg-curve} 中三幕完整的统计学剧情——它把第 5 章与第 6 章的每一块理论都变成了眼见为实的数字：
\begin{enumerate}
  \item \textbf{裸 MLE 过拟合（$\lambda = 0.00$）}：虽然训练集准确率高达 $95.83\%$、训练损失一路俯冲到 $0.149$，但模型拼命放大了 28 维噪声特征的权重（$\norm{w}_2$ 膨胀至 $5.825$，呼应命题~\ref{prop:ch6-weight-diverge}！），导致验证集准确率暴跌至 $78.89\%$，验证 BCE 更是走出刺眼的\textbf{U 形反弹曲线}恶化至 $0.7219$！
  \item \textbf{最优偏差-方差甜点（$\lambda = 0.10$）}：高斯先验 MAP 把权重范数温和压制在 $1.309$，虽然训练准确率略降到 $91.67\%$（主动接受微小偏差），却把验证准确率一口气拉升 $+6.11\%$ 至 $\mathbf{85.00\%}$（180 条样本答对 $73 + 80 = 153$ 条！），验证 BCE 骤降至 $\mathbf{0.4060}$，$F_1 = \mathbf{0.8439}$（实测 ROC-AUC $= \mathbf{0.9160}$）！更重要的是，把这 $153/180$ 条验证命中记录接回第 5 章第 5.4 节的 Bootstrap 流水线，我们不仅报出点估计 $85.00\%$，还能诚实报告其 \textbf{95\% Bootstrap 置信区间 $[80.00\%,\; 90.00\%]$}（$\widehat{\mathrm{SE}} = 2.63\%$）！
  \item \textbf{过强惩罚欠拟合（$\lambda = 5.00$）}：把 $\lambda$ 拧到 $5.00$ 时，先验把权重强行捏扁到 $\norm{w}_2 = 0.093$，连前 2 维真实信号也被过度压缩，验证准确率回落至 $81.67\%$、验证 BCE 升回 $0.6563$——这就是高偏差导致的\textbf{欠拟合（Underfitting）}！
\end{enumerate}

\subsection{全卷收束与第 3 卷预告}
\label{subsec:ch6-wrapup}

恭喜你！至此你已经完整通关《本科人工智能七卷本自学讲义》的\textbf{第一阶段（第 1 卷 + 第 2 卷）数理与工程底座}：
\begin{itemize}
  \item 你不再把损失函数当成黑盒魔法——你知道最小二乘是高斯 MLE、岭回归与 Lasso 是高斯/拉普拉斯 MAP、二元与多分类交叉熵是 Bernoulli/Categorical MLE 与最小 KL 散度；
  \item 你不再怕高维矩阵求导与数值爆炸——你会用迹技巧推梯度、用 \texttt{check\_gradient} 把梯度 Bug 扼杀在 $10^{-11}$ 精度下、用 Log-Sum-Exp 护航指数运算；
  \item 你拥有了严苛的实验科学素养——你识破了数据泄漏的 80\% 纯噪声骗局、看穿了 96\% 负类的准确率陷阱，并会用 Bootstrap 置信区间克制地解读每一次实验涨落。
\end{itemize}

\begin{insightBox}{迈向第 3 卷：《经典人工智能：搜索、博弈、概率图与序列决策》}
手握第 1--2 卷的硬核数学武器，接下来第 3 卷将带你正式踏入人工智能的核心推理与决策殿堂（对标 Berkeley CS188 / Stanford CS221）：
\begin{enumerate}
  \item \textbf{启发式搜索与对抗博弈}：从 BFS/Dijkstra 跃升到 $A^*$ 启发式可采纳性证明、Alpha-Beta 剪枝与蒙特卡洛树搜索（MCTS，直接动用第 2 卷第 3 章的 Hoeffding 界构造 UCB 公式！）；
  \item \textbf{概率图模型（PGM）}：将第 2 卷第 1 章的\textbf{条件独立、d-分离与 Explaining Away} 扩展为完整的贝叶斯网络、变量消元精确推断与吉布斯/重要性采样近似推断；
  \item \textbf{马尔可夫决策过程（MDP）与强化学习底座}：用动态规划（值迭代、策略迭代）与时序差分学习解决序贯决策，为你后续攀登第 4 卷经典机器学习、第 5 卷深度学习与第 6 卷大模型/强化学习前沿铺设坚不可摧的阶梯！
\end{enumerate}
\end{insightBox}

% ====================================================================
\section{本章习题}
\label{sec:ch6-exercises}

\exTier{A. 概念辨析与反例查错}{严防正则化、凸优化与评测工程的高频失误}
\begin{enumerate}[label=\textbf{6.\arabic*}]
  \item \textbf{判断正误并纠错}：请判断以下三句话的真假，并分别用本章定理或具体反例说明原因：
  (a) “在做带 $L_2$ 正则化的线性回归或逻辑回归时，应该把偏置项 $b$ 也放进 $\frac{\lambda}{2}(\norm{w}_2^2 + b^2)$ 一起惩罚，正则化效果更彻底。”
  (b) “只要训练数据完全线性可分，不加正则化的逻辑回归经梯度下降充分训练后一定会收敛到唯一的有限全局最优权重 $w^*$。”
  (c) “在仅有 1\% 欺诈交易的金融反欺诈任务上，某分类器在测试集上取得了 99.0\% 的准确率，说明它已经达到极高商用水平。”

  \item \textbf{Softmax 数值溢出与数据泄漏双排查}：某同学编写了一套多分类流水线，代码包含如下两段实现：
  \begin{enumerate}[label=(\roman*)]
    \item 前向与损失：\texttt{P = np.exp(Z) /} \texttt{np.sum(np.exp(Z),} \texttt{axis=1,} \texttt{keepdims=True)}，随后计算损失 \texttt{loss =} \texttt{-np.mean(}\allowbreak\texttt{np.log(}\allowbreak\texttt{P[np.arange(n), y]))}；
    \item 预处理与划分：先在全量数据 \texttt{X\_all} 上执行 \texttt{X\_all =} \texttt{(X\_all - X\_all.mean(0))} \texttt{/ X\_all.std(0)} 并挑选与 \texttt{y\_all} 相关性最高的 20 维特征，再切分为 70\% 训练集与 30\% 测试集。
  \end{enumerate}
  请分别指出两段代码会在什么具体场景下崩盘或产生严重幻觉，并写出修复后的正确 NumPy 写法。
\end{enumerate}

\exTier{B. 具体小例手算}{算通逻辑回归单步更新、线性边界偏移与混淆矩阵}
\begin{enumerate}[label=\textbf{6.\arabic*}, resume]
  \item \textbf{逻辑回归一步梯度更新与远点拖拽验算}：
  (a) 设单条二维样本 $x = (1, -2)^T$，标签 $y = 1$，当前参数 $w_0 = (0, 0)^T$、$b_0 = 0$，$L_2$ 系数 $\lambda = 0.2$，学习率 $\eta = 0.5$。手算初始 logit $z_0$、预测概率 $\hat{p}_0$、带正则项的 BCE 损失值 $\mathcal{L}(w_0, b_0)$、梯度 $(\nabla_w \mathcal{L}, \partial\mathcal{L}/\partial b)$，以及一步梯度下降后的 $(w_1, b_1)$ 与新预测概率 $\hat{p}_1$；
  (b) 验算例~\ref{ex:ch6-ols-drag}：给定 7 个点 $(-2,0), (-1,0), (0,0), (1,1), (2,1), (3,1), (30,1)$，手算样本均值 $\bar{x}, \bar{y}$、离差和 $\sum (x_i - \bar{x})^2$ 与 $\sum (x_i-\bar{x})(y_i-\bar{y})$，验证最小二乘斜率 $w = 15/668$、截距 $b = 311/668$，且 $0.5$ 阈值边界确为 $x_1^* = 23/15$。

  \item \textbf{不平衡混淆矩阵指标手算与阈值权衡}：某罕见病筛查模型在 $N=200$ 名受检者（其中 20 名真实患病）上测得混淆矩阵：$\mathrm{TP}=16,\; \mathrm{FP}=4,\; \mathrm{FN}=4,\; \mathrm{TN}=176$。
  (a) 手算 Accuracy、Precision、Recall、FPR 与 $F_1$ 分数；
  (b) 与始终预测健康的“偷懒基线”（$\hat{y}\equiv 0$）对比 Accuracy 与 $F_1$；
  (c) 若将分类阈值调低使混淆矩阵变为 $\mathrm{TP}=19,\; \mathrm{FP}=19,\; \mathrm{FN}=1,\; \mathrm{TN}=161$，重算 Precision、Recall 与 $F_1$，并分析在“癌症初筛”与“自动封号系统”两种场景下应如何抉择阈值。
\end{enumerate}

\exTier{C. AI 核心推演与实用证明}{夯实 Softmax 梯度推导与逻辑回归安全步长界}
\begin{enumerate}[label=\textbf{6.\arabic*}, resume]
  \item \textbf{Softmax 平移不变性、交叉熵梯度与 $K=2$ 退化证明}：对 $K$ 分类样本，logits $z \in \R^K$，Softmax 概率 $\hat{p}_k = e^{z_k} / \sum_{c=1}^K e^{z_c}$，one-hot 标签 $y \in \{0,1\}^K$，交叉熵损失 $\ell(z, y) = -\sum_{k=1}^K y_k \ln \hat{p}_k$。
  (a) 证明对任意常数 $c \in \R$ 恒有 $\mathrm{softmax}(z + c\mathbf{1}) = \mathrm{softmax}(z)$（即 Log-Sum-Exp 减最大值的代数依据）；
  (b) 严格求导证明 $\frac{\partial \ell}{\partial z_j} = \hat{p}_j - y_j$（对任意 $j=1,\dots,K$）；
  (c) 说明当 $K=2$ 时，将 Softmax 回归的两组参数 $(w_1, b_1), (w_2, b_2)$ 平移消去一组（令 $w_0' = 0, b_0' = 0$，记 $w = w_1 - w_0,\; b = b_1 - b_0$）即精确退化为二元逻辑回归。
  \exHint{(b) 直接把 $\ell(z,y)$ 化简为 $\ln\big(\sum_{c=1}^K e^{z_c}\big) - \sum_{k=1}^K y_k z_k$ 再对 $z_j$ 求偏导，两行即可证明！}

  \item \textbf{Sigmoid + MSE 非凸区间与逻辑回归 Hessian 安全学习率上界}：
  (a) 完整验证式~\eqref{eq:ch6-sigmse-d2}：当 $y=1$、$p=\sigma(z)$ 时，$\ell_{\text{MSE}}(z) = \frac{1}{2}(p-1)^2$ 的二阶导数为 $\ell''_{\text{MSE}}(z) = p(1-p)^2(3p-1)$，并由此解出严格凹区间为 $z < -\ln 2$；
  (b) 对带 $L_2$ 正则的无偏置二元逻辑回归 $\mathcal{L}(w) = \frac{1}{N}\sum_{i=1}^N [\ln(1+e^{x_i^T w}) - y_i x_i^T w] + \frac{\lambda}{2}\norm{w}_2^2$，利用基本不等式 $p(1-p) \le \frac{1}{4}$ 证明其 Hessian 处处满足矩阵不等式
  \[
    \lambda I_d \;\preceq\; H(w) \;\preceq\; \left(\frac{1}{4N}\lambda_{\max}(X^T X) + \lambda\right) I_d,
  \]
  并结合第 1 卷第 6 章的平滑凸函数步长界，写出保证全批量梯度下降稳定收敛的安全学习率上界 $\eta_{\max}$。
  \exHint{(b) 对任意单位向量 $v$，用 $0 < \hat{p}_i(1-\hat{p}_i) \le \frac{1}{4}$ 放缩二次型 $v^T H(w) v = \frac{1}{N}\sum_i \hat{p}_i(1-\hat{p}_i)(x_i^T v)^2 + \lambda$，再结合第 1 卷第 3 章 Rayleigh 商上界 $v^T(X^T X)v \le \lambda_{\max}(X^T X)$！}
\end{enumerate}

\exTier{D. 可自检编程实验}{动手交付多分类 Softmax 回归、数据泄漏复现与 ROC-AUC}
\begin{enumerate}[label=\textbf{6.\arabic*}, resume]
  \item \textbf{纯 NumPy 多分类 Softmax 回归 + Log-Sum-Exp + 有限差分梯度校验}：固定随机种子 \texttt{2026}，生成 $K=3$ 类二维高斯混合数据（每类 $100$ 条，均值分别为 $(-2,0), (2,0), (0, 2.5)$，协方差均为 $0.8 I_2$），按 $70\%/30\%$ 分层划分训练/验证集并严格仅在训练集上标准化。手写带 Log-Sum-Exp 稳定化与 $L_2$ 正则（$\lambda = 0.05$）的 \texttt{SoftmaxRegressionL2} 类。\\
  \textbf{自检标准}：
  (a) 在迷你张量上调用中心有限差分梯度校验（$\epsilon = 10^{-5}$），断言对 $(W, b)$ 的整体解析梯度相对误差 $< 10^{-9}$；
  (b) 对极端 logits \texttt{z = np.array([1000.0, 1001.0, 999.0])} 调用你的 \texttt{softmax} 与 \texttt{logsumexp}，断言输出不含 \texttt{nan}/\texttt{inf} 且概率和与 $1.0$ 误差 $< 10^{-12}$；
  (c) 在 $90$ 条验证集上分类准确率 $> 90\%$（参考值约 $92.22\%$）。
  \exHint{将矩阵 $W$ 与偏置 $b$ 展平拼为一维向量 \texttt{theta =} \texttt{np.concatenate(}\allowbreak\texttt{[W.ravel(), b])}，传入第 1 卷的 \texttt{check\_gradient} 即可一键校验全部多分类参数梯度！}

  \item \textbf{复现数据泄漏灾难并手写 ROC-AUC 阈值扫描}：
  \begin{enumerate}[label=(\alph*)]
    \item 复现第 6.4 节的纯噪声实验（$N=60, D=500$，独立重复 50 次，40 训练 / 20 测试，按相关性选 Top-10 特征），对比泄漏版与干净版的平均测试准确率；
    \item 复用第 6.5 节的 30 维多尺度逻辑回归数据与最优模型（$\lambda = 0.10$），按预测概率从高到低排序，用纯 NumPy 手写梯形积分（\texttt{np.trapezoid} 或手动相邻差分累加）计算验证集上的精确 \textbf{ROC-AUC}，并打印阈值 $\tau \in \{0.2, 0.3, 0.4, 0.5, 0.6, 0.7, 0.8\}$ 下的 Precision、Recall 与 $F_1$ 权衡表。
  \end{enumerate}
  \textbf{自检标准}：(a) 断言泄漏版 50 次平均测试准确率 $> 75\%$（参考值 $80.40\%$），而干净版落在 $[45\%, 55\%]$ 内（参考值 $49.50\%$）；(b) 断言验证集 ROC-AUC 落在 $[0.90, 0.95]$ 之间（参考值 $0.9160$），且在 $\tau = 0.5$ 时准确率 $= 85.00\%$、$F_1 = 0.8439$。
  \exHint{手算 ROC-AUC 最稳健的纯 NumPy 写法：按 \texttt{order = np.argsort(-scores)} 降序排列真实标签 \texttt{y\_s = y\_va[order]}，用 \texttt{np.cumsum} 一次算出所有阈值点的 \texttt{TPR} 与 \texttt{FPR}，首端拼 $0$ 后用梯形面积公式求和即可！}
\end{enumerate}


% ==================== Back Matter / Appendix ====================
\appendix
% !TeX root = ../main.tex
\ctexset{
  chapter={
    name={附录\ ,},
    number=\Alph{chapter},
    aftername={\hfill{\fontsize{38}{38}\selectfont\sffamily\bfseries\color{tealMist}\thechapter}\par\vspace{-4mm}}
  }
}
\renewcommand{\chaptermark}[1]{\markboth{附录 \thechapter\quad #1}{}}

\chapter[全卷 48 道分层习题完整精解]{全卷 48 道分层习题完整精解}
\label{app:solutions}

\chapEpigraph{如果你想学会游泳，就必须下水；如果你想成为解题高手，就必须去解题。}{乔治·波利亚（George P\'olya），《怎样解题》}

\begin{chapIntro}{本附录使用建议}
本附录给出第 1--6 章全部 48 道分层习题的完整解答：
\begin{itemize}
  \item \textbf{A 层（概念辨析与反例查错）}：先给出明确正误判定，再给出最小显式反例或故障根因与修复代码；
  \item \textbf{B 层（小例手算与数值直觉）}：保留精确分数与四位小数，全部中间算术均已由 Python 脚本交叉核验；
  \item \textbf{C 层（实用推导与性质证明）}：给出零跳步的完整证明，并标注其在后续卷次（搜索、决策树、VAE、大模型对齐）的直接落点；
  \item \textbf{D 层（编程实战与可自检实验）}：给出纯 NumPy 核心实现逻辑、固定随机种子 \texttt{2026} 下的实测参考输出与断言自检说明。
\end{itemize}
建议先合上附录独立作答与编码，再逐项对照推演细节与数值断言。
\end{chapIntro}

% =========================================================================
\section{第 1 章习题精解（1.1--1.8）}
% =========================================================================

\subsection*{习题 1.1\quad 独立、互斥与条件独立辨析}
\begin{enumerate}[label=(\alph*), leftmargin=1.8em]
  \item \textbf{错误。} 若 $A \cap B = \varnothing$，则 $\Prob(A \cap B) = 0$；但已知 $\Prob(A) > 0$ 且 $\Prob(B) > 0$，故 $\Prob(A)\Prob(B) > 0 \neq \Prob(A \cap B)$。事实上，正概率的互斥事件恰恰是\textbf{极度负相关}的（$A$ 发生即保证 $B$ 绝不发生）。\textbf{最小反例}：抛一枚均匀硬币，$A = \{\text{正面}\}$，$B = \{\text{反面}\}$，$\Prob(A)=\Prob(B)=0.5$，但 $\Prob(A\cap B)=0 \neq 0.25$。
  \item \textbf{错误。} 边际独立不能推出条件独立（典型机制为汇合结构 / Explaining Away）。\textbf{最小反例}：设 $X, Y \stackrel{\text{i.i.d.}}{\sim} \mathrm{Bernoulli}(0.5)$，令 $Z = X \oplus Y$（异或，或令 $Z = X + Y$）。显然 $X \indep Y$；但一旦给定 $Z = 1$，只要知道 $X = 1$ 就必然确定 $Y = 0$，即 $\Prob(Y=0 \mid X=1, Z=1) = 1 \neq \Prob(Y=0 \mid Z=1) = 0.5$，故 $X \not\indep Y \mid Z$。
  \item \textbf{错误。} 概率密度函数 $p_X(x)$ 只需满足非负性 $p_X(x) \ge 0$ 与总积分为 $1$（$\int_{-\infty}^{+\infty} p_X(x)\,dx = 1$），单点高度完全可以大于 $1$。\textbf{最小反例}：均匀分布 $X \sim \mathrm{Uniform}(0, 0.5)$，在 $x \in [0, 0.5]$ 上 $p_X(x) = 2 > 1$。
  \item \textbf{正确。} 期望的线性性 $\E[X + Y] = \E[X] + \E[Y]$ 直接来源于求和/积分运算的线性性，\textbf{完全不要求} $X$ 与 $Y$ 独立或不相关，只要 $\E[|X|] < \infty$ 且 $\E[|Y|] < \infty$ 即恒成立。
\end{enumerate}

\subsection*{习题 1.2\quad 协方差为零与线性相关性陷阱}
\begin{enumerate}[label=(\arabic*), leftmargin=1.8em]
  \item \textbf{验证 $\Cov(X, Y) = 0$ 但 $X \not\indep Y$}：随机变量 $X$ 等概率取 $\{-2, 0, 2\}$，对应的 $Y = |X|$ 分别取 $\{2, 0, 2\}$，积 $XY$ 分别取 $\{-4, 0, 4\}$。因此：
  \[
    \E[X] = \frac{-2 + 0 + 2}{3} = 0, \qquad
    \E[Y] = \frac{2 + 0 + 2}{3} = \frac{4}{3}, \qquad
    \E[XY] = \frac{-4 + 0 + 4}{3} = 0.
  \]
  于是 $\Cov(X, Y) = \E[XY] - \E[X]\E[Y] = 0 - 0\cdot \frac{4}{3} = 0$。然而，$Y$ 是 $X$ 的确定性非线性函数，显然不独立：例如 $\Prob(X = 0, Y = 2) = 0$，但 $\Prob(X = 0)\Prob(Y = 2) = \frac{1}{3} \cdot \frac{2}{3} = \frac{2}{9} \neq 0$。
  \item \textbf{为何 $\Var(X+Y) = \Var(X) + \Var(Y)$ 不能推出独立}：因为一般恒等式 $\Var(X+Y) = \Var(X) + \Var(Y) + 2\Cov(X, Y)$ 表明，方差可加性等价于 $\Cov(X, Y) = 0$（线性不相关），而非独立。直接沿用上例验证：
  \begin{align*}
    \Var(X) &= \frac{(-2)^2 + 0^2 + 2^2}{3} = \frac{8}{3}, \qquad
    \Var(Y) = \E[Y^2] - (\E[Y])^2 = \frac{8}{3} - \frac{16}{9} = \frac{8}{9}, \\
    X + Y &\in \{0, 0, 4\} \implies \E[X+Y] = \frac{4}{3},\quad
    \Var(X+Y) = \frac{0^2 + 0^2 + 4^2}{3} - \frac{16}{9} = \frac{32}{9} = \frac{8}{3} + \frac{8}{9}.
  \end{align*}
  尽管 $\Var(X+Y) = \Var(X) + \Var(Y)$ 精确成立，$X$ 与 $Y$ 仍强依赖。
\end{enumerate}

\subsection*{习题 1.3\quad 两厂芯片溯源与检测后验}
记 $H_A, H_B$ 为芯片来自 A、B 厂，$E$ 为“芯片是次品”，$A$ 为“检测仪报警”。已知 $\Prob(H_A)=0.70$、$\Prob(H_B)=0.30$、$\Prob(E\mid H_A)=0.01$、$\Prob(E\mid H_B)=0.04$。
\begin{enumerate}[label=(\arabic*), leftmargin=1.8em]
  \item 由全概率公式，随机一颗芯片为次品的总概率为：
  \[
    \Prob(E) = \Prob(H_A)\Prob(E\mid H_A) + \Prob(H_B)\Prob(E\mid H_B)
             = 0.70 \times 0.01 + 0.30 \times 0.04 = 0.007 + 0.012 = 0.019.
  \]
  由 Bayes 定理，次品来自 B 厂的后验概率为：
  \[
    \Prob(H_B \mid E) = \frac{\Prob(H_B)\Prob(E\mid H_B)}{\Prob(E)}
                      = \frac{0.012}{0.019} = \frac{12}{19} \approx 0.6316 \;(63.16\%).
  \]
  \item 已知 $\Prob(A\mid E)=0.95$、$\Prob(A\mid \neg E)=0.02$，且由 (1) 得 $\Prob(E)=0.019$、$\Prob(\neg E)=0.981$。再次使用全概率公式求触发报警的总概率：
  \[
    \Prob(A) = \Prob(A\mid E)\Prob(E) + \Prob(A\mid \neg E)\Prob(\neg E)
             = 0.95 \times 0.019 + 0.02 \times 0.981 = 0.01805 + 0.01962 = 0.03767.
  \]
  报警芯片确实为次品的后验概率为：
  \[
    \Prob(E \mid A) = \frac{0.01805}{0.03767} = \frac{1805}{3767} \approx 0.4792 \;(47.92\%).
  \]
  注意：尽管检测仪灵敏度高达 $95\%$、误报率仅 $2\%$，但由于芯片整体次品基率仅 $1.9\%$，报警芯片中仍有超过一半是良品（假阳性），这是基率效应的典型体现。
\end{enumerate}

\subsection*{习题 1.4\quad 草地湿润中的 Explaining Away 精确手算}
因为 $R \indep S$，四组联合概率 $\Prob(R=r, S=s, W=1) = \Prob(R=r)\Prob(S=s)\Prob(W=1\mid R=r, S=s)$ 分别为：
\begin{align*}
  \Prob(R=1, S=1, W=1) &= 0.2 \times 0.1 \times 0.99 = 0.0198, \\
  \Prob(R=1, S=0, W=1) &= 0.2 \times 0.9 \times 0.80 = 0.1440, \\
  \Prob(R=0, S=1, W=1) &= 0.8 \times 0.1 \times 0.90 = 0.0720, \\
  \Prob(R=0, S=0, W=1) &= 0.8 \times 0.9 \times 0.00 = 0.0000.
\end{align*}
\begin{enumerate}[label=(\arabic*), leftmargin=1.8em]
  \item 总概率 $\Prob(W=1) = 0.0198 + 0.1440 + 0.0720 + 0 = 0.2358 = \frac{1179}{5000}$。因此仅观测到草地湿润时，洒水器开启的后验概率为：
  \[
    \Prob(S=1 \mid W=1) = \frac{0.0198 + 0.0720}{0.2358} = \frac{0.0918}{0.2358} = \frac{51}{131} \approx 0.3893 \;(38.93\%).
  \]
  \item 进一步确认昨晚下过雨（$R=1$）后，在 $R=1, W=1$ 条件下的总权重为 $0.0198 + 0.1440 = 0.1638$，故：
  \[
    \Prob(S=1 \mid W=1, R=1) = \frac{0.0198}{0.0198 + 0.1440} = \frac{0.0198}{0.1638} = \frac{11}{91} \approx 0.1209 \;(12.09\%).
  \]
  \item \textbf{变化方向与物理含义}：$\Prob(S=1)=10.00\% \xrightarrow{\text{观测 }W=1} 38.93\% \xrightarrow{\text{再观测 }R=1} 12.09\%$。看到草地湿润大幅抬高了“洒水器开了”的怀疑度；但一旦得知昨晚下过雨，“下雨”已经充分解释了草地湿润，因而将洒水器开启的后验概率压回至接近先验水平——这就是贝叶斯网络汇合结构（Collider）中的\textbf{辩解效应（Explaining Away）}。
\end{enumerate}

\subsection*{习题 1.5\quad 全方差定律（Law of Total Variance）证明}
\begin{proof}
由方差定义 $\Var(X) = \E[X^2] - (\E[X])^2$，并在等式右侧同时加上和减去 $\E\bigl[(\E[X \mid Y])^2\bigr]$：
\[
  \Var(X) = \Bigl(\E[X^2] - \E\bigl[(\E[X \mid Y])^2\bigr]\Bigr)
          + \Bigl(\E\bigl[(\E[X \mid Y])^2\bigr] - (\E[X])^2\Bigr).
\]
分别化简两项：
\begin{enumerate}[label=(\arabic*), leftmargin=1.8em]
  \item 对第一项，由全期望定律有 $\E[X^2] = \E\bigl[\E[X^2 \mid Y]\bigr]$，再由期望线性性合并：
  \[
    \E[X^2] - \E\bigl[(\E[X \mid Y])^2\bigr]
    = \E\Bigl[\E[X^2 \mid Y] - (\E[X \mid Y])^2\Bigr]
    = \E\bigl[\Var(X \mid Y)\bigr],
  \]
  其中利用了条件方差展开 $\Var(X \mid Y) = \E[(X - \E[X\mid Y])^2 \mid Y] = \E[X^2 \mid Y] - (\E[X\mid Y])^2$。
  \item 对第二项，记随机变量 $U := \E[X \mid Y]$，由全期望定律得 $\E[U] = \E[\E[X \mid Y]] = \E[X]$，因此：
  \[
    \E\bigl[(\E[X \mid Y])^2\bigr] - (\E[X])^2 = \E[U^2] - (\E[U])^2 = \Var(U) = \Var\bigl(\E[X \mid Y]\bigr).
  \]
\end{enumerate}
将两项相加即得 $\Var(X) = \E[\Var(X \mid Y)] + \Var(\E[X \mid Y])$。
\end{proof}
\textbf{直观解释}：若将 $Y$ 视为类别/分组标签，则 $\E[\Var(X \mid Y)]$ 度量\textbf{类内平均残差方差}（知道 $Y$ 之后仍无法消除的剩余随机性），$\Var(\E[X \mid Y])$ 度量\textbf{类间条件均值的方差}（能被 $Y$ 解释掉的那部分变化量）；总方差恰好等于二者之和。

\subsection*{习题 1.6\quad Log-Odds 形式下的朴素贝叶斯线性分类器本质}
\begin{proof}
由 Bayes 定理与条件独立假设 \eqref{eq:ch1-naive-bayes}，两类后验概率之比消去了相同的证据分母 $\Prob(x)$：
\[
  \frac{\Prob(y=1 \mid x)}{\Prob(y=0 \mid x)}
  = \frac{\Prob(y=1)\prod_{j=1}^d \Prob(x_j \mid y=1)}{\Prob(y=0)\prod_{j=1}^d \Prob(x_j \mid y=0)}
  = \frac{\pi}{1 - \pi} \prod_{j=1}^d \frac{p_j^{x_j}(1 - p_j)^{1 - x_j}}{q_j^{x_j}(1 - q_j)^{1 - x_j}}.
\]
两边取自然对数，并将含 $x_j$ 的项与不含 $x_j$ 的常数项分离：
\begin{align*}
  \log \frac{\Prob(y=1 \mid x)}{\Prob(y=0 \mid x)}
  &= \log \frac{\pi}{1 - \pi}
     + \sum_{j=1}^d \Bigl[x_j \log \frac{p_j}{q_j} + (1 - x_j)\log \frac{1 - p_j}{1 - q_j}\Bigr] \\
  &= \sum_{j=1}^d x_j \log \frac{p_j(1 - q_j)}{q_j(1 - p_j)}
     + \Bigl(\log \frac{\pi}{1 - \pi} + \sum_{j=1}^d \log \frac{1 - p_j}{1 - q_j}\Bigr)
   = w^T x + b.
\end{align*}
其中权重向量 $w \in \R^d$ 与偏置 $b \in \R$ 的显式表达式为：
\[
  w_j = \log \frac{p_j(1 - q_j)}{q_j(1 - p_j)} \quad (j=1,\dots,d), \qquad
  b = \log \frac{\pi}{1 - \pi} + \sum_{j=1}^d \log \frac{1 - p_j}{1 - q_j}.
\]
由此立得 $\Prob(y=1 \mid x) = \sigma(w^T x + b)$，说明二值朴素贝叶斯（生成式模型）与第 6 章的逻辑回归（判别式模型）共享完全相同的线性对数几率决策边界，区别仅在于前者通过估计联合分布间接推出 $(w, b)$，后者直接极小化交叉熵拟合 $(w, b)$。
\end{proof}

\subsection*{习题 1.7\quad 高维对数空间朴素贝叶斯与 Laplace 平滑实验}
\textbf{核心纯 NumPy 实现逻辑}：将特征矩阵 $X \in \{0,1\}^{N\times d}$ 的对数联合概率写成矩阵乘法 $X \log P^T + (1 - X)\log(1 - P)^T + \log\pi$，再用 Log-Sum-Exp 归一化避免 $d=4000$ 维连乘下溢：
\begin{lstlisting}[language=Python]
def fit_nb(X, y, alpha=1.0):
    pi = np.array([np.mean(y == 0), np.mean(y == 1)])
    p = np.vstack([
        (X[y == c].sum(axis=0) + alpha) / (np.sum(y == c) + 2.0 * alpha)
        for c in (0, 1)
    ])
    return pi, p

def predict_proba_log(X, pi, p):
    lj = np.log(pi)[None, :] + X @ np.log(p).T + (1.0 - X) @ np.log(1.0 - p).T
    m = np.max(lj, axis=1, keepdims=True)
    log_z = m + np.log(np.sum(np.exp(lj - m), axis=1, keepdims=True))
    return np.exp(lj - log_z)
\end{lstlisting}
\textbf{验证输出与断言自检（\texttt{default\_rng(2026)}）}：
\begin{itemize}
  \item \textbf{直接概率连乘}：在 $d=4000$ 维下，每条样本约含 200 个出现词与 3800 个未出现词，连乘量级约 $0.05^{200} \cdot 0.95^{3800} \approx 10^{-345}$，低于 IEEE 754 float64 最小正数（$\approx 2.2 \times 10^{-308}$），实测 $95.0\%$ 的测试样本两类分子同时下溢为 \texttt{0.0}，无法计算后验。
  \item \textbf{$\alpha=0$（无平滑）}：实测训练集存在未见词（$\hat{p}_{jc}=0$），对数空间触发 $\log 0 = -\infty$（28 处 \texttt{-inf}），在 $0 \cdot (-\infty)$ 运算中产生大量 \texttt{NaN}。
  \item \textbf{$\alpha=1.0$（Laplace 平滑 + 对数空间）}：零 \texttt{NaN}，测试集准确率达 \textbf{95.25\%}（$>90\%$），所有样本后验概率和与 $1$ 的最大偏差为 $5.68 \times 10^{-14} < 10^{-12}$，全部自检断言通过。
\end{itemize}

\subsection*{习题 1.8\quad 蒙特卡洛验证 Explaining Away 与条件协方差}
\textbf{核心纯 NumPy 实现逻辑}：
\begin{lstlisting}[language=Python]
rng = np.random.default_rng(2026)
N = 200_000
R = (rng.random(N) < 0.2).astype(int)
S = (rng.random(N) < 0.1).astype(int)
pw = np.where((R == 1) & (S == 1), 0.99,
     np.where((R == 1) & (S == 0), 0.80,
     np.where((R == 0) & (S == 1), 0.90, 0.00)))
W = (rng.random(N) < pw).astype(int)
corr_marg = np.corrcoef(R, S)[0, 1]
ps_w = S[W == 1].mean()
ps_wr = S[(W == 1) & (R == 1)].mean()
corr_cond = np.corrcoef(R[W == 1], S[W == 1])[0, 1]
\end{lstlisting}
\textbf{验证输出与断言自检（\texttt{default\_rng(2026)}）}：
\begin{itemize}
  \item 边际相关系数 $\widehat{\mathrm{Corr}}(R, S) = +0.00321$（$|r| < 0.01$，确认边际独立）；
  \item $\hat{\Prob}(S=1 \mid W=1) = 0.38885$（与理论值 $51/131 \approx 0.38931$ 偏差 $0.00047 < 0.005$）；
  \item $\hat{\Prob}(S=1 \mid W=1, R=1) = 0.12341$（与理论值 $11/91 \approx 0.12088$ 偏差 $0.00253 < 0.005$）；
  \item 条件相关系数 $\widehat{\mathrm{Corr}}(R, S \mid W=1) = -0.82621$（理论精确值为 $-0.83037$，稳稳落入 $[-0.86, -0.80]$），定量证明观测共同结果会诱发强烈的条件负相关。
\end{itemize}

% =========================================================================
\section{第 2 章习题精解（2.1--2.8）}
% =========================================================================

\subsection*{习题 2.1\quad 核心分布与多元高斯概念辨析}
\begin{enumerate}[label=(\alph*), leftmargin=1.8em]
  \item \textbf{错误。} 两个边际分布皆为标准正态，不保证联合分布为二维高斯。\textbf{经典反例}：设 $X_1 \sim \mathcal{N}(0, 1)$，独立取对称符号 $S \in \{-1, +1\}$（等概率 $0.5$），令 $X_2 = S X_1$。显然 $X_2$ 边际分布亦为 $\mathcal{N}(0, 1)$，但 $(X_1, X_2)$ 的概率质量全部集中在两条对角线 $x_2 = \pm x_1$ 上，且 $X_1 + X_2$ 在 $0$ 点具有 $0.5$ 的离散点质量，绝非正态分布。
  \item \textbf{正确。} 对 $X \sim \mathrm{Pois}(\lambda)$，由幂级数展开直接验证 $\E[X] = \lambda$ 且 $\E[X(X-1)] = \lambda^2$，故 $\Var(X) = \E[X(X-1)] + \E[X] - (\E[X])^2 = \lambda^2 + \lambda - \lambda^2 = \lambda = \E[X]$。
  \item \textbf{正确。} 对\textbf{联合高斯}向量，零协方差即意味着协方差矩阵为对角阵 $\Sigma = \diag(\sigma_1^2, \sigma_2^2)$，联合密度函数精确分解为两个一维高斯密度的乘积 $p(x_1, x_2) = p(x_1)p(x_2)$，故 $X_1 \indep X_2$。（注意：离开“联合高斯”前提该结论失效，如习题 1.2。）
  \item \textbf{错误。} 由多元高斯条件分布定理（定理 \ref{thm:ch2-cond-mvg}），条件协方差矩阵为 Schur 补 $\Sigma_{1\mid 2} = \Sigma_{11} - \Sigma_{12}\Sigma_{22}^{-1}\Sigma_{21}$，其表达式\textbf{完全不依赖于观测值 $x_2$}，仅由先验协方差结构决定（且恒满足 $\Sigma_{1\mid 2} \preceq \Sigma_{11}$）。
\end{enumerate}

\subsection*{习题 2.2\quad 非法协方差阵与欧氏 vs 马氏距离反直觉查错}
\begin{enumerate}[label=(\arabic*), leftmargin=1.8em]
  \item \textbf{两种方法识破非法协方差阵 $\Sigma = \begin{psmallmatrix} 1.0 & 1.2 \\ 1.2 & 1.0 \end{psmallmatrix}$}：
  \begin{itemize}
    \item \textbf{特征值法}：特征多项式 $(1-\lambda)^2 - 1.2^2 = 0$，解得 $\lambda_1 = 2.2 > 0$、$\lambda_2 = -0.2 < 0$（同时 $\det\Sigma = -0.44 < 0$），故 $\Sigma$ 不是半正定矩阵（例如取 $v = (1,-1)^T$ 得 $v^T\Sigma v = -0.4 < 0$），违背协方差矩阵必满足 $\Sigma \succeq 0$ 的充要条件。
    \item \textbf{相关系数法}：由定义 $\rho_{12} = \frac{\Sigma_{12}}{\sqrt{\Sigma_{11}\Sigma_{22}}} = \frac{1.2}{\sqrt{1.0\times 1.0}} = 1.2 > 1$，违背 Cauchy--Schwarz 不等式 $|\rho_{12}| \le 1$。
  \end{itemize}
  \item \textbf{欧氏距离 vs 马氏距离对比}：对 $\mathcal{N}(0, \diag(16, 1))$，各轴标准差分别为 $\sigma_1 = 4, \sigma_2 = 1$，$\det\Sigma = 16$：
  \begin{itemize}
    \item 点 $a = (4, 0)^T$：欧氏距离 $\|a\|_2 = 4$；马氏距离 $d_M(a, 0) = \sqrt{4^2/16 + 0^2/1} = 1$；概率密度 $p(a) = \frac{1}{2\pi\cdot 4}e^{-1/2} = \frac{e^{-0.5}}{8\pi} \approx 0.02413$。
    \item 点 $b = (0, 2)^T$：欧氏距离 $\|b\|_2 = 2$；马氏距离 $d_M(b, 0) = \sqrt{0^2/16 + 2^2/1} = 2$；概率密度 $p(b) = \frac{1}{2\pi\cdot 4}e^{-2} = \frac{e^{-2}}{8\pi} \approx 0.005385$。
  \end{itemize}
  \textbf{原理解释}：虽然 $\|b\|_2 = 2 < \|a\|_2 = 4$，但第 1 维自然波动幅度很大（$\sigma_1=4$），$a$ 仅偏离 $1$ 个标准差；第 2 维非常紧凑（$\sigma_2=1$），$b$ 已偏离 $2$ 个标准差，导致密度比 $p(a)/p(b) = e^{1.5} \approx 4.482$。马氏距离按各主轴方差自动归一化，因此才是高斯密度等高线的真实几何尺子。
\end{enumerate}

\subsection*{习题 2.3\quad 二项评测概率与 Beta 先验手算}
\begin{enumerate}[label=(\arabic*), leftmargin=1.8em]
  \item $S_5 \sim \mathrm{Bin}(5, 0.8)$：
  \[
    \E[S_5] = np = 5 \times 0.8 = 4.0, \qquad
    \Var(S_5) = np(1-p) = 5 \times 0.8 \times 0.2 = 0.8.
  \]
  至少答对 4 题的精确概率为：
  \begin{align*}
    \Prob(S_5 \ge 4) &= \Prob(S_5 = 4) + \Prob(S_5 = 5)
    = \binom{5}{4}(0.8)^4(0.2)^1 + \binom{5}{5}(0.8)^5 \\
    &= 5 \times 0.4096 \times 0.2 + 0.32768 = 0.4096 + 0.32768 = 0.73728 = \frac{2304}{3125}.
  \end{align*}
  \item $X \sim \mathrm{Beta}(3, 2)$（$\alpha = 3, \beta = 2$）：
  \begin{itemize}
    \item \textbf{归一化常数}：$B(3, 2) = \frac{\Gamma(3)\Gamma(2)}{\Gamma(5)} = \frac{2!\cdot 1!}{4!} = \frac{2}{24} = \frac{1}{12}$。
    \item \textbf{密度函数}：$p(x) = \frac{1}{B(3,2)} x^{3-1}(1-x)^{2-1} = 12x^2(1 - x),\; x \in [0, 1]$（其余处为 $0$）。
    \item \textbf{均值与方差}：$\E[X] = \frac{\alpha}{\alpha+\beta} = \frac{3}{5} = 0.6$；$\Var(X) = \frac{\alpha\beta}{(\alpha+\beta)^2(\alpha+\beta+1)} = \frac{6}{25\times 6} = \frac{1}{25} = 0.04$。
    \item \textbf{众数}：令 $p'(x) = 12(2x - 3x^2) = 12x(2 - 3x) = 0$，在 $(0,1)$ 内唯一极大值点为 $x^* = \frac{\alpha-1}{\alpha+\beta-2} = \frac{2}{3} \approx 0.6667$。
    \item \textbf{先验伪计数解释}：与均匀先验 $\mathrm{Beta}(1, 1)$ 相比，$\mathrm{Beta}(3, 2)$ 相当于在看到真实数据前，额外注入了 $\alpha - 1 = 2$ 次正面伪计数与 $\beta - 1 = 1$ 次反面伪计数（共 $3$ 次伪观测，也常以总等效样本量 $\alpha+\beta=5$ 刻画先验强度）。
  \end{itemize}
\end{enumerate}

\subsection*{习题 2.4\quad 2D 多元高斯椭球主轴与条件分布精确手算}
已知 $\mu = \begin{pmatrix} 0 \\ 1 \end{pmatrix}$，$\Sigma = \begin{pmatrix} 5 & 2 \\ 2 & 2 \end{pmatrix}$。
\begin{enumerate}[label=(\arabic*), leftmargin=1.8em]
  \item \textbf{特征分解与等密度椭圆半轴}：特征方程为
  \[
    \det(\Sigma - \lambda I) = (5 - \lambda)(2 - \lambda) - 4 = \lambda^2 - 7\lambda + 6 = (\lambda - 6)(\lambda - 1) = 0,
  \]
  解得特征值 $\lambda_1 = 6,\; \lambda_2 = 1$。
  \begin{itemize}
    \item 对 $\lambda_1 = 6$：$(\Sigma - 6I)u_1 = \begin{psmallmatrix} -1 & 2 \\ 2 & -4 \end{psmallmatrix}u_1 = 0$，得单位特征向量 $u_1 = \frac{1}{\sqrt{5}}\begin{pmatrix} 2 \\ 1 \end{pmatrix}$；
    \item 对 $\lambda_2 = 1$：$(\Sigma - I)u_2 = \begin{psmallmatrix} 4 & 2 \\ 2 & 1 \end{psmallmatrix}u_2 = 0$，得单位特征向量 $u_2 = \frac{1}{\sqrt{5}}\begin{pmatrix} -1 \\ 2 \end{pmatrix}$（满足 $u_1^T u_2 = 0$）。
  \end{itemize}
  等密度椭圆 $d_M^2(x, \mu) = 1$ 的长半轴长为 $\sqrt{\lambda_1} = \sqrt{6} \approx 2.4495$（沿 $u_1$ 方向），短半轴长为 $\sqrt{\lambda_2} = 1$（沿 $u_2$ 方向）。
  \item \textbf{平方马氏距离}：$\det\Sigma = 5\times 2 - 2^2 = 6$，故 $\Sigma^{-1} = \frac{1}{6}\begin{pmatrix} 2 & -2 \\ -2 & 5 \end{pmatrix}$。对 $x = (2, 3)^T$，偏差向量 $\delta = x - \mu = (2, 2)^T$，于是：
  \[
    d_M^2(x, \mu) = \delta^T \Sigma^{-1} \delta
    = \frac{1}{6}\begin{pmatrix} 2 & 2 \end{pmatrix}
      \begin{pmatrix} 2 & -2 \\ -2 & 5 \end{pmatrix}
      \begin{pmatrix} 2 \\ 2 \end{pmatrix}
    = \frac{1}{6}\begin{pmatrix} 2 & 2 \end{pmatrix}\begin{pmatrix} 0 \\ 6 \end{pmatrix}
    = \frac{12}{6} = 2.
  \]
  （亦可在特征基下核验：$u_1^T\delta = \frac{6}{\sqrt{5}},\; u_2^T\delta = \frac{2}{\sqrt{5}}$，故 $d_M^2 = \frac{36/5}{6} + \frac{4/5}{1} = \frac{6}{5} + \frac{4}{5} = 2$。）
  \item \textbf{条件分布 $X_1 \mid (X_2 = 3)$}：由定理 \ref{thm:ch2-cond-mvg}，
  \[
    \mu_{1\mid 2} = \mu_1 + \frac{\Sigma_{12}}{\Sigma_{22}}(3 - \mu_2) = 0 + \frac{2}{2}(3 - 1) = 2, \qquad
    \sigma_{1\mid 2}^2 = \Sigma_{11} - \frac{\Sigma_{12}^2}{\Sigma_{22}} = 5 - \frac{2^2}{2} = 3.
  \]
  故 $X_1 \mid (X_2 = 3) \sim \mathcal{N}(2, 3)$。观测 $X_2$ 将 $X_1$ 的方差从 $5$ 降至 $3$，方差缩减比例为 $\frac{5 - 3}{5} = 40\%$（恰好等于相关系数平方 $\rho_{12}^2 = \frac{4}{10} = 0.4$）。
\end{enumerate}

\subsection*{习题 2.5\quad 独立高斯之和与样本均值的 $1/\sqrt{n}$ 标准误}
\begin{proof}
\begin{enumerate}[label=(\arabic*), leftmargin=1.8em]
  \item 因 $X_1 \sim \mathcal{N}(\mu_1, \Sigma_1)$ 与 $X_2 \sim \mathcal{N}(\mu_2, \Sigma_2)$ 相互独立，其联合密度等于边际密度之积，直接合并指数项即可见拼接向量 $U := \begin{pmatrix} X_1 \\ X_2 \end{pmatrix} \in \R^{2d}$ 服从联合高斯分布：
  \[
    U = \begin{pmatrix} X_1 \\ X_2 \end{pmatrix}
    \sim \mathcal{N}\!\left(\begin{pmatrix} \mu_1 \\ \mu_2 \end{pmatrix},\;
         \begin{pmatrix} \Sigma_1 & 0 \\ 0 & \Sigma_2 \end{pmatrix}\right).
  \]
  选取线性变换矩阵 $A = \begin{pmatrix} I_d & I_d \end{pmatrix} \in \R^{d \times 2d}$，则 $X_1 + X_2 = A U$。由多元高斯仿射封闭性定理（定理 \ref{thm:ch2-affine}），立得：
  \[
    X_1 + X_2 \sim \mathcal{N}\!\left(A\begin{pmatrix} \mu_1 \\ \mu_2 \end{pmatrix},\;
      A\begin{pmatrix} \Sigma_1 & 0 \\ 0 & \Sigma_2 \end{pmatrix}A^T\right)
    = \mathcal{N}(\mu_1 + \mu_2,\; \Sigma_1 + \Sigma_2).
  \]
  \item 对 $n$ 个 i.i.d.\ 样本 $X_1, \dots, X_n \sim \mathcal{N}(\mu, \sigma^2)$，由 (1) 数学归纳法可知 $\sum_{i=1}^n X_i \sim \mathcal{N}(n\mu, n\sigma^2)$；再施加标量线性变换 $\frac{1}{n}$，得：
  \[
    \bar{X}_n = \frac{1}{n}\sum_{i=1}^n X_i \sim \mathcal{N}\!\left(\mu,\; \frac{\sigma^2}{n}\right).
  \]
  样本均值的标准误为 $\mathrm{SE}(\bar{X}_n) = \frac{\sigma}{\sqrt{n}}$。要令标准误减半（$\frac{\sigma}{\sqrt{n'}} = \frac{1}{2}\frac{\sigma}{\sqrt{n}}$），必须满足 $\sqrt{n'} = 2\sqrt{n} \iff n' = 4n$，这就是统计学与蒙特卡洛中“精度翻倍需 $4$ 倍样本量”的根源。
\end{enumerate}
\end{proof}

\subsection*{习题 2.6\quad 白化等价性与高维高斯薄球壳现象}
\begin{proof}
\begin{enumerate}[label=(\arabic*), leftmargin=1.8em]
  \item 由实对称正定矩阵谱分解 $\Sigma = Q\Lambda Q^T$，取唯一对称正定平方根逆 $\Sigma^{-1/2} = Q\Lambda^{-1/2}Q^T$，满足 $(\Sigma^{-1/2})^T = \Sigma^{-1/2}$ 且 $\Sigma^{-1/2}\Sigma\Sigma^{-1/2} = I_d$。对 $Z = \Sigma^{-1/2}(X - \mu)$ 应用仿射封闭性定理（定理 \ref{thm:ch2-affine}）：
  \[
    \E[Z] = \Sigma^{-1/2}(\mu - \mu) = 0, \qquad
    \Cov(Z) = \Sigma^{-1/2}\Sigma (\Sigma^{-1/2})^T = I_d,
  \]
  故 $Z \sim \mathcal{N}(0, I_d)$，各分量 $Z_1, \dots, Z_d \stackrel{\text{i.i.d.}}{\sim} \mathcal{N}(0, 1)$。进一步展开平方马氏距离：
  \[
    d_M^2(X, \mu) = (X - \mu)^T \Sigma^{-1} (X - \mu)
    = (X - \mu)^T \Sigma^{-1/2}\Sigma^{-1/2} (X - \mu)
    = Z^T Z = \|Z\|_2^2 = \sum_{i=1}^d Z_i^2.
  \]
  \item 对每个标准正态分量 $Z_i \sim \mathcal{N}(0, 1)$，有 $\E[Z_i^2] = \Var(Z_i) = 1$ 且 $\E[Z_i^4] = 3$，从而
  \[
    \Var(Z_i^2) = \E[Z_i^4] - (\E[Z_i^2])^2 = 3 - 1^2 = 2.
  \]
  因为 $Z_1, \dots, Z_d$ 相互独立，由期望线性性与独立方差可加性得：
  \[
    \E\bigl[d_M^2(X, \mu)\bigr] = \sum_{i=1}^d \E[Z_i^2] = d, \qquad
    \Var\bigl(d_M^2(X, \mu)\bigr) = \sum_{i=1}^d \Var(Z_i^2) = 2d.
  \]
  \item \textbf{高维薄球壳直觉}：当 $d = 1{,}000$ 时，平方马氏半径的均值为 $1000$，标准差仅为 $\sqrt{2d} \approx 44.72$（相对波动仅 $4.5\%$），对应马氏半径 $d_M(X, \mu) \approx \sqrt{d} \approx 31.62$，且相对厚度随 $O(1/\sqrt{d})$ 趋近于零！虽然密度函数在原点 $\mu$ 处最高，但高维球壳体积以 $r^{d-1}$ 爆炸式增长，密度衰减与体积膨胀的乘积在 $r \approx \sqrt{d}$ 处形成极窄尖峰——这就是高维高斯“中心空心、质量全在薄球壳”的几何真相（也是第 5 卷 VAE 隐空间先验匹配的核心直觉）。
\end{enumerate}
\end{proof}

\subsection*{习题 2.7\quad 3D 高斯 Cholesky 采样、白化还原与多元条件切片验证}
\textbf{核心纯 NumPy 实现逻辑}：
\begin{lstlisting}[language=Python]
rng = np.random.default_rng(2026)
mu = np.array([1.0, -1.0, 2.0])
Sigma = np.array([[4.0, 1.2, -0.8], [1.2, 2.0, 0.6], [-0.8, 0.6, 1.5]])
L = np.linalg.cholesky(Sigma)
X = mu + rng.normal(size=(200_000, 3)) @ L.T
W = np.linalg.solve(L, (X - mu).T).T
dm2 = np.sum(W ** 2, axis=1)
# Conditional X1 | (X2=0.0, X3=2.5) via Schur complement
S12, S22 = Sigma[0:1, 1:3], Sigma[1:3, 1:3]
mu_cond = float(mu[0] + (S12 @ np.linalg.solve(S22, np.array([1.0, 0.5]))).item())
var_cond = float(Sigma[0, 0] - (S12 @ np.linalg.solve(S22, S12.T)).item())
mask = (np.abs(X[:, 1] - 0.0) < 0.15) & (np.abs(X[:, 2] - 2.5) < 0.15)
\end{lstlisting}
\textbf{验证输出与断言自检（\texttt{default\_rng(2026)}）}：
\begin{itemize}
  \item 经验均值最大绝对误差 $0.00358$，经验协方差最大绝对误差 $0.01032 < 0.03$；
  \item 白化样本 $W$ 经验协方差与 $I_3$ 最大绝对误差 $0.00512 < 0.03$；平方马氏距离样本均值 $2.9976$（理论 $d=3$，落入 $[2.95, 3.05]$）、方差 $5.9849$（理论 $2d=6$，落入 $[5.85, 6.15]$）；
  \item 理论条件参数精确值为 $\mu_{1\mid 2,3} = \frac{47}{33} \approx 1.4242$、$\sigma_{1\mid 2,3}^2 = \frac{373}{165} \approx 2.2606$；切片内共落入 $1{,}289$ 条样本，经验均值 $1.3256$（偏差 $0.0986 < 0.15$）、经验方差 $2.1718$（偏差 $0.0888 < 0.15$），全部断言通过。
\end{itemize}

\subsection*{习题 2.8\quad 欧氏距离 vs 马氏距离的隐蔽异常检测实战}
\textbf{核心纯 NumPy 实现逻辑}（严格禁止 \texttt{np.linalg.inv}，仅用 Cholesky + 三角求解）：
\begin{lstlisting}[language=Python]
rng = np.random.default_rng(2026)
Sigma = np.array([[4.0, 1.8], [1.8, 1.0]])
L = np.linalg.cholesky(Sigma)
X_norm = rng.normal(size=(2000, 2)) @ L.T
eigvals, eigvecs = np.linalg.eigh(Sigma)
u_min = eigvecs[:, 0]
signs = np.where(np.arange(20) % 2 == 0, 1.0, -1.0)
X_anom = signs[:, None] * 1.8 * u_min[None, :] + rng.normal(0, 0.05, size=(20, 2))
X_all = np.vstack([X_norm, X_anom])
d_euc = np.linalg.norm(X_all, axis=1)
Y = np.linalg.solve(L, X_all.T)
d_mah = np.sqrt(np.sum(Y ** 2, axis=0))
\end{lstlisting}
\textbf{验证输出与断言自检（\texttt{default\_rng(2026)}）}：
\begin{itemize}
  \item $\Sigma$ 特征值为 $\lambda_{\min} \approx 0.15693$、$\lambda_{\max} \approx 4.84307$（正常长轴标准差 $\sqrt{\lambda_{\max}} \approx 2.2007$）；
  \item 沿短轴 $1.8\,u_{\min}$ 注入的隐蔽异常点欧氏距离仅约 $1.8$，淹没在正常云内部，故欧氏距离 Recall@25 为 \textbf{0/20}（$0\%$）；
  \item 其理论马氏距离高达 $1.8 / \sqrt{\lambda_{\min}} \approx 4.5439$，远超正常点 $99.9\%$ 轮廓，故马氏距离 Recall@25 精准命中 \textbf{20/20}（$100\%$）。
\end{itemize}

% =========================================================================
\section{第 3 章习题精解（3.1--3.8）}
% =========================================================================

\subsection*{习题 3.1\quad 极限、浓度不等式与采样核心概念辨析}
\begin{enumerate}[label=(\alph*), leftmargin=1.8em]
  \item \textbf{错误。} Markov 不等式严格要求随机变量\textbf{几乎必然非负（$X \ge 0$）}，否则负值会把 $\E[X]$ 拉低甚至导致右边远小于真实尾概率。\textbf{最小反例}：设 $\Prob(X = 1) = 0.8$、$\Prob(X = -3) = 0.2$，则 $\E[X] = 0.8 - 0.6 = 0.2 > 0$；取 $a = 1$ 时，左侧 $\Prob(X \ge 1) = 0.8$，但右侧 $\E[X]/a = 0.2 < 0.8$，不等式反号！
  \item \textbf{正确。} 函数 $\varphi_1(x) = \ln x$ 在 $(0, +\infty)$ 上满足 $\varphi_1''(x) = -1/x^2 < 0$（严格凹），$\varphi_2(x) = 1/x$ 满足 $\varphi_2''(x) = 2/x^3 > 0$（严格凸）。因 $X > 0$ 且不几乎必然为常数，由严格 Jensen 不等式立得 $\E[\ln X] < \ln\E[X]$ 且 $\E[1/X] > 1/\E[X]$。
  \item \textbf{错误。} 经典中心极限定理（CLT）要求单样本具有\textbf{有限方差（$0 < \sigma^2 < \infty$）}。\textbf{经典反例}：标准 Cauchy 分布 $X_i \stackrel{\text{i.i.d.}}{\sim} \mathrm{Cauchy}(0, 1)$（密度 $p(x) = \frac{1}{\pi(1+x^2)}$），其均值与方差均不存在；无论 $n$ 多大，样本均值 $\bar{X}_n$ 的精确分布始终仍是 $\mathrm{Cauchy}(0, 1)$，永不趋近正态。
  \item \textbf{正确。} 重要性采样依赖恒等式 $\int p(x)f(x)\,dx = \int q(x)\frac{p(x)f(x)}{q(x)}\,dx$，仅在支集覆盖条件 $\{x: p(x)f(x)\neq 0\} \subseteq \{x: q(x)>0\}$ 下成立。若在某个正概率区域 $S$ 上 $q(x) = 0$ 而 $p(x)f(x) \neq 0$，采样点永远落不进 $S$，估计量收敛到截断积分 $\int_{S^c} p(x)f(x)\,dx$，存在无法通过增大 $M$ 消除的系统偏差。
\end{enumerate}

\subsection*{习题 3.2\quad 重要性采样“无偏但方差无穷”反例}
\begin{enumerate}[label=(\arabic*), leftmargin=1.8em]
  \item 目标密度 $p(x) = \frac{1}{2\sqrt{2\pi}}\exp(-x^2/8)$，提议密度 $q(x) = \frac{1}{\sqrt{2\pi}}\exp(-x^2/2)$，故权重为：
  \[
    w(x) = \frac{p(x)}{q(x)} = \frac{1}{2}\exp\!\left(\frac{3}{8}x^2\right).
  \]
  直接验证期望：
  \[
    \E_{X\sim q}[w(X)] = \int_{-\infty}^{+\infty} q(x)\frac{p(x)}{q(x)}\,dx = \int_{-\infty}^{+\infty} p(x)\,dx = 1,
  \]
  因此对常数函数 $f(x) = 1$ 严格无偏。
  \item 计算权重二阶矩：
  \[
    \E_{X\sim q}\bigl[w(X)^2\bigr]
    = \int_{-\infty}^{+\infty} \frac{p(x)^2}{q(x)}\,dx
    = \int_{-\infty}^{+\infty} \frac{\frac{1}{8\pi}e^{-x^2/4}}{\frac{1}{\sqrt{2\pi}}e^{-x^2/2}}\,dx
    = \frac{1}{4\sqrt{2\pi}} \int_{-\infty}^{+\infty} \exp\!\left(+\frac{x^2}{4}\right)dx = +\infty.
  \]
  因为被积函数在 $|x| \to \infty$ 时呈超指数发散，故 $\Var_q(w(X)) = \E_q[w(X)^2] - 1 = +\infty$。
  \item \textbf{尾部设计原则}：为保证 $\E_q[w^2] = \int p(x)\frac{p(x)}{q(x)}\,dx < \infty$，在尾部区域 $|x|\to\infty$ 必须避免权重 $p(x)/q(x)$ 失控增长；工程上应保证提议分布 $q$ 的尾部\textbf{至少与目标 $p$ 一样厚}（例如对正态目标选更大方差的高斯或重尾 Student-$t$ 分布），使 $w(x)$ 有界或缓增。
\end{enumerate}

\subsection*{习题 3.3\quad 三阶梯浓度不等式正面交锋}
已知 $n = 100$，$\mu = \E[\bar{X}_{100}] = 0.5$，$\Var(\bar{X}_{100}) = 0.0025$，且 $X_i \in [0, 1]$。
\begin{enumerate}[label=(\arabic*), leftmargin=1.8em]
  \item \textbf{Markov 不等式（单侧）}：因 $\bar{X}_{100} \ge 0$，
  \[
    \Prob(\bar{X}_{100} \ge 0.65) \le \frac{\E[\bar{X}_{100}]}{0.65} = \frac{0.5}{0.65} = \frac{10}{13} \approx 0.7692.
  \]
  \item \textbf{Chebyshev 不等式（双侧）}：取 $\varepsilon = 0.15$，
  \[
    \Prob(|\bar{X}_{100} - 0.5| \ge 0.15) \le \frac{\Var(\bar{X}_{100})}{0.15^2} = \frac{0.0025}{0.0225} = \frac{1}{9} \approx 0.1111.
  \]
  \item \textbf{Hoeffding 不等式}：因每个变量范围 $b_i - a_i = 1$，指数项 $2n\varepsilon^2 = 2 \times 100 \times 0.15^2 = 4.5$：
  \begin{align*}
    \text{单侧：}\quad &\Prob(\bar{X}_{100} \ge 0.65) \le e^{-4.5} \approx 0.0111, \\
    \text{双侧：}\quad &\Prob(|\bar{X}_{100} - 0.5| \ge 0.15) \le 2e^{-4.5} \approx 0.0222.
  \end{align*}
\end{enumerate}
\textbf{松紧对比}：单侧上界从 Markov 的 $0.7692$ 骤降至 Hoeffding 的 $0.0111$（精确二项尾概率约为 $0.00176$）；双侧上界从 Chebyshev 的 $0.1111$ 压低 5 倍至 Hoeffding 的 $0.0222$，体现“利用的信息越多（仅均值 $\to$ 均值+方差 $\to$ 独立+有界），浓度界越紧”的三阶梯跃迁。

\subsection*{习题 3.4\quad AI 评测测试集规模设计：Hoeffding vs CLT}
给定允许误差 $\varepsilon = 0.02$（$\varepsilon^2 = 0.0004$）。
\begin{enumerate}[label=(\arabic*), leftmargin=1.8em]
  \item \textbf{Hoeffding 最小样本量}：令 $2\exp(-2n\varepsilon^2) \le \delta$，解得 $n \ge \frac{1}{2\varepsilon^2}\ln\frac{2}{\delta} = 1250\ln\frac{2}{\delta}$：
  \begin{itemize}
    \item $95\%$ 置信度（$\delta = 0.05$）：$n \ge 1250 \ln 40 \approx 1250 \times 3.68888 = 4611.10 \implies n_{\min} = \mathbf{4{,}612}$；
    \item $99\%$ 置信度（$\delta = 0.01$）：$n \ge 1250 \ln 200 \approx 1250 \times 5.29832 = 6622.90 \implies n_{\min} = \mathbf{6{,}623}$。
  \end{itemize}
  \item \textbf{CLT 渐近 $95\%$ 样本量（$p \approx 0.90$）}：由 $1.96\sqrt{\frac{p(1-p)}{n}} \le 0.02$ 解得：
  \[
    n_{\mathrm{CLT}} \ge \frac{1.96^2 \times 0.90 \times 0.10}{0.02^2}
    = \frac{3.8416 \times 0.09}{0.0004} = 864.36 \implies n_{\mathrm{CLT}} = \mathbf{865}.
  \]
  \item \textbf{相差逾 5 倍的两大原因}：(i) \textbf{方差信息}：Hoeffding 在最坏方差 $p(1-p) \le 0.25$（$p=0.5$）下做无分布保守估计，而 $p=0.90$ 的真实方差 $0.09$ 不到最坏情形的一半；(ii) \textbf{保证强度}：Hoeffding 提供对\textbf{任意有限 $n$} 严格成立的非渐近硬保证（包含额外的松弛系数），而 CLT 依赖 $n\to\infty$ 的渐近正态近似，紧得多但在极小样本或极端概率（如 $p=0.999$）下会失真。
\end{enumerate}

\subsection*{习题 3.5\quad Softplus 凸性与加权算术-几何平均不等式}
\begin{proof}
\begin{enumerate}[label=(\arabic*), leftmargin=1.8em]
  \item 对 $f(x) = \ln(1 + e^x)$ 求一阶与二阶导数：
  \[
    f'(x) = \frac{e^x}{1 + e^x} = \frac{1}{1 + e^{-x}} = \sigma(x), \qquad
    f''(x) = \sigma'(x) = \sigma(x)\bigl(1 - \sigma(x)\bigr).
  \]
  对任意 $x \in \R$ 均有 $\sigma(x) \in (0, 1)$，故 $f''(x) > 0$ 处处成立，$f$ 在 $\R$ 上\textbf{严格凸}。直接应用 Jensen 不等式（定理 \ref{thm:ch3-jensen}），即得对任意有限期望随机变量 $X$ 恒有 $\E[\ln(1 + e^X)] \ge \ln(1 + e^{\E[X]})$。
  \item 构造取值于 $\{a_1, \dots, a_k\}$ 的离散随机变量 $X$，使其概率质量函数为 $\Prob(X = a_i) = w_i$（合法因 $w_i \ge 0$ 且 $\sum_{i=1}^k w_i = 1$）。对严格凹函数 $\varphi(t) = \ln t$（$t > 0$）应用 Jensen 不等式 $\E[\ln X] \le \ln(\E[X])$：
  \[
    \sum_{i=1}^k w_i \ln a_i \le \ln\!\left(\sum_{i=1}^k w_i a_i\right)
    \iff \ln\!\left(\prod_{i=1}^k a_i^{w_i}\right) \le \ln\!\left(\sum_{i=1}^k w_i a_i\right).
  \]
  两边施加严格递增的指数函数 $\exp(\cdot)$，立得加权 AM--GM 不等式 $\prod_{i=1}^k a_i^{w_i} \le \sum_{i=1}^k w_i a_i$。因 $\ln t$ 严格凹，\textbf{等号成立的充要条件}是 $X$ 几乎必然为常数，即所有具有正权重（$w_i > 0$）的 $a_i$ 全都相等。
\end{enumerate}
\end{proof}

\subsection*{习题 3.6\quad 重要性采样的理论最优提议分布与“零方差悖论”}
\begin{proof}
\begin{enumerate}[label=(\arabic*), leftmargin=1.8em]
  \item 当 $q$ 覆盖 $p$ 支集时，$\E_q[Y] = \int q(x) f(x)\frac{p(x)}{q(x)}\,dx = \int f(x)p(x)\,dx = I$。由方差展开公式：
  \[
    \Var_q(Y) = \E_q[Y^2] - (\E_q[Y])^2 = \E_q\!\left[\frac{f(X)^2 p(X)^2}{q(X)^2}\right] - I^2.
  \]
  \item 对凸函数 $\varphi(y) = y^2$ 应用 Jensen 不等式，立得 $\E_q[Y^2] \ge (\E_q[Y])^2 = I^2$，即 $\Var_q(Y) \ge 0$。当取 $q^*(x) = \frac{f(x)p(x)}{I}$ 时（显然 $q^*(x) > 0$ 且 $\int q^*(x)\,dx = 1$，为合法密度），对任意采样点 $X \sim q^*$ 均有
  \[
    Y = f(X)\frac{p(X)}{q^*(X)} = f(X)\frac{p(X)}{f(X)p(X)/I} \equiv I\quad (\text{常数}),
  \]
  故 $\Var_{q^*}(Y) = \Var(I) = 0$。反之，若 $\Var_q(Y) = 0$，则 $f(X)p(X)/q(X)$ 几乎处处等于常数 $I$，必有 $q(x) = f(x)p(x)/I$ 几乎处处成立。
  \item \textbf{“零方差悖论”与工程指导原则}：构造 $q^*(x)$ 恰好需要知道归一化分母 $I = \int f(t)p(t)\,dt$——而 $I$ 正是我们费尽心思想要估计的未知积分！尽管无法直接采样 $q^*$，它却指明了黄金设计原则：\textbf{好的提议分布 $q(x)$ 应当尽量与 $|f(x)|p(x)$ 的形状成正比}，把采样预算集中在被积函数与目标密度乘积最大的区域（而非仅盯着 $p(x)$ 最大的区域），这正是习题 3.8 中把正态中心平移到 $4$ 的原因。
\end{enumerate}
\end{proof}

\subsection*{习题 3.7\quad 浓度不等式实测对比与右偏分布 CLT 演化实验}
\textbf{核心纯 NumPy 实现逻辑}：
\begin{lstlisting}[language=Python]
rng = np.random.default_rng(2026)
for n in [20, 60, 100, 200, 300]:
    means = rng.binomial(n, 0.5, size=20_000) / n
    emp = np.mean(np.abs(means - 0.5) >= 0.1)
    cheb = min(1.0, 0.25 / (n * 0.01))
    hoeff = 2.0 * np.exp(-2.0 * n * 0.01)
for n in [1, 5, 30, 100]:
    z = np.sqrt(n) * (rng.exponential(1.0, size=(20_000, n)).mean(axis=1) - 1.0)
    skew, cov95 = np.mean(z ** 3), np.mean(np.abs(z) <= 1.96)
\end{lstlisting}
\textbf{验证输出与断言自检（\texttt{default\_rng(2026)}）}：
\begin{itemize}
  \item \textbf{浓度上界对比}：在 $n = 20, 60, 100, 200, 300$ 时，Chebyshev 上界依次为 $1.00000$、$0.41667$、$0.25000$、$0.12500$、$\mathbf{0.08333}$；Hoeffding 上界依次为 $1.34064$、$0.60239$、$0.27067$、$0.03663$、$\mathbf{0.00496}$；实测尾概率依次为 $0.26135$、$0.09330$、$0.03435$、$0.00360$、$\mathbf{0.00040}$（$n=300$ 落入 $[0.0001, 0.0010]$，且 $n\ge 200$ 后 Hoeffding 指数级胜出）。
  \item \textbf{右偏 $\mathrm{Exp}(1)$ 的 CLT 演化}：随 $n = 1, 5, 30, 100$ 增大，样本偏度依次为 $1.8253$、$0.9472$、$0.3186$、$0.2652$（沿理论 $2/\sqrt{n}$ 快速趋零），$[-1.96, 1.96]$ 覆盖率依次为 $0.9480$、$0.9554$、$0.9536$、$0.9507$，稳定收敛于 $0.950 \pm 0.005$。
\end{itemize}

\subsection*{习题 3.8\quad 罕见尾事件重要性采样与权重塌缩 ESS 诊断}
\textbf{核心纯 NumPy 实现逻辑}：
\begin{lstlisting}[language=Python]
# (1) P(Z > 4) with q = N(4, 1): w(z) = p(z)/q(z) = exp(-4z + 8)
z_is = rng.normal(4.0, 1.0, size=20_000)
p_is = np.mean((z_is > 4.0) * np.exp(-4.0 * z_is + 8.0))
# (2) E_{N(0, 4)}[X^2] with q = N(0, q_std^2) and ESS/M
x = rng.normal(0.0, q_std, size=10_000)
log_w = (-0.5*(x/2.0)**2 - np.log(2.0)) - (-0.5*(x/q_std)**2 - np.log(q_std))
w = np.exp(log_w)
est, ess_ratio = np.mean(w * x**2), (w.sum()**2) / (len(w) * np.sum(w**2))
\end{lstlisting}
\textbf{验证输出与断言自检（\texttt{default\_rng(2026)}，50 次独立重复）}：
\begin{itemize}
  \item \textbf{罕见尾事件}（真值 $3.1671\times 10^{-5}$）：朴素 MC 有 $46\%$ 的轮次直接输出 $0$（均值 $4.00\times 10^{-5}$，标准差高达 $4.29\times 10^{-5} > 3\times 10^{-5}$）；平移提议 $q = \mathcal{N}(4, 1)$ 的 50 轮均值为 $3.1710 \times 10^{-5}$，单轮标准差仅 $5.14\times 10^{-7} < 7\times 10^{-7}$，平均相对误差仅 \textbf{1.31\%}（方差缩减约 7000 倍！）。
  \item \textbf{ESS 权重塌缩诊断}（真值 $4.0$）：厚尾提议 $q_2 = \mathcal{N}(0, 3^2)$ 均值为 $4.0052$、标准差仅 $0.0293 < 0.05$、平均 $\mathrm{ESS}/M = \mathbf{0.8314} > 0.75$；瘦尾提议 $q_1 = \mathcal{N}(0, 1)$ 均值偏低且剧烈抖动（$3.7059 \pm 2.2287$，标准差 $> 0.5$），平均 $\mathrm{ESS}/M$ 塌缩至 \textbf{0.1190}（$< 0.25$），直观验证了习题 3.2 的无穷方差灾难。
\end{itemize}

% =========================================================================
\section{第 4 章习题精解（4.1--4.8）}
% =========================================================================

\subsection*{习题 4.1\quad 信息论核心概念辨析}
\begin{enumerate}[label=(\alph*), leftmargin=1.8em]
  \item \textbf{错误。} KL 散度不满足对称性，也不满足三角不等式，因此不是数学意义上的距离度量。\textbf{反例}：由习题 4.2(1)，对 $P=(0.9,0.1)$ 与 $Q=(0.5,0.5)$，有 $\DKL(P\|Q)\approx 0.3681 \neq \DKL(Q\|P)\approx 0.5108$；取 $R=(0.7,0.3)$，还可验证 $\DKL(P\|Q) > \DKL(P\|R) + \DKL(R\|Q)$ 亦可发生（例如取更极端的 $P=(0.99,0.01), R=(0.5,0.5), Q=(0.01,0.99)$ 时 $\DKL(P\|Q)\approx 4.50 > \DKL(P\|R)+\DKL(R\|Q)\approx 0.64 + 1.61 = 2.25$）。
  \item \textbf{正确。} 离散分布满足 $0 \le p(x) \le 1 \implies -\log p(x) \ge 0$，故 Shannon 熵恒非负；但连续密度 $p(x)$ 可大于 $1$，令 $-\ln p(x) < 0$。\textbf{实例}：均匀分布 $X \sim \mathrm{Uniform}(0, a)$ 的微分熵 $h(X) = \ln a$，当 $a = 0.5$ 时 $h(X) = -\ln 2 \approx -0.6931 < 0$（对高斯 $\mathcal{N}(0, \sigma^2)$，$h(X) = \frac{1}{2}\ln(2\pi e\sigma^2)$，当 $\sigma$ 足够小时亦为负）。
  \item \textbf{错误。} $H(Y \mid X) = \sum_{x} P(X=x) H(Y \mid X=x) \le H(Y)$ 仅在\textbf{对所有 $x$ 取平均后}成立，并不保证每一个具体取值 $X = x_0$ 处的逐点条件熵都更小。\textbf{最小反例}：设 $P(X=0,Y=0)=0.90,\; P(X=0,Y=1)=0,\; P(X=1,Y=0)=0.05,\; P(X=1,Y=1)=0.05$。此时边缘分布 $P(Y=1)=0.05$，边缘熵 $H(Y) = H_2(0.05) \approx 0.2864\;\text{bit}$；但一旦观测到罕见值 $X = 1$，条件分布为 $P(Y\mid X=1) = (0.5, 0.5)$，逐点条件熵飙升至 $H(Y \mid X=1) = 1.0\;\text{bit} > H(Y)$。
  \item \textbf{正确。} 由换底公式 $\log_2 p = \frac{\ln p}{\ln 2}$，有 $\bar{H}_{\mathrm{bit}} = \frac{\bar{H}_{\mathrm{nat}}}{\ln 2}$，于是：
  \[
    2^{\bar{H}_{\mathrm{bit}}} = \exp\!\bigl(\bar{H}_{\mathrm{bit}}\cdot \ln 2\bigr)
    = \exp\!\bigl(\bar{H}_{\mathrm{nat}}\bigr) = \mathrm{PPL}.
  \]
  困惑度本质上是预测概率几何平均的倒数 $\bigl(\prod_{t=1}^T q(w_t\mid w_{<t})\bigr)^{-1/T}$，与所用对数底完全无关。
\end{enumerate}

\subsection*{习题 4.2\quad KL 散度非对称与支集爆炸反例}
\begin{enumerate}[label=(\arabic*), leftmargin=1.8em]
  \item 取 $P = (0.9, 0.1)$ 与 $Q = (0.5, 0.5)$，使用自然对数 $\ln$：
  \begin{align*}
    \DKL(P\|Q) &= 0.9\ln\frac{0.9}{0.5} + 0.1\ln\frac{0.1}{0.5}
                = 0.9\ln 1.8 - 0.1\ln 5 \approx 0.52901 - 0.16094 = \mathbf{0.3681}\;\text{nat}, \\
    \DKL(Q\|P) &= 0.5\ln\frac{0.5}{0.9} + 0.5\ln\frac{0.5}{0.1}
                = 0.5\ln\frac{25}{9} = \ln\frac{5}{3} \approx \mathbf{0.5108}\;\text{nat}.
  \end{align*}
  两者不等的直觉原因：$\DKL(P\|Q)$ 以 $P$ 为权重，把 $90\%$ 权重压在 $x=0$（比值仅 $1.8$），只在 $x=1$ 处赋予 $10\%$ 权重；而 $\DKL(Q\|P)$ 以均匀分布 $Q$ 为权重，在 $x=1$ 处以 $50\%$ 的高权重遭遇 $5$ 倍概率低估（$\ln 5$），因而惩罚更重。
  \item \textbf{有限 vs 无穷的反例}：取 $P = (1, 0)$ 与 $Q = (0.5, 0.5)$（约定 $0\ln 0 = 0$）。则：
  \[
    \DKL(P\|Q) = 1\cdot\ln\frac{1}{0.5} + 0 = \ln 2 \approx 0.6931 < \infty,
    \qquad
    \DKL(Q\|P) = 0.5\ln\frac{0.5}{1} + 0.5\ln\frac{0.5}{0} = +\infty.
  \]
  \textbf{物理含义}：在正向散度 $\DKL(P\|Q)$ 中，凡是 $P(x) > 0$ 的地方必须保证 $Q(x) > 0$ 否则惩罚爆炸（\textbf{零回避 / Mode-Covering}，这里 $Q$ 覆盖了 $P$ 的全部支集所以有限）；而在反向散度 $\DKL(Q\|P)$ 中，凡是 $P(x) = 0$ 的地方必须强迫 $Q(x) = 0$（\textbf{零强制 / Mode-Seeking}），本例中 $P(1)=0$ 但 $Q(1)=0.5>0$，故反向 KL 立刻发散到 $+\infty$。
\end{enumerate}

\subsection*{习题 4.3\quad $2\times 2$ 联合分布全套信息量手算}
由联合分布表立得两个边际分布：
\[
  P_X(0) = \tfrac{1}{2} + 0 = \tfrac{1}{2},\quad P_X(1) = \tfrac{1}{4} + \tfrac{1}{4} = \tfrac{1}{2}; \qquad
  P_Y(0) = \tfrac{1}{2} + \tfrac{1}{4} = \tfrac{3}{4},\quad P_Y(1) = 0 + \tfrac{1}{4} = \tfrac{1}{4}.
\]
逐项以 $\log_2$（单位 bit）手算（记 $0\log_2 0 = 0$）：
\begin{align*}
  H(X) &= -\frac{1}{2}\log_2\frac{1}{2} - \frac{1}{2}\log_2\frac{1}{2} = \mathbf{1.0000}\;\text{bit}, \\
  H(Y) &= -\frac{3}{4}\log_2\frac{3}{4} - \frac{1}{4}\log_2\frac{1}{4}
        = 2 - \frac{3}{4}\log_2 3 \approx \mathbf{0.8113}\;\text{bit}, \\
  H(X, Y) &= -\frac{1}{2}\log_2\frac{1}{2} - 0 - \frac{1}{4}\log_2\frac{1}{4} - \frac{1}{4}\log_2\frac{1}{4}
           = \frac{1}{2} + \frac{1}{2} + \frac{1}{2} = \frac{3}{2} = \mathbf{1.5000}\;\text{bit}.
\end{align*}
直接由条件分布算条件熵：给定 $X=0$ 时 $Y$ 必然为 $0$（$H(Y\mid X=0)=0$）；给定 $X=1$ 时 $P(Y\mid X=1)=(\frac{1}{2}, \frac{1}{2})$（$H(Y\mid X=1)=1$），故：
\[
  H(Y \mid X) = \frac{1}{2}\cdot 0 + \frac{1}{2}\cdot 1 = \frac{1}{2} = \mathbf{0.5000}\;\text{bit}.
\]
给定 $Y=0$ 时 $P(X\mid Y=0)=(\frac{2}{3}, \frac{1}{3})$；给定 $Y=1$ 时 $X$ 必然为 $1$（$H(X\mid Y=1)=0$），故：
\[
  H(X \mid Y) = \frac{3}{4}\!\left(-\frac{2}{3}\log_2\frac{2}{3} - \frac{1}{3}\log_2\frac{1}{3}\right) + \frac{1}{4}\cdot 0
              = \frac{3}{4}\log_2 3 - \frac{1}{2} \approx \mathbf{0.6887}\;\text{bit}.
\]
互信息：
\[
  I(X; Y) = H(Y) - H(Y \mid X) = \left(2 - \frac{3}{4}\log_2 3\right) - \frac{1}{2}
          = \frac{3}{2} - \frac{3}{4}\log_2 3 \approx \mathbf{0.3113}\;\text{bit}.
\]
\textbf{核验链式法则与对称性}：
\begin{itemize}
  \item $H(X) + H(Y\mid X) = 1 + 0.5 = 1.5 = H(X, Y)$，且 $H(Y) + H(X\mid Y) = (2 - \frac{3}{4}\log_2 3) + (\frac{3}{4}\log_2 3 - \frac{1}{2}) = 1.5 = H(X, Y)$；
  \item $H(X) - H(X\mid Y) = 1 - (\frac{3}{4}\log_2 3 - \frac{1}{2}) = \frac{3}{2} - \frac{3}{4}\log_2 3 = H(Y) - H(Y\mid X) = I(X; Y)$。
\end{itemize}

\subsection*{习题 4.4\quad 语言模型困惑度与一维高斯 KL 手算}
\begin{enumerate}[label=(\arabic*), leftmargin=1.8em]
  \item 4 个词的条件概率分别为 $2^{-1},\; 2^{-2},\; 2^{-3},\; 2^{-1}$，对应自信息（bit）依次为 $1, 2, 3, 1$：
  \begin{align*}
    \bar{H}_{\mathrm{bit}} &= \frac{1 + 2 + 3 + 1}{4} = \frac{7}{4} = \mathbf{1.7500}\;\text{bit}, \\
    \bar{H}_{\mathrm{nat}} &= \frac{7}{4}\ln 2 \approx 1.75 \times 0.69315 = \mathbf{1.2130}\;\text{nat}, \\
    \mathrm{PPL} &= 2^{7/4} = e^{\bar{H}_{\mathrm{nat}}} \approx \mathbf{3.3636}.
  \end{align*}
  \item 代入一维高斯 KL 闭式公式（命题 \ref{prop:ch4-gaussian-kl}）$\DKL(\mathcal{N}(\mu_1,\sigma_1^2)\|\mathcal{N}(\mu_2,\sigma_2^2)) = \ln\frac{\sigma_2}{\sigma_1} + \frac{\sigma_1^2 + (\mu_1-\mu_2)^2}{2\sigma_2^2} - \frac{1}{2}$：
  \begin{align*}
    \DKL\bigl(\mathcal{N}(1, 2^2) \,\|\, \mathcal{N}(0, 1)\bigr)
    &= \ln\frac{1}{2} + \frac{4 + 1^2}{2\times 1} - \frac{1}{2}
     = 2 - \ln 2 \approx \mathbf{1.3069}\;\text{nat}, \\
    \DKL\bigl(\mathcal{N}(0, 1) \,\|\, \mathcal{N}(1, 2^2)\bigr)
    &= \ln\frac{2}{1} + \frac{1 + (-1)^2}{2\times 4} - \frac{1}{2}
     = \ln 2 - \frac{1}{4} \approx \mathbf{0.4431}\;\text{nat}.
  \end{align*}
  \textbf{代价不对称解读}：用窄分布 $\mathcal{N}(0,1)$ 逼近宽分布 $\mathcal{N}(1, 2^2)$（前式）会漏掉宽分布尾部质量，二次惩罚除以窄方差 $2\sigma_2^2=2$ 被急剧放大（$1.3069$）；反过来用宽分布去罩住窄分布（后式），尾部已被充分覆盖，代价不到前者的三分之一（$0.4431$）。
\end{enumerate}

\subsection*{习题 4.5\quad 单高斯正向 KL 极小化的精确矩匹配定理}
\begin{proof}
\begin{enumerate}[label=(\arabic*), leftmargin=1.8em]
  \item 将正向 KL 散度拆成负微分熵与交叉熵之和：
  \[
    \DKL\bigl(P \,\|\, \mathcal{N}(\mu, \sigma^2)\bigr)
    = -h(P) + \E_{X\sim P}\bigl[-\ln q_{\mu,\sigma^2}(X)\bigr].
  \]
  因为 $-h(P)$ 与 $(\mu, \sigma^2)$ 无关，极小化正向 KL 等价于极小化交叉熵 $J(\mu, \sigma^2)$：
  \begin{align*}
    J(\mu, \sigma^2)
    &= \E_{X\sim P}\!\left[\frac{1}{2}\ln(2\pi\sigma^2) + \frac{(X - \mu)^2}{2\sigma^2}\right] \\
    &= \frac{1}{2}\ln(2\pi\sigma^2) + \frac{\E_P[(X - \mu_P + \mu_P - \mu)^2]}{2\sigma^2} \\
    &= \frac{1}{2}\ln(2\pi\sigma^2) + \frac{\sigma_P^2 + (\mu_P - \mu)^2}{2\sigma^2},
  \end{align*}
  其中交叉项 $2(\mu_P - \mu)\E_P[X - \mu_P] = 0$ 精确消去。
  \item 先对 $\mu$ 求偏导：
  \[
    \frac{\partial J}{\partial \mu} = -\frac{\mu_P - \mu}{\sigma^2} = 0 \implies \mu^* = \mu_P.
  \]
  事实上，对任意固定的 $\sigma^2 > 0$，$(\mu_P - \mu)^2 \ge 0$ 均在 $\mu^* = \mu_P$ 处唯一取到全局最小 $0$。将 $\mu = \mu_P$ 代回并记 $v := \sigma^2 > 0$，目标函数化为一元函数 $g(v) = \frac{1}{2}\ln(2\pi v) + \frac{\sigma_P^2}{2v}$：
  \[
    g'(v) = \frac{1}{2v} - \frac{\sigma_P^2}{2v^2} = \frac{v - \sigma_P^2}{2v^2}.
  \]
  在 $(0, \sigma_P^2)$ 上 $g'(v) < 0$ 严格递减，在 $(\sigma_P^2, +\infty)$ 上 $g'(v) > 0$ 严格递增，故 $v^* = (\sigma^*)^2 = \sigma_P^2$ 为唯一全局极小值点。证毕。
  \item 对双峰混合 $P = 0.5\mathcal{N}(-2, 0.6^2) + 0.5\mathcal{N}(2, 0.6^2)$，由全期望与二阶矩可加性立得：
  \begin{align*}
    \mu_{\mathrm{fwd}} &= \E_P[X] = 0.5(-2) + 0.5(2) = \mathbf{0}, \\
    \sigma_{\mathrm{fwd}}^2 &= \E_P[X^2] - \mu_{\mathrm{fwd}}^2
     = 0.5\bigl(0.6^2 + (-2)^2\bigr) + 0.5\bigl(0.6^2 + 2^2\bigr) = 0.36 + 4 = \mathbf{4.36},
  \end{align*}
  故 $\sigma_{\mathrm{fwd}} = \sqrt{4.36} \approx \mathbf{2.0881}$。
\end{enumerate}
\end{proof}

\subsection*{习题 4.6\quad 对称且有界的 Jensen--Shannon 散度}
\begin{proof}
\begin{enumerate}[label=(\arabic*), leftmargin=1.8em]
  \item 混合分布 $M = \frac{1}{2}(P + Q)$ 关于 $P, Q$ 对称，交换 $P, Q$ 仅交换两项顺序，故 $\mathrm{JSD}(P\|Q) = \mathrm{JSD}(Q\|P)$。由 Gibbs 不等式，$\DKL(P\|M) \ge 0$ 且 $\DKL(Q\|M) \ge 0$，故 $\mathrm{JSD}(P\|Q) \ge 0$。等号成立当且仅当 $\DKL(P\|M) = 0$ 且 $\DKL(Q\|M) = 0$，即 $P = M$ 且 $Q = M$，等价于 $P = Q$。
  \item 由全概率公式，$X$ 的边际分布恰好为 $P_X(x) = \frac{1}{2}P(x) + \frac{1}{2}Q(x) = M(x)$。条件熵为 $H(X \mid Z) = \frac{1}{2}H(X\mid Z=0) + \frac{1}{2}H(X\mid Z=1) = \frac{1}{2}H(P) + \frac{1}{2}H(Q)$。因此互信息满足：
  \begin{align*}
    I(X; Z) &= H(X) - H(X \mid Z) = H(M) - \frac{1}{2}H(P) - \frac{1}{2}H(Q) \\
    &= -\sum_x \frac{P(x)+Q(x)}{2}\ln M(x) + \frac{1}{2}\sum_x P(x)\ln P(x) + \frac{1}{2}\sum_x Q(x)\ln Q(x) \\
    &= \frac{1}{2}\sum_x P(x)\ln\frac{P(x)}{M(x)} + \frac{1}{2}\sum_x Q(x)\ln\frac{Q(x)}{M(x)}
     = \mathrm{JSD}(P \,\|\, Q).
  \end{align*}
  （连续情形将求和替换为积分完全同构。）
  \item 利用互信息对称性与离散熵非负性 $H(Z \mid X) \ge 0$：
  \[
    \mathrm{JSD}(P \,\|\, Q) = I(X; Z) = H(Z) - H(Z \mid X) \le H(Z) = \ln 2 \quad (\text{或 } 1\;\text{bit}).
  \]
  即使 $P, Q$ 支集完全不相交，仍有 $M(x) \ge \frac{1}{2}P(x)$ 且 $M(x) \ge \frac{1}{2}Q(x)$，此时由 $X$ 可 100\% 判别 $Z$（$H(Z\mid X)=0$），$\mathrm{JSD}(P\|Q)$ 精确达到上界 $\ln 2$ 而绝不发散。
\end{enumerate}
\end{proof}

\subsection*{习题 4.7\quad 安全信息论工具箱与决策树根节点信息增益选特征}
\textbf{核心纯 NumPy 实现逻辑}（用布尔掩码 \texttt{p > 0} 安全规避 $0\log 0$ 产生 \texttt{NaN}）：
\begin{lstlisting}[language=Python]
def entropy(p, base=2.0):
    p = np.asarray(p, dtype=float)
    m = p > 0
    return float(-np.sum(p[m] * np.log(p[m])) / np.log(base))

def kl_div(p, q, base=2.0):
    p, q = np.asarray(p, dtype=float), np.asarray(q, dtype=float)
    m = p > 0
    if np.any(q[m] <= 0):
        return np.inf
    return float(np.sum(p[m] * np.log(p[m] / q[m])) / np.log(base))

def info_gain_bits(y, feat):
    hy = entropy([np.mean(y == 0), np.mean(y == 1)], base=2.0)
    hy_f = sum(
        np.mean(feat == v) * entropy([np.mean(y[feat == v] == 0),
                                      np.mean(y[feat == v] == 1)], base=2.0)
        for v in (0, 1) if np.any(feat == v)
    )
    return hy - hy_f
\end{lstlisting}
\textbf{验证输出与断言自检}：
\begin{itemize}
  \item 第 (1) 问：100 对含零离散分布恒等式最大误差 $< 10^{-15} < 10^{-12}$，最小 KL 值 $\ge -10^{-15}$，零 \texttt{NaN}。
  \item 第 (2) 问：$H(y) = H_2(4/8) = 1.0000\;\text{bit}$；各特征精确经验信息增益：
  \begin{itemize}
    \item $F_1$（$F_1=1$ 中 3 正 1 负，$F_1=0$ 中 1 正 3 负）：$\mathrm{IG}(y, F_1) = 1 - H_2(0.75) = \mathbf{0.1887}\;\text{bit}$；
    \item $F_2$（$F_2=1, 0$ 中各 2 正 2 负，与标签完全独立）：$\mathrm{IG}(y, F_2) = 1 - 1 = \mathbf{0.0000}\;\text{bit}$；
    \item $F_3$（$F_3=1$ 中 4 正 1 负，$F_3=0$ 中 0 正 3 负）：$\mathrm{IG}(y, F_3) = 1 - \frac{5}{8}H_2(0.8) - \frac{3}{8}\cdot 0 = \mathbf{0.5488}\;\text{bit}$。
  \end{itemize}
  故决策树根节点最优分裂特征为 $\mathbf{F_3}$。
\end{itemize}

\subsection*{习题 4.8\quad 双峰分布正向 vs 反向 KL 数值寻优复现}
\textbf{核心纯 NumPy 实现逻辑}（在网格上直接评估精确高斯密度 $q(x;\mu,\sigma)$ 并用黎曼求和近似积分）：
\begin{lstlisting}[language=Python]
x = np.linspace(-6.0, 6.0, 601)
dx = x[1] - x[0]
npdf = lambda z, m, s: np.exp(-0.5*((z - m)/s)**2) / (np.sqrt(2*np.pi)*s)
p = 0.5 * npdf(x, -2.0, 0.6) + 0.5 * npdf(x, 2.0, 0.6)
p /= np.sum(p) * dx
for m in np.linspace(-2.5, 2.5, 101):
    for s in np.linspace(0.3, 2.8, 101):
        q = npdf(x, m, s)
        fwd = float(np.sum(p * np.log((p + 1e-15) / (q + 1e-15))) * dx)
        rev = float(np.sum(q * np.log((q + 1e-15) / (p + 1e-15))) * dx)
\end{lstlisting}
\textbf{验证输出与断言自检}：
\begin{itemize}
  \item \textbf{正向 KL 网格最优}：$\mu_{\mathrm{fwd}} = 0.000,\; \sigma_{\mathrm{fwd}} = 2.100$（与习题 4.5 精确解 $\sqrt{4.36}\approx 2.088$ 最邻近的步长 $0.025$ 格点），对应 $\DKL(P\|Q) = \mathbf{0.5552}$、$\DKL(Q\|P) = \mathbf{1.1806}$；
  \item \textbf{反向 KL 网格最优}：$\mu_{\mathrm{rev}} = -2.000$（由对称性 $+2.000$ 同样最优）、$\sigma_{\mathrm{rev}} = 0.600$，精准锁死单个峰，对应 $\DKL(Q\|P) = \mathbf{0.6919} \approx \ln 2$、而正向散度飙升至 $\DKL(P\|Q) = \mathbf{10.3062} > 9.5$。与表 \ref{tab:ch4-kl-bimodal} 及全部自检标准高度吻合。
\end{itemize}

% =========================================================================
\section{第 5 章习题精解（5.1--5.8）}
% =========================================================================

\subsection*{习题 5.1\quad 统计估计与置信区间判断正误并纠错}
\begin{enumerate}[label=(\alph*), leftmargin=1.8em]
  \item \textbf{错误。} 由偏差-方差分解（定理 \ref{thm:ch5-mse-decomp}），$\mathrm{MSE}(\hat{\theta}) = \Bias(\hat{\theta})^2 + \Var(\hat{\theta})$。无偏仅令第一项为 $0$，但若方差过大，总 MSE 反而劣于引入微小偏差换取大幅方差缩减的收缩估计量。\textbf{具体反例}：正态样本方差估计（见习题 5.7，$n=5,\;\sigma^2=4$），无偏 $S^2$ 的 $\mathrm{MSE}=8.0$，而有偏 MLE $\hat{\sigma}_n^2$ 的 $\mathrm{MSE}=5.76$、收缩估计 $\tilde{\sigma}^2$ 的 $\mathrm{MSE}=5.333$，均显著低于无偏估计量；另见习题 5.5 的均值收缩定理。
  \item \textbf{错误。} MLE 只具有\textbf{渐近无偏性}，有限样本下经常存在偏差。\textbf{两个经典反例}：(i) 正态总体方差的 MLE $\hat{\sigma}_n^2 = \frac{1}{n}\sum (X_i - \bar{X})^2$ 满足 $\E[\hat{\sigma}_n^2] = \frac{n-1}{n}\sigma^2 < \sigma^2$（系统性低估）；(ii) $\mathrm{Uniform}(0, \theta)$ 的 MLE $\hat{\theta}_{\mathrm{MLE}} = \max_i X_i$ 满足 $\E[\hat{\theta}_{\mathrm{MLE}}] = \frac{n}{n+1}\theta < \theta$（见习题 5.3(b)）。
  \item \textbf{错误。} 在 $n = 180$ 条验证样本上，$85\%$ 附近的单点准确率标准误约为 $\sqrt{0.85\times 0.15 / 180} \approx 2.66\%$，95\% 置信区间宽度约 $\pm 5.2\%$。两模型 $1.0\%$ 的差距（仅差约 2 道题）完全落在随机抽样波动区间内，必须使用配对 Bootstrap（如习题 5.8）或 McNemar 检验核查差值置信区间是否排除 $0$，绝不能仅凭点估计差 1 个百分点下结论。
\end{enumerate}

\subsection*{习题 5.2\quad Bootstrap 失效反例诊断}
\begin{enumerate}[label=(\alph*), leftmargin=1.8em]
  \item \textbf{3 次全正面 $\{1, 1, 1\}$ 上的 Bootstrap 坍缩}：由于原始经验分布 $\hat{F}_n$ 在 $x=1$ 处赋予了 $100\%$ 概率质量，无论运行多少轮有放回重采样，每一批 $B = 5{,}000$ 份重采样样本全都是 $\{1, 1, 1\}$，样本均值恒等于 $1.0$。因此百分位 Bootstrap 得到的 95\% 置信区间退化为单点区间 $\mathbf{[1.0, 1.0]}$（宽度为 $0$）！这严重违背常识——仅抛 3 次硬币根本不足以 95\% 确信这枚硬币绝不会出现反面。\\
  \textbf{Beta(2, 2) 先验的化解机制}：引入温和的对称先验 $p \sim \mathrm{Beta}(2, 2)$ 后，吸收 3 次全正面得到光滑后验 $p \mid \mathcal{D} \sim \mathrm{Beta}(5, 2)$，其 MAP 估计回拉至 $\frac{5-1}{5+2-2} = 0.80$、后验均值 $\frac{5}{7} \approx 0.7143$，且后验密度在 $(0, 1)$ 上保留了合理的尾部不确定性，彻底避免零样本支集塌缩。
  \item \textbf{均匀分布最大值 $T = \max_i X_i$ 的 Bootstrap 失效}：设原样本最大值为 $M_n = \max_{1\le i\le n} X_i$。由于任何有放回重采样都只能从 $\{X_1, \dots, X_n\}$ 中抽取元素，重采样最大值 $M_n^*$ 满足两个病态性质：(i) 恒有 $M_n^* \le M_n$，永远无法越过样本上界去探索 $(M_n, \theta]$；(ii) $M_n^*$ 恰好等于原最大值 $M_n$ 的概率为“至少抽中一次 $M_n$”的概率：
  \[
    \Prob^*(M_n^* = M_n) = 1 - \left(1 - \frac{1}{n}\right)^n \xrightarrow{n\to\infty} 1 - e^{-1} \approx 63.21\%!
  \]
  而真实抽样分布 $M_n/\theta$ 是连续分布，命中单点 $\theta$ 的概率为 $0$。Bootstrap 在端点处出现了 $63.2\%$ 的离散点质量，无法逼近真实边界极值分布，导致置信区间系统性失效。
\end{enumerate}

\subsection*{习题 5.3\quad 指数分布与均匀分布的 MLE 手算}
\begin{enumerate}[label=(\alph*), leftmargin=1.8em]
  \item \textbf{指数分布 $\mathrm{Exp}(\lambda)$ 的 MLE}：对数似然函数为
  \[
    \ell(\lambda) = \sum_{i=1}^n \ln\bigl(\lambda e^{-\lambda x_i}\bigr) = n\ln\lambda - \lambda\sum_{i=1}^n x_i.
  \]
  令一阶导数为零：$\ell'(\lambda) = \frac{n}{\lambda} - \sum_{i=1}^n x_i = 0 \implies \hat{\lambda}_{\mathrm{MLE}} = \frac{n}{\sum_{i=1}^n x_i} = \frac{1}{\bar{x}}$。又 $\ell''(\lambda) = -\frac{n}{\lambda^2} < 0$，故为唯一全局极大值点。\\
  当 $n=4$、观测为 $\{1, 2, 3, 2\}$ 时，$\bar{x} = \frac{1+2+3+2}{4} = 2.0$，故 $\hat{\lambda}_{\mathrm{MLE}} = \frac{1}{2} = \mathbf{0.5}$。
  \item \textbf{均匀分布 $\mathrm{Uniform}(0, \theta)$ 的 MLE}：联合似然函数为
  \[
    L(\theta) = \prod_{i=1}^n \frac{1}{\theta}\mathbf{1}_{[0,\theta]}(x_i)
              = \frac{1}{\theta^n}\mathbf{1}_{\{\theta \ge \max_i x_i\}}.
  \]
  \textbf{为何不能令 $\ell'(\theta) = 0$}：当 $\theta < M_n := \max_i x_i$ 时 $L(\theta) = 0$；当 $\theta \ge M_n$ 时，$L(\theta) = \theta^{-n}$ 是严格单调递减函数，其导数 $\ell'(\theta) = -n/\theta$ 处处严格为负、永不为零，极大值落在参数定义域的\textbf{边界点}（支集随待估参数 $\theta$ 变化，违反经典微积分内点极值条件），因此直接得：
  \[
    \hat{\theta}_{\mathrm{MLE}} = M_n = \max_{1\le i\le n} X_i.
  \]
  由给定密度 $f_{M_n}(t) = \frac{n t^{n-1}}{\theta^n}$（$0 \le t \le \theta$）计算期望：
  \[
    \E[\hat{\theta}_{\mathrm{MLE}}] = \int_0^\theta t\cdot\frac{n t^{n-1}}{\theta^n}\,dt
    = \frac{n}{\theta^n}\left[\frac{t^{n+1}}{n+1}\right]_0^\theta = \frac{n}{n+1}\theta.
  \]
  当 $n = 3$ 时，$\E[\hat{\theta}_{\mathrm{MLE}}] = \mathbf{\frac{3}{4}\theta}$，偏差为 $\Bias(\hat{\theta}_{\mathrm{MLE}}) = \frac{3}{4}\theta - \theta = \mathbf{-\frac{1}{4}\theta}$（系统性低估真值 $25\%$）。
\end{enumerate}

\subsection*{习题 5.4\quad Beta-Bernoulli 共轭的顺序更新性质}
已知先验 $p \sim \mathrm{Beta}(\alpha_0, \beta_0) = \mathrm{Beta}(2, 8)$。
\begin{enumerate}[label=(\alph*), leftmargin=1.8em]
  \item \textbf{第一批数据后}（$n_1 = 10$ 次曝光，$k_1 = 4$ 次点击，$n_1 - k_1 = 6$ 次未点击）：
  \begin{itemize}
    \item 后验分布：$p \mid \mathcal{D}_1 \sim \mathrm{Beta}(2 + 4,\; 8 + 6) = \mathbf{\mathrm{Beta}(6, 14)}$；
    \item $\hat{p}_{\mathrm{MLE}}^{(1)} = \frac{4}{10} = \mathbf{0.4000}$；
    \item $\hat{p}_{\mathrm{MAP}}^{(1)} = \frac{6 - 1}{6 + 14 - 2} = \frac{5}{18} \approx \mathbf{0.2778}$；
    \item 全贝叶斯后验均值：$\E[p \mid \mathcal{D}_1] = \frac{6}{6 + 14} = \frac{6}{20} = \mathbf{0.3000}$。
  \end{itemize}
  \item \textbf{顺序更新等价性验证}：
  \begin{itemize}
    \item \emph{分两批顺序更新}：以 $\mathrm{Beta}(6, 14)$ 为新先验，吸收第二批 $n_2 = 90$ 次曝光中的 $k_2 = 26$ 次点击与 $64$ 次未点击，得后验 $\mathrm{Beta}(6 + 26,\; 14 + 64) = \mathbf{\mathrm{Beta}(32, 78)}$；
    \item \emph{一次性合并更新}：以原始先验 $\mathrm{Beta}(2, 8)$ 吸收全部 $n = 100$ 次曝光中的 $k = 30$ 次点击与 $70$ 次未点击，得后验 $\mathrm{Beta}(2 + 30,\; 8 + 70) = \mathbf{\mathrm{Beta}(32, 78)}$，两者\textbf{完全一致}（充分统计量与共轭闭包的直接体现）。
  \end{itemize}
  在 $n = 100$ 时：$\hat{p}_{\mathrm{MLE}}^{(2)} = \frac{30}{100} = \mathbf{0.3000}$，$\hat{p}_{\mathrm{MAP}}^{(2)} = \frac{32 - 1}{32 + 78 - 2} = \frac{31}{108} \approx \mathbf{0.2870}$，后验均值 $\frac{32}{110} \approx \mathbf{0.2909}$。\\
  \textbf{MLE 与 MAP 差距变化}：$n=10$ 时 $|\text{MLE} - \text{MAP}| = 0.4000 - 0.2778 = \mathbf{0.1222}$；$n=100$ 时迅速缩窄近 10 倍至 $0.3000 - 0.2870 = \mathbf{0.0130}$。说明数据稀缺时先验强力纠偏，而随样本量增大，似然逐渐主导后验，MAP 向 MLE 收敛。
\end{enumerate}

\subsection*{习题 5.5\quad 收缩估计量以微偏换低 MSE（正则化前奏）}
\begin{proof}
\begin{enumerate}[label=(\arabic*), leftmargin=1.8em]
  \item 已知 $\E[\bar{X}] = \mu$、$\Var(\bar{X}) = \sigma^2/n$。对 $\hat{\mu}_\alpha = \alpha \bar{X}$：
  \[
    \Bias(\hat{\mu}_\alpha) = \E[\alpha\bar{X}] - \mu = (\alpha - 1)\mu, \qquad
    \Var(\hat{\mu}_\alpha) = \alpha^2 \Var(\bar{X}) = \alpha^2 \frac{\sigma^2}{n}.
  \]
  由均方误差分解定理（定理 \ref{thm:ch5-mse-decomp}）：
  \[
    \mathrm{MSE}(\hat{\mu}_\alpha) = \Bias(\hat{\mu}_\alpha)^2 + \Var(\hat{\mu}_\alpha)
    = (\alpha - 1)^2 \mu^2 + \alpha^2 \frac{\sigma^2}{n}.
  \]
  \item $f(\alpha) := \mathrm{MSE}(\hat{\mu}_\alpha)$ 是关于 $\alpha$ 的二次函数，二次项系数 $\mu^2 + \sigma^2/n > 0$（开口严格向上）。令一阶导数为零：
  \[
    f'(\alpha) = 2(\alpha - 1)\mu^2 + 2\alpha\frac{\sigma^2}{n} = 0
    \implies \alpha^* = \frac{\mu^2}{\mu^2 + \sigma^2/n} = \frac{n\mu^2}{n\mu^2 + \sigma^2}.
  \]
  因 $\sigma^2 > 0$ 且 $\mu \neq 0$，显然严格有 $0 < \alpha^* < 1$。又因 $\alpha = 1$ 对应无偏估计 $\bar{X}$ 且 $\alpha^* \neq 1$ 是严格凸二次函数的唯一全局极小点，必严格满足：
  \[
    \mathrm{MSE}(\hat{\mu}_{\alpha^*}) = \frac{\mu^2 \cdot \sigma^2/n}{\mu^2 + \sigma^2/n}
    = \alpha^* \cdot \frac{\sigma^2}{n} < \frac{\sigma^2}{n} = \mathrm{MSE}(\bar{X}).
  \]
  \item \textbf{与高斯先验 MAP 的联系}：对比式 \eqref{eq:ch5-gauss-conj}，当置零均值先验 $\mu \sim \mathcal{N}(0, \tau^2)$ 时，MAP 估计恰好是线性收缩 $\hat{\mu}_{\mathrm{MAP}} = \frac{n}{n + \sigma^2/\tau^2}\bar{X}$；其收缩因子 $\frac{n\tau^2}{n\tau^2 + \sigma^2}$ 与上述最优因子 $\frac{n\mu^2}{n\mu^2 + \sigma^2}$ 形式完全一致（先验方差 $\tau^2$ 正是对未知真值信号强度 $\mu^2$ 的替代）。这揭示了 $L_2$ 正则化（岭回归）提升泛化的本质：主动接受微小向零偏差，换取更大幅度的方差下降。
\end{enumerate}
\end{proof}

\subsection*{习题 5.6\quad 泊松-伽马共轭与多元高斯均值 MAP 推导}
\begin{proof}
\begin{enumerate}[label=(\arabic*), leftmargin=1.8em]
  \item \textbf{泊松-伽马共轭}：合并先验与 $n$ 条独立泊松观测的似然，仅保留含 $\lambda$ 的因子：
  \begin{align*}
    p(\lambda \mid x_1, \dots, x_n)
    &\propto p(\lambda)\prod_{i=1}^n p(x_i; \lambda)
     \propto \Bigl(\lambda^{\alpha - 1} e^{-\beta\lambda}\Bigr)
             \prod_{i=1}^n \Bigl(e^{-\lambda}\lambda^{x_i}\Bigr) \\
    &= \lambda^{\bigl(\alpha + \sum_{i=1}^n x_i\bigr) - 1} \exp\!\bigl(-(\beta + n)\lambda\bigr).
  \end{align*}
  对照 Gamma 密度核 $\lambda^{\alpha'-1}e^{-\beta'\lambda}$，立得后验仍为 Gamma 分布：
  \[
    \lambda \mid \mathcal{D} \;\sim\; \mathrm{Gamma}\!\left(\alpha',\; \beta'\right), \qquad
    \alpha' = \alpha + \sum_{i=1}^n x_i,\quad \beta' = \beta + n.
  \]
  由此后验均值为 $\E[\lambda \mid \mathcal{D}] = \frac{\alpha'}{\beta'} = \frac{\alpha + \sum_{i=1}^n x_i}{\beta + n}$；对对数后验 $(\alpha' - 1)\ln\lambda - \beta'\lambda$ 求导置零，当 $\alpha' = \alpha + \sum x_i > 1$ 时唯一 MAP 为 $\hat{\lambda}_{\mathrm{MAP}} = \frac{\alpha + \sum_{i=1}^n x_i - 1}{\beta + n}$。
  \item \textbf{多元高斯均值 MAP}：写出对数后验并忽略与 $\mu$ 无关的归一化常数：
  \[
    \ln p(\mu \mid \mathcal{D})
    = -\frac{1}{2\sigma^2}\sum_{i=1}^n \|x_i - \mu\|_2^2 - \frac{1}{2\tau^2}\|\mu\|_2^2 + \mathrm{const}.
  \]
  对向量 $\mu \in \R^d$ 求梯度并令其为零（注意 Hessian 为 $-(\frac{n}{\sigma^2} + \frac{1}{\tau^2})I_d \prec 0$，保证唯一全局极大）：
  \[
    \nabla_\mu \ln p(\mu \mid \mathcal{D})
    = \frac{1}{\sigma^2}\sum_{i=1}^n (x_i - \mu) - \frac{1}{\tau^2}\mu = 0
    \implies \left(\frac{n}{\sigma^2} + \frac{1}{\tau^2}\right)\mu = \frac{n}{\sigma^2}\bar{x},
  \]
  两边同乘 $\sigma^2$ 并整理即得 $\hat{\mu}_{\mathrm{MAP}} = \frac{n}{n + \sigma^2/\tau^2}\bar{x}$。
\end{enumerate}
\end{proof}

\subsection*{习题 5.7\quad 蒙特卡洛实测三种方差估计量的偏差与 MSE}
\textbf{理论推导速核}：因 $SS := \sum_{i=1}^n (X_i - \bar{X})^2$ 满足 $SS/\sigma^2 \sim \chi^2_{n-1}$（自由度 $n-1=4$，均值 $4$，方差 $8$），对一般形式 $c\cdot SS$ 有 $\E = 16c,\; \Var = 128c^2$，代入 $c \in \{1/4, 1/5, 1/6\}$ 立得精确理论 MSE：$\mathrm{MSE}(S^2)=8.000$、$\mathrm{MSE}(\hat{\sigma}_n^2)=144/25=5.760$、$\mathrm{MSE}(\tilde{\sigma}^2)=16/3\approx 5.333$。\\
\textbf{核心纯 NumPy 实现逻辑}：
\begin{lstlisting}[language=Python]
rng = np.random.default_rng(2026)
X = rng.normal(0.0, 2.0, size=(20_000, 5))
ss = np.sum((X - X.mean(axis=1, keepdims=True)) ** 2, axis=1)
s_unb, s_mle, s_shr = ss / 4.0, ss / 5.0, ss / 6.0
for name, est in [("S^2", s_unb), ("MLE", s_mle), ("Shrink", s_shr)]:
    print(name, est.mean(), est.mean() - 4.0, np.mean((est - 4.0) ** 2))
\end{lstlisting}
\textbf{验证输出与断言自检（\texttt{default\_rng(2026)}）}：
\begin{itemize}
  \item 无偏 $S^2$（$/4$）：均值 $3.9855$（误差 $0.0145 < 0.08$），偏差 $-0.0145$，经验 MSE $= \mathbf{7.9151}$（理论 $8.000$）；
  \item MLE $\hat{\sigma}_n^2$（$/5$）：均值 $3.1884$（距 $3.2$ 误差 $0.0116 < 0.08$），偏差 $-0.8116$，经验 MSE $= \mathbf{5.7242}$（理论 $5.760$）；
  \item 收缩版 $\tilde{\sigma}^2$（$/6$）：均值 $2.6570$（理论 $2.6667$），偏差 $-1.3430$，经验 MSE $= \mathbf{5.3214}$（理论 $5.333$）。
\end{itemize}
严格满足 $\mathrm{MSE}(\tilde{\sigma}^2) < \mathrm{MSE}(\hat{\sigma}_n^2) < \mathrm{MSE}(S^2)$，全部断言通过。

\subsection*{习题 5.8\quad 实现配对 Bootstrap 比较两个分类器}
\textbf{核心纯 NumPy 实现逻辑}：
\begin{lstlisting}[language=Python]
def paired_bootstrap(both_ok, only_a, only_b, both_bad, B=10_000, seed=2026):
    rng = np.random.default_rng(seed)
    hits_a = np.array([1]*both_ok + [1]*only_a + [0]*only_b + [0]*both_bad)
    hits_b = np.array([1]*both_ok + [0]*only_a + [1]*only_b + [0]*both_bad)
    idx = rng.integers(0, len(hits_a), size=(B, len(hits_a)))
    diffs = (hits_b[idx].mean(axis=1) - hits_a[idx].mean(axis=1)) * 100.0
    lo, hi = np.percentile(diffs, [2.5, 97.5])
    return float(lo), float(hi), float(np.mean(diffs > 0))
\end{lstlisting}
\textbf{验证输出与断言自检（\texttt{default\_rng(2026)}，$B=10{,}000$）}：
\begin{itemize}
  \item \textbf{情形 (a) 微弱差异}（点估计 $\Delta = +1.11\%$）：配对 95\% CI 为 $\mathbf{[-1.67\%, +3.89\%]}$，跨越 $0$，且 $\Prob^*(\Delta > 0) = 0.7278$，统计上\textbf{不能}宣称模型 B 显著优于模型 A；
  \item \textbf{情形 (b) 显著差异}（点估计 $\Delta = +6.67\%$）：配对 95\% CI 为 $\mathbf{[+2.78\%, +11.11\%]}$，下界严格大于 $0$，且 $\Prob^*(\Delta > 0) = \mathbf{0.9989} > 0.99$，可确认模型 B 具有统计显著优势。
\end{itemize}

% =========================================================================
\section{第 6 章习题精解（6.1--6.8）}
% =========================================================================

\subsection*{习题 6.1\quad 逻辑回归与评测工程判断正误并纠错}
\begin{enumerate}[label=(\alph*), leftmargin=1.8em]
  \item \textbf{错误。} 工程与理论中通常\textbf{不惩罚偏置项 $b$}。原因是偏置项仅控制整体输出基线（例如回归中的目标均值 $\bar{y}$，或分类中的类别先验对数几率），不增加模型随特征变化的复杂度；若惩罚 $b$，在目标值整体平移（如房价从万元改千元）或类别极度不平衡时，会把模型基准强行拉向 $0$，破坏平移不变性并引入无谓的系统偏差。
  \item \textbf{错误。} 当训练数据完全线性可分时，存在某个方向 $u$ 使全部样本满足 $(2y_i - 1)x_i^T u > 0$。沿射线 $w(t) = t u$（$t \to +\infty$），不加正则的交叉熵损失严格单调趋近于 $0$ 但永不可达（$\inf_w \mathcal{L}(w) = 0$，无有限极小值点），因此梯度下降会使 $\|w\|_2 \to +\infty$ 发生\textbf{权重发散}而非收敛到有限最优解！加入任意微小 $L_2$ 正则项（$\lambda > 0$）才能使目标严格强凸并保证唯一有限最优解。
  \item \textbf{错误。} 在仅有 1\% 欺诈交易的极度不平衡数据上，一个什么都不学、永远预测“正常交易（$\hat{y}\equiv 0$）”的废柴模型就能拿到 \textbf{99.0\% 的准确率}，而它对真正欺诈的召回率 Recall $= 0\%$、$F_1 = 0$，毫无商用价值。必须结合 Precision、Recall、$F_1$ 或 PR-AUC/ROC-AUC 综合评估。
\end{enumerate}

\subsection*{习题 6.2\quad Softmax 数值溢出与数据泄漏双排查}
\begin{enumerate}[label=(\roman*), leftmargin=1.8em]
  \item \textbf{Softmax 与交叉熵的数值崩盘}：
  \begin{itemize}
    \item \textbf{崩盘场景}：只要某条样本的 logit 较大（如 float64 下 $z_{ik} > 710$），\texttt{np.exp(Z)} 立刻上溢为 \texttt{inf}，进而 \texttt{inf / inf} 变为 \texttt{NaN}；反之，若预测正确类的概率 $\hat{p}_{i, y_i}$ 因极端错分下溢为 \texttt{0.0}，随后 \texttt{np.log(0.0)} 产生 \texttt{-inf} 甚至 \texttt{NaN}，直接炸毁训练梯度。
    \item \textbf{修复写法}（利用平移不变性与 Log-Sum-Exp 直接在 log 空间算交叉熵）：
  \end{itemize}
\begin{lstlisting}[language=Python]
Z_shift = Z - np.max(Z, axis=1, keepdims=True)
lse = np.max(Z, axis=1) + np.log(np.sum(np.exp(Z_shift), axis=1))
loss = float(np.mean(lse - Z[np.arange(len(y)), y]))
P = np.exp(Z_shift) / np.sum(np.exp(Z_shift), axis=1, keepdims=True)
\end{lstlisting}
  \item \textbf{全量预处理与特征选择导致的数据泄漏}：
  \begin{itemize}
    \item \textbf{幻觉场景}：先在全量 \texttt{X\_all} 上计算均值/标准差并利用全量标签 \texttt{y\_all} 挑选 Top-20 相关特征，等价于让模型在训练前偷看了测试集分布与标签！在高维小样本下，即使所有特征全是纯随机噪声，被挑出的 20 维特征也会与测试标签存在虚假相关，制造高达 $80\%$ 以上的虚假高准确率幻觉（见习题 6.8(a)）。
    \item \textbf{修复写法}（严格\textbf{先切分，再仅用训练集拟合标准化参数与特征索引}）：
  \end{itemize}
\begin{lstlisting}[language=Python]
# 1. Split first
X_tr_raw, y_tr = X_all[tr_idx], y_all[tr_idx]
X_te_raw, y_te = X_all[te_idx], y_all[te_idx]
# 2. Fit scaler & feature selector ONLY on training fold
mu, sd = X_tr_raw.mean(axis=0), X_tr_raw.std(axis=0) + 1e-8
X_tr_std, X_te_std = (X_tr_raw - mu) / sd, (X_te_raw - mu) / sd
corrs_tr = np.abs(X_tr_std.T @ (y_tr - y_tr.mean()))
top20 = np.argsort(corrs_tr)[-20:]
X_tr_final, X_te_final = X_tr_std[:, top20], X_te_std[:, top20]
\end{lstlisting}
\end{enumerate}

\subsection*{习题 6.3\quad 逻辑回归一步梯度更新与远点拖拽验算}
\begin{enumerate}[label=(\alph*), leftmargin=1.8em]
  \item \textbf{单步梯度下降精确手算}：已知 $x=(1, -2)^T,\; y=1,\; w_0=(0,0)^T,\; b_0=0,\; \lambda=0.2,\; \eta=0.5$。
  \begin{itemize}
    \item 初始前向：$z_0 = x^T w_0 + b_0 = \mathbf{0}$，$\hat{p}_0 = \sigma(0) = \mathbf{0.5}$；
    \item 带正则 BCE 损失：$\mathcal{L}(w_0, b_0) = \ln(1 + e^0) - 1\cdot 0 + \frac{0.2}{2}\|0\|_2^2 = \mathbf{\ln 2 \approx 0.6931}$；
    \item 梯度计算（不惩罚偏置 $b$）：
    \[
      \nabla_w \mathcal{L} = (\hat{p}_0 - y)x + \lambda w_0 = (-0.5)\begin{pmatrix} 1 \\ -2 \end{pmatrix} + 0 = \begin{pmatrix} \mathbf{-0.5} \\ \mathbf{1.0} \end{pmatrix}, \qquad
      \frac{\partial \mathcal{L}}{\partial b} = \hat{p}_0 - y = \mathbf{-0.5};
    \]
    \item 一步更新（$\eta = 0.5$）：
    \[
      w_1 = w_0 - 0.5\begin{pmatrix} -0.5 \\ 1.0 \end{pmatrix} = \begin{pmatrix} \mathbf{0.25} \\ \mathbf{-0.50} \end{pmatrix}, \qquad
      b_1 = 0 - 0.5(-0.5) = \mathbf{0.25};
    \]
    \item 更新后预测：$z_1 = x^T w_1 + b_1 = 1(0.25) + (-2)(-0.50) + 0.25 = \mathbf{1.50}$，新预测概率 $\hat{p}_1 = \sigma(1.50) = \frac{1}{1 + e^{-1.5}} \approx \mathbf{0.8176}$（朝真实标签 $y=1$ 大幅推进）。
  \end{itemize}
  \item \textbf{7 点最小二乘远点拖拽精确验算}：数据点为 $(-2,0), (-1,0), (0,0), (1,1), (2,1), (3,1), (30,1)$，$n=7$。
  \begin{align*}
    \sum_{i=1}^7 x_i &= -2 - 1 + 0 + 1 + 2 + 3 + 30 = 33 \implies \bar{x} = \mathbf{\frac{33}{7}}, \\
    \sum_{i=1}^7 y_i &= 0 + 0 + 0 + 1 + 1 + 1 + 1 = 4 \implies \bar{y} = \mathbf{\frac{4}{7}}.
  \end{align*}
  利用平方和分解 $\sum x_i^2 = 4 + 1 + 0 + 1 + 4 + 9 + 900 = 919$ 与 $\sum x_i y_i = 1 + 2 + 3 + 30 = 36$：
  \begin{align*}
    S_{xx} &= \sum_{i=1}^7 (x_i - \bar{x})^2 = \sum x_i^2 - n\bar{x}^2
            = 919 - 7\cdot\frac{1089}{49} = \frac{6433 - 1089}{7} = \mathbf{\frac{5344}{7}}, \\
    S_{xy} &= \sum_{i=1}^7 (x_i - \bar{x})(y_i - \bar{y}) = \sum x_i y_i - n\bar{x}\bar{y}
            = 36 - 7\cdot\frac{33}{7}\cdot\frac{4}{7} = \frac{252 - 132}{7} = \mathbf{\frac{120}{7}}.
  \end{align*}
  由此斜率与截距精确为：
  \[
    w = \frac{S_{xy}}{S_{xx}} = \frac{120}{5344} = \mathbf{\frac{15}{668}}, \qquad
    b = \bar{y} - w\bar{x} = \frac{4}{7} - \frac{15}{668}\cdot\frac{33}{7}
      = \frac{2672 - 495}{7\times 668} = \frac{2177}{7\times 668} = \mathbf{\frac{311}{668}}.
  \]
  令 $w x_1^* + b = \frac{1}{2} = \frac{334}{668}$，解得 $\frac{15}{668}x_1^* = \frac{23}{668} \implies x_1^* = \mathbf{\frac{23}{15} \approx 1.5333}$，导致正类点 $(1, 1)$ 被右移后的阈值边界错误判负！
\end{enumerate}

\subsection*{习题 6.4\quad 不平衡混淆矩阵指标手算与阈值权衡}
已知总样本 $N = 200$，正类 $P = 20$，负类 $N_- = 180$。
\begin{enumerate}[label=(\alph*), leftmargin=1.8em]
  \item 当 $\mathrm{TP}=16,\; \mathrm{FP}=4,\; \mathrm{FN}=4,\; \mathrm{TN}=176$ 时：
  \begin{align*}
    \mathrm{Accuracy} &= \frac{16 + 176}{200} = \mathbf{96.00\%}, &
    \mathrm{Precision} &= \frac{16}{16 + 4} = \mathbf{80.00\%}, \\
    \mathrm{Recall} &= \frac{16}{16 + 4} = \mathbf{80.00\%}, &
    \mathrm{FPR} &= \frac{4}{4 + 176} = \frac{1}{45} \approx \mathbf{2.22\%}, \\
    F_1 &= \frac{2\times 0.80\times 0.80}{0.80 + 0.80} = \mathbf{0.8000}.
  \end{align*}
  \item \textbf{对比偷懒基线（$\hat{y}\equiv 0$）}：$\mathrm{TP}=\mathrm{FP}=0,\; \mathrm{FN}=20,\; \mathrm{TN}=180$，其 $\mathrm{Accuracy} = 180/200 = \mathbf{90.00\%}$（看似只低 6 个百分点），但 $\mathrm{Recall} = 0\%$、$F_1 = \mathbf{0.0000}$，彻底暴露准确率在不平衡场景下的钝化失真。
  \item \textbf{调低阈值后}（$\mathrm{TP}=19,\; \mathrm{FP}=19,\; \mathrm{FN}=1,\; \mathrm{TN}=161$）：
  \[
    \mathrm{Precision} = \frac{19}{38} = \mathbf{50.00\%}, \qquad
    \mathrm{Recall} = \frac{19}{20} = \mathbf{95.00\%}, \qquad
    F_1 = \frac{2\times 0.50\times 0.95}{0.50 + 0.95} = \frac{19}{29} \approx \mathbf{0.6552}.
  \]
  \textbf{场景决策原则}：(i) \textbf{癌症初筛}中漏诊（FN）的代价是危及生命，误报（FP）仅需复查确认，应\textbf{调低阈值}优先拉高 Recall（如选 $95\%$ 召回方案）；(ii) \textbf{自动封号系统}中误封正常用户（FP）严重损害口碑与客诉成本，应\textbf{调高阈值}优先守住高 Precision（如 $80\%$ 甚至 $95\%+$ 精确率，可疑边缘账号转人工复审）。
\end{enumerate}

\subsection*{习题 6.5\quad Softmax 平移不变性、交叉熵梯度与 $K=2$ 退化证明}
\begin{proof}
\begin{enumerate}[label=(\alph*), leftmargin=1.8em]
  \item \textbf{平移不变性}：对任意常数 $c \in \R$ 与任意类别 $k \in \{1,\dots,K\}$：
  \[
    \mathrm{softmax}(z + c\mathbf{1})_k
    = \frac{e^{z_k + c}}{\sum_{m=1}^K e^{z_m + c}}
    = \frac{e^c e^{z_k}}{e^c \sum_{m=1}^K e^{z_m}}
    = \frac{e^{z_k}}{\sum_{m=1}^K e^{z_m}} = \mathrm{softmax}(z)_k.
  \]
  取 $c = -\max_m z_m$ 即得到数值稳定的 Log-Sum-Exp 与防上溢 Softmax。
  \item \textbf{交叉熵梯度}：因 one-hot 标签满足 $\sum_{k=1}^K y_k = 1$，将损失函数化简为：
  \[
    \ell(z, y) = -\sum_{k=1}^K y_k \ln\frac{e^{z_k}}{\sum_{c=1}^K e^{z_c}}
    = \left(\sum_{k=1}^K y_k\right)\ln\!\left(\sum_{c=1}^K e^{z_c}\right) - \sum_{k=1}^K y_k z_k
    = \ln\!\left(\sum_{c=1}^K e^{z_c}\right) - \sum_{k=1}^K y_k z_k.
  \]
  对任意 $z_j$ 直接求偏导，立得极为优雅的“预测减标签”形式：
  \[
    \frac{\partial \ell}{\partial z_j}
    = \frac{e^{z_j}}{\sum_{c=1}^K e^{z_c}} - y_j = \hat{p}_j - y_j.
  \]
  \item \textbf{$K=2$ 精确退化为二元逻辑回归}：设二分类（类别 $0, 1$）的 Softmax 参数为 $(w_0, b_0)$ 与 $(w_1, b_1)$，由 (a) 的平移不变性，将两个 logit 同时减去 $w_0^T x + b_0$，并记 $w := w_1 - w_0,\; b := b_1 - b_0$，则类别 $1$ 的后验概率为：
  \[
    \hat{p}_1 = \frac{e^{w_1^T x + b_1}}{e^{w_0^T x + b_0} + e^{w_1^T x + b_1}}
              = \frac{e^{(w_1-w_0)^T x + (b_1-b_0)}}{1 + e^{(w_1-w_0)^T x + (b_1-b_0)}}
              = \frac{1}{1 + e^{-(w^T x + b)}} = \sigma(w^T x + b),
  \]
  相应地 $\hat{p}_0 = 1 - \hat{p}_1$，且多分类交叉熵精确化为二元 BCE 损失。
\end{enumerate}
\end{proof}

\subsection*{习题 6.6\quad Sigmoid + MSE 非凸区间与逻辑回归 Hessian 安全学习率上界}
\begin{proof}
\begin{enumerate}[label=(\alph*), leftmargin=1.8em]
  \item 记 $p = \sigma(z) \in (0, 1)$，利用 $p' = p(1 - p)$，对 $\ell_{\mathrm{MSE}}(z) = \frac{1}{2}(p - 1)^2$ 连续求导：
  \begin{align*}
    \ell'_{\mathrm{MSE}}(z) &= (p - 1)p' = -p(1 - p)^2, \\
    \ell''_{\mathrm{MSE}}(z) &= -\bigl[p'(1 - p)^2 + p\cdot 2(1 - p)(-p')\bigr]
    = -p'(1 - p)\bigl[(1 - p) - 2p\bigr] \\
    &= p(1 - p)^2(3p - 1).
  \end{align*}
  因为 $p(1 - p)^2 > 0$ 处处恒正，$\ell''_{\mathrm{MSE}}(z)$ 的符号完全由 $(3p - 1)$ 决定：
  \[
    \ell''_{\mathrm{MSE}}(z) < 0 \iff 3p - 1 < 0 \iff p = \frac{1}{1 + e^{-z}} < \frac{1}{3}
    \iff e^{-z} > 2 \iff z < -\ln 2.
  \]
  由此证得当 $y = 1$ 且模型严重错判（$z < -\ln 2 \approx -0.6931$）时，Sigmoid + MSE 损失处于严格凹区且梯度趋近于零，产生梯度消失与非凸平台。
  \item 对带 $L_2$ 正则的无偏置逻辑回归 $\mathcal{L}(w) = \frac{1}{N}\sum_{i=1}^N [\ln(1+e^{x_i^T w}) - y_i x_i^T w] + \frac{\lambda}{2}\|w\|_2^2$，其梯度为 $\nabla \mathcal{L}(w) = \frac{1}{N}\sum_{i=1}^N (\hat{p}_i - y_i)x_i + \lambda w$，Hessian 矩阵为：
  \[
    H(w) = \frac{1}{N}\sum_{i=1}^N \hat{p}_i(1 - \hat{p}_i)\,x_i x_i^T + \lambda I_d.
  \]
  任取单位向量 $v \in \R^d$（$\|v\|_2 = 1$），计算二次型 $v^T H(w) v$：
  \[
    v^T H(w) v = \frac{1}{N}\sum_{i=1}^N \hat{p}_i(1 - \hat{p}_i)(x_i^T v)^2 + \lambda.
  \]
  由基本不等式 $0 < \hat{p}_i(1 - \hat{p}_i) \le \frac{1}{4}$（当 $\hat{p}_i = 1/2$ 取等）与第 1 卷 Rayleigh 商上界 $\sum_{i=1}^N (x_i^T v)^2 = v^T (X^T X) v \le \lambda_{\max}(X^T X)$：
  \[
    \lambda \;\le\; v^T H(w) v \;\le\; \frac{1}{4N}v^T (X^T X) v + \lambda \;\le\; \frac{1}{4N}\lambda_{\max}(X^T X) + \lambda.
  \]
  因 $v$ 任意，即得矩阵不等式 $\lambda I_d \preceq H(w) \preceq L I_d$，其中全局平滑常数为 $L := \frac{1}{4N}\lambda_{\max}(X^T X) + \lambda$。由第 1 卷第 6 章 $L$-平滑凸函数梯度下降收敛定理，保证全批量梯度下降单调下降并稳定收敛的安全学习率上界为：
  \[
    0 < \eta < \eta_{\max} = \frac{2}{L} = \frac{2}{\frac{1}{4N}\lambda_{\max}(X^T X) + \lambda}
    \quad \Bigl(\text{工程推荐保守步长 }\eta \le \frac{1}{L}\Bigr).
  \]
\end{enumerate}
\end{proof}

\subsection*{习题 6.7\quad 纯 NumPy 多分类 Softmax 回归 + Log-Sum-Exp + 有限差分梯度校验}
\textbf{核心纯 NumPy 实现逻辑}：
\begin{lstlisting}[language=Python]
def logsumexp(Z, axis=-1):
    m = np.max(Z, axis=axis, keepdims=True)
    return np.squeeze(m + np.log(np.sum(np.exp(Z - m), axis=axis, keepdims=True)), axis=axis)

def softmax(Z):
    e = np.exp(Z - np.max(Z, axis=-1, keepdims=True))
    return e / np.sum(e, axis=-1, keepdims=True)

def loss_and_grad(theta, X, y, K, lam):
    n, d = X.shape
    W, b = theta[:d*K].reshape(d, K), theta[d*K:]
    Z = X @ W + b
    loss = float(np.mean(logsumexp(Z, axis=1) - Z[np.arange(n), y]) + 0.5*lam*np.sum(W*W))
    P = softmax(Z)
    P[np.arange(n), y] -= 1.0
    return loss, np.concatenate([((X.T @ P) / n + lam * W).ravel(), P.mean(axis=0)])
\end{lstlisting}
\textbf{验证输出与断言自检（\texttt{default\_rng(2026)}，$\lambda = 0.05$，800 轮 GD，$\eta = 0.25$）}：
\begin{itemize}
  \item \textbf{有限差分梯度校验}（$\epsilon = 10^{-5}$）：整体参数 $(W, b)$ 相对误差 $= \mathbf{1.76 \times 10^{-11}} < 10^{-9}$；
  \item \textbf{极端 logits 稳定性测试}：\texttt{[1000.0, 1001.0, 999.0]} 下 \texttt{softmax} 输出 \texttt{[0.2447, 0.6652, 0.0900]}（和与 $1.0$ 误差 $< 10^{-15}$），\texttt{logsumexp} 输出 $1001.4076$，零 \texttt{NaN}/\texttt{inf}；
  \item \textbf{验证集分类准确率}：在 90 条分层验证集上精确命中 \textbf{92.22\%}（$83/90 > 90\%$），全部断言通过。
\end{itemize}

\subsection*{习题 6.8\quad 复现数据泄漏灾难并手写 ROC-AUC 阈值扫描}
\textbf{核心纯 NumPy ROC-AUC 梯形积分实现逻辑}：
\begin{lstlisting}[language=Python]
order = np.argsort(-scores)
y_s = y_va[order]
tps, fps = np.cumsum(y_s == 1), np.cumsum(y_s == 0)
tpr = np.concatenate([[0.0], tps / tps[-1]])
fpr = np.concatenate([[0.0], fps / fps[-1]])
roc_auc = float(np.sum((fpr[1:] - fpr[:-1]) * (tpr[1:] + tpr[:-1]) * 0.5))
\end{lstlisting}
\textbf{验证输出与断言自检}：
\begin{enumerate}[label=(\alph*), leftmargin=1.8em]
  \item \textbf{纯噪声数据泄漏灾难}（$N=60, D=500$，50 次重复，真实 Bayes 上界 $50\%$）：泄漏版平均测试准确率高达 \textbf{80.40\%}（$>75\%$），干净版平均测试准确率回落至 \textbf{49.50\%}（落入 $[45\%, 55\%]$），无可辩驳地证明“测试前全量选特征”凭空伪造了 30 个百分点的幻觉性能。
\par\pagebreak[4]
  \item \textbf{30 维逻辑回归（$\lambda = 0.10$）验证集 ROC-AUC 与阈值权衡实测}：验证集精确 $\mathrm{ROC\text{-}AUC} = \mathbf{0.9160}$（落入 $[0.90, 0.95]$）；各阈值下的 Precision、Recall 与 $F_1$ 实测如下表所示（$\tau=0.5$ 时精确命中 $\mathrm{Acc}=85.00\%,\; F_1=0.8439$，而 $\tau=0.4$ 处 $F_1$ 略高至 $0.8482$）：
\end{enumerate}

\begin{center}
\small
\begin{tabular}{@{}ccccc@{}}
\toprule
分类阈值 $\tau$ & Accuracy & Precision & Recall & $F_1$ 分数 \\
\midrule
$0.2$ & $70.00\%$ & $62.94\%$ & $98.90\%$ & $0.7692$ \\
$0.3$ & $80.56\%$ & $72.95\%$ & $97.80\%$ & $0.8357$ \\
$0.4$ & $83.89\%$ & $81.00\%$ & $89.01\%$ & $\mathbf{0.8482}$ \\
$0.5$ & $\mathbf{85.00\%}$ & $89.02\%$ & $80.22\%$ & $0.8439$ \\
$0.6$ & $80.56\%$ & $90.00\%$ & $69.23\%$ & $0.7826$ \\
$0.7$ & $74.44\%$ & $\mathbf{92.45\%}$ & $53.85\%$ & $0.6806$ \\
$0.8$ & $67.22\%$ & $92.11\%$ & $38.46\%$ & $0.5426$ \\
\bottomrule
\end{tabular}
\end{center}


\end{document}
